\documentclass[12pt,twoside, openany]{book}

% IMPORTANT: Compile with XeLaTeX on Overleaf

% For book dimensions
\usepackage[
  paperwidth=6in,
  paperheight=9in,
  inner=0.8in,      % Margin near the spine
  outer=0.6in,      % Margin near the edge
  top=0.8in,
  bottom=0.8in,
  includeheadfoot   % Ensures headers/footers stay inside the margins
]{geometry}

% Prevent text overflow
\sloppy

\usepackage[strings]{underscore}

% Dotted leaders for chapters
\usepackage{tocloft}
\renewcommand{\cftdotsep}{1.5}
\renewcommand{\cftchapleader}{\cftdotfill{\cftdotsep}}

% \renewcommand{\cftchapfont}{\normalfont}
\renewcommand{\cftchappagefont}{\normalfont}

% Indentation settings
\cftsetindents{chapter}{0em}{2em}

% For no Κεφάλαιο Ν before the title
\usepackage{titlesec}

\titleformat{\chapter}
  {\normalfont\LARGE\filcenter} 
  {\thechapter}
  {1em}
  {}

\titlespacing*{\chapter}
  {0pt}   % left margin
  {1cm}   % space before chapter title
  {4.5cm}   % space after chapter title

% Language support
\usepackage{polyglossia}
\setdefaultlanguage{greek}
\setotherlanguage{english}

% Font support 
\setmainfont{Georgia}
\newfontfamily\greekfont{Georgia}
\newfontfamily\englishfont{Georgia}

% Typography improvements
\usepackage{microtype}
\usepackage{setspace}
\usepackage{csquotes}
\usepackage{emptypage}

\begin{document}

\frontmatter

\title{
\Huge ΣΠΑΡΚ ΘΟΤΣ\\[0.5em]
\large\itshape Μια ιστορία για υπεραναλυτές\\
\normalsize\itshape (A story for overthinkers)
}
\author{Γιώργος Τόλης}
\date{}
\mainmatter


\maketitle

\cleardoublepage

\tableofcontents

\cleardoublepage

\chapter{}

«Δεν ξέρω πλέον ποιος είμαι, δεν ξέρω πλέον ούτε πού είμαι, αλλά σήμερα θα φανερωθεί η αλήθεια». Αυτές ήταν οι τελευταίες λέξεις που άκουσα από τον πατέρα μου εκείνη την μέρα. Όταν όμως είσαι οκτώ χρονών δεν μπορείς να αποκρυπτογραφήσεις τέτοιες προτάσεις. Τι πάει να πει «Δεν ξέρω πού είμαι»; Ακόμη χειρότερα, τι πάει να πει «Δεν ξέρω ποιος είμαι»; Εγώ, ας πούμε, είμαι ο Σπαρκ Θοτς, δεν έχω κάποια αμφιβολία για αυτό. Όσον αφορά το που είμαι, μπορείς να πεις ότι είναι λίγο κακός στην γεωγραφία, αλλά τουλάχιστον στην πόλη που μένουμε ξέρω πάντα που είμαι, που βρίσκομαι.

Η εικόνα μου όμως για αυτές τις ερωτήσεις σύντομα θα άλλαζε το βράδυ της ίδιας μέρας που έμαθα τα τραγικά νέα. Ο πατέρας μου είχε πεθάνει το πρωί κατά την διάρκεια της αποστολής του με την Ειδική Υπηρεσία. Η υπηρεσία είχε στείλει έναν πράκτορα στο σπίτι μας για να με ενημερώσει για τις εξελίξεις εκείνης της μέρας, αλλά φαινόταν ότι δεν είχε καμία εμπειρία με παιδιά. Δεν ξέρω τι πίστευε ότι θα γινόταν αν έλεγε ξερά σε ένα οκτάχρονο μοναχοπαίδι όπως εγώ πως ο πατέρας του απεβίωσε, αλλά μια ήρεμη αντίδραση σίγουρα δεν ήταν στην λίστα του. Η αντίδραση μου στα νέα την στιγμή εκείνη ήταν για κάποιο λόγο να πω «Σας ευχαριστώ πολύ. Καλή συνέχεια...» συνοδευόμενος από ένα δάκρυ να κυλάει από το δεξί μου μάτι. Όταν τελικά είδα πως ο πράκτορας δεν ήθελε κάτι παραπάνω από εμένα, έκλεισα την πόρτα και έπεσα στο πάτωμα σκεπτόμενος τις τελευταίες λέξεις που είχα ακούσει από τον πατέρα μου. «Δεν ξέρω ποιος είμαι.», «Δεν ξέρω πού είμαι.» και «Σήμερα θα φανερωθεί η αλήθεια». Ποια αλήθεια κρύβεται στην τελευταία πρόταση; Τι διαφορετικό είχε αυτή η τελευταία υπόθεση του πατέρα μου με την Ειδική Υπηρεσία που τον είχε οδηγήσει να ξεστομίζει τέτοιες φιλοσοφίες; 

Η μητέρα μου ήταν σε άλλη πόλη λόγω μιας δικής της αποστολής με την Ειδική Υπηρεσία, οπότε δεν είχα κάποιον να μιλήσω εκείνο το βράδυ. Ήμουν ολομόναχος στο σπίτι και κοίταζα απλά το ταβάνι με απόγνωση, χωρίς να μπορούν να φύγουν από το μυαλό μου οι λέξεις του πατέρα μου. Το ήξερα ότι μετά από λίγο καιρό ίσως θα τις είχα ξεχάσει τελείως, αλλά εκείνη την νύχτα δεν μπόρεσα να κοιμηθώ λόγω των σκέψεων που με περιέβαλαν σχετικά με τον θάνατο του πατέρα μου, την κατάσταση της μητέρας μου και, το κυριότερο, την χρησιμότητα μου σε αυτή την οικογένεια. Μέχρι στιγμής δεν έχω καταφέρει κάτι στη ζωή μου. Δεν ξέρω καν αν έπρεπε να έχω καταφέρει κάτι. Η ρουτίνα μου είναι απλή: σχολείο, διάβασμα, ύπνος. Προσφέρω κάτι στην οικογένεια μου; ή είμαι απλά ένα βάρος για τους γονείς μου; Μάλλον είναι πολύ μικρός ακόμη για να κάνω τέτοιες ερωτήσεις, αλλά εκείνο το βράδυ ένιωθα πως όντως ήμουν άχρηστος, πως δεν έχω λόγο ύπαρξης. 

Το μόνο καθημερινό κίνητρο που είχα ήταν η ευχαρίστηση που ένιωθα καθώς παίζαμε παιχνίδια με τον πατέρα μου πριν πάμε για ύπνο. Ο πατέρας μου είχε πιο "σταθερό", θα έλεγε κανείς, πρόγραμμα στην Ειδική Υπηρεσία σε σχέση με την μητέρα μου, γιατί είχαν τελείως άλλες ειδικότητες. Ο πατέρας μου ήταν κυρίως ανακριτής, αλλά λόγω των ικανοτήτων του μπορούσε να πηγαίνει και σε αποστολές με την ομάδα του. Η μητέρα μου από την άλλη είναι κατάσκοπος. Από μικρή, μου είχε πει, της άρεσε να μαθαίνει τα πάντα για κάποιον, το οποίο κατά την γνώμη είναι κάπως παράξενο, αλλά για τις δουλειές της Ειδικής Υπηρεσίας χρειάζονται αρκετοί κατάσκοποι έτσι ώστε να αντλούν πληροφορίες για τους εγκληματίες που προσπαθούν να συλλάβουν. Το μόνο αρνητικό που έχει αυτή η ειδικότητα είναι το ωράριο. Η μητέρα μου δεν βρίσκεται πολλές ώρες στο σπίτι. Τις περισσότερες ώρες βρίσκεται σε πόστο κατασκοπείας ή στα κεντρικά της Ειδικής Υπηρεσίας, όπου αναλύει τις πληροφορίες που έμαθε για τους υπόπτους που κατασκόπευε. Λιγότερο επικίνδυνη ειδικότητα κατά την γνώμη μου, σε σχέση με τις αποστολές που πήγαινε ο πατέρας μου, αλλά δεν παύει να έχει και αυτή κίνδυνο, τον κίνδυνο της ανακάλυψης της και στην συνέχεια ποιος ξέρει τι. Αν ο ύποπτος που παρακολουθεί αποδεικνύεται όντως ένοχος και προσπαθήσει να την απαγάγει; Αν δεν θέλει να της αφήσει καν την ευκαιρία να γυρίσει πίσω στην Ειδική Υπηρεσία και την σκοτώσει επιτόπου; Το οκτάχρονο μυαλό δεν μπορεί να κρατηθεί από τέτοιου είδους σκέψεις όταν μόλις έχει μάθει ότι ο ένας του γονιός δεν είναι πλέον ζωντανός. Ευτυχώς για μένα, το επόμενο πρωί η μητέρα μου άνοιξε την πόρτα και με είδε ξαπλωμένο στο πάτωμα και έπεσε κατευθείαν πάνω μου, χωρίς να ακουστεί ούτε μια λέξη. Εκείνη η στιγμή ήταν μόνο συναίσθημα. Δύο άνθρωποι σε απόγνωση να κλαίνε στο πάτωμα και να μην μπορούν να σκεφτούν τίποτα άλλο εκτός από τον τρίτο αυτόν άνθρωπο, που για πρώτη φορά έλειπε από το σπίτι. 

Η μητέρα μου, αφότου ηρέμησε λίγο, με σήκωσε από το πάτωμα, κάθισε στα γόνατα της και το πρόσωπο της έκανε μια έκφραση σαν να ήθελε να πει κάτι, αλλά το στόμα της δεν μπορούσε να κουνηθεί. Τα μάτια της ορθάνοιχτα με κοιτούσαν, όχι πλέον με βλέμμα απόγνωσης, αλλά με ένα βλέμμα που μου μετέφερε τον φόβο της. Εκείνη την στιγμή πήγα να την ακουμπήσω και εγώ στους ώμους, αλλά εκείνη μου έδιωξε τα χέρια και έπεσε στα πόδια μου κλαίγοντας για ακόμη μια φορά. Μου φάνηκε λογικό εκείνη την στιγμή, γιατί προφανώς είχε πεθάνει ένας από τους πιο σημαντικούς άντρες στη ζωή της, αν όχι ο πιο σημαντικός, οπότε προσπάθησα και εγώ να κάνω το πιο απλό που μπορούσα να κάνω εκείνη τη στιγμή, να την καθησυχάσω.

Όπως ήταν έτσι καμπουριασμενη και έμοιαζε με χελώνα που έχει κρυφτεί μέσα στο κέλυφος της, έβαλα το χέρι μου στην πλάτη της και της είπα «Όλα θα πάνε καλά», που ήταν το μόνο πράγμα ο οκτάχρονος εαυτός μου μπορούσε να σκεφτεί υπό την πίεση που βρισκόταν. Εκείνη την στιγμή, χωρίς η μητέρα μου να κάνει κάποια κίνηση προς το μέρος μου και χωρίς να σταματήσει να κλαίει, της άκουσα να λέει:

\textasciitilde Πρέπει να ζήσεις Σπαρκ... πρέπει να ζήσεις...\textasciitilde

«Μην ανησυχείς μητέρα, θα ζήσω», της απάντησα χαμογελώντας.

Η μητέρα μου όταν της απάντησα παραξενεύτηκε, κοίταξε προς το μέρος μου, χαμογέλασε και είπε:

«Το είπα δυνατά αυτό ε; Δεν πειράζει. Έλα εδώ να σε κάνω μια αγκαλιά ακόμη.»

Καθώς την έκανα αγκαλιά, προσπάθησα να αδειάσω το μυαλό μου από τις αρνητικές σκέψεις αλλά δεν μπορούσα. Το μόνο πράγμα που μπορούσα να σκεφτώ εκείνη τη στιγμή ήταν ο πατέρας μου και οι τελευταίες του λέξεις.

«Κοίτα την ώρα, έχει πάει 6:30, θες να κοιμηθείς καθόλου ή θα πας σχολείο άυπνος;», μου είπε η μητέρα μου επειδή αισθάνθηκε πως είχα αρχίσει να κάνω πάλι αρνητικές σκέψεις.

«Αφού δεν προλαβαίνω να κοιμηθώ ούτε να προσπαθούσα. Δεν αξίζει. Μπορείς να μου φτιάξεις πρωινό;» της είπα με έναν υποτονικό τόνο και κάθισα στο τραπέζι.

Η μητέρα μου δεν ξόδεψε καθόλου χρόνο και πήγε στην κουζίνα. Σπάνια την έβλεπα να μαγειρεύει, καθώς έλειπε συνέχεια από το σπίτι. Ο πατέρας μου δεν ήταν και από τους καλύτερους μάγειρες, αλλά τουλάχιστον είχα φαγητό κάθε μέρα στο σπίτι, οπότε δεν μπορώ να παραπονεθώ. Η μητέρα μου δεν ξέρω πραγματικά τι έτρωγε κατά της διάρκεια της κατασκοπείας της, ή αν έτρωγε καν. Δεν μπορώ να φανταστώ κάποιο χρονικό διάστημα που να είναι τόσο χαλαρή ώστε να κάνει διάλειμμα για φαγητό, ή τουλάχιστον κάποιο σνακ. Ακόμη χειρότερα, ποτέ κοιμάται; Γιατί προκύπτει πάλι το ίδιο πρόβλημα. Πότε υπάρχει χρόνος στο πόστο της για διάλειμμα; Είναι πραγματικά ηρωίδα η μητέρα μου, δεν ξέρω πως τα καταφέρνει.

«Μαμά!»

«Ναι Σπαρκ;» απάντησε χωρίς να γυρίσει.

«Τρως καθόλου;» την ρώτησα από περιέργεια.

«Προφανώς παιδί μου. Αν δεν έτρωγα δεν θα μπορούσα να είμαι εδώ αυτή τη στιγμή.» μου απάντησε γελώντας. «Τι σε έπιασε τώρα; Θέλεις να μάθεις μήπως και αν κοιμάμαι καθόλου; Χαχα, δεν χρειάζεται να ανησυχείς για μένα Σπαρκ. Από εδώ και πέρα θα είμαι πιο κοντά σου από ποτέ, οπότε μην αγχώνεσαι.» μου είπε και γύρισε να πάρει δύο αυγά από το ψυγείο.

«Εντάξει» είπα ενω τράβηξα την καρέκλα για να κάτσω στο τραπέζι.

\chapter{}

Δεν θυμάμαι πολλά άλλα πράγματα από εκείνο το πρωί με την μητέρα μου. Το μόνο που θυμάμαι είναι να ετοιμάζομαι να φύγω για το σχολείο θλιμμένος και εκείνη να μου λέει.

«Σπαρκ θελω από εδώ και πέρα να χαμογελάς. Ήμασταν άτυχοι για μία μέρα. Δεν χρειάζεται αυτή η ατυχία να μας χαλάσει όλη την υπόλοιπη ζωή.» 

Εκείνη τη στιγμή δεν απλά της χαμογέλασα και έφυγα από το σπίτι, αλλά με το που έκλεισα την πόρτα, το χαμόγελο αυτό ξεθώριασε. Δεν άκουσα την μητέρα μου. Η θλίψη δεν με άφηνε να χαρώ ούτε στιγμή.

Η υποτονικότητα μου αρχικά νόμιζα ότι πέρασε απαραίτητη στο σχολείο. Περνούσαν οι ώρες και οι συμμαθητές μου συμπεριφέρονταν λες και ήταν μια κανονική μέρα. Ο φίλος μου ο Βασι ειδικά δεν έλεγε να σταματήσει να μιλάει σήμερα. Δεν έδινα πολύ σημασία σε αυτά που έλεγε γιατί είχα αλλά πράγματα στο μυαλό μου αλλά ό,τι και να ήταν θα έπρεπε να είναι πολύ σημαντικό για να μου μιλάει ακόμη και μέσα στην τάξη για αυτό. Ο Βασι γενικά είναι ένας πολύ εξωστρεφής άνθρωπος ο οποίος με είχε βρει την πρώτη μέρα της σχολικής μας... "καριέρας"; "σταδιοδρομίας"; Δεν ξέρω πραγματικά, αλλά πιστεύω άρχισε να με κάνει παρέα γιατί είμαι αρκετά εσωστρεφής άνθρωπος και αυτό είναι ό,τι καλύτερο για τον Βασι, καθώς αυτός μπορεί να μιλάει όσο θέλει χωρίς να τον διακόπτουν. Εμένα αυτό μου άρεσε, είχα τουλάχιστον μία παρέα στο σχολείο. Μερικές φορές ήταν αστείος, μερικές φορές απλά νάρκισσος, αλλά ποτέ μα ποτέ αρνητικός προς εμένα. Δεν με είχε πειράξει ποτέ, εκτός από φιλικά πειράγματα που προφανώς υπάρχουν σε όλες τις φιλίες. 

Αυτό που μου άρεσε ιδιαίτερα στον Βασι ήταν ότι με υπερασπιζόταν σε καταστάσεις που ο εσωστρεφής εαυτός μου φοβόταν να μιλήσει. Λόγω του Βασι, κάνεις στο σχολείο δεν μας ενοχλούσε, ούτε τολμούσε να μας ενοχλήσει, γιατί όλοι ήξεραν ότι ο Βασι είχε έναν ιδιαίτερο τρόπο διαχείρισης των καταστάσεων και μπορούσε να κερδίσει κάθε ντιμπέιτ ή να αποκρούσει κάθε προσβολή χωρίς να προσβάλει κάποιον πίσω. Αυτή η ασφάλεια και η γαλήνη που μου πρόσφερε ο Βασι στον χώρο του σχολείου είναι πραγματικά κάτι να είμαι ευγνώμων. 

Η υποτονικότητα μου δεν έλεγε να φύγει εκείνη τη μέρα. Βεβαίως, όσο είμαι με το Βασι γελάω, χαίρομαι και δεν ανησυχώ, αλλά δεν μπορώ να ξεχάσω έτσι απλά ότι ο πατέρας μου δεν είναι πλέον μαζί μου. Μου ήρθαν πάλι κάποια στιγμή οι λέξεις του πατέρα μου ενώ κάναμε μάθημα και ξανάρχισα τους συνειρμούς για το "ποιος είμαι" και το "που είμαι". Ακόμη δεν μπορούσα να καταλάβω την δυσκολία του πατέρα μου πάνω σε αυτό το ζήτημα. Μάλλον είναι μικρός ακόμη και δεν ξέρω πως δουλεύει η ζωή. Μάλλον όταν μεγαλώσω θα μπορέσω να βρω απάντηση. Αυτή τη στιγμή όχι μόνο δεν μπορώ να σκεφτώ καθαρά, αλλά δεν μπορώ να παρακολουθήσω και στο μάθημα. Έχει περάσει όλη η σχολική μέρα και δεν έχω μάθει τίποτα καινούργιο, όχι γιατί δεν έχω κάνει μάθημα, αλλά γιατί έκανα σκέψεις μόνος μου κοιτώντας τους τοίχους της αίθουσας. Αν και ωραίοι οι τοίχοι, τι θα πω στην μητέρα μου όταν γυρίσω σπίτι; Ότι δεν άκουσα την συμβουλή της και ήμουν θλιμμένος όλη μέρα; Ότι δεν κάναμε καθόλου μάθημα όλη μέρα; Το δεύτερο δεν πρόκειται να το πιστέψει. Τουλάχιστον για μια φορά μετά από πολύ καιρό θα ήταν η μητέρα μου που θα με υποδεχόταν σπίτι αντί για τον πατέρα μου. Θέλω να δω και τι φαγητό θα έχει φτιάξει, αν έχει φτιάξει κάτι, καταλαβαίνω ότι η κατάσταση είναι δύσκολη και για τους δυο μας, αλλά αφού μου έφτιαξε πρωινό, γιατί να μην έχει φτιάξει και μεσημεριανό;

Το τελευταίο κουδούνι ήχησε. Όλοι σηκώθηκαν να επιστρέψουν στα σπίτια τους. Όταν πήγα να σηκωθώ, είδα την δασκάλα να με πλησιάζει και να μου κάνει νόημα να μείνω στην τάξη.

«Είσαι καλά Σπαρκ; Σε παρατηρώ όλη την μέρα και φαίνεται να έχεις ένα υποτονικό ύφος. Έχεις κάτι στο μυαλό σου;»

«Τι σημαίνει "Δεν ξέρω ποιος είμαι";» την ρώτησα χωρίς να την κοιτάξω ή να σηκωθώ από την θέση μου.

«Τι φιλοσοφική ερώτηση είναι αυτή;» είπε η δασκάλα προβληματισμένη. «Ο καθένας μπορεί να μάθει ποιος είναι μέσα από τις εμπειρίες του. Δεν υπάρχει κάποιος αυστηρός ορισμός για το "ποιος είναι" κάποιος.» είπε η δασκάλα και χαμογέλασε

Την κοίταξα για λίγο και ξανακοίταξα το θρανίο μου.

«Πιστεύετε... ότι είμαι χρήσιμος;»

«Έχεις γίνει φιλόσοφος σήμερα;» με ρώτησε η δασκάλα γελώντας. «Όλοι οι άνθρωποι είναι χρήσιμοι με τον δικό τους τρόπο. Ας πούμε κοίτα...»

Η δασκάλα λύγισε τα γόνατα της για να έρθει στην δίπλα μου και μου έδειξε το χέρι της. Έστρεψε την προσοχή μου στην παλάμη της και με μια κίνηση του χεριού της, μια μικρή φωτιά εμφανίστηκε πάνω από την παλάμη της δασκάλας.

«Εγώ ας πούμε, έχω μαγεία Φωτιάς. Μπορώ να βγάλω φωτιά από τα χέρια μου. Βέβαια...» είπε η δασκάλα και έσβησε γρήγορα την φωτιά από το χερι της, τινάζοντας το λίγο λες και είχε χτυπήσει τον καρπό της, «...δεν μπορώ να τη χρησιμοποιήσω για πολύ ώρα, γιατί καίγομαι εύκολα. Υπάρχουν άλλοι άνθρωποι, όμως, στην Ειδική Υπηρεσία ειδικά, οι οποίοι έχουν και αυτοί μαγεία Φωτιάς, αλλά δεν καίγονται τόσο εύκολα γιατί έχουν γεννηθεί με φυσική αντοχή στα μειονεκτήματα αυτής της μαγείας. Αυτό με κάνει εμένα άχρηστη; Όχι βέβαια. Η χρησιμότητα μου στην κοινωνία είναι να μεταδίδω γνώσεις, πράγμα που δεν έχει απαραίτητα να κάνει με την μαγεία μου. Δεν χρειάζεται να είσαι ανασφαλής για την δυνατότητα να χρησιμοποιήσεις την μαγεία σου αν σκέφτεσαι αυτό. Όλοι έχουν δυσκολίες με την μαγεία τους, ακόμη και οι ενήλικες από ότι είδες. Από όσο ξέρω, όμως, οι γονείς σου είναι κορυφαίοι πράκτορες της Ειδικής Υπηρεσίας, έτσι δεν είναι; Σίγουρα θα μπορέσουν να σε βοηθήσουν να εξελιχθείς.»

Την κοίταξα πάλι για λίγο και ξανακοίταξα το θρανίο μου. Αυτή την φορά, ένα δάκρυ μου ξέφυγε και κύλησε στο αριστερό μου μάγουλο.

«Γιατί υπάρχει η μαγεία;» ρώτησα την δασκάλα προσπαθώντας να κρατήσω τα συναισθήματα μου για τον εαυτό μου, καθώς τα λόγια της δασκάλας μου θύμισαν πάλι τον πατέρα μου. «Γιατί κανείς δεν αναρωτιέται τον λόγο ύπαρξης της μαγείας; Κάνει τον κόσμο πολύ πιο περίπλοκο από ότι χρειάζεται»

«Η μαγεία, Σπαρκ, πράγματι είναι από τα μεγαλύτερα μυστήρια του κόσμου μας. Το μόνο που ξέρουμε για αυτή είναι πως γύρω στα 200 με 220 χρόνια πριν άρχισαν να εμφανίζονται τα πρώτα σημάδια μαγείας στους ανθρώπους. Κανείς δεν ξέρω από πού προήλθε ή γιατί υπάρχει. Αυτός είναι και ο λόγος που ονομάστηκε μαγεία. Ο κάθε άνθρωπος πλέον μπορεί να χρησιμοποιήσει την δική του μαγεία. Βέβαια, μέσα σε αυτά τα 200 χρόνια, πολλοί λίγοι καταφέρνουν να χρησιμοποιούν την μαγεία τους χωρίς να τους κρατάει κάτι πίσω. Ας πούμε, η μαγεία Φωτιάς που έχω εγώ είμαι από φύσης πολύ καταστροφική, όμως το μειονεκτήματα ότι ο χρήστης της τραυματίζεται, ή τουλάχιστον τρέχει το ρίσκο να τραυματιστεί, κάθε φορά που την χρησιμοποιεί κρατάει τους περισσότερους χρήστες αυτής της μαγείας πίσω από την πλήρη ας το πούμε δύναμη που θα νόμιζε ότι θα είχε κάποιος. 

Υπάρχουν βέβαια άνθρωποι, όπως σου είπα και πριν, που έχουν γενετικές "ιδιοτροπίες" ας το πούμε. Αυτοί οι άνθρωποι μπορούν να χρησιμοποιήσουν την μαγεία τους με πολύ μεγαλύτερη... ένταση; Δεν ξέρω πως να το εξηγήσω τώρα, αλλά αυτοί οι άνθρωποι συνήθως έχουν τρεις επιλογές. Η μία είναι να ζήσουν την ζωή τους όπως θα την ζούσαν χωρίς να έχουν το γενετικό πλεονέκτημα. Συνήθως όμως όταν ένας άνθρωπος έχει την πλήρη δύναμη μιας μαγείας δύσκολα αντιστέκεται στην ιδέα να την χρησιμοποιεί, οπότε οι περισσότεροι κάνουν μια από τις ακόλουθες επιλογές. Η δεύτερη επιλογή, που είναι και η πιο συχνή είναι, αφορά την Ειδική Υπηρεσία. Η συντριπτική πλειοψηφία των πρακτόρων της Ειδικής Υπηρεσίας είναι χρήστες μαγείας με γενετικά πλεονεκτήματα. Με την βοήθεια αυτών, οι διαφορές υποθέσεις που διερευνούνται από την Ειδική Υπηρεσία λύνονται με πολύ γρήγορους ρυθμούς. 

Η τρίτη επιλογή είναι... να χρησιμοποιήσει κάποιος την μαγεία για τον εαυτό του. Υπάρχουν κακόβουλοι άνθρωποι που χρησιμοποιούν τις μαγείες τους για σκοπούς που κανείς δεν θα φανταζόταν πριν την ύπαρξη της μαγείας. Είναι πραγματικά...»

Το τηλέφωνό της δασκάλας χτύπησε. Δεν πρόλαβα να δω το όνομα που έγραφε αλλά από την έκφραση της κατάλαβα πως μάλλον είχε αργήσει για κάποιο ραντεβού.

«Λοιπόν Σπαρκ πολύ φλυαρησα. Ήθελα να σου πω συμπερασματικά να μην αγχώνεσαι. Όλα θα πάνε καλά. Ο,τι και να περνάς να ξέρεις πως αν έχεις ανθρώπους δίπλα σου θα μπορέσεις να επανέλθεις.» μου είπε η δασκάλα και με συνόδευσε έξω από την τάξη. Εκεί  βρισκόταν ο Βασι και με περίμενε.

«Επιτέλους τελειώσατε. Τι συζητάγατε τόση ώρα θα ήθελα να ήξερα. Κάτι για μαγεία άκουσα στο τέλος που άρχισα να βαριέμαι να κάθομαι και άρχισα να κρυφακουω. Εξελίχθηκε η μαγεία σου; Πες μου ότι μεγάλωσαν οι κεραυνοί σου. Πλάκα θα είχε. Θα μπορούσαμε επιτέλους να παλέψουμε κάποια στιγμή. Τώρα δεν μπορώ να σταματήσω να το φαντάζομαι...» άρχισε να μονολογεί κλασικά ο Βασι καθώς η δασκάλα μας άφησε μόνους να γυρίσουμε στα σπίτια μας.

Καθώς περπατάμε ο Βασι περιέγραφε, με λίγες υπερβολές, το πως θα εξελισσόταν η μάχη μας στο μυαλό του. Εγώ βέβαια δεν τον άκουγα γιατί σκεφτόμουν τον πατέρα μου. Από αυτόν κληρονόμησα την μαγεία Κεραυνού, ή τουλάχιστον έτσι νομίζω πως δουλεύει η μαγεία. Βέβαια ο πατέρας μου είχε "γενετικά πλεονεκτήματα" όπως τα ονομάτισε η δασκάλα και μπορούσε άνετα να χρησιμοποιεί τους κεραυνούς του χωρίς να τον κρατάει κάτι πίσω. Εγώ βέβαια μέχρι τώρα δεν μπορώ ούτε να σοκαρω κάποιων με τους κεραυνούς από τα χέρια μου. Ούτε να φορτίσω το κινητό μου δεν μπορώ, που θα έλεγε κάποιος είναι ένα  πλεονέκτημα της μαγείας μου. 

Όλοι λένε πως η εξέλιξη μου είναι αναπόφευκτη λόγω του πατέρα μου, όμως εγώ αμφιβάλω. Και αν ο πατέρας μου ήταν όντως όσο καλός όσο λένε με τους Κεραυνούς του, γιατί πέθανε; Δεν μπορούσε να βγει από εκείνη την κατάσταση που βρέθηκε με την μαγεία του; Αν όμως η μαγεία του πατέρα μου ήταν το μόνο δυνατό του σημείο δεν θα με προβληματίζει τόσο. Ο πατέρας μου ήταν φημισμένος ως ένας από τους πιο έξυπνους ανακριτές της Ειδικής Υπηρεσίας. Γιατί δεν μπορούσε να χρησιμοποιήσει αυτή την εξυπνάδα εκεί που βρέθηκε; Τι είδους αποστολή ήταν; Πραγματικά δεν μπορώ να σκεφτώ κατάσταση που ο πατέρας μου να βρίσκεται σε μειονεκτική θέση. Θέλω πραγματικά να πάω ο ίδιος στην Ειδική Υπηρεσία και να τους ρωτήσω.

«...και μετά ΜΠΟΥΜ. Πιάνουμε τους κακουλιδες με τις μαγείες μας και επιστρέφουμε στην Ειδική Υπηρεσία ως ήρωες. Τι λες; Θα γίνουμε πράκτορες μαζί;» είπε ο Βασι.

«Πώς προέκυψε αυτό;» τον ρώτησα, αφού είχα αγνοήσει όλη την ιστορία που είχε αναπτύξει καθώς περπατούσαμε.

«Θες να γίνουμε πράκτορες μαζί όταν μεγαλώσουμε ή όχι; Θα έχει πλάκα. Τι λες;» ρώτησε πάλι ο Βασι απευθύνοντας μου τον λόγο, κάτι που δεν έκανε συχνά.

«Θα... το σκεφτώ» του απάντησα, αφήνοντας τον στο ενδιάμεσο της κατάφασης και της άρνησης.

«Σκέψου πόσο ωραία θα είναι. Θα πηγαίνουμε ταξίδια σε άλλες πόλεις. Θα γνωρίζουμε νέους ανθρώπους. Θα...» άρχισε να μονολογεί πάλι ο Βασι. Δεν γνωρίζω αν αυτά που έλεγε είναι σωστά αλλά αν όντως η δουλειά στην Ειδική Υπηρεσία είναι όπως την περιγράφει, ίσως να μην είναι και τόσο άσχημα τελικά. 

Οι γονείς μου δεν έμπαιναν ποτέ σε λεπτομέρειες σχετικά με την δουλειά τους λόγω της φύσης της. Οι υποθέσεις της Ειδικής Υπηρεσίας είναι υποτίθεται κρυφές για να έχουν το στοιχείο της έκπληξης, ή τουλάχιστον έτσι το είχα εγώ στο μυαλό μου. Αν η πρόοδος της Ειδικής Υπηρεσίας για κάθε υπόθεση ξέφευγε κάπως από τα κεντρικά και το μάθαινε κάποιος "ένοχος", προφανώς θα άλλαζε τα σχέδια του αντίστοιχα. Καταλαβαίνω λοιπόν γιατί οι γονείς μου δεν συζητούσαν ποτέ με εμένα για την δουλειά τους. Από την άλλη, όταν αρχίζεις και μιλάς στον εαυτό σου λέγοντας "Δεν ξέρω ποιος είμαι" και "δεν ξέρω που είμαι", ενώ το παιδί σου είναι δίπλα, προφανώς εγώ θα οδηγηθεί σε σκέψεις. Πώς δημιουργείται αυτή η αμφιβολία σε έναν άνθρωπο; Μήπως η Ειδική Υπηρεσία 

«Σπαρκ φτάσαμε στο φανάρι που χωρίζονται οι δρόμοι μας. Ορεβουάρ!... ή τι λένε οι Γάλλοι...»

«Καλό μεσημέρι Βασι...» του είπα και κοίταξα κάτω.

«Θα γίνουμε πράκτορες έτσι;» με ξαναρωτησε ο Βασι χαμογελώντας.

«Μπορεί...» του απάντησα για να σταματήσει να με ρωτάει.

«Αυτό ήθελα να ακούσω. Άντε Τσάο!» είπε ο Βασι και έστριψε δεξιά για να πάει στο σπίτι του όσο εγώ περίμενα το φανάρι.

Όταν έγινε πράσινο το φανάρι, συνέχισα το περπάτημα και άρχισα να αναρωτιέμαι αν όντως μπορούσα να γίνω πράκτορας της Ειδικής Υπηρεσίας. Αφού υπάρχουν και πράκτορες χωρίς γενετικά πλεονεκτήματα, γιατί να μην μπορώ και γω να είμαι ένας από αυτούς; Μπορεί να μην είμαι ο αρχηγός κάποιας ομάδας, αλλά αν η δουλειά στην Ειδική Υπηρεσία σημαίνει ταξίδια και γνωριμίες, γιατί όχι; Αν η δουλειά στην Ειδική Υπηρεσία είναι σαν τις αστυνομικές ταινίες που εξερευνούν μυστήρια και καταλήγουν στο αποκορύφωμα με μια τελική μάχη μεταξύ "καλών" και "κακών" θα είχε όντως πλάκα. Τουλάχιστον αν δεν είμαι εγώ από αυτούς που μάχονται, καθώς δεν μπορώ να χρησιμοποιήσω την μαγεία μου για τέτοιους σκοπούς. Οι κεραυνοί μου είναι πολύ μικροί και αδύναμοι για να τραυματιστεί κάποιος από αυτούς. Εντάξει, δεν με πειράζει να μην λαμβάνω μέρος σε αυτές τις μάχες, αν γίνονται καν μάχες. Μπορεί η Ειδική Υπηρεσία να αιφνιδιάζει τους ενόχους και να τους συλλαμβάνει χωρίς να ανάβουν τα αίματα. Ποιος ξέρει. Μόνο ένας τρόπος υπάρχει να μάθει κανείς.

Ως πράκτορας θα μπορούσα επίσης να βρω τον λόγο που πέθανε ο πατέρας μου, καθώς ο ευγενικός πράκτορας που είχε έρθει σπίτι μου εκείνο το βράδυ δε μου έδωσε παραπάνω πληροφορίες. Αν ρωτήσω έτσι στα ξερά δεν νομίζω να απαντούσε κάνεις. Ένας νέος πράκτορας να ζητάει πληροφορίες για τον θάνατο ενός κορυφαίου πράκτορα είναι από οποία οπτική και να το σκεφτεί κανείς λίγο ύποπτο. Θα χρειαστεί να κερδίσω πρώτα την εκτίμηση των ανθρώπων μέσα στην Ειδική Υπηρεσία. Για κάτσε... η μαμα δεν μπορεί να μάθει για ποιο λόγο πέθανε ο μπαμπας; Αν δεν ξέρει ήδη προφανώς. Αν ξέρει όμως γιατί δεν το μοιράζεται μαζί μου; Δεν μπορώ να καταλάβω γιατί όλα στην Ειδική Υπηρεσία είναι τόσο κρυφά. Δεν χρειάζεται κιόλας να είναι όλα κρυφά. Δεν να κάνουν και εξαίρεση για την περίπτωσή μου; 

Όσο σκεφτόμουν έφτασα ασυνήθιστα γρήγορα στο σπίτι. Άνοιξα την πόρτα και κοίταξα την γωνία που αφήνουμε τα παπούτσια μας. Ένα ζευγάρι προφανώς έλειπε. Κάθισα λίγο παραπάνω αφότου έκλεισα την πόρτα κοιτώντας την γωνία αυτή που δεν θα ήταν ποτέ ξανά ίδια. Ένα ακόμη δάκρυ κύλησε στο μάγουλο μου, το δεξί αυτή τη φορά. Άκουσα βήματα να έρχονται από την κουζίνα σαν κάποιος να έτρεχε.

«Σπαρκ!» είπε η μητέρα μου όταν με είδε στην πόρτα όρθιο κοιτώντας την γωνία των παπουτσιών μας. Έτρεξε προς το μέρος μου και με αγκάλιασε.

\textasciitilde Άργησες και με αγχωσες. Ξέρεις πως αισθάνεται μια γυναίκα όταν παθαίνει κάτι εκτός απροόπτου και προσπαθεί να αποφύγει την επανεμφάνιση του ίδιου γεγονότος. Πραγματικά...\textasciitilde 

«Μαμά είμαι καλά. Μην αγχώνεσαι.»

«Τι λες Σπαρκ; Δεν σου είπα τίποτα. Μόνο το όνομα σου φώναξα καθώς ερχόμουν. Ήμουν... χαρούμενη που σε είδα δύο φορές μέσα σε μια μέρα, αυτό ήταν όλο» είπε η μητέρα μου καθώς χαλάρωνε την αγκαλιά της και ακούμπησε τα χέρια της στους ώμους μου. «Έλα τώρα να φας τίποτα. Πρέπει να πεινάς. Βγάλε τα παπούτσια σου και έλα.»

\textasciitilde Όχι ότι πρόλαβα να φτιάξω κάτι καλό, έφτιαξα κάτι πρόχειρο με όσα είχε ο πατέρας σου στην κουζίνα\textasciitilde 

«Δεν πειράζει μαμά. Ό,τι και να έφτιαξες θα το φάω.» είπα για να την καθησυχάσω, αφού ξέρω ότι είχε να μαγειρέψει αρκετό καιρό και πιθανότατα ήταν αγχωμενη.

«Τι είπες Σπαρκ;» ρώτησε προβληματισμένη η μητέρα μου όσο έβγαζε τα χέρια της από τους ώμους μου. 

«Είπα κάτι λάθος;» της απάντησα χαμογελώντας. Εκείνη την στιγμή η απάντηση της με έκανε να νομίζω ότι την πρόσβαλα με τα λόγια μου.

«Τίποτα. Δεν πειράζει. Πρέπει να είναι η φαντασία μου...» μου αποκρίθηκε η μητέρα μου και πήγε να ετοιμάσει το τραπέζι. «Άσε τα πράγματα σου και έλα να φάμε.»

\chapter{}

«Βάλατε όλα τα πράγματα στο αυτοκίνητο;»

«Ναι αδερφέ, εδώ και κανένα μισάωρο. Περίμενε λίγο να τελειώσουμε αυτό τον γύρο και θα φύγουμε.»

«Πάλι πόκερ παίζεται; Κανόνισε αδελφέ να πορώσεις τα παιδιά μου με αυτή τη χαζομάρα και πηγαίνουν στο καζίνο.»

«Σιγά τα αυγά. Δεν πρόκειται να πάθουν κάτι τα παιδιά σου, είναι έξυπνα αγόρια. Εξάλλου, είστε εικοσάρηδες και οι δύο τώρα, έτσι δεν είναι;»

«Θείε τελείωνε τη φλυαρία και παίξε για να τελειώνουμε.»

«Α ναι συγγνώμη. Πάω πάσο. Θέλω να δω και το τελευταίο φύλο πριν σας κερδίσω. Ρίξε την τελευταία κάρτα Ρομπ.»

«Άσσος μπαστούνι. Για ποντάρετε τώρα.»

«Ποντάρω άλλα 20.»

«Άλλα 20; Με διέλυσες μικρέ. Είπαμε να παίξουμε με μικρά ποσά αλλά τώρα νομίζω ότι έχεις καλό χέρι. Βλέπω κάτω τρία μπαστούνια μπαστούνια, οπότε αν έχεις και εσύ άλλα δύο σίγουρα κερδίζεις την κέντα μου, οπότε έχω μικρές πιθανότητες να κερδίσω, άρα...»

«Άρα περιμένεις θείε, δεν είναι η σειρά σου. Είμαστε τρία άτομα που παίζουμε εκτός του Ρομπ που είναι ντίλερ.»

«Α συγγνώμη, παρακαλώ παίξε.»

«Ταιριάζω τα 20 σου αδερφέ. Τώρα παίζεις θείε.»

«100»

«Τι 100; Αφού πριν έλεγες ότι θα έχανες;»

«Παίξτε μικροί. Ο πατέρας σας βιάζεται δεν βλέπετε.»

«Με την ησυχία σας παιδιά, δεν βιάζομαι. Ίσα ίσα μου αρέσει να βλέπω τον αδελφό μου να χάνει.»

«Σιγά μην χάσω.»

«Ταιριάζω τα 100 σου θείε. Νομίζω ότι λες αλήθεια και πιστεύεις ότι οι πιθανότητες είναι με το μέρος σου, αλλά ατύχησες.»

«Εγώ αποσύρομαι. Δεν θα παίξω για τόσα πολλά λεφτά. Κρίμα τα 20 που πόνταρα πριν τον θείο όμως.»

«Δεν πειράζει μικρέ. Κοίτα πως θα κερδίσουμε τον αδελφό σου τώρα. Ποντάρω άλλα 400»

«Θείε είπαμε μικρά ποσά. Τι νομίζεις ότι κάνεις τώρα; Θες να πάμε ταξίδι και να μην έχω φράγκο;»

«Καλά καλά μην με βρίσεις κιόλας. Για τα 100 όμως δεν είπες τίποτα; Ό,τι σε συμφέρει κάνεις μικρέ.»

«Εγώ λέω να μην ποντάρουμε άλλο. 100 είναι ήδη πολλά.»

«Α τώρα τα αναφέρεις τα 100; Γιατί; Φοβάσαι μην χάσεις; Μικρός είσαι ακόμη δεν πειράζει. Από τα λάθη μαθαίνεις. Ποντάρω 200 αντί για 400 γιατί είμαι καλός άνθρωπος.»

«Εγώ δεν ποντάρω είπα πάνω από 100.»

«Αποσύρσου τότε.»

«Όχι.»

«Φοβάσαι;»

«Είπα τους λόγους μου. Πες του τίποτα ρε Ρομπ. Δικός σου πατέρας είναι.»

«Ναι μπαμπά έχει δίκιο. Δεν αξίζει να ποντάρουμε τόσα λεφτά πριν αρχίσουμε το ταξίδι. Γυρίστε τις κάρτες σας.»

«Όχι δεν τις γυρνάω. Πόνταρε 200 ή αποσύρσου.»

«Έλα αδελφέ νομίζω δεν χρειάζεται να μακρυγορείς άλλο. Μπλοφάρεις και φαίνεται από μίλια μακριά.»

«Που ξέρεις ότι δεν έχω κέντα χρώμα. Εξάλλου τα μπαστούνια κάτω είναι άσσος, δυάρι και πέντε.»

«Επειδή δεν είσαι καλός ψεύτης. Για αυτό ξέρω ότι μπλοφάρεις. Και επειδή είδα τις κάρτες σου...»

«Αμάν ρε Ριντ. Πάντα σε μπελάδες με βάζεις»

«Οπότε θείε δεν έχεις τίποτα;»

«Ναι δεν έχω. Άντε πάρε τα 100 και πάμε. Δεν ξαναπαίζω άλλο πόκερ σήμερα.»

«Φοβάσαι;»

«Άσε μας ρε Ριντ. Πάμε στο αμάξι τώρα γιατί νευρίασα με εσένα.»

«Ποιος θα οδηγήσει, αν επιτρέπεται;»

«Εγώ προφανώς. Αφού ο αδελφός μου είναι νευριασμένος και εσείς είστε μικροί ακόμη για τέτοιες διαδρομές.»

«Δεν είμαι μικρός πλέον μπαμπά.»

«Τουλάχιστον φέρεστε σαν μικροί...»

«Τι πάει να πει αυτό;»

«Το γεγονός ότι δεν το καταλαβαίνεις δεν βοηθάει την θέση σου.»

«Άσε με τουλάχιστον να κάτσω μπροστά. Παρακαλώ. Παρακαλώ πολύ.»

«Ο θείος είπαμε θα κάτσει μπροστά για να με βοηθάει. Εσείς καθίστε πίσω και απολαύστε την διαδρομή. Αν θες έχω κάτι μικρά αστυνομικά βιβλία στο πίσω κάθισμα να πάρεις καμιά ματιά, γιατί θα είμαστε 5 ώρες στον δρόμο.»

«Σιγά μην διαβάσω εγώ βιβλίο. Ειδικά και μέσα στο αυτοκίνητο που θα κουνάει κιόλας; Άστο καλύτερα. Δεν με ξέρεις καθόλου καλά μπαμπά.»

«Ναι πλάκα έκανα καλέ. Περισσότερο για τον αδελφό σου το είπα, που του αρέσει να διαβάζει και να αφαιρείται. Κέρδισε και στον τελευταίο γύρο τον θειο του, οπότε πρέπει να είναι η τυχερή του σήμερα. Καλά δεν λέω μικρέ;»

«Παρακαλώ; Αφαιρέθηκα πάλι συγγνώμη, δεν άκουσα τι λέγατε...»

«Έλα Ριντ τελειώνετε. Θα αργήσουμε.»

«Κοίτα ποιος μιλάει. Τέλος πάντων, πάμε παιδιά.»

\chapter{}

«Σπαρκ! Ξέρεις που έβαζε ο μπαμπας σου τη σπάτουλα;»

«Πρέπει να είναι σε αυτό το ντουλάπι αριστερά από τον φούρνο.» είπα στην μητέρα μου χωρίς να σηκωθώ από το τραπέζι που διάβαζα.

«Μα καλά τρεις μέρες που είμαι εδώ έχω ανοίξει όλα τα ντουλάπια και δεν την έχω δει. Πάλι καλά που έχεις καλή μνήμη παιδί μου. Έλα εδώ να σε ζουλήξω» είπε η μητέρα μου και άρχισε προς το μέρος μου με την σπάτουλα.

«Μη μαμά, διαβάζω» της είπα και την έδιωξα με το χέρι μου.

«Καλά δεν πειράζει, πάω να συνεχίσω να μαγειρεύω» είπε η μητέρα μου και γύρισε στην κουζίνα.

Δεν διάβαζα και κάτι δύσκολο εκείνη τη στιγμή, αλλά ήταν από εκείνη την μέρα που είχα μάθει τα νέα για τον πατέρα μου δεν είχα προσέξει ούτε δευτερόλεπτο στο σχολείο, οπότε ήθελα να έχω λίγη ησυχία για να αφομοιώσω το υλικό που είχα μπροστά μου.

Μα καλά όμως, γιατί στο σχολείο δεν κάνουμε κάτι χρήσιμο; Δηλαδή, την ιστορία ας πούμε την θεωρώ άχρηστη. Τι με ενδιαφέρει τι έκαναν κάποιοι άνθρωποι χιλιάδες χρόνια πριν από εμένα και γιατί πρέπει να το μάθω απ’ έξω; Δεν μας μαθαίνουν καν σύγχρονη ιστορία, τουλάχιστον από την ανακάλυψη της μαγείας και μετά. Αυτό θα με ενδιέφερε όντως. Είχα ρωτήσει πριν κάτι βδομάδες όμως πριν πεθάνει ο πατέρας μου την δασκάλα αν η ιστορία σε μεγαλύτερες τάξεις θα περιέχει και ιστορία της μαγείας, αλλά αυτή μου απάντησε πως δεν ξέρουμε αρκετά για την μαγεία ακόμη, παρόλο που έχουν περάσει 200 χρόνια.

Η δασκάλα λέει συνέχεια πως για να ξέρουμε ποιοι είμαστε, πως αναπτύχθηκε η κοινωνία και προς τα που πρέπει να κινηθούμε σαν κοινωνία, πρέπει να ξέρουμε ιστορία. Δεν κατάλαβα ποτέ το γιατί, αλλά για να το λέει η δασκάλα κάτι θα ξέρει. Ποιοι είμαστε... «Δεν ξέρω ποιος είμαι»... Λες ο πατέρας μου να έπαθε κάτι σε σχέση με την ιστορία; Μπα δεν νομίζω. Αφού πράκτορας ήταν, τι την χρειάζονται αυτοί την ιστορία οι πράκτορες;

Τι χρειάζεται να ξέρει ένας πράκτορας πλάκα πλάκα. Να γιατί το σχολείο δεν μας προετοιμάζει. Εγώ πως θα βρω δουλειά ξέροντας το πυθαγόρειο θεώρημα και να χωρίζω τις προτάσεις σε υποκείμενο, αντικείμενο και ρήμα; Γιατί δεν υπάρχουν μαθήματα που να σχετίζονται με πραγματικά επαγκέλματα; Και ας μη μιλήσω για το ότι δεν διδάσκουν τίποτα σε σχέση με την μαγεία. Είμαι σίγουρος πως αν ρωτήσω την δασκάλα θα μου πει πάλι «...επειδή δεν ξέρουμε αρκετά για την μαγεία, παρόλο που έχουν περάσει 200 χρόνια...». 

Πάλι αφαιρέθηκα. Τώρα που ήμουν; Α ναι...

«Σπαρκ τρώμε βραδινό. Κάνε λίγο χώρο να βάλουμε τα πιάτα. Αν θες μετά θα σε βοηθήσω και να διαβάσεις για να τελειώσεις νωρίς.»

«Ευχαριστώ μαμά» είπα και μετακίνησα τα βιβλία μου στην τρίτη καρέκλα του τραπεζιού, στην οποία δεν καθόταν κάποιος πλέον. 

«Τι διαβάζεις τόση ώρα θα ήθελα να ήξερα. Δεν είσαι και τόσο μεγάλος για να χρειάζεται να αφιερώνεις τόσες ώρες στο διάβασμα.»

«Δεν είμαι μεγάλος;»

«Όχι, εννοώ ότι τα μαθήματα σου τώρα θα έπρεπε να είναι σχετικά εύκολα. Αν κάθεσαι πέντε ή έξι ώρες τώρα και διαβάζεις, σε λίγα χρόνια τι θα κάνεις; Θα ξενυχτάς;»

«Δεν... ξέρω...» είπα και έφαγα την πρώτη μπουκιά.

Η μητέρα μου έφερε και το δικό της πιάτο και έκατσε απέναντι μου για να αρχίσει να τρώει. Το σπίτι είναι πάντα ήσυχο όταν τρώμε, είτε ήταν με τον πατέρα μου, είτε ήταν με τη μητέρα μου. Μακάρι να είχε τέτοια ησυχία και στο σχολείο. Όλη μέρα, σε κάθε μάθημα, ακούγεται βαβούρα. Δεν το αντέχω καθόλου. Είναι σαν να με ακολουθεί ένα σμήνος από έντομα. Λες να υπάρχει μαγεία Ησυχίας; Πλάκα θα είχε. Να μπορούσα με μία σκέψη να τους κάνω όλους να ησυχάσουν και να μην ακούω τίποτα εκτός από τις σκέψεις μου. Ίσως υπάρχει κάποιος στην Ειδική Υπηρεσία με τέτοια μαγεία. Αλλά σε τι θα χρησίμευε στην Ειδική Υπηρεσία μία τέτοια μαγεία; Ίσως σε κατασκοπίες; Μπορεί.

«Μαμά»

«Ναι Σπαρκ;»

«Υπάρχει πράκτορας με μαγεία Ησυχίας;»

«Ησυχίας; Δεν νομίζω να έχω ακούσει τέτοιο πράγμα. Έχω ακούσει όμως πως με μαγεία Ανέμου μπορούν μερικοί πράκτορες να αλλάζουν την κατεύθυνση των ηχητικών κυμάτων, δημιουργώντας έτσι ήσυχες περιοχές.»

«Πολύ περίεργο αυτό. Γιατί να τα κάνουν όλα αυτά με την μαγεία Ανέμου και να μην μπουν κατευθείαν στη δράση;»

«Δεν είναι όλα δράση Σπαρκ. Ο Βασι στα είπε αυτά; Καλά μπορεί να είναι και η ειδικότητα μου, αλλά εγώ έχω παλέψει συνολικά δύο φορές για να ξεφύγω ή να πιάσω κάποιον, αλλά έχουμε και ειδικούς στις μάχες, οπότε και σε κίνδυνο να είμαστε, έρχονται και μας βοηθάνε οι συνάδελφοι μας. Επίσης, δεν είναι καθόλου περίεργο. Είναι αντιθέτως πολύ δημιουργικό κατά την γνώμη μου. Η δημιουργικότητα να ξέρεις ανταμείβεται στην Ειδική Υπηρεσία. Υπάρχουν πράκτορες που δεν είναι σας εμένα ή τον πατέρα σου με την μαγεία, δηλαδή...»

«Δεν έχουν γενετικά πλεονεκτήματα;» συμπλήρωσα την πρόταση της.

«Μπορείς να το πεις και έτσι. Δεν ξέρω πως λέγεται αλλά τέλος πάντων. Αυτοί οι πράκτορες χρησιμοποιούν με δημιουργικούς τρόπους τις μαγείες τους και βοηθάνε σε κάθε υπόθεση με τον δικό τους τρόπο. Ας πούμε, εκτός από αυτό με τα ηχητικά κύματα που σου είπα, ένας συνεργάτης μου έχει μαγεία Αντικατάστασης, αλλά δεν είναι γενετικά πλεονεκτήματα όπως το λες εσύ. Το μόνο που μπορεί ουσιαστικά να κάνει είναι να αλλάξει θέση σε δύο αντικείμενα ή ανθρώπους που είναι εντός του οπτικού του πεδίου, και αυτό πειραματικά έχουμε δει πως μπορεί να το κάνει μία φορά κάθε τέσσερις περίπου ώρες. Σε συνδυασμό όμως με την δική μου μαγεία Αλλαγής Όψης, μπορώ να υποδυθώ έναν “συνεργάτη” του υπόπτου που κατασκοπεύουμε και ο συνάδελφος μου να αλλάξει τη θέση μου με αυτή του ”συνεργάτη” του υπόπτου, κάνοντας τη δουλειά μου πολύ πιο εύκολη. Όσον αφορά τον “συνεργάτη”, συνήθως έχουμε μαζί μας και έναν άλλο συνάδελφο που είναι ειδικός στις μάχες ο οποίος ηρεμεί την κατάσταση χωρίς να χρειάζεται ο “αδύναμος” μας συνάδελφος να μπει σε κίνδυνο.»

«Γιατί ξαφνικά τώρα άρχισες να μιλάς για την Ειδική Υπηρεσία;» ρώτησα την μητέρα μου επειδή παραξενεύτηκα που άρχισε να μιλάει φαινομενικά μόνη της για την δουλειά της.

«Δεν σου αρέσει;» μου απάντησε

«Όχι, μου αρέσει. Απλά... δεν μου έχετε ξαναμιλήσει ποτέ για την δουλειά σας στην Ειδική Υπηρεσία.»

«Δεν νομίζω ότι μας έχεις ρωτήσει και ποτέ. Εγώ νόμιζα πως δεν σε ενδιέφερε η Ειδική Υπηρεσία, για αυτό δεν άνοιγα ποτέ συζήτηση για τέτοια θέματα.»

«Με ενδιαφέρει, απλά... δεν ξέρω αν μπορώ να γίνω πράκτορας.»

«Μη το λες. Σου είπα, δεν χρειάζεται να έχεις “γενετικές ανωμαλίες”, ή όπως το είπες. Μπορείς να γίνεις ανακριτής για παράδειγμα. Οι ανακριτές είναι η ειδικότητα με τους περισσότερους “αδύναμους” χρήστες μαγείας, γιατί είναι ξεκάθαρα ψυχολογική δουλειά. Θες να γίνεις πράκτορας Σπαρκ;» με ρώτησε η μητέρα μου καθώς σηκώθηκε από το τραπέζι για να πάρει τα πιάτα μας και να τα πάει στον νιπτήρα. 

«Μου το πρότεινε ο Βασι, αλλά μάλλον το είπε για πλάκα...»

«Που ξέρεις ότι δεν το εννοούσε; Δεν ξέρω αν το βλέπεις καθόλου εσύ, αλλά οι κεραυνοί σου είναι λίγο ασταθής από τότε που μας άφησε ο πατέρας σου.»

«Τι εννοείς;»

«Κοιτάξου λίγο στον καθρέφτη και έρχομαι. Δώσε μου δύο λεπτά να πλύνω τα πιάτα μας και έρχομαι.»

Πηγαίνοντας στον καθρέφτη αναρωτήθηκα τι εννοούσε η μητέρα μου. “Οι κεραυνοί σου είναι ασταθείς”. Τι μπορεί να σημαίνει αυτό; Μήπως μπορώ πλέον να δημιουργώ μεγάλους κεραυνούς όπως ο πατέρας μου; Δεν νομίζω. Δεν έχω γενετικά πλεονεκτήματα, αν είχα θα φαινόταν από πιο μικρή ηλικία. Αλλά είναι κάπως περίεργο και οι δύο γονείς μου να έχουν γεννηθεί με φυσική αντοχή στα μειονεκτήματα των μαγειών τους και εγώ όχι. Μήπως είμαι υιοθετημένος; Αποκλείεται, αφού έχω την μαγεία του πατέρα μου και η μητέρα μου σίγουρα ξέρει αν με γέννησε αυτή η όχι, οπότε δεν νομίζω να είναι αυτό.

Κοιτώντας τον καθρέφτη το είδα. Γύρω από το σώμα μου υπήρχε μία μικρή στρώση από κεραυνούς που αναβόσβηναν λες και ήταν χριστουγεννιάτικα φωτάκια που κάνουν κύκλους γύρω από το σώμα μου. Ο ρυθμός τους βέβαια ήταν λίγο ασυγχρόνιστος, κάποιοι κεραυνοί αναβόσβηναν λίγο πιο γρήγορα από κάποιους άλλους με αποτέλεσμα όσο πιο πολύ ώρα τους κοιτούσα να άλλαζε το μοτίβο της σειράς με την οποία αναβοσβήνουν. Δεν ξέρω και αν θεωρούνται κεραυνοί βέβαια, αυτά τα μικροσκοπικά “ηλεκτρικά κύματα” θα αποκαλούσε κανείς αν τα έβλεπε. Πιο πολύ με σπίθες παρά με κεραυνούς μοιάζουν, έτσι μικρούλες και τυχαίες που φαίνονται.

«Είδες;» είπε η μητέρα σκύβοντας δίπλα μου και ακουμπώντας το ένα της χέρι στον ώμο μου. \textasciitilde Μα καλά, δεν έχεις κοιταχτεί καθόλου στον καθρέφτη τρεις μέρες τώρα;\textasciitilde 

«Έχω κοιταχτεί, αλλά δεν το είχα προσέξει» της απάντησα.

«Τι είπες Σπαρκ;» με κοίταξε η μητέρα μου προβληματισμένη.

«Ότι έχω κοιταχτεί...»

«Ναι αυτό το άκουσα, αλλά που κολλάει αυτό;»

«Εσύ δεν ρώτησες αν έχω κοιταχτεί στον καθρέφτη;»

«Δεν θυμάμαι να ρωτάω κάτι. Τουλάχιστον όχι φωναχτά.»

«Α... πρέπει να ήταν η φαντασία μου...» είπα και κοίταξα το πάτωμα.

«Δεν πειράζει Σπαρκ, και εγώ λέω κάποιες φορές άσχετα πράγματα. Μπορεί να πήρες και από μένα κάτι τελικά» είπε η μητέρα μου και μου χαμογέλασε. «Κοίτα όμως τώρα. Αυτοί οι τυχαίοι κεραυνοί γύρω από το σώμα σου, ακίνδυνοι μεν, δεν είναι φυσιολογικοί. Συνήθως όποτε κάποιου η μαγεία αρχίζει και γίνεται ανεξέλεγκτη, αυτό σημαίνει πως είναι συναισθηματικά ασταθής, ή τουλάχιστον αυτό έχω παρατηρήσει. Για αυτό σου είχα πει εκείνη τη μέρα να χαμογελάς.» \textasciitilde Βέβαια, στη δική σου περίπτωση είναι τελείως ακίνδυνη γιατί οι κεραυνοί σου δεν μπορούν να προκαλέσουν ζημιά σε κάποιον ή κάτι όποτε δεν χρειάζεται να ασχοληθούμε με αυτόν τον τομέα.\textasciitilde  «Για πες μου, πώς είσαι συναισθηματικά τις τελευταίες μέρες; Είναι κάτι που σε προβληματίζει εκτός από τον θάνατο του μπαμπά;»

«Δεν... με προβληματίζει κάτι άλλο. Μόνο ο μπαμπάς, αλλά...»

\textasciitilde Αλλά;\textasciitilde 

«Οι τελευταίες λέξεις που τον άκουσα να λέει εκείνο το πρωί ήταν κάτι σαν “Δεν ξέρω ποιος είμαι, δεν ξέρω που είμαι” και κάτι άλλα που δεν τα θυμάμαι ακριβώς. Τι σημαίνει αυτό; Τι είχε περάσει ο μπαμπάς στην Ειδική Υπηρεσία και έλεγε τέτοιες φιλοσοφίες; Γιατί, αν ήταν και αυτός συναισθηματικά ασταθής, δεν είχε και αυτός κεραυνούς γύρω του; Θα μου πεις, αυτός μπορεί να τους ελέγξει πιο καλά από εμένα τους κεραυνούς του, αλλά και πάλι δεν θα έπρεπε να υπάρχουν σημάδια; Είναι σαν...»

\textasciitilde Σπαρκ, δεν ήξερα ότι σε προβληματίζει τόσο η κατάσταση του πατέρα σου.\textasciitilde 

«Προφανώς και με προβληματίζει η κατάσταση του πατέρα μου. Έγιναν όλα πάρα πολύ ξαφνικά και δεν μπορώ να καταλάβω το γιατί.»

Η μητέρα μου με κοίταξε με σοβαρό ύφος, κουνώντας το κεφάλι της λίγο προς τα αριστερά. 

«Σε ποιον απάντησες τώρα;» ρώτησε με σοβαρό τόνο.

«Σε... εσένα;» της απάντησα προβληματισμένος εγώ αυτή τη φορά.

«Αχ δεν ξέρω πολύ φυσική, ούτε βιολογία, αλλά θέλω να δοκιμάσω κάτι» είπε η μητέρα μου και πήγε στην κουζίνα να πάρει κάποια πράγματα, αφήνοντας με μπροστά στον καθρέφτη μόνο μου για λίγη ώρα. 

Όταν γύρισε, είχε ένα χαμόγελο στο πρόσωπο της λες και είχε κερδίσει κάποιο μεγάλο χρηματικό έπαθλο. Έσκηψε πάλι δίπλα μου και μου αποκρίθηκε.

«Τι έχω πίσω από την πλάτη μου;»

«Που θες να ξέρω;» της απάντησα, χωρίς να καταλάβω την σημασία της ερώτησης. Η μητέρα μου σήκωσε το δεξί της χέρι που είχε πίσω από την πλάτη της, ρίχνοντας κάποια πράγματα πράγματα στο πάτωμα, ένα εκ των οποίων ήταν ένα μαχαίρι το οποίο χτύπησε στο πλακάκι και έκανε μία μικρή ρωγμή στο πάτωμα.

«Ουπς. Δεν πειράζει, θα το φτιάξουμε άλλη ώρα» είπε η μητέρα μου και ακούμπησε το χέρι της στον ώμο μου. «Τώρα τι έχω πίσω από την πλάτη μου;»\textasciitilde ...μανταρίνι, πιρούνι, μανταρίνι, πιρούνι, μανταρίνι...\textasciitilde 

«Τι μανταρίνι πιρούνι; Μαμά είσαι καλά;»

«Μπίνγκο» είπε η μητέρα μου και μου έδειξε το αριστερό της χέρι στο οποίο κρατούσε ένα μανταρίνι και ένα πιρούνι. Χαρούμενη με την απάντηση μου, μάζεψε και τα υπόλοιπα πράγματα που είχαν πέσει και τα πήγε στην κουζίνα. 

«Τι σημαίνει μπίνγκο; Δεν καταλαβαίνω τι προσπαθείς να κάνεις μαμά;» την ρώτησα ακολουθώντας την στην κουζίνα.

«Έλα ρε Σπαρκ. Πριν λίγο για αυτό ακριβώς το πράγμα σου μιλούσα.»

«Ποιο πράγμα;»

«Για την δημιουργικότητα» είπε η μητέρα μου καθώς άφηνε το μαχαίρι και το πιρούνι μέσα στο συρτάρι. «Είσαι μικρός ακόμη οπότε μπορεί να μην το καταλάβεις, αλλά κάτσε να σου εξηγήσω τι παρατήρησα με πιο ευθύ τρόπο.»

Η μητέρα μου πήγε πάλι στο τραπέζι. Την ακολούθησα και κάθισα δίπλα της. Που κολλάει η δημιουργικότητα; Τι σχέση έχουν οι κεραυνοί μου με τη δημιουργικότητα;

«Λοιπόν Σπαρκ. Ας αρχίσουμε από τα βασικά... περίπου. Δεν είμαι και τελείως σίγουρη, αλλά νομίζω πως τα νεύρα του σώματος μας στέλνουν ηλεκτρικά σήματα στον εγκέφαλο.»

«Ωραία. Αλλά αυτό...»

«Περίμενε. Εσύ από μικρός δεν μπορούσες να φτιάξεις μεγάλους, σχετικά, κεραυνούς, μόνο αυτές τις μικρές σπίθες που είδες στον καθρέφτη.»

«Το ξέρω, αλλά που...»

«Περίμενε. Θέλω να ακούσεις όλο αυτό που θέλω να πω και μετά να μου πεις ό,τι θες. Λοιπόν, που ήμουνα; Α ναι. Παρόλο όμως που δεν μπορείς να παράγεις μεγάλους, δεν παύει να ισχύει ότι έχεις μαγεία κεραυνών. Απ’ ότι φαίνεται, ακόμη και αυτούς τους μικρούς κεραυνούς μπορείς να τους ελέγχεις και...» η μητέρα μου σταμάτησε για λίγο να μιλάει, συνεχίζοντας όμως να με κοιτάζει επίμονα.

«Και;»

«Ναι αυτό είναι. Δώσε μου το χέρι σου» είπε η μητέρα μου και άπλωσε το δεξί της χέρι περιμένοντας να της δώσω και γω το δικό μου. Της έδωσα το χέρι μου και την κοίταξα στα μάτια.

\textasciitilde Και να διαβάζεις τις σκέψεις άλλων\textasciitilde\  είπε η μητέρα μου χωρίς να κουνήσει τα χείλια της. Έμεινα άφωνος.

«Πώς; Εγώ το έκανα αυτό; Μαμά βοήθεια, δεν καταλαβαίνω.»

«Η θεωρία μου είναι πως αγγίζοντας με, μπορείς να οδηγήσεις κάπως τους “κεραυνούς” σου μέσα από τα νεύρα μου και να μου διαβάσεις τις σκέψεις. Ίσως δημιουργείς κάπως ηλεκτρικό πεδίο ή κάτι, πραγματικά δεν ξέρω φυσική, για αυτό το λέω μόνο σαν θεωρία. Αυτό όμως δεν αναιρεί τη χρησιμότητα αυτής της ικανότητας σου. Δεν ξέρουμε ακόμη πως ακριβώς το κάνεις, αλλά μπορείς να τη χρησιμοποιήσεις αν θες να μπεις στην Ειδική Υπηρεσία, ειδικά ως ανακριτής.» είπε η μητέρα μου και μου χαμογέλασε πλατιά, λες και προσκαλούσε να μπω στην Ειδική Υπηρεσία από οκτώ χρονών παιδί.

«Θα... το σκεφτώ. Είμαι μικρός ακόμη.» της απάντησα σκήβοντας το κεφάλι.

Η μητέρα μου δεν μου απάντησε. Κατάλαβε ότι μάλλον ήθελα λίγο χρόνο να επεξεργαστώ την κατάσταση. Πήρα τα βιβλία μου και πήγα στο δωμάτιο μου. Ηλεκτρικά σήματα; Διάβασμα σκέψεων; Πώς γίνεται από το πουθενά να ανακαλύπτουμε κάτι που με την πρώτη ματιά δεν έχει καμία σχέση με την μαγεία μου, αλλά να συνδέεται μάλλον άμεσα και ίσως είναι ικανότητα κλειδί για την υπόλοιπη ζωή μου;

Ανακριτής ε; Καλό ακούγεται. Ούτε μάχες, ούτε κίνδυνος. Μόνο υπό έλεγχο συζητήσεις με υπόπτους στα κεντρικά της Ειδικής Υπηρεσίας. Μπορεί να μην χρειάζεται καν συζήτηση βέβαια, γιατί αν μπορώ να διαβάζω τις σκέψεις των υπόπτων, τότε η δουλειά μου θα είναι πολύ εύκολη. Απ’ ότι κατάλαβα, όμως, μπορώ να διαβάσω μόνο τις τρέχον σκέψεις κάποιου, όχι και τις αναμνήσεις του, οπότε ίσως τελικά χρειάζεται λίγη συζήτηση με τους υπόπτους, αλλά δεν πειράζει. Από την άλλη, δεν θα πηγαίνω ταξίδια, όπως έλεγε ο Βασι. Αλλά τι ταξίδια να κάνουν οι πράκτορες τώρα που το σκέφτομαι. Πολύ αγχωτικό θα είναι γιατί λογικά θα ακολουθούν κάποια σχέδια για να ολοκληρώσουν αποστολές, ή κάτι τέτοιο. Και αν κάνω κάτι λάθος και πάει όλο το σχέδιο στράφι; Άστο καλύτερα. Πιο ωραία θα είναι να δουλεύω μόνος μου σαν ανακριτής και να βρίσκω μόνος μου τι να κάνω και πότε να το κάνω. Ίσως η Ειδική Υπηρεσία δεν είναι και όσο κακή δουλειά όσο την φανταζόμουν. Ειδικά με το διάβασμα σκέψεων θα είναι παιχνιδάκι. Θα τελειώνω τις ανακρίσεις μου μέσα σε μερικά λεπτά και την υπόλοιπη μέρα θα κάθομαι και θα χαλαρώνω. Είμαι χαρούμενος μόνο και μόνο που το σκέφτομαι.

Βγήκα από το δωμάτιο μου βιαστικά και πήγα στην μητέρα μου χαρούμενος.

«Τι έγινε Σπαρκ. Οι κεραυνοί σου έφυγαν. Είσαι καλύτερα τώρα;» ρώτησε η μητέρα μου χαρούμενη.

«Θέλω να μπω στην Ειδική Υπηρεσία.»

\chapter{}

Το τελευταίο κουδούνι ήχησε και σήμερα. Όλοι σηκώθηκαν για να επιστρέψουν σπίτια τους. Βαρετή μέρα σήμερα. Δεν μάθαμε κάτι καινούργιο. Δηλαδή μάθαμε... αλλά τίποτα τόσο σημαντικό... νομίζω... Δεν πειράζει όμως, γιατί σήμερα θέλω να μιλήσω στον Βασι για την νέα μου ικανότητα. Ελπίζω να ενθουσιαστεί όσο και εγώ όταν κατάλαβα τι ακριβώς μπορώ να κάνω. Μακάρι να με ακούσει βέβαια, γιατί ποτέ δεν με ακούει. Ή τουλάχιστον έτσι νομίζω. Μπορεί να κάνω λάθος, όπως έκανα με την μητέρα μου και την Ειδική Υπηρεσία.

«Σπαρκ πάμε;» με ρώτησε ο Βασι.

«Ναι» του απάντησα. «Όσο περπατάμε θέλω να σου πω και κάτι...»

«Ωραία αλλά γρήγορα. Θέλω να πάω γρήγορα σπίτι γιατί η μαμά μου είπε πως σήμερα θα έχουμε μπριζόλα και μου έχει ανοίξει η όρεξη.»

«Δώσε μου το χέρι σου» του είπα καθώς βγαίναμε από τους τοίχους του σχολείου.

«Γιατί; Τι θες να κάνεις; Μη μου πεις ότι θα με σοκάρεις με τους μεγάλους σου κεραυνούς. Ωχ, φοβάμαι τώρα. Σε παρακαλώ να είσαι προσεκτικός, γιατί θέλω πολύ να φάω μπριζόλα. Μην με...»

«Μη μιλάς για λίγο» τον διέκοψα και του έκανα νεύμα να μου δώσει το χέρι του. Αυτός τελικά δέχτηκε και μου έδωσε σιγά σιγά το χέρι του, τραβόντας το πίσω δύο φορές για να τσεκάρει αν θα τον σοκάρω με τους υποτιθέμενους μεγάλους μου κεραυνούς. «Μη τραβιέσαι πίσω. Πρέπει να με ακουμπάς για να δουλέψει αυτό που θέλω να κάνω.»

«Καλά» είπε ο Βασι και μου έδωσε το χέρι του κλείνοντας ταυτόχρονα τα μάτια του φοβισμένος. \textasciitilde ...μη με σοκάρεις, μη με σοκάρεις, μη με...\textasciitilde 

«Δεν θα σε σοκάρω είπαμε. Τώρα σκέψου κάτι άσχετο χωρίς να μιλήσεις.» του είπα κοιτώντας τον στα μάτια.

\textasciitilde Κάτι άσχετο, ε; Η μπριζόλα δεν είναι άσχετη γιατί την ανέφερα πιο πριν. Θα σκεφτώ το αντίθετο της μπριζόλας τότε. Τι μπορεί να είναι αυτό; Λαχανικά; Όχι δεν είναι τα λαχανικά, αφού τρώμε σαλάτα μαζί με την μπριζόλα. Τι άλλο θα μπορούσε να είναι το αντίθετο της μπριζόλας;\textasciitilde 

«Σταμάτα να προσπαθείς να βρεις το αντίθετο της μπριζόλας και σκέψου κάτι.»

«Πώς ήξερες πως ψάχνω να βρω αντίθετο... Αυτό είναι το κόλπο σου; Μαντεύεις τι μπορώ να σκεφτώ; Ε μα προφανώς και μπορείς να μαντέψεις τι μπορώ να σκεφτώ, αφού κάνουμε παρέα τόσο καιρό, είναι σαν να είμαστε πρακτικά αδέλφια. Το χέρι που κολλάει όμως;»

«Δεν “μαντεύω”, διαβάζω τις σκέψεις σου χρησιμοποιώντας τη μαγεία μου.» του απάντησα κοιτώντας τον επίμονα για να δω την αντίδραση του.

«Α... εντάξει...» μου απάντησε και σταμάτησε να μου κρατάει το χέρι. «Αυτό είναι όλο;»

«Τι εννοείς “αυτό είναι όλο”; Δεν είσαι καθόλου χαρούμενος για μένα; Που επιτέλους καταφέρνω να κάνω κάτι χρήσιμο με την μαγεία μου;»

«Δεν καταλαβαίνω που ακριβώς μπορεί να χρησιμεύσει αυτή η “δεξιότητα” σου, αλλά αν είσαι εσύ χαρούμενος, θα είμαι και εγώ χαρούμενος» μου είπε ο Βασι και μου χαμογέλασε.

«Δεν το εννοούσα έτσι. Εννοώ πως τώρα θα μπορέσουμε να πάμε επιτέλους και οι δύο στην Ειδική Υπηρεσία. Εσύ προφανώς μπορείς να μπεις εύκολα γιατί έχεις γενετική αντοχή στα μειονεκτήματα της μαγείας σου, αλλά εγώ βρήκα επιτέλους τι μπορώ να κάνω για να γίνω πράκτορας της Ειδικής Υπηρεσίας.»

«Και οι κεραυνοί σου; Δεν μεγάλωσαν; Δεν θες να γίνεις ειδικός στις μάχες μαζί μου;»

«Δεν λέω ότι δεν θέλω να γίνω ειδικός στις μάχες, απλά... δεν μπορώ... γιατί οι κεραυνοί μου δεν “μεγάλωσαν” όπως λες κι εσύ. Κατέληξες μόνος σου σε αυτό το συμπέρασμα...» του είπα κοιτώντας τον λίγο εκνευρισμένος. «Μάλλον θα γίνω... ανακριτής»

«Και οι μάχες που λέγαμε; Δεν θα πηγαίνουμε μαζί σε αποστολές; Σε ταξίδια; Προτιμάς δηλαδή να κάθεσαι σε ένα δωμάτιο όλη μέρα και να “διαβάζεις τις σκέψεις” εγκληματιών; Αυτό λες;» άρχισε να μου φωνάζει ο Βασι κουνώντας εκφραστικά τα χέρια του σαν Ιταλός.

«Δεν... είπα αυτό. Απλά πιστεύω πως είναι καλύτερο για μένα, που δεν έχω γενετική αντοχή, να γίνω... ανακριτής...» είπα σκύβοντας το κεφάλι μου γιατί αγρίεψα τον μόνο μου φίλο και αγχώθηκα μην τσακωθούμε και χωριστούμε. «...αυτό είναι όλο.»

«Δεν μου αρέσει η επιλογή σου. Αλλά ποιος είμαι εγώ να σου πω τι να κάνεις. Πάντως να ξέρεις, αν ποτέ θες κάποιον ειδικό στις μάχες, εγώ θα είμαι εκεί να σε βοηθήσω σαν μεγάλος αδελφός.» μου είπε και μου χαμογέλασε. Πώς το κάνει πάντα αυτό; Αλλάζει συναισθήματα και κάνει λες και δεν έγινε τίποτα. Με είδε που αγχώθηκα και είπε να αποκλιμακώσει την κατάσταση; 

«Τι σκέφτεσαι Σπαρκ; Όλα καλά;» με ρώτησε ο Βασι

«Ναι όλα καλά. Πρώτη φορά με ρωτάς τι σκέφτομαι. Σε... ευχαριστώ που ρώτησες.»

«Δεν σε ρωτάω πρώτη φορά. Απλά όταν σε ρωτάω συνήθως είσαι πιο αφαιρεμένος και με αγνοείς.»

«Α... συγγνώμη...»

«Δεν χρειάζεται να λες συγγνώμη κάθε φορά που κάνεις κάτι ρε Σπαρκ. Φίλοι είμαστε. Καταλαβαίνω τι θες να μου πεις με μία ματιά, τουλάχιστον έτσι νομίζω τις περισσότερες φορές.»

«Α... συγγνώμη... α... ναι δεν...» είπα και ο Βασι άρχισε να γελάει, προφανώς με το “λάθος” μου.

«Δεν πειράζει, και μένα μου ξεφεύγουν λέξεις κατά λάθος.» είπε και με χτύπησε στον ώμο. «Κοίτα, φτάσαμε στο φανάρι. Βλέπεις πόσο γρήγορα περνάει η ώρα όταν περνάς καλά;»

Δεν του απάντησα λεκτικά. Τον κοίταξα απλά και χαμογέλασα.

«Σου ξαναλέω, ανακριτή Σπαρκ, αν θελήσεις ποτέ κάποιον ειδικό στις μάχες, είμαι ένα τηλεφώνημα μακριά.» μου είπε ο Βασι καθώς γύρναγε δεξιά για να πάει σπίτι του.

«Είμαστε ακόμη οκτώ χρονών, δεν ξέρω γιατί επιμένεις τόσο πολύ. Και αν δεν μας δεχτούν στην Ειδική Υπηρεσία όταν κάνουμε αίτηση; Που την βρίσκεις τόση αισιοδοξία;»

«Εσύ που την βρίσκεις τόση απαισιοδοξία; Πλάκα κάνω. Δεν θέλω να σε πειράξω ξανά σήμερα γιατί σε βλέπω είσαι λίγο ντεμί. Αουφίνταρζεν τώρα. Τα λέμε την Δευτέρα.»

«Καλή συνέχεια Βασι.» του είπα και γύρισα προς το φανάρι που ήταν κόκκινο. 

“Ειδικός στις μάχες”; Ναι καλά. Άμα κάποιος ενήλικας του επιτεθεί αυτή τη στιγμή, ότι μαγεία και να έχει, ο Βασι δεν θα μπορέσει να κάνει τίποτα. Σίγουρα ειρωνεύεται λίγο τον εαυτό του, δε γίνεται να έχει τόση αυτοπεποίθηση. Τέλος πάντων, δεν πρέπει να αργήσω πάλι, γιατί την προηγούμενη φορά η μαμά είχε μεικτά συναισθήματα. Δεν καταλαβαίνω όμως γιατί επιμένει τόσο, είμαστε ακόμη οκτώ χρονών, έχουμε όλη τη ζωή μπροστά μας. Μπορεί να αλλάξουν τα πράγματα, ποιος ξέρει, και να γίνει ο Βασι ανακριτής και εγώ ειδικός στις μάχες. Πλάκα θα είχε.

\chapter{}

«Πότε πέρασε ο καιρός ρε Σπαρκ. Την μία είμαστε στο σχολείο, την άλλη είμαστε είκοσι και συζητάμε σε μία καφετέρια τα προβλήματα μας. Δηλαδή εγώ σου λέω τα προβλήματα μου, δεν ξέρω αν αυτή τη στιγμή με   ακούς ή έχεις αφαιρεθεί πάλι.»

«Όχι σε ακούω Βασι, απλά παράλληλα σκέφτομαι τι θα γίνει το απόγευμα. Αγχώνομαι πολύ. Αν δεν μπορέσω να επιδείξω την ικανότητα μου σωστά στον αξιολογητή που θα μας εξετάσει; Αν πάλι μου δώσουν ειδικότητα που δεν μου ταιριάζει; Αν...»

«Δεν χρειάζεται να αγχώνεσαι ρε φίλε τόσο πολύ. Οι αρνητικές σκέψεις οδηγούν σε αυτοεκπληρούμενη προφητεία. Χαλάρωσε λίγο. Τόσο καιρό που σε ξέρω όλο αρνητικές σκέψεις είσαι, έτσι νομίζω τουλάχιστον. Από τη στιγμή που κάναμε τις αιτήσεις μας στην Ειδική Υπηρεσία και μας κάλεσαν για αξιολόγηση, αυτό σημαίνει πως τους ενδιαφέρουν οι ικανότητες μας. Εξάλλου, εσένα θα σε ξέρουν ήδη λόγω των γονιών σου. Άλλος ένας λόγος να μην αγχώνεσαι... ή μάλλον όχι, γράψε λάθος, μπορεί να έχουν υψηλές προσδοκίες λόγω των γονιών σου. Αλλα ποιος νοιάζεται; Θα πάμε χαλαρά να επιδείξουμε τις μαγείες μας, ή τα “κόλπα” μας στην δική σου περίπτωση, και αν μας δεχτούν, μας δέχτηκαν. Δύσκολα, όμως, να ξέρεις σε απορρίπτουν αν σε έχουν ήδη καλέσει για αξιολόγηση, η οποία δεν είναι τίποτα παραπάνω από μία συνέντευξη. Δεν ξέρω γιατί την αποκαλούν “Αξιολόγηση”. Ίσως για να ακούγεται πιο κλασάτο, που δεν ακούγεται, αλλά τέλος πάντων. Ίσως γιατί αξιολογούν την μαγεία μας και μας δίνουν κατάλληλη ειδικότητα; Μόνο αυτό μπορώ να σκεφτώ. Α μπορεί...» άρχισε να μονολογεί πάλι ο Βασι.

Χαλαρά, ε; Δίκιο έχει. Τόσα χρόνια που προπονώ το διάβασμα σκέψεων δεν μου έχει προκύψει κάτι παράξενο. Δεν ξέρω καν γιατί αγχώνομαι. Ίσως γιατί είναι η πρώτη μου δουλειά; Ίσως γιατί θα έχουν υψηλές προσδοκίες για μένα, όπως είπε ο Βασι; Αλλά πως ακριβώς μπορεί να πάει στραβά. Τώρα που το σκέφτομαι δεν θα χρειαστεί να κάνω και κάτι περίπλοκο. Απλά θα πω στον αξιολογητή ότι έχω μαγεία Κεραυνών, αλλά μπορώ να διαβάζω σκέψεις ξέρω γω και πιθανότατα θα μου δώσει την ειδικότητα που μου ταιριάζει, ανακριτής δηλαδή. Αν του το αναφέρω κιόλας μπορεί να τον οδηγήσω και σε αυτή την επιλογή. 

Α ναι, και ο πατέρας μου ήταν ανακριτής. Πώς δεν το έχω σκεφτεί τόσο καιρό. Καλά, ο πατέρας μου δεν ήταν κλασικός ανακριτής βέβαια, πήγαινε σε αποστολές με την ομάδα του, από ότι είχα καταλάβει, και έπαιρνε μέρος και σε μάχες. Λογικό μου ακούγεται, αφού είχε και αυτός μαγεία κεραυνών. Δεν κατάλαβα ποτέ όμως πως προέκυψε η ειδικότητα του ανακριτή στην περίπτωση του πατέρα μου. Ίσως ήξερε να ψυχολογεί τους απέναντι του. Μπορεί να χρησιμοποιούσε τους κεραυνούς του για να προκαλέσει ηλεκτροσόκ στους υπόπτους μέχρι να ομολογήσουν. Πιο βίαιο μεν, αλλά αν δούλευε, ποιος είμαι εγώ να το κρίνω;

«...και η κοπέλα μου δε συμφώνησε, αλλά ρε φίλε τι να κάνω; Από μικρός θέλω να μπω στην Ειδική Υπηρεσία. Θα αλλάξω το όνειρο μου για μία τυχαία που γνώρισα πριν δύο χρόνια; Στο κάτω κάτω, αν εν τέλη δεν μου αρέσει, απλά παραιτούμαι και αρχίζω άλλη δουλειά. Δεν θα πεθάνω κιόλας. Δεν συμφωνείς Σπαρκ;»

«Όχι και τυχαία...»

«Ναι εντάξει, αλλά δεν θα αλλάξω γνώμη έτσι απλά επειδή μου είπε η κοπέλα μου όχι. Καλά που να καταλαβαίνεις εσύ που δεν έχεις βρει ακόμη κοπέλα. Εσύ πλάκα πλάκα γιατί θες να μπεις στην Ειδική Υπηρεσία;» με ρώτησε ο Βασι και ήπιε την τελευταία γουλιά από τον καφέ του.

«Έχω τους λόγους μου» είπα προσπαθώντας να φανώ μυστήριος, αλλά μου φαίνεται ακούστηκα περισσότερο περίεργος παρά μυστήριος, καθώς όταν το είπα ο Βασι έσκασε στα γέλια.

«Εντάξει κύριε Σπαρκ. Κράτα τους λόγους σου μυστικούς. Θα μου τους πεις όταν νιώσεις έτοιμος, το ξέρω ότι θα μου τους πεις.» είπε ο Βασι και κάθισε πίσω στην καρέκλα του. Με το που ακούμπησε την πλάτη του όμως στην πλάτη της καρέκλας, τον πήρε τηλέφωνο η κοπέλα του. «Ούτε ένα λεπτό δεν μπορώ να κάτσω ρε φίλε. Έρχομαι σε λίγο, Σπαρκ.» είπε και σηκώθηκε για να απαντήσει στο τηλεφώνημα.

Εγώ συνέχισα να κοιτάω έξω από το τζάμι του καφενείου. Δεν υπήρχαν και πολλοί περαστικοί, που μου φάνηκε παράξενο γιατί ήταν 3:27 το μεσημέρι στο κέντρο της πόλης, ημέρα Πέμπτη. Ίσως οι περισσότεροι βρίσκονταν στις δουλειές τους. Ίσως έτυχε. Ότι και να συνέβαινε πάντως ήταν σίγουρα παράξενο. Μπορεί να ήταν απλά η φαντασία μου. 

Καθώς κοιτούσα έξω από το παράθυρο, είδα μία κοπέλα με καστανά μαλλιά πιασμένα σε κοτσίδα και ένα κλειδί αυτοκινήτου στο δεξί της χέρι που το κουνούσε σαν να ήταν μπρελόκ να πλησιάζει τον Βασι και να τον ρωτάει κάτι. Αυτός στην αρχή φάνηκε σαν να την απορρίπτει, κάνοντας νεύμα με το χέρι του, αλλά αυτή επέμενε και του έδειξε κάτι στο κινητό της. Ο Βασι στη συνέχεια της έδειξε προς μία κατεύθυνση και αυτή τον ευχαρίστησε και έφυγε. Δεν ήταν σίγουρα η κοπέλα του Βασι, καθώς με την κοπέλα του μιλούσε στο κινητό. Ποια να ήταν άραγε; Τι ήθελε από τον Βασι;

«Ποια ήταν αυτή η κοπέλα, Βασι;» τον ρώτησα όταν μπήκε πάλι μέσα στο καφενείο.

«Που θες να ξέρω. Απλά ήρθε και με ρώτησε αν ξέρω που είναι τα κεντρικά της Ειδικής Υπηρεσίας γιατί λέει τελείωσε η μπαταρία του κινητού της και δεν ήξερε που να πάει. Τυχερή ήταν που τα κεντρικά είναι δύο στενά από το καφενείο» είπε ο Βασι με νευριασμένο τόνο.

«Μόνο αυτό; Νόμιζα ότι θα ήταν κάτι παραπάνω...»

«Άσε μας ρε Σπαρκ. Δεν θέλω να εκνευριστώ περισσότερο» είπε ο Βασι και ακούμπησε την πλάτη του στην πλάτη της καρέκλας, βγάζοντας μία ανάσα ανακούφισης. Εγώ συνέχισα να κοιτάζω έξω από το τζάμι του καφενείου.

«Ξέρεις Σπαρκ, μπορεί να μην μιλάμε τόσο συχνά όσο τώρα όταν γίνουμε και οι δύο πράκτορες στην Ειδική Υπηρεσία. Τα προγράμματα μας θα είναι τελείως γεμάτα κάθε μέρα, οπότε θέλω για μία στιγμή την προσοχή σου.»

Γύρισα και τον κοίταξα. Είχε σκύψει προς το τραπέζι και είχε του τεντωμένο λες και ήθελε να κάνουμε χειραψία.

«Υποσχέσου μου πως ό,τι και να γίνει, θα μιλάμε τουλάχιστον τηλεφωνικός» είπε ο Βασι και απλωσε παραπάνω το χέρι του δείχνοντας μου πως όντως ήθελε να κάνουμε κάποιου είδους χειραψία.

«Αυτό είναι όλο;» του είπα δίνοντας του διστακτικά το χέρι μου.

«Ναι αυτό είναι όλο.» είπε ο Βασι και χαμογέλασε ενώ έπιασε το χέρι μου. \textasciitilde Δεν θέλω να σε χάσω Σπαρκ. Είσαι ο μόνος μου φίλος. Όλοι οι άλλοι με αποφεύγουν γιατί είμαι υπερκοινωνικός. Αν χάσω εσένα θα είμαι...\textasciitilde 

«Δεν χαθούμε Βασι. Μην κάνεις τέτοιες σκέψεις, όπως λες και συ. Αυτοεκ... πως το είχες πει; Κάτι για αυτοεκπληρούμενη προφητεία;»

«Ξέχασα ότι διαβάζεις σκέψεις...» είπε ο Βασι και κύλησε ένα δάκρυ στο αριστερό του μάγουλο. «Σε ευχαριστώ πολύ Σπαρκ. Πραγματικά μου έφτιαξες τη διάθεση.»

«Με την κοπέλα σου τι έγινε;» τον ρώτησα από απορία γιατί ήταν αρκετή ώρα έξω από το καφενείο και μίλαγε.

«Τα ίδια με πριν, μου λέει να μη πάω στην Ειδική Υπηρεσία. Μην ασχολείσαι. Δεν πρόκειται να βγάλεις άκρη.» μου είπε και έκανε νεύμα σε έναν υπάλληλο για να πληρώσουμε. «Δεν σε πειράζει να πληρώνουμε σιγά σιγά; Έχει πάει 4 και οι συνεντεύξεις μας είναι γύρω στις 5 και οι δύο, οπότε μπορούμε να πάμε να περπατήσουμε λίγο εδώ γύρω για να χαλαρώσουν τα πνεύματα.»

«Ναι, γιατί όχι.» του απάντησα και άφησα το ποσό που αντιστοιχούσε στον καφέ μου στο τραπέζι.

«Όχι όχι, σε παρακαλώ. Θέλω να σε κεράσω μία φορά» είπε ο Βασι και πλήρωσε κατευθείαν και τους δύο καφέδες μας όταν ήρθε ο υπάλληλος.

«Σε ευχαριστώ Βασι. Αν και... δεν χρειαζόταν.»

«Τι σημαίνει δεν χρειαζόταν; Φίλοι είμαστε. Μία σε κερνάω εγώ, μία με κερνάς εσύ. Απλά αυτή τη φορά ένιωσα πως πρέπει να κεράσω και το έκανα»

«Τι σημαίνει το ένιωσα; Πώς ξέρεις ότι “έπρεπε” να κεράσεις εσύ;»

«Αυτό είναι ένα μειονέκτημα σου Σπαρκ. Δεν μπορείς να σκεφτείς με συναισθήματα. Είσαι πολύ κλειστός. Δεν ξέρω καν αν έχεις συναισθήματα. Σε βλέπω σχεδόν κάθε μέρα και αισθάνομαι ότι μιλάω με ρομπότ. Και όποτε σε ρωτάω “τι σκέφτεσαι;” ή “πώς είσαι;” μου λες απλά όλα καλά και συνεχίζεις ότι έκανες χωρίς να με αφήσεις να σε βοηθήσω»

Έσκυψα το κεφάλι γιατί ο υψωμένος τόνος του Βασι με άγχωσε.

«Δεν στο λέω για να τσακωθούμε. Στο λέω σαν φιλική συμβουλή. Άρχισε να εκφράζεσαι περισσότερο. Τώρα που θα πας στην Ειδική Υπηρεσία και θα χρειάζεται να συνεργάζεσαι με άλλους πράκτορες, θα είσαι μουγκός και εκεί; Πάρε πρωτοβουλίες. Μίλα. Όπως μιλάς, μερικές φορές, σε μένα.»

«Σε ευχαριστώ για τη συμβουλή Βασι. Ας σηκωθούμε τώρα. Πάμε να περπατήσουμε.» είπα και σηκώθηκα από την καρέκλα μου

«Έτσι μπράβο.» είπε ο Βασι χαμογελώντας και με ακολούθησε.

Καθώς βγαίναμε από το καφενείο, ένας λίγο ψηλότερος από εμάς άντρας,  λίγο μεγαλύτερος ηλικιακά, με περίεργη φράντζα και ζακέτα μιας ποδοσφαιρικής ομάδας, μας πλησίασε από πίσω.

«Συγγνώμη καλοί μου κύριοι. Μήπως γνωρίζετε που βρίσκεται αυτή η διεύθυνση;» ρώτησε ο άντρας και μας έδειξε στο κινητό του την διεύθυνση της Ειδικής Υπηρεσίας.

«Είσαι ο δεύτερος που με ρωτάει σήμερα» πετάχτηκε ο Βασι. «Είναι κάποιου είδους τεστ; Αν ναι, μπορείτε να το κόψετε, περάσαμε.»

«Όχι καλέ μου κύριε. Βλέπετε, έχω μια συνέντευξη για την πρώτη μου δουλειά στις 4:45, οπότε χρειάζεται να παρευρεθώ όσο το δυνατόν νωρίτερα» είπε ο ψιλός άντρας, καθώς έβαζε το κινητό στην τσέπη του.

«Μήπως αυτή η “δουλειά” που λέτε σχετίζεται με την Ειδική Υπηρεσία;» ρώτησε ο Βασι και τον κοίταξε επίμονα.

«Ναι! Πώς το ξέρατε;» είπε ο άντρας έκπληκτος.

«Αφού μου ζητήσατε...»

«Ναι το ξέρω, πλάκα κάνω.» είπε ο άντρας γελώντας. «Απλά ύλπιζα να με βοηθούσατε κιόλας, γιατί βλέπετε βιάζομαι λίγο.»

«Κάτσε λίγο. Γιατί ήρθες και μας ρώτησες οδηγίες για το πώς να πας στην Ειδική Υπηρεσία ενώ έχεις κινητό; Δεν μπορείς να το ψάξεις μόνος σου;» ρώτησε ο Βασι, κάνοντας μισό βήμα πίσω. Άρχισα να παρατηρώ λίγο παραπάνω τις κινήσεις του Βασι και μου φαίνεται πως νομίζει ότι ο άντρας είναι λίγο ύποπτος.

«Σας άκουσα να μιλάτε για την Ειδική Υπηρεσία λίγο πιο πριν και ύλπιζα να κάνω καινούργιους φίλους, αυτό είναι όλο. Το όνομα μου είναι Κουπ παρεμπιπτόντως.» είπε ο άντρας με την περίεργη φράντζα και έδωσε το χέρι του στον Βασι.

«Πες το από την αρχή ρε Κουπ. Εγώ είμαι ο Βασι και αυτός εδώ είναι ο Σπαρκ.» είπε ο Βασι δίνοντας το χέρι του στον άντρα.

«Χάρηκα πολύ Σπαρκ.»

«Και γω...» του είπα δίνοντας επίσης το χέρι μου.

\textasciitilde Δεν νιώθω πολύ δυνατή μαγεία από αυτόν εδώ. Σίγουρα είναι στην Ειδική Υπηρεσία;\textasciitilde  διάβασα τις σκέψεις του Κουπ, αλλά δεν το σχολίασα.

«Ας πηγαίνω λοιπόν.» είπε ο Κουπ, φτιάχνοντας τη ζακέτα του και τινάζοντας λίγο το σήμα της ομάδας του.

«Γιατί δεν έρχεσαι μαζί μας;» τον σταμάτησε ο Βασι βάζοντας το χέρι του στο στήθος του Κουπ.«Βλέπεις έχουμε και εμείς αξιολόγηση στην Ειδική Υπηρεσία γύρω στις 5:00, οπότε μπορούμε να αρχίσουμε να πηγαίνουμε προς τα εκεί όλοι μαζί.»

«Βλέπετε βιάζομαι αρκετά, οπότε ίσως χρειαστεί να πάω μόνος μου...»

«Σπαρκ τι λες; Πάμε με τον Κουπ στην Ειδική Υπηρεσία; Δεν πειράζει να περιμένουμε και λίγο για την αξιολόγηση μας μέσα στα κεντρικά.» μου είπε ο Βασι κρατώντας το χέρι του στο στήθος του Κουπ και γυρνώντας το κεφάλι του αριστερά για να με κοιτάξει.

«Ναι, γιατί όχι.» του απάντησα και γύρισα για να παρατηρήσω την αντίδραση του Κουπ. Αυτός είχε σκύψει το κεφάλι του και είχε παρατήσει την προσπάθεια του να ξεφύγει από τον Βασι.

«Πολύ ωραία, λοιπόν, κύριοι.  Ας ξεκινήσουμε όμως γιατί δεν θέλω να αργήσω» είπε ο Κουπ, διώχνοντας το χέρι του Βασι εκνευρισμένος από το στήθος του.

«Σιγά καλέ. Ακόμη 4:07 είναι η ώρα και τα κεντρικά της Ειδικής Υπηρεσίας είναι 5 λεπτά περπάτημα από εδώ.» του απάντησε ο Βασι και του χαμογέλασε ειρωνικά. «Αλλά αν επιμένεις τόσο πολύ»

«5 λεπτά;» ρώτησα τον Βασι. «Αφού πριν είπες είναι δύο στενά από το καφενείο.»

«Λεπτομέρειες...» είπε ο Βασι και μου έγνεψε το κεφάλι λες και μου έλεγε να παρατηρήσω παραπάνω τον Κουπ, ο οποίος φαινόταν εκνευρισμένος με την ειρωνεία του Βασι. «Πάμε Κουπ. Δεν σε βλέπω καλά. Μπορεί να είναι το στυλ σου να φτάνεις πολύ νωρίς σε συνεντεύξεις. Δεν θα συνεχίσω να σε κρατάω πίσω.»

«Σας ευχαριστώ καλοί μου κύριοι...»

«Και μη μας λες κύριους. Είμαστε μικρότεροι από εσένα.» τον έκοψε ο Βασι.

«Σε αυτό κάνετε λάθος. Βλέπετε είμαι 22 και σεις φαίνεστε...»

«20» τον ξανα έκοψε ο Βασι. «Άρα εσείς κάνετε λάθος καλέ μου κύριε. Πάμε λοιπόν. Και να μας μιλάς στον ενικό από εδώ και πέρα.» του είπε ο Βασι και ξεκίνησε να περπατάει προς τα κεντρικά της Ειδικής Υπηρεσίας.

Άφησα τον Κουπ να περάσει μπροστά μου και συνέχισα από πίσω του. Αρκετά ύποπτος έχω να πω. Μπορεί απλά να ήθελε να κάνει φίλους πριν μπει στην Ειδική Υπηρεσία, ποιος ξέρει. Οι σκέψεις του όμως δεν με πείθουν για την αθωότητα του. Τι εννοεί όμως “δεν νιώθω δυνατή μαγεία από αυτόν εδώ”; Πως ξέρει πως είμαι “αδύναμος” με τους κεραυνούς μου; Αφού δεν αναφέρθηκε ούτε μία φορά η μαγεία στη συζήτηση μας. Μήπως η μαγεία του σχετίζεται με την αξιολόγηση της δύναμης των ανθρώπων που έχει απέναντι του;

«Κουπ... να σου κάνω μία ερώτηση;» ρώτησα τον Κουπ ακουμπόντας με το αριστερό μου χέρι τον ώμο του.

«Παρακαλώ Σπαρκ. Ρώτα ό,τι θες.» μου απάντησε, μειώνοντας λίγο ταχύτητα στο βηματισμό του για να έρθει δίπλα μου. Προσπάθησα να κρατήσω το χέρι μου πάνω στον ώμο του, αλλά μου έπεσε λίγο και το κράτησα στην πλάτη του για να μην προκαλέσω υποψία.

«Τι μαγεία έχεις, αν επιτρέπεται;» τον ρώτησα. Για κάποιο λόγο δεν σκεφτότανε κάτι εκείνη την ώρα.

«Μαγεία Ελαστικότητας. Μπορώ να τεντώσω τα άκρα μου μέχρι και 100 μέτρα. Καλά, 50 μέτρα, γιατί μετά τα 50 συνήθως είναι επίπονο. Αλλά δεν χρειάζεται κιόλας. 100 μέτρα είναι υπερβολή, απλά είναι η μεγαλύτερη απόσταση που έχει μετρηθεί σε άνθρωπο με την μαγεία μου, για αυτό λέω συνήθως 100 μέτρα για να εντυπωσιάζω τον κόσμο. Παρόλο που...»

«Πολυλογάδες φτάσαμε.» είπε ο Βασι, διακόπτοντας πάλι τον Κουπ και ανοίγοντας μας την πόρτα.

\chapter{}

«Πιάσε μου μία το νερό μου αδελφέ»

«Δεν μπορώ, έχω βάλει τα όλα τα νερά στο πορτμπαγκάζ.»

«Και δεν σκέφτηκες ότι μπορεί να διψάσει κάνεις μέσα στις 5 ώρες που θα είμαστε στο αυτοκίνητο;»

«Θα κάνουμε διάλειμμα ρε Ριντ στα μισά της διαδρομής. Μην αγχώνεσαι. Θα πάρουμε τότε τα νερά.»

«Ναι, αλλά αρχίσαμε πριν 20 λεπτά. Εγώ τι πρέπει να κάνω που διψάω τώρα; Να περιμένω άλλες δύο ώρες;»

«Έπρεπε να είχες πιεί νερό πριν φύγουμε. Αντί να μας κοιτάς σαν μούμια να παίζουμε πόκερ, θα μπορούσες απλά να είχες βάλει ένα ποτήρι νερό να πιεις.»

«Μη με εκνευρίζεις άλλο. Θέλω να είμαι ήρεμος όταν οδηγώ με πάνω από 100 χιλιόμετρα. Πες μου τουλάχιστον αν έβαλες το κολατσιό μας κάπου που να μπορούμε να το φτάσουμε.»

«Ναι, τα έχει ο Ρομπ. Μη μου πεις ότι πεινάς τώρα;»

«Σου είπα μη με φουντώνεις άλλο. Θέλω ηρεμία.»

«Καλά εντάξει, μη φωνάζεις κιόλας.»

«Μπαμπά χαλάρωσε και εσύ λίγο, αφού βλέπεις ότι ο θείος ο Ριντ οδηγάει.»

«Καλά Ρομπ. Όλοι εναντίον μου είστε σήμερα. Να δούμε τις επόμενες μέρες πως θα είστε, γιατί αν συνεχίσει αυτό το πράγμα, θα γυρίσω πίσω μόνος.»

«Και πως ακριβώς σκοπεύεις να το κάνεις αυτό αδελφέ;»

«Έχω τα κονέ μου εγώ, μην αγχώνεσαι.»

«Δεν αμφιβάλω...»

«Να βάλω τουλάχιστον μουσική; Ή θα είμαστε μουγκά σε όλη τη διαδρομή;»

«Βάλε, ξέρω γω. Περισσότερο με ενοχλεί η φωνή σου, παρά η μουσική.»

«Παιδιά τι σταθμό λέτε να βάλω;»

«Βρήκες και εσύ τώρα ανθρώπους να ρωτήσεις θείε. Αφού δεν ακούμε ράδιο εμείς. Πώς θες να γνωρίζουμε τους σταθμούς που υπάρχουν;»

«Έχει δίκιο μπαμπά. Βάλε ό,τι θες τώρα και αν δεν αρέσει σε κανέναν απλά θα σου πούμε να το αλλάξεις.»

«Έχεις δίκιο. Εσείς η νέα γενιά με τα ίντερνετ και τις εφαρμογές σας που να ξέρετε πως είναι το ράδιο. Βάζω στοίχημα πως αν σας έδινα ένα παλιό ράδιο δεν θα μπορείτε καν να το ανοίξετε.»

«Σιγά θείε, πόσο δύσκολο μπορεί να είναι;»

«Δεν είναι, απλά δεν είστε εξοικειωμένοι με τέτοιες παλιομοδίτικες συσκευές, οπότε θα σας πάρει αρκετή ώρα. Έχετε συνηθίσει να τα έχετε όλα έτοιμα. Ανοίγετε απλά μία εφαρμογή και της λέτε “παίξε αυτό” ή “βρες μου αυτό” και αφήνετε το μπλιμπλίκι να σας κάνει τη ζωή εύκολη.»

«Δεν είναι καλό αυτό θείε; Γιατί να θέλουμε να κάνουμε τη ζωή μας δύσκολη, αν υπάρχει πιο εύκολος τρόπος;»

«Για μένα δεν είναι εύκολος αυτός ο τρόπος. Εσείς που έχετε μεγαλώσει με ένα κινητό στο χέρι προφανώς το βρίσκεται εύκολο. Εμείς είμαστε σε μία ηλικία που δύσκολα μαθαίνουμε καινούργια πράγματα. Για αυτό και βλέπεις τους περισσότερους 50χρονους και 60χρονους να μην έχουν καν κινητό αφής, που εσείς το θεωρείτε δεδομένο. Είναι γιατί δεν μπορούν να προσαρμοστούν με τις εξελίξεις όσο γρήγορα όσο εσείς οι νέοι.»

«Εγώ έχω κινητό αφής»

«Ποιος σου μίλησε εσένα Ριντ. Μη προσπαθείς να με κάνεις να φανώ ότι λέω κάτι που δεν ισχύει.»

«Ναι αλλά εγώ έχω κινητό αφής. Είμαι μικρότερος από 59 δηλαδή;»

«Όχι δεν εννοώ αυτό ρε αδελφέ. Είπα πως οι “περισσότεροι” 50χρονοι δεν έχουν κινητό αφής. Πάλι όμως εκεί εσύ, να περάσει το δικό σου.»

«Εγώ τι σου φταίω που μου έδωσαν στη δουλειά μου κινητό αφής δωρεάν. Δεν ξέρω καν να το χρησιμοποιώ, απλά λέω ότι το έχω και χρησιμοποιώ το παλιό μου κινητό με τα κουμπάκια.»

«Άρα αποδεικνύεις το επιχείρημα μου.»

«Πάνω κάτω...»

«Όχι πάνω κάτω. Με φουντώνεις εσύ τώρα. Παιδιά μήπως θέλει κανείς να αλλάξουμε θέση μόλις σταματήσουμε;»

«Εγώ όχι θείε. Θέλω να είμαι στο αριστερό μέρος του αμαξιού. Ρομπ εσύ;»

«Μπα, καλά είμαι εδώ στο κέντρο. Μου αρέσει να βλέπω τον δρόμο μπροστά. Με ηρεμεί έτσι όπως παίρνει τις στροφές ο θείος ο Ριντ»

«Εσύ μικρέ πίσω μου θες αλλαγή;»

«Συγγνώμη τι λέμε; Δεν άκουγα... διάβαζα...»

«Δεν πειράζει, άστο. Συνέχισε να διαβάζεις.»

\chapter{}

Όταν μπήκαμε μέσα στα κεντρικά της Ειδικής Υπηρεσίας, φάνηκε λες και το κτίριο ξαφνικά μεγάλωσε. Απ’ έξω φαινόταν σαν μία απλή επταόροφη πολυκατοικία. Το μόνο που την ξεχώριζε από τις υπόλοιπες πολυκατοικίες ήταν η γυάλινη περιστρεφόμενη πόρτα, αντί για κανονική ξύλινη πόρτα, που περιβαλλόταν από μία τζαμαρία. Μπαίνοντας μέσα όμως, οι τοίχοι απομακρύνθηκαν, το ταβάνι ψήλωσε και η γραμματέας, που απ’ έξω από το κτίριο φαινόταν λες και είναι δίπλα στην πόρτα, τώρα ήταν αρκετά μέτρα μακριά μας. Έστρεψα το βλέμμα μου στον Βασι και τον Κουπ και είχαν και οι δύο παρόμοια αντίδραση με εμένα.

«Ελάτε καλέ, μην κάθεστε μπροστά στην πόρτα» μας φώναξε η γραμματέας. 

Προχωρήσαμε λοιπόν προς την γραμματέα, παραμένοντας έκπληκτοι με την κατάσταση που βρισκόμασταν. Απεριόριστοι πράκτορες, που δεν μπορούσες να τους διακρίνεις απ’ έξω από τη τζαμαρία, πήγαιναν πέρα δώθε στο κτίριο, μερικοί φαινόντουσαν σαν να βιάζονται, άλλοι συζητούσαν μεταξύ τους περπατώντας, άλλοι κάθονταν μόνοι τους και έπιναν καφέ, όλοι τους όμως φορούσαν παρόμοια ενδυμασία: ένα καφέ σακάκι, μερικοί το είχαν κουμπωμένο, άλλοι ανοιχτό, μία ζώνη με ένα πιστόλι στη δεξιά μεριά, και ένα παντελόνι που σε άλλοι το ταίριαζαν με το σακάκι, άλλοι όχι και τόσο. Μάλλον το σακάκι και η ζώνη ήταν τα σημαντικά μέρη της ενδυμασίας και από μέσα οι πράκτορες φορούσαν διαφορετικά ρούχα, σε περίπτωση ίσως που χρειαστεί να είναι πιο ευέλικτοι και το σακάκι δεν τους βοηθάει σε αυτό. Είδα εξάλλου μερικούς πράκτορες που έχουν ανοιχτό το σακάκι τους να φοράνε άσχετα ρούχα κάτω από αυτό. Ένας φορούσε την ίδια ποδοσφαιρική μπλούζα με τον Κουπ και όταν περάσαμε από δίπλα του μας χαμογέλασε.

Φτάνοντας στην γραμματέα, η ώρα είχε πάει 4:10. Μπροστά μας ήταν ένας πράκτορας που μας ήταν γυρισμένος πλάτη και έγραφε κάτι στο κινητό του. Ο Κουπ χάζευε ακόμη τους τοίχους και το ταβάνι. Ο Βασι που είχε ξεπεράσει το σοκ πήρε πρωτοβουλία και του μίλησε.

«Συγγνώμη, είστε στην ουρά;»

«Α, όχι. Κάτι σημείωνα, περάστε.» είπε και πήγε λίγα μέτρα δεξιά, αφήνοντας μας χώρο για να μιλήσουμε με τη γραμματέα.

«Μη φεύγεις. Θα σε χρειαστώ σε λίγο» είπε η γραμματέας και έβγαλε μία μικροσκοπική καρέκλα από την τσέπη της, την πέταξε στον πράκτορα και αυτός την ακούμπησε κάτω. Η καρέκλα τότε άλλαξε χρώμα, από καφέ σε μωβ, και το μέγεθος της πολλαπλασιάστηκε. Από ενός περίπου κυβικού εκατοστού που ήταν πριν, τώρα ήταν μια φυσιολογική καρέκλα. Η καρέκλα άλλαξε άλλη μία φορά, από μωβ τώρα πίσω στο καφέ, και ο πράκτορας κάθισε στην καρέκλα, χωρίς να έχει κάποια αντίδραση και συνέχισε να γράφει στο κινητό του. Οι τρεις μας είχαμε μείνει άφωνοι.

«Λοιπόν, σας παρακαλώ πείτε μου τα ονόματα σας και σε τι μπορώ να σας βοηθήσω» είπε η γραμματέας και πάτησε κάποια κουμπιά στον υπολογιστή της.

«Εσείς το κάνετε αυτό;» ρώτησε ο Βασι με έναν τόνο έκπληξης που δεν ταίριαζε στον χαρακτήρα του.

«Ποιο;» ρώτησε η γραμματέας και συνέχισε να πατάει πράγματα στον υπολογιστή της, λες και δεν ήταν προφανές για τι πράγμα μιλούσε ο Βασι.

Κοιταχτήκαμε μεταξύ μας, γιατί νομίζαμε πως κάναμε εμείς κάποιο λάθος. Είδα τον Βασι να κάνει κάποιες κινήσεις με τα χέρια του, σαν να ήθελε να υπολογίσει μαθηματικά πως μπορεί να συμβαίνει κάτι τέτοιο, αλλά δεν μπορούσε να καταλάβει.

«Α, μάλλον εννοείς τη μαγεία μου. Ναι εγώ το κάνω» μας απάντησε αφότου μας είδε να προβληματιζόμαστε η γραμματέας. «Συγγνώμη, δεν κατάλαβα. Λοιπόν, ονόματα.»

«Δεν θα μας εξηγήσετε πως το κάνετε;» ρώτησε ο Κουπ.

«Αν σας το εξηγήσω, θα χαλάσει η μαγεία του όλου πράγματος. Εξάλλου δεν είστε πράκτορες. Δεν έχω υποχρέωση να σας εξηγώ πως δουλεύουν τα πράγματα εντός της Ειδικής Υπηρεσίας σε ξένους. Πείτε μου τα ονόματα σας τώρα ή τουλάχιστον πείτε μου πώς μπορώ να σας βοηθήσω.»

«Είμαστε εδώ για αξιολόγηση.» είπε ο Βασι, ακουμπόντας το αριστερό του χέρι στο γραφείο.

«Α. Πάρα πολύ ωραία. Άρα εσείς πρέπει να είστε οι κύριοι Κουπ, Σπαρκ και Βασι. Νωρίς ήρθατε, προς τι η βιασύνη;» ρώτησε κοιτώντας τον Βασι ο οποίος ήταν ο πιο κοντινός της.

«Μην κοιτάτε εμένα. Ο ψηλός θα σας δώσει εξηγήσεις.» είπε ο Βασι και έδειξε τον Κουπ.

«Απλά δεν ήθελα να αργήσω, αυτό είναι όλο» είπε ο Κουπ και χαμογέλασε στον Βασι και τη γραμματέα.

«Ούτε εμείς καταλαβαίνουμε τη βιασύνη του, απλά είπαμε να τον συνοδεύσουμε, αφού έχουμε αξιολόγηση ακριβώς μετά από αυτόν» είπε ο Βασι στην γραμματέα, θέλοντας να δικαιολογήσει την ώρα που φτάσαμε.

«Δεν χρειάζεται να μου απολογείσαι. Αν αργούσατε θα υπήρχε πρόβλημα. Να δούμε τώρα...» είπε η γραμματέας και κοίταξε τον υπολογιστή της. «...ωραία. Είστε έτοιμοι. Πάρτε αυτά τα χαρτιά, ο καθένας το δικό του, μη τα ανακατέψετε» συνέχισε η γραμματέας και μας έδωσε μερικά χαρτιά με τα ονόματα μας στην κορυφή και από κάτω ερωτήσεις λες και ήταν διαγώνισμα. «Μην ασχοληθείτε με αυτά τα χαρτιά, είναι για τον αξιολογητή σας. Απλά του τα δίνεται όταν μπείτε στο δωμάτιο. Ο συνάδελφος μου από εδώ θα σας πάει στον χώρο αναμονής» είπε κάνοντας νεύμα προς τον πράκτορα που καθόταν στην καρέκλα που είχε αλλάξει μαγικά μέγεθος.

«Και τους τρεις;» ρώτησε αυτός σταματώντας να γράφει και γυρίζοντας εκνευρισμένος το βλέμμα του προς την γραμματέα.

«Γιατί; Σε κούρασε η κοπελίτσα που ήρθε πριν και δεν μπορείς να οδηγήσεις άλλους καλεσμένους στον χώρο αναμονής;» είπε η γραμματέας με ειρωνικό τόνο.

«Όχι βέβαια. Απλά... ξαφνιάστηκα. Ξεκινήστε μία να περπατάτε.» είπε ο πράκτορας, αλλά παρέμεινε στην καρέκλα του.

«Προς τα που;» ρώτησε ο Βασι. «Δεν μας βοηθάτε και πολύ.»

«Ως προς κάποια κατεύθυνση, τώρα αυτό σε μάρανε;» του απάντησε ο πράκτορας.

Ο Βασι, λίγο εκνευρισμένος, άρχισε να τρέχει προς τον πράκτορας, λες και ήθελε να ξεκινήσει κάποια μάχη, αλλά ένα βέλος εμφανίστηκε στα πόδια του. Ο πράκτορας που καθόταν έκανε μια κίνηση προς τα αριστερά με τον δείκτη του δεξιού του χεριού και ο Βασι εκσφενδονίστηκε στον διάδρομο που βρισκόταν στα αριστερά του πράκτορα που καθόταν. Ο Κουπ και εγώ κοιταχτήκαμε χωρίς να βγάλουμε άχνα, καθώς είχαμε συγκλονιστεί με την επίδειξη μαγείας που μας είχαν κάνει η γραμματέας και ο πράκτορας. Μετά από λίγο ακούστηκε ένας ήχος λες και ο Βασι χτύπησε σε κάποιο τραπέζι και έσπασε κάποιο βάζο.

«Εσείς θα περπατήσετε ή θα καθήσετε εκεί;» μας ρώτησε ο πράκτορας.

Αρχίσαμε λοιπόν να περπατάμε προς τον πράκτορα και εμφανίστηκαν δύο βέλη, ένα ανάμεσα στα πόδια του καθενός, και αρχίσαμε να αιωρούμαστε, αλλά η ταχύτητα με την οποία μετακινούμασταν έμεινε ίδια με την ταχύτητα με την οποία περπατούσαμε. Κοίταξα τον πράκτορα που καθόταν και έκανε την ίδια κίνηση με τον δείκτη του αριστερού του χεριού και τα βέλη άλλαξαν κατεύθυνση και αρχίσαμε να μετακινούμαστε προς τον ίδιο διάδρομο που εκσφενδονίστηκε ο Βασι, αλλά με την ίδια ταχύτητα που περπατούσαμε.

«Μην ανησυχείτε για τον φίλο σας. Το πολύ πολύ να έσπασε κανένα βάζο, αλλά έχουμε ανθρώπους με μαγεία Επισκευής, οπότε δεν χρειάζεται να πληρώσει τίποτα» είπε ο πράκτορας καθώς σηκωνόταν από την καρέκλα. Εμείς συνεχίσαμε να μετακινούμαστε με σταθερή ταχύτητα προς τον διάδρομο. «Θα τους συνοδέψω μέχρι ένα σημείο και μετά θα πάω στην ομάδα μου. Καλή συνέχεια.» είπε ο πράκτορας στην γραμματέα και γύρισε προς το μέρος μας. 

«Μπορείς να μου πετάξεις μία πίσω την καρέκλα;» του είπε η γραμματέας δείχνοντας την καρέκλα του που είχε συρρικνωθεί πάλι στο αρχικό της μέγεθος. Ο πράκτορας έσκυψε, έπιασε την καρέκλα και την πέταξε προς την γραμματέα. Από την καρέκλα εμφανίστηκε ένα βέλος και η τροχιά την οποία ακολουθούσε η μικροσκοπική καρέκλα άλλαζε με βάση τις κινήσεις του χεριού του πράκτορα. Η καρέκλα έφτασε τελικά στην τσέπη της γραμματέα, από την οποία είχε βγει, και ο πράκτορας της χαμογέλασε και την χερέτησε.

Ο πράκτορας που είχε μείνει λίγο πίσω λόγω της παράκλησης που του είχε κάνει η γραμματέας, έτρεξε λίγο για να έρθει δίπλα μας. Δεν μας μίλησε, απλά συνέχισε να περπατάει δίπλα μας ενώ εμείς αιωρούμασταν. Στους τοίχους του διαδρόμου, ενώ προχωρούσαμε, παρατηρούσα τις κορνίζες που ήταν κρεμασμένες. Όλες είχαν έναν πράκτορα και στην κάτω δεξιά γωνία είχαν ένα όνομα και μία ειδικότητα. Τι είδους πράκτορες είναι αυτοί; Γιατί έχουν κορνιζες κρεμασμένες; Αν γίνω και γω κορυφαίος πράκτορας θα κερδίσω και εγώ κορνίζα; Παρατηρώντας λίγο παραπάνω όμως άρχισα να βλέπω και νεότερα πρόσωπα, όχι μόνο 40χρονες και 50χρονους. Άρα το κριτήριο για να πάρεις κορνίζα μάλλον δεν είναι η σπουδαιότητα; Μπορεί από την άλλη να είναι, αλλά αυτοί οι νέοι να είναι τόσο πετυχημένοι που να έχουν κερδίσει κορνίζα στα κεντρικά της Ειδικής Υπηρεσίας.

Όσο προχωρούσαμε όμως, όλο και περισσότερες κορνίζες με νέα πρόσωπα υπήρχαν. Πόσοι πετυχημένοι νέοι υπάρχουν στην Ειδική Υπηρεσία; Μήπως όταν γίνεσαι πράκτορας απλά κρεμάνε μία κορνίζα;

«Να σε ρωτήσω κάτι;» είπα στον πράκτορα που μας συνόδευε.

«Τι θες;» μου είπε χωρίς να γυρίσει να με κοιτάξει.

«Εσύ έχεις κορνίζα;»

Δεν μου απάντησε.

«Πώς παίρνω κορνίζα;» τον ξαναρωτησα.

Δεν μου απάντησε.

«Δεν ξέρεις, ε;»

«Δεν είναι ότι δεν ξέρω, αλλά θα μάθεις σύντομα. Αν δεν πας στην πρώτη σου αποστολή δεν μπορείς να μάθεις εκ πείρας.» μου απάντησε υποτονικά χωρίς πάλι να με κοιτάξει.

«Ατύχησες. Εγώ θέλω να γίνω ανακριτής, δεν θα πηγαίνω σε αποστολές. Μπορείς να μου πεις τώρα πως παίρνω κορνίζα;» τον ρώτησα με έναν τόνο που ταίριαζε περισσότερο στην συμπεριφορά του Βασι, παρά στη δική μου

«Ποιος σου είπε πως οι ανακριτές δεν πάνε σε αποστολές;» με ρώτησε ο πράκτορας

«Η μητέρα μου;» του απάντησα

«Τι σχέση έχει η μητέρα σου με την Ειδική Υπηρεσία;» με ρώτησε ο πράκτορας εκνευρισμένος.

«Είναι... κατάσκοπος...» του απάντησα λίγο φοβισμένος, σκύβοντας το κεφάλι μου.

Ο πράκτορας γύρισε αργά και με κοίταξε για λίγο.

«Ήξερα ότι κάτι μου θύμιζε κάτι η φάτσα σου» μου είπε ο πράκτορας και μου χαμογέλασε. «Αν περάσεις την αξιολόγηση σου, μπορεί να είμαστε και στην ίδια ομάδα, αφού είπες θες να γίνεις ανακριτής.»

«Γιατί το λες αυτό; Και πώς και οι ανακριτές πηγαίνουν σε αποστολές; Νόμιζα πως η δουλειά του ανακριτή είναι απλά να... “ανακρίνει” υπόπτους. Τι παραπάνω κάνει ένας απλός ανακριτής; Ειδικά κάποιος “αδύναμος” ανακριτής.»

«”Αδύναμος”; Δεν σε καταλαβαίνω.» μου απάντησε ο πράκτορας.

«Ξέρεις... χωρίς... πως να στο πω τώρα... “γενετικά πλεονεκτήματα“» του είπα, προσπαθώντας να του εξηγήσω τι εννοούσα.

«Μάλλον εννοείς πράκτορες κλάσης Β. Ναι, σίγουρα αυτό εννοείς.» μου απάντησε ο πράκτορας και γύρισε το κεφάλι του πάλι προς τα μπροστά.

«Κλάσης Β; Υπάρχουν και κλάσης Α; Τι ονομασία είναι αυτή;» τον ρώτησα από περιέργεια.

«Στην κλάση Β εντάσσονται όλοι οι άνθρωποι, στην πράκτορες ή μη, οι οποίοι είναι “αδύναμοι”, όπως λες εσύ, με την μαγεία τους. Στην κλάση Α αντιθέτως εντάσσονται οι... “δυνατοί”; Έτσι θα τους έλεγες; Καταλαβαίνεις πάντως τι εννοώ. Εγώ, για παράδειγμα, είμαι κλάσης Α. Εκπλήσσομαι πάντως που δεν ξέρεις αυτή τη βασική ορολογία. Καλά θα μου πεις στο σχολείο δεν αναφέρεται τίποτα για την μαγεία, αλλά η μητέρα σου είναι κορυφαία κατάσκοπος. Δεν σου το έχει αναφέρει ποτέ αυτό;»

«Δεν... όχι δεν μου το έχει αναφέρει» του απάντησα απογοητευμένος.

«Δεν πειράζει, για αυτό είσαι εδώ, για να μάθεις.»

Όταν το είπε αυτό ο πράκτορας που μας συνόδευε, άκουσα στο βάθος ένα μουρμουρητό που έμοιαζε με τη φωνή του Βασι. Καθώς πλησιάζαμε, οι υποψίες μου επιβεβαιώθηκαν. Ο διάδρομος, μετά από αρκετή ώρα... αιώρησης, είχε μία πόρτα στα δεξιά μας, η οποία ήταν ανοικτή και από μέσα ακουγόταν η φωνή του Βασι. Μπαίνοντας μέσα στο δωμάτιο, είδαμε τον Βασι ξαπλωμένο στο πάτωμα μαζί με δύο κοπέλες. Η μία ήταν καστανομάλλα με κυματιστό μαλλί και ενδυμασία νοσοκόμας, η οποία καθόταν στα γόνατα δίπλα από τον Βασι κρατώντας τις παλάμες της λίγα εκατοστά πάνω του, από τις οποίες έλαμπε ένα είδος κίτρινου φωτός. Η άλλη κοπέλα ήταν λίγο πιο δίπλα καθισμένη σε μία καρέκλα δίπλα σε μία δεύτερη πόρτα και σχολίαζε σιγανόφωνα αυτά που μουρμούριζε ο Βασι. Ήταν η κοπέλα με την κοτσίδα που είχε μιλήσει στον Βασι έξω από το καφενείο. Τι σύμπτωση και αυτή. Πρώτα ο Κουπ και μετά αυτή. Σε λίγο μπορεί να πετύχουμε και τη μητέρα μου στην Ειδική Υπηρεσία. Μπα, αυτό είναι αδύνατο, αφού έχει φύγει πάλι για κάποιες μέρες για μία αποστολή, οπότε δεν θα είναι στην πόλη μάλλον.

«Λοιπόν, εγώ σας αφήνω εδώ.» είπε ο πράκτορας που μας συνόδευε και μας άφησε να περπατήσουμε μόνοι μας.

«Όχι δεν θα πας πουθενά» είπε ο Βασι και σηκώθηκε από το πάτωμα στο οποίο είχε ξαπλώσει. «Έχουμε μία μπίσνα να τελειώσουμε εμείς οι δύο. Νομίζεις πως μπορείς με τραυματίσεις πριν την αξιολόγηση μου και να μην υπάρχουν συνέπειες; Για έλα δω...» είπε ο Βασι και έτρεξε πάλι προς το μέρος του πράκτορα με επιθετικές προθέσεις.

Ο πράκτορας, απογοητευμένος από την συμπεριφορά του Βασι, απλά έδειξε με το δάχτυλο του το πάτωμα και ο Βασι φάνηκε να γλιστράει και να πέφτει με τα μούτρα στο πάτωμα.

«Για αυτό έχουμε ανθρώπους σαν την Ιφι εδώ. Αν τραυματίζεσαι, απλά πες το στις νοσοκόμες και θα σε θεραπεύσουν με τις μαγείες τους.» είπε ο πράκτορας και έφυγε από το δωμάτιο.

«Καθίστε λίγο ακόμη κύριε να σας θεραπεύσω» είπε η νοσοκόμα και γύρισε τον Βασι ανάσκελα. «Ηρεμήστε για λίγο και θα είστε μία χαρά στο άψε σβήσε»

«Αυτό είπες και την προηγούμενη φορά, αλλά περίμενα κανένα πεντάλεπτο στο πάτωμα και δεν είχες τελειώσει» είπε νευριασμένος ο Βασι.

«...σιγά μην περίμενες και δεκάλεπτο...» σχολίασε η κοπέλα που καθόταν στην γωνία.

«Λοιπόν...» ακούστηκε μία ανδρική φωνή από μέσα από το δίπλα δωμάτιο και άνοιξε η δεύτερη πόρτα του δωματίου στην οποία καθόταν δίπλα η κοπέλα που σχολίαζε. «Α, εσύ πρέπει να είσαι το ραντεβού για τις 4:30.» την ρώτησε ένας άντρας με περιποιημένο μουστάκι που βγήκε μέσα από εκείνο το δωμάτιο. Γύρισε και κοίταξε προς το μέρος μας. «Και εσείς... πρέπει να είστε τα υπόλοιπα ραντεβού; Για να δω λίγο...» κοίταξε το ρολόι που βρισκόταν στον τοίχο. «4:19. Προς τι η βιασύνη ρε παιδιά; Γιατί είστε όλοι τόσο νωρίς;»

«Θέλετε να φύγουμε και να ξανάρθουμε;» είπε ειρωνικά ο Βασι ξαπλωμένος στο πάτωμα.

«Όχι, ίσα ίσα μου κάνετε τη δουλειά πιο εύκολη. Θα τελειώσουμε πιο γρήγορα έτσι. Παρακαλώ, κυρία μου, περάστε, αν είστε έτοιμη, για να ξεκινήσουμε την αξιολόγησή σας.» είπε ο άντρας με το μουστάκι και η κοπέλα που καθόταν στη γωνία σηκώθηκε και πήγε στο άλλο δωμάτιο. «Οι υπόλοιποι... καθίστε, θα έρθει και σας η ώρα σας» είπε χιουμοριστικά ο άντρας και έκλεισε την πόρτα.

\chapter{}

«Επόμενος!» ακούστηκε η φωνή του άντρα με το περιποιημένο μουστάκι από το δίπλα δωμάτιο, ανοίγοντας ίσα ίσα την πόρτα.

«Τόσο γρήγορα;» απάντησε ο Βασι που ήταν ακόμη ξαπλωμένος στο πάτωμα.

«Ναι, τόσο γρήγορα. Ελάτε τώρα γρήγορα, γιατί έχουμε και άλλες δουλείες» του απάντησε ο άντρας από μέσα από το δωμάτιο.

«Κουπ θα πας εσύ; Αφού, έτσι κι αλλιώς, εσύ είχες ραντεβού πριν από εμάς» πρότεινα στον Κουπ.

Ο Κουπ δεν μου απάντησε ευθέως, απλά σηκώθηκε, μου χαμογέλασε και μπήκε μέσα στο δωμάτιο. Η κοπέλα με την κοτσίδα που μπήκε πριν από εμάς δεν βγήκε ποτέ από το δωμάτιο. Λες να έχουν και άλλα δωμάτια που συνδέονται αλυσιδωτά όπως ο χώρος αναμονής που βρισκόμασταν και το δίπλα δωμάτιο; Ίσως υπάρχει κάποιου είδους μαγεία τηλεμεταφοράς και ο άντρας με το μουστάκι την μετέφερε σε κάποιο άλλο δωμάτιο στην Ειδική Υπηρεσία; Πολύ περίεργο πάντως.

Επίσης, γιατί ο διάδρομος που ακολουθήσαμε για να έρθουμε ως εδώ ήταν τόσο μακρύς; Όταν στρίψαμε για να μπούμε στον χώρο αναμονής δεν είχαμε φτάσει καν στο τέλος του διαδρόμου. Πόσο μακριά λες να εκτείνεται αυτός ο διάδρομος; Και οι κορνίζες; Προφανώς, για να γράφουν και ειδικότητες δίπλα από τα ονόματα του καθενός, αυτό σημαίνει ότι όλα τα εικονιζόμενα πρόσωπα είναι πράκτορες. Πώς μπορώ να πάρω και εγώ κορνίζα άραγε; Γιατί με ενδιαφέρει καν τόσο πολύ να “κερδίσω” κορνίζα; Μάλλον απλά μου φάνηκε εντυπωσιακό σε συνδυασμό με την ποσότητα τους και του μήκους του διαδρόμου.

«Είστε έτοιμος κύριε» είπε η νοσοκόμα που ήταν πάνω από τον Βασι.

«Επιτέλους» είπε ο Βασι και σηκώθηκε. «Ήθελα να ‘ξερα τι κάνεις τόση ώρα. Είμαι ξαπλωμένος στο πάτωμα και συ απλά κάθεσαι από πάνω μου και μου λες να μην κουνιέμαι. Πραγματικά...»

«Επόμενος!» ακούστηκε η φωνή του άντρα με το μουστάκι καθώς άνοιξε η πόρτα του δίπλα δωματίου. Ο Κουπ βγήκε από το δωμάτιο με ένα χαμόγελο στο πρόσωπό του και, χωρίς να μας πει τίποτα, προχώρησε στον διάδρομο και έσπευσε προς την είσοδο του κτιρίου.

«Σπαρκ πήγαινε. Άσε με λίγο να συνέλθω και θα μπω μετά από σένα.» μου είπε ο Βασι και κάθισε στη καρέκλα δίπλα μου.

Σηκώθηκα λοιπόν και άλλαξα δωμάτιο. Μέσα στο δωμάτιο, αν μπορεί κανείς να αποτελέσει αυτόν τον χώρο δωμάτιο, ήταν ένα γραφείο με τρεις καρέκλες, μία πίσω και δύο μπροστά από το γραφείο. Στο γραφείο αυτό καθόταν ο άντρας με το μουστάκι που με φώναξε, αλλά και η κοπέλα με την κοτσίδα που είχε μπει πριν από εμάς. Αυτό που μου έκανε εντύπωση όμως δεν ήταν το γραφείο, αλλά ο χώρος που ήταν αντικριστά από το γραφείο. Δέντρα. Πάρα πολλά δέντρα. Μία μικρή λίμνη λίγο πιο μακριά. Φωτιές ακόμη πιο μακριά. Στο βάθος κάτι μικρά κτίρια σαν μία πόλη που είχε μόνο τρεις δρόμους. Τα κτίρια αυτά φαίνονταν να βρίσκονται τόσο μακριά όσο το μέγεθος του διαδρόμου. Λες για αυτό να είχαν τόσο μακρύ διάδρομο; Για να μπορέσουν να έχουν ένα τέτοιο “δωμάτιο”; Το ταβάνι ήταν ξαφνικά υπερβολικά ψιλό και οι τοίχοι ήταν σαν να είχαν μεγαλώσει πάλι. Το μόνο που με χώριζε από τα δέντρα, την λίμνη, τις φωτιές και τα κτίρια ήταν ένα τεράστιο τζάμι με μία μικρή γυάλινη πόρτα κοντά στην πόρτα από την οποία μπήκα.

«Καθίστε κύριε Σπαρκ Θοτς» μου είπε ο άντρας με το μουστάκι, και αυτό έκανα. «Λέγομαι Νικ και είμαι ο αρχηγός της Ειδικής Υπηρεσίας» συνέχισε και μου έδωσε το χέρι του.

«Χάρηκα για την γνωριμία» του απάντησα και του έδωσα το χέρι μου.

\textasciitilde Δεν αισθάνομαι τόσο δυνατή μαγεία από σένα, αλλά είναι οριακά κλάσης Β. Για να δούμε όμως αν η αίτηση του έλεγε την αλήθεια...\textasciitilde «Λοιπόν καθίστε κύριε και δώστε μου το χαρτί σας» είπε ο Νικ και καθώς μου άφηνε το χέρι.

«Ορίστε» του απάντησα και του έδωσα το χαρτί που μας είχε δώσει η γραμματέας.

«Πολύ ωραία, λοιπόν...» είπε ο Νικ και άρχισε να γράφει στο χαρτί κάποια πράγματα.

Όσο σημείωνε ο Νικ, έστρεψα το βλέμμα προς την κοπέλα με την κοτσίδα που καθόταν πλέον απέναντι μου. Μουρμούριζε στον εαυτό της κάποιες λέξεις που δεν μπορούσα να καταλάβω, λες και ήταν τελείως άλλη γλώσσα, και έπαιζε με τα δάχτυλα της. Ταυτόχρονα κούναγε τα πόδια της πάνω κάτω λες και ήταν αγχωμένη, αλλά η έκφραση της και ο τόνος της φωνής της ήταν περισσότερο εκνευρισμένος παρά νευριασμένος. 

«Τι κοιτάς;» μου είπε η κοπέλα χωρίς να με κοιτάξει καν.

«Τί... ποτα. Απλά... με παραξένεψε που δεν βγήκες  πριν από αυτό το δωμάτιο, ενώ ο φίλος μας βγήκε με το που τελείωσε» της απάντησα, προσπαθώντας να εξηγήσω τον εαυτό μου.

«Α» απάντησε και συνέχισε να μουρμουράει περίεργες λέξεις παίζοντας με τα δάχτυλα της.

«Ωραία κύριε Σπαρκ, είμαστε σχεδόν έτοιμοι. Θέλω μόνο να μου πείτε τι σκέφτομαι αυτή τη στιγμή» μου είπε ο Νικ κοιτώντας με επίμονα.

«Α θέλετε να σας κάνω επίδειξη της ικανότητας μου; Λυπάμαι, αλλά πρέπει να έχω σωματική επαφή με τον στόχο μου για να...»

«Γιατί; Εσύ δεν το ελέγχεις;» με έκοψε ο Νικ.

«Ναι... αλλά για να περάσουν οι κεραυνοί από το σώμα μου στο σώμα του στόχου μου και το αντίθετο, πρέπει να έχω σωματική επαφή. Αυτό ήθελα να σας...»

«Κεραυνοί;» προβληματίστηκε ο άντρας.

«Ναι. Έχω μαγεία Κεραυνών. Δεν το έγραψε η αίτηση μου;» απάντησα αγχωμένος μήπως είχα ξεχάσει να γράψω τι μαγεία είχα στην αίτηση μου.

«Ναι... δικό μου λάθος. Δεν τις διαβάζω πολύ συγκεντρωμένος τις αιτήσεις. Απλά βλέπω ποιοι έχουν ωραίες ικανότητες και τους καλώ για αξιολόγηση» μου απάντησε ο Νικ αφότου κοίταξε την αίτηση μου στον υπολογιστή του. «Και πάλι όμως. Δεν θα έπρεπε να μπορείς να ελέγξεις τους κεραυνούς σου έτσι ώστε να πάνε από τα δικά μου νεύρα στα δικά σου; Τουλάχιστον έτσι το περιγράφεις στην αίτηση σου αν καταλαβαίνω καλά.»

«Να σας πως την αλήθεια δεν το έχω δοκιμάσει και ποτέ. Είχα στο μυαλό μου πως η σωματική επαφή ήταν απαραίτητη για να...»

«Δε με ενδιαφέρει τι νόμιζες και τι όχι. Έλα διάβασε τις σκέψεις μου για να τελειώνουμε.» με έκοψε πάλι ο Νικ και με ξανακοίταξε επίμονα.

Δεν ήξερα τι να κάνω. Είχα πάρα πολύ καιρό να προσπαθήσω να ελέγξω τους κεραυνούς μου συνειδητά. Συνήθως το διάβασμα σκέψεων μου έβγαινε σχεδόν αυτόματα με τη σωματική επαφή, για αυτό εξάλλου το ανακάλυψε και η μητέρα μου. 

Τέντωσα το χέρι μου προς το κεφάλι του Νικ και προσπάθησα να βγάλω μερικούς κεραυνούς, αλλά μάταια. Προσπάθησα άλλη μία και αυτή τη φορά μπόρεσα να βγάλω κάποιους κεραυνούς από τα δάχτυλα μου, οι οποίοι έφτασαν στον Νικ και γύρισαν πίσω στο χέρι μου. \textasciitilde ...ζίνο, καζίνο, καζίνο, καζίνο...\textasciitilde\  ήταν αυτό που σκεφτόταν εκείνη τη στιγμή.

«Καζίνο;» του είπα χαρούμενος τραβόντας γρήγορα το χέρι μου πίσω.

«Ωραία, τώρα κάντο χωρίς το χέρι σου. Αν πας σε κάποια αποστολή θα έχεις τεντωμένο το χέρι σου προς τον ύποπτο; Κινεί υποψίες, ελπίζω να το καταλαβαίνεις.» μου είπε ο Νικ.

«Δεν... είμαι σίγουρος αν μπορώ να το κάνω...» του απάντησα σκύβοντας το κεφάλι μου απογοητευμένος.

«Έχεις προσπαθήσει;»

«Ό... χι...»

«Έ τότε γιατί τα παρατάς;» μου είπε ο Νικ και μου χαμογέλασε. «Έλα τι περιμένεις, δοκίμασε»

Έδεσα τα χέρια μου πίσω από την πλάτη μου, για να δείξω στον Νικ ότι δεν θα προσπαθούσα να τα χρησιμοποιήσω. Έκλεισα τα μάτια μου και συγκεντρωθηκα στην εργασία που μου είχε δοθεί, να διαβάσω τις σκέψεις του χωρίς να χρησιμοποιήσω τα χέρια μου. Προσπάθησα να βγάλω κεραυνούς από κάποιο σημείο του σώματος μου εκτός από τα άκρα μου, αλλά δεν μπορούσα. Προσπάθησα να συγκεντρωθώ παραπάνω, αλλά μάταια. Δεν μπορούσα να κάνω τίποτα. Αυτό ήταν, είχα αποτύχει. Άνοιξα τα μάτια μου και έλυσα τα χέρια μου. Προσπάθησα μία τελευταία φορά, χωρίς να μπορέσω να βγάλω κάποιον κεραυνό από το σώμα μου, αλλά άκουσα \textasciitilde ...άσος, άσος, άσος....\textasciitilde 

«Άσος;»

«Πάρα πολύ ωραία. Άλλη μία φορά»

Πώς το έκανα; Διάβασα όντως τις σκέψεις του χωρίς να βγάλω κεραυνούς από το σώμα μου; Κοίταξα τον Νικ στα μάτια και προσπάθησα ξανά.

\textasciitilde ...μπαστούνι, πέντε μπαστούνι, πέντε μπα....\textasciitilde 

«Πέντε μπαστούνι;»

«Είδες; Τελικά μπορείς. Να ξέρεις, μπορεί να μη βλέπεις τους κεραυνούς σου, αλλά προφανώς όποτε συγκεντρώνεσαι και προσπαθείς να χρησιμοποιήσεις τη μαγεία σου, πάντα κάτι βγαίνει από το σώμα σου, ορατό ή μη. Οι κεραυνοί σου πιθανότατα αυτές τις δύο φορές μπορεί να μην ήταν ορατοί, αλλά αφού διάβασες τις σκέψεις μου σίγουρα υπήρχαν.» μου είπε ο Νικ και σημείωσε κάτι στο τέλος του χαρτιού, αφήνοντας την υπογραφή του. «Δεν νομίζω να χρειαστεί να χρησιμοποιήσουμε το υπόλοιπο δωμάτιο για την αξιολόγηση σου. Οπότε...»

«Τι είναι αυτό το “δωμάτιο”;» τον διέκοψα.

«Αίθουσα εξάσκησης» μου απάντησε ο Νικ. «Εδώ εξασκούνται οι συνάδελφοι μας που είναι ειδικοί στις μάχες.»

«Και το γυαλί είναι για να σας προστατεύει από τον κίνδυνο;» τον ρώτησα από περιέργεια.

«Όχι καλέ, δεν είναι δικό μου αυτό το γραφείο, η γραμματέας μου το δίνει όποτε έχω αξιολογήσεις»

«Α, νόμιζα πως ο χώρος εδώ ήταν δικός σας...»

«Πολλά πρέπει να νομίζεις που δεν ισχύουν. Ας πούμε, σου είπα πριν ότι με λένε Νικ. Στην πραγματικότητα με λένε Μαλχη»

«Μαλχη;»

«Ναι. Σιγά μην έδινα το πραγματικό μου όνομα σε κάποιον ξένο που υπάρχει πιθανότητα να μην τον ξαναδώ ποτέ. Οι πράκτορες της Ειδικής Υπηρεσίας πρέπει να έχουν ένα στοιχείο μυστηρίου Σπαρκ. Θα το καταλάβεις αυτό όσο συνεχίσεις να δουλεύεις εδώ.»

«Άρα πέρασα την αξιολόγηση;» τον ρώτησα χαρούμενος.

«Βεβαίως. Όσον αφορά την ειδικότητα σου τώρα...»

«Ανακριτής;»

«Σαν να διάβασες τις σκέψεις μου» μου είπε ο Μαλχη και γέλασε μόνος του με το αστείο του. «Δεν πιστεύω όμως να νομίζεις όμως πως οι ανακριτές στην Ειδική Υπηρεσία δεν πηγαίνουν σε αποστολές; Γιατί οι περισσότεροι ξένοι έχουν αυτή τη πεποίθηση.»

«Όχι. Γιατί να το πιστεύω αυτό;» του απάντησα, λέγοντας προφανώς ψέματα.

«Ωραία, λοιπόν. Πάρε το σακάκι και τη ζώνη με το όπλο σου και συνόδευσε παρακαλώ την Μπλαιρ στο τέλος του διαδρόμου. Εκεί θα βρείτε τρεις πόρτες. Μπείτε στην αριστερή και γνωρίστε την ομάδα σας και θα έρθω και εγώ σε λίγο να...»

«Κάτσε τι;» τον διέκοψε γιατί είχα μπερδευτεί.

«Τι έγινε;» με ρώτησε προβληματισμένος.

«Τίποτα... απλά μου φάνηκαν πολλά...» του απάντησα παίρνοντας το σακάκι και την ζώνη με το όπλο από τα χέρια του και τα φόρεσα.

«Τι πολλά; Δεν είπα τίποτα περίπλοκο. Απλά μόλις βγείτε πηγαίνεται δεξιά, περπατήστε όλο τον διάδρομο και πηγαίνετε στην ομάδα σας» μου είπε ο Μαλχη δείχνοντας και τις κατευθύνσεις με τα χέρια του.

«Ε... ντάξει»

«Πάμε» είπε η κοπέλα και βγήκε από το δωμάτιο.

«Επόμενος!» φώναξε ο Μαλχη. «Τρέχα να την προλάβεις, ώστε να πάτε μαζί στην ομάδα σας.»

«Εντάξει» του απάντησα και βιάστηκα να βγω από το δωμάτιο. Είδα τον Βασι να κάθετε ακόμη στη καρέκλα που καθόταν πριν μπω στην “αίθουσα εξάσκησης”, όπως την αποκάλεσε ο Μαλχη

«Σπαρκ. Πώς πήγε;» με ρώτησε ο Βασι.

«Επόμενος είπα!» φώναξε πάλι ο Μαλχη.

«Καλύτερα να πηγαίνεις, θα μιλήσουμε μετά.» του είπα και έτρεξα να προλάβω την κοπέλα. Είδα την έκφραση του προσώπου του Βασι να αλλάζει από χαρούμενη σε προβληματισμένη πριν συνεχίσω να τρέχω.

«Μη βρίσω το σπανάκι σου, είπα επόμενος!» άκουσα πάλι τον Μαλχη.

«Τώρα έρχομαι, μη φωνάζεις» άκουσα τον Βασι να απαντάει καθώς απομακρυνόμουν.

Όταν τελικά έφτασα την κοπέλα, είχα λαχανιάσει. Προσπάθησα να την σταματήσω για λίγο, αλλά αυτή συνέχισε να περπατάει με τον γρήγορο σχετικά ρυθμό της.

«Περίμενε... λίγο...» είπα βαριανασαίνοντας.

«Και ποιος είσαι εσύ να μου πεις να περιμένω» μου απάντησε χωρίς να σταματήσει. 

«Καλά...» είπα και έτρεξα λίγο ακόμη για να φτάσω δίπλα της.

Παρατήρησα ότι φορούσε και αυτή παρόμοιο σακάκι και ζώνη με εμένα. Είχε τα χέρια της στις τσέπες του σακακιού και μουρμουρούσε ακόμη περίεργες λέξεις, αυτή την φορά λίγο πιο μελωδικά.

«Μπλαιρ, ε;» την ρώτησα. Αυτή αναστέναξε.

«Ναι» μου απάντησε χωρίς να με κοιτάξει.

«Αν επιτρέπεται, τι είναι αυτές οι λέξεις που λες τόση ώρα;» την ρώτησα από περιέργεια.

«Σολφέζ» μου απάντησε μονολεκτικά.

«Τι σημαίνει αυτό;» την ξαναρωτησα. Αυτή σταμάτησε να μουρμουρίζει για λίγο και με κοίταξε.

«Νότες» είπε, χωρίς να εξηγήσει την απάντηση της και γύρισε το βλέμμα της πάλι ευθεία.

«Τι εννοείς;» την ρώτησα άλλη μία φορά, επειδή δυσκολευόμουν να καταλάβω τις μονολεκτικές τις απαντήσεις.

«Αν είναι να συνεχίζεις να με ρωτάς χαζές ερωτήσεις να μου το πεις τώρα, γιατί δεν θα τα πάμε καλά στην ίδια ομάδα» είπε εκνευρισμένη και με χτύπησε στον σβέρκο.

«Άουτς... δεν χρειαζόταν αυτό.»

«Εγώ έκρινα ότι χρειαζόταν. Έχεις κάποια άλλη ερώτηση;» με ρώτησε με το χέρι της οπλισμένο.

«Όχι...» της απάντησα και έσκυψα το κεφάλι.

«Πάμε γρήγορα τότε, γιατί με περιμένει κάποιος.»

\chapter{}

Μπαίνοντας στο δωμάτιο που μας είχε πει ο Μαλχη είδαμε τρεις άντρες, δύο νέους και έναν λίγο μεγαλύτερης ηλικίας θα έλεγε κανείς. Ο ένας νέος ήταν ο πράκτορας που είχε συνοδεύσει εμένα και τον Ρομπ μέχρι τον χώρο αναμονής με την μαγεία του. Μόλις με είδε γέλασε και συνέχισε να γράφει στο κινητό του. Η Μπλαιρ με το που μπήκε στο δωμάτιο έτρεξε προς το μέρος του μεγάλου άντρα και τον αγκάλιασε. Είχε σταματήσει να μουρμουρίζει πλέον και η έκφραση του προσώπου της ήταν πλέον χαρούμενη αντί για σοβαρή.

«Ήρθα Χερμπ. Επιτέλους θα είμαστε μαζί.» είπε στον μεγάλο ηλικιακά άντρα με έναν τόνο ανακούφισης.

«Ναι μην ανησυχείς Μπλαιρ. Δεν θα χρειαστεί να ανησυχήσεις πότε ξανά» της απάντησε ο άντρας, που μάλλον λεγόταν Χερμπ, και την αγκάλιασε και αυτός.

Κοίταξα λίγο πιο προσεκτικά το δωμάτιο. Ήταν ένα απλό δωμάτιο γραφείου, με τριθέσιους καναπέδες στον δεξί και τον αριστερό τοίχο και μία βιβλιοθήκη που μάλλον περιείχε φακέλους με πληροφορίες πίσω από το γραφείο. Πάνω στο γραφείο υπήρχε ένα ταμπελάκι με το όνομα μου. Το πλησίασα και το σήκωσα για να το δω καλύτερα.

«”Σπαρκ Θοτς”» μου είπε ο άντρας που αγκάλιαζε την Μπλαιρ. «Συγχαρητήρια στην νέα σας θέση ως ανακριτής. Το όνομα μου είναι Χερμπ. Χάρηκα για την γνωριμία.»

«Και εγώ χάρηκα, αλλά πως... γιατί είχατε το όνομα μου στο γραφείο ήδη έτοιμο; Πώς ξέρατε...»

«Ο Μαλχη είναι πίσω απ’ όλα. Δεν υπήρχε περίπτωση να μην περάσεις την αξιολόγηση με τέτοιες ικανότητες που έχεις. Σε χρειαζόμαστε , βλέπεις, για τις αποστολές μας.» μου απάντησε ο Χερμπ. «Το γραφείο είναι δικό σου, και το δωμάτιο αυτό είναι της ομάδας μας. Εδώ θα συζητάμε τις υποθέσεις που θα αναλαμβάνουμε σαν ομάδα. Τώρα περιμένουμε τον Μαλχη να τελειώσει τις αξιολογήσεις και να έρθει να σας εξηγήσει την υπόθεση που μελετάμε ήδη οι τρεις μας, για να αρχίσετε να βοηθάτε και εσείς από αύριο.» είπε ο Χερμπ και άφησε από την αγκαλιά του την Μπλαιρ.

«Από αύριο;» ρώτησε η Μπλαιρ.

«Βεβαίως, αφού σήμερα είναι Πέμπτη. Αν μπορούσαμε θα ξεκινούσαμε και σήμερα, αλλά δεν βολεύει και τόσο» της απάντησε ο Χερμπ χαμογελώντας ζεστά στο τέλος της πρότασης του. «Μπορώ να σας κάνω μία ερώτηση κύριε Θοτς;»

«Παρακαλώ. Ρωτήστε ό,τι θέλετε.»

«Σχετίζεστε κάπως με τον σπουδαίο ανακριτή Θοτς, που έχει και αυτός μαγεία Κεραυνών;»

Δεν του απάντησα. Συνέχισα να κοιτάζω το ταμπελάκι με το όνομα μου.

«Όχι;»

«Είναι... ήταν... ο πατέρας μου»

«Ω, συγγνώμη. Δεν ήθελα να σας φέρω σε δύσκολη θέση. Βλέπετε, ήταν παλιός μου συνάδελφος. Αυτό εδώ ήταν το γραφείο του. Είχαμε πάει σε πολλές αποστολές μαζί, αλλά στην τελευταία του αποστολή επέμενε ότι έπρεπε να πάει μόνος του. Κάτι έλεγε για αλήθειες και ότι ήθελε να βρει τον εαυτό του, δεν είχα καταλάβει να σου πω την αλήθεια, αλλά μάλλον το γεγονός ότι πήγε μόνος του του στοίχισε την ζωή. Μπορείτε να δείτε πάντως την κορνίζα του απ’ έξω από το δωμάτιο αν...»

Έτρεξα γρήγορα έξω από το δωμάτιο. Πώς το είχα παρατηρήσει; Το πρόσωπο του πατέρα μου ακριβώς δίπλα από την πόρτα κορνιζαρισμένο. “...δεν ξέρω ποιος είμαι...δεν ξέρω που είμαι...” μπορούσα να ακούσω την εικόνα να λέει, καθώς αυτές οι λέξεις μου είχαν μείνει στο μυαλό τα τελευταία 12 χρόνια. Κάτω δεξιά έγραφε προφανώς το όνομα του και την ειδικότητα του: ανακριτής.

«Τι κοιτάς;» άκουσα την φωνή του Μαλχη. Γύρισα το βλέμμα μου στα αριστερά και τον είδα να με πλησιάζει φτιάχνοντας το μουστάκι του. Δεν του απάντησα. Έστρεψα το βλέμμα μου πάλι στον πατέρα μου.

«Αν επικεντρώνεσαι στους νεκρούς, δεν θα μπορέσεις να συνεχίσεις την δουλειά στην Ειδική Υπηρεσία. Σε κάθε υπόθεση όλοι ρισκάρουμε τις ζωές μας για το ευρύτερο καλό, την ασφάλεια των πολιτών» είπε ο Μαλχη και ακούμπησε το χέρι του στον ώμο μου. \textasciitilde Τελείωνε να κοιτάς τον πατέρα σου και πάμε μέσα\textasciitilde 

«Δεν μπορούσες να είσαι λίγο πιο διακριτικός;» του είπα χωρίς να τον κοιτάξω καθώς ένα δάκρυ κύλησε στο μάγουλο.

«Όχι δεν μπορούσα, αυτό είναι το στυλ μου. Και η Μπλαιρ έτσι είναι, οπότε αν θες να συνεχίσεις να δουλεύεις μαζί μας πρέπει να το συνηθίσεις...» \textasciitilde ...κλαψιάρη\textasciitilde 

«Καλά» είπα και έδιωξα το χέρι του από τον ώμο μου.

«Στο ξαναλέω πάντως, μη δεθείς πολύ με κανέναν. Όλοι θα πεθάνουν, σύντομα ή όχι. Μπορεί να πεθάνω και γω, ποιος ξέρει. Οι ζωές μας είναι σε κίνδυνο από την στιγμή που δεχόμαστε αυτή τη δουλειά. Οπότε πάρτο σαν φιλική συμβουλή αυτό. Καλά δεν σου λέω να γίνεις ρομπότ, αλλά κομμένο το κλάψιμο» μου είπε και προχώρησε προς το δωμάτιο.

«Ε... ντάξει...» του απάντησα.

«Επίσης, άλλη μία συμβουλή, μη διαβάζεις τις σκέψεις των συναδέλφων σου, τουλάχιστον στις αποστολές. Χάνεις χρόνο, αφού είμαστε όλοι με το μέρος σου» μου είπε σταματώντας λίγο πριν την πόρτα.

«Ε...ντάξει» του απάντησα και μπήκαμε μαζί στο δωμάτιο που ήταν η υπόλοιπη ομάδα. Όλοι σταμάτησαν ό,τι έκαναν και κοίταξαν τον Μαλχη. 

«Λοιπόν. Συστηθήκατε ή πρέπει να το κάνω εγώ και αυτό;» είπε ο Μαλχη στρέφοντας το βλέμμα του σε μένα.

«Μόνο με τον Χερμπ μιλήσαμε, με τους άλλους δεν...»

«Πω, πάλι εγώ έχω τον κλήρο. Λοιπόν. Χερμπ. Ειδικός στις μάχες, μαγεία Δύναμης, κλάσης Α.» είπε ο Μαλχη δείχνοντας τον Χερμπ.

«Μαγεία Δύναμης; Τι σημαίνει αυτό;» ρώτησα τον Χερμπ από περιέργεια, γιατί το όνομα της μαγείας του μου φάνηκε περίεργο.

«Μπορώ να αυξάνω και να μειώνω το μέγεθος των μυών μου κάνοντας με πιο δυνατό και καθώς μεγαλώνουν οι μυς μου μεγαλώνει και το σώμα μου. Προφανώς υπάρχει και ένα όριο στο πόσο μεγάλους μπορώ να κάνω τους μυς μου πριν κοπούν, αλλά αυτό το όριο είναι πολύ μεγαλύτερο από τους μπόντι μπίλντερς που έχεις ακουστά, αλλιώς δεν θα ήταν ιδιαίτερη η μαγεία μου»

«Σωστά... Και εσύ;» κοίταξα τον πράκτορα που μας είχε συνοδεύσει μέχρι τον χώρο αναμονής.

«Ρομπ, κατάσκοπος. Η μαγεία μου όπως είδες και πριν είναι μαγεία Τροχιάς. Μπορώ να αλλάξω την τροχιά οποιουδήποτε αντικειμένου, έμβιου ή μη, ή ακόμη και την τροχιά κάποιας μαγείας. Βέβαια δεν μπορώ να αλλάξω την ταχύτητα του αντικειμένου, για αυτό σας είχα πει να αρχίσετε να περπατάτε πριν σας σηκώσω και σας συνοδεύσω μέχρι τον χώρο αναμονής.»

«Α και γιατί χρειαζόταν να το κάνεις αυτό; Λες και  δεν μπορούσαμε να περπατήσουμε χωρίς εσένα...» τον ρώτησε η Μπλαιρ που μάλλον είχε πέσει και αυτή θύμα της μαγείας του Ρομπ.

«Είναι ένα κόλπο που μου ζητάει να κάνω η γραμματέας στους νέους για να δει την αντίδραση τους»

«Άρα όντως απλά έκανες επίδειξη της μαγείας σας στην είσοδο χωρίς λόγο» είπε εκνευρισμένη η Μπλαιρ.

«Δεν μπορώ να το αρνηθω, αλλά δεν μπορώ να πω και πως ήταν “επίδειξη”. Απλά σας μετέφερα μέχρι τη μέση του διαδρόμου»

«Καλά καλά. Ένα άγγιγμα να σου έκανα και το δαχτυλάκι σου ούτε χαλίκι να κουνήσει» είπε η Μπλαιρ γυρνώντας πλάτη στον Ρομπ και τινάζοντας την κοτσίδα της.

«Γιατί καλέ; Έχεις μήπως κάποιου είδους “άντι-μαγείας”;» την ρώτησε ειρωνικά ο τρίτος άντρας.

«Μπίνγκο» είπε ο Χερμπ.

«Ναι έχω άντι-μαγεία. Σε σχέση με εσάς τους υπόλοιπους, δεν μπορώ να πετάω φωτιές από τα χέρια μου ή να αλλάζω μαγικά την κατεύθυνση ενός αντικειμένου...»

«...τροχιά....» την διέκοψε ο Ρομπ.

«Τροχιά, ξετροχιά. Η “μαγεία” μου εμένα είναι να ουδετεροποιώ, ή στην γλώσσα σας να “αχρηστεύω”, άλλες μαγείες, είτε ακουμπόντας τον χρήστη, είτε μέσω κάποιου αντικειμένου που του έχω δώσει λίγη από την αύρα της μαγείας μου» είπε η Μπλαιρ και κάθισε στον καναπέ δίπλα στον τρίτο άντρα που την σχολίασε προ ολίγου. «Εσύ τώρα τι πρόβλημα έχεις και κάθεσαι μόνος σου στον καναπέ;» είπε και τον άρχισε να τον αγγίζει επανειλημμένα στον ώμο.

«Σταμάτα, σκέφτομαι» είπε ο άντρας.

«Τι σκέφτεσαι καλέ. Αφού ο ανακριτής μας είναι ο Σπαρκ και ο κατάσκοπος μας είναι ο κύριος τροχιάς. Εσύ λογικά είσαι ειδικός στις μάχες σαν εμένα και τον Χερμπ, δεν χρειάζεται να σκέφτεσαι και πολύ...» του είπε η Μπλαιρ συνεχίζοντας να τον πειράζει στον ώμο.

«Δεν μπορώ να έχω προσωπικά προβλήματα δηλαδή;» της απάντησε ο άντρας διώχνοντας το χέρι της από πάνω του.

«Ηρεμήστε! Φυτ, συστήσου στην Μπλαιρ και τον Σπαρκ.» είπε ο Μαλχη δείχνοντας μου τον άντρα που καθόταν δίπλα στην Μπλαιρ.

«Με λένε Φυτ, από ότι καταλάβατε, είμαι ειδικός στις μάχες...»

«Άρα το βρήκα!» τον διέκοψε η Μπλαιρ.

«Ναι, αφού κάθε ομάδα έχει έναν ανακριτή, έναν κατάσκοπο και τρεις ειδικούς στις μάχες, δεν ήταν και τόσο δύσκολο...» την ειρωνεύτηκε ο Φυτ.

«Μπορείτε να σταματήσετε να τσακώνεστε και να τελειώνετε;» είπε εκνευρισμένος ο Μαλχη. «Μαγεία και κλάση γρήγορα, τελείωνε.»

«Έχω μαγεία...»

«Δύναμης. Το πέτυχα;» είπε η Μπλαιρ. «Σε λένε Φυτ, οπότε τι άλλη μαγεία θα μπορούσες να έχεις;»

«...Φύσης... έχω μαγεία Φύσης...» είπε ο Φυτ κοιτώντας περίεργα την Μπλαιρ. «Μπορώ κυρίως να ελέγχω και να παράγω ρίζες δέντρων, τις οποίες χρησιμοποιώ στις μάχες.»

«Είναι πολύ καταστροφική μαγεία, όσο απλή και να ακούγεται» συμπλήρωσε ο Χερμπ. «Αν δείτε πως την χρησιμοποιεί στις μάχες, ειδικά όταν χρησιμοποιεί τις μυτερές του ρίζες. Ο Φυτ είναι από τους πιο δυνατούς μας πράκτορες κλάσης Α.»

«Υπερβάλεις Χερμπ...» σχολίασε ο Φυτ.

«Καθόλου. Απλά είσαι μετριόφρον. Να ξέρεις, χωρίς εσένα μπορεί και να μην ζούσαμε εμείς οι τρεις αυτή τη στιγμή. Καλά δεν λέω Ρομπ;»

«Δεν ξέρω αν δεν θα “ζούσαμε”, αλλά ναι ο Φυτ είναι σίγουρα πιο δυνατός από εμάς και μας έχει κάνει την ζωή πιο εύκολη από τότε που εντάχθηκε στην ομάδα» απάντησε ο Ρομπ και χαμογέλασε στον Φυτ.

«Ωραία, τέλος τα είπατε. Η ώρα έχει πάει 4:50 και έλεγα να πάω πιο νωρίς σήμερα στο καζίνο αφού τελείωσα τις αξιολογήσεις, αλλά απ’ ότι φαίνεται θα χρειαστεί να κάτσουμε πολύ ώρα εδώ αν συνεχίζεται να σχολιάζεται ο ένας τον άλλον...» είπε ο Μαλχη και σταύρωσε τα χέρια του.

«Και το σχέδιο;» τον ρώτησε ο Χερμπ. «Δεν θα το εξηγήσεις στην Μπλαιρ και τον Σπαρκ;»

«Πω, το είχα ξεχάσει. Ευχαριστώ ρε Χερμπ. Θα σε κεράσω μία μπύρα στο καζίνο. Θα έρθεις έτσι;» είπε ο Μαλχη και του έκλεισε το μάτι.

«Βεβαίως, περίμενα να τελειώσεις τις αξιολογήσεις για να πάμε» του απάντησε ο Χερμπ χαμογελώντας.

«Ωραία. Λοιπόν...»

«Καθίστε, εγώ δεν συστήθηκα» διέκοψα τον Μαλχη

«Δεν χρειάζεται. Όλοι εδώ ξέρουν, με τον ένα ή με τον άλλο τρόπο ότι διαβάζεις σκέψεις, τους το είπα το μεσημέρι πριν πάω στην αίθουσα εξάσκησης και προετοιμαστώ για τις αξιολογήσεις.» μου είπε ο Μαλχη.

«Εξάλλου το σχέδιο μας εξελίσσεται γύρω από εσένα, οπότε όφειλα να τους το αναφέρω.»

«Ε...ντάξει...» του απάντησα. «Μπορώ τουλάχιστον να κάνω ένα τηλεφώνημα πριν αρχίσεις να μας εξηγείς το σχέδιο;»

«Κάνε ότι θες, απλά μην αργήσεις. Έχεις δύο λεπτά, πήγαινε.» μου είπε εκνευρισμένος ο Μαλχη.

Έτρεξα για ακόμη μια φορά έξω από το δωμάτιο, αυτή τη φορά όχι για να δω τον πατέρα μου, αλλά για να πάρω τηλέφωνο τον Βασι. “...αύρα της μαγείας μου...”, τι σημαίνει αυτό; Επίσης, πώς ο Κουπ και ο Μαλχη κατάλαβαν πως είμαι “κλάσης Β” μόνο με μία χειραψία; Δεν καταλαβαίνω. Δεν έχω χρόνο για σκέψεις τώρα, πρέπει να πάρω τηλέφωνο τον Βασι να του μιλήσω έστω για λίγο. 

Έβγαλα το κινητό μου από την τσέπη μου και προσπάθησα να τον καλέσω. Αυτός όμως δεν μου απάντησε. Τον πήρα τηλέφωνο και δεύτερη φορά, αλλά μια απ’ τα ίδια. Τρίτη φορά, τίποτα. Δεν είχα χρόνο να πάω μέχρι την είσοδο για να ελέγξω αν με περιμένει εκεί ο Βασι, οπότε του έστειλα ένα μήνυμα “Πάρε με σε 20”, υπολογίζοντας περίπου πόση ώρα θα χρειαζόταν ο βιαστικός Μαλχη να μας εξηγήσει το σχέδιο.Καθώς έκλεινα το κινητό μου είδα την νοσοκόμα που είχε θεραπεύσει τον Βασι πριν την αξιολόγηση μας να έρχεται προς το μέρος μου. 

«Εσείς πρέπει να είστε... ο κύριος Θοτς... Με λένε Ιφι, χάρηκα για την γνωριμία» είπε η νοσοκόμα και μου έδωσε το χέρι της. Το βλέμμα ήταν στραμμένο στο πάτωμα και φαινόταν λες και ντρεπόταν να μου μιλήσει.

«Χάρηκα, Ιφι. Να σε ρωτήσω κάτι;»

«Ναι, ό,τι θέλετε. Πείτε μου.»

«Μήπως είδατε τον άντρα που θεραπεύσατε στον χώρο αναμονής καθώς ερχόσασταν;»

«Ό... χι... τουλάχιστον δεν τον πρόσεξα...»

«Δεν πειράζει, σας ευχαριστώ. Συγγνώμη τώρα, πρέπει να πηγαίνω στο γραφείο μου» της είπα και γύρισα πίσω στο δωμάτιο της ομάδας μας, ή πλέον “το γραφείο μου”.

«Ούτε ένα λεπτό δεν πέρασε» μου είπε ειρωνικά ο Μαλχη.

«Δεν πειράζει. Πες μας τώρα αυτά που θέλεις να μας πεις και άσε με εμένα.» του απάντησα. Κοίταξα πίσω μου και είδα την Ιφι να κλείνει την πόρτα προσπαθώντας λίγο να κρυφοκοιταξει και να κρυφακούσει τι λέγαμε.

«Πολύ ωραία λοιπόν. Ακούστε προσεκτικά όλοι, γιατί θα σας δώσω και πληροφορίες που δεν σας ήταν γνωστές πριν...»

\chapter{}

«Έρχεσαι συχνά εδώ;»

«Μόνο τις Κυριακές. Γιατί ρωτάς;»

«Τίποτα, απλά δεν σε έχω ξαναδεί ποτέ σε αυτό το μπαρ»

«Αποκλείεται, αφού έρχομαι εδώ... τέλος πάντων»

«Γιατί σταμάτησες;»

«Δεν είναι τίποτα...» \textasciitilde δεν πρέπει να μάθει γιατί είμαι εδώ. Που ξέρω γω αν είναι πράκτορας ή όχι και θέλει να με καλοπιάσει κερνόντας με μπύρα;\textasciitilde 

«Θα πιεις καθόλου; Δεν σε κέρασα για να κάθεσαι να με κοιτάς»

«Ναι γιατί όχι...» \textasciitilde ...αλλά μετά πρέπει να φύγω γρήγορα. Η δουλειά μου δεν μπορεί να σταματήσει για κανέναν. Αν τους αφήσω θα πάθουν στερητικό σύνδρομο\textasciitilde 

«Να σε ρωτήσω φίλε, τι δουλειά κάνεις;»

«Είμαι...»\textasciitilde ...τι να του πω τώρα...\textasciitilde 

«Ναι;»

«Έμπορος φαρμάκων.»

«Ά, μεταφέρεις τα απαραίτητα από τη μία πόλη σε άλλη, ε; Ή τα φτιάχνεις και τα πουλάς;»

«Το πρώτο θα μπορούσε να πει κανείς περιγράφει πιο σωστά τη δουλειά μου»\textasciitilde ...σίγουρα αυτός είναι πράκτορας, πρέπει να σηκωθώ να φύγω...\textasciitilde «Λοιπόν φίλε, δεν θα πιω άλλο. Ευχαριστώ για το ωραίο βράδυ, αλλά θα πρέπει να πηγαίνω»

«Εντάξει φίλε μου, καλό βράδυ και καλή τύχη» του είπα και έσκυψα γρήγορα να ειδοποιήσω τους υπόλοιπους. Αυτός με είδε που άνοιξα το κινητό μου και μάλλον ανακουφίστηκε που με ξεφορτώθηκε.

«Θα έρθεις με εμάς τώρα» άκουσα την Μπλαιρ να του λέει.

«Ποια νομίζεις ότι είσαι;» της απάντησε ο “φίλος" μου.

«Μπλαιρ Γκαμ, πράκτορας της Ειδικής Υπηρεσίας. Σε παρακαλώ φόρα αυτές τις χειροπέδες και πάμε στο βανάκι»

Ο “φίλος” μου άρχισε να τρέχει προς την είσοδο του μπαρ, αλλά εκεί βρήκε τις ρίζες του Φυτ να του μπλοκάρουν την απόδραση. Εκνευρισμένος γύρισε και πήγε να πυροβολήσει την Μπλαιρ, που περπάταγε προς το μέρος του, η τροχιά της σφαίρας όμως άλλαξε τελευταία στιγμή και το πόδι του θύτη κατέληξε να γίνεται θύμα της σφαίρας. Η Μπλαιρ του έβαλε τις χειροπέδες με τη βία και τον πήγε στο βανάκι όπου περίμενε ο Χερμπ με άλλους τρεις ύποπτους.

«Άλλος ένας μας έμεινε για σήμερα» μου είπε ο Ρομπ που βρισκόταν πίσω από το μπαρ κάνοντας τον μπάρμαν.

«Ναι, και ο μόνος κλάσης Α. Ας ελπίσουμε να μην υπάρξουν φασαρίες» είπε ο Φυτ και κάθισε δίπλα μου.  «Ρομπ, είπες πως ο τελευταίος συνήθως έρχεται μετά τις 10:00 το βράδυ και είναι ακόμη 9:47; Έχουμε οπότε χρόνο. Ας ελπίζουμε τουλάχιστον πως δεν ήταν στη γειτονιά, γιατί θα είχε ακούσει σίγουρα το πυροβολητό»

«Πολύ καλό σχέδιο πάντως, μπράβο σε όποιον το οργάνωσε. Πώς ξέρουμε όμως με σιγουριά πως όλοι οι ναρκομανείς και οι έμποροι ναρκωτικών έρχονται σε αυτό το μπαρ;» ρώτησα από περιέργεια, γιατί όλα τα πήγαιναν με το μέρος μας και άρχισα να έχω υποψίες.

«Α αυτό το ξέρουμε λόγω του Ρομπ. Τον τελευταίο μήνα έχει μελετήσει όλους τους ναρκομανείς που είχαμε στους φακέλους στα κεντρικά» είπε ο Φυτ χτυπώντας τον Ρομπ στον ώμο.

«Εντάξει, όχι και όλους... όσους μπόρεσα.» είπε ο Ρομπ. «Οι έρευνα μου έδειξε πως αυτό το τετράγωνο που βρισκόμαστε είναι το μέρος που μαζεύονται όλοι οι ναρκομανείς για να προμηθευτούν. Υπέθεσα λοιπόν πως οι προμηθευτές τους θα θέλουν να ξεκουραστούν κάπου κοντά και άρχισα να κρατάω φωτογραφίες τους και να τους ακολουθώ όποτε έβρισκα την ευκαιρία και όλοι προμηθευτές έρχονταν για κάποιο λόγο σε αυτό το μπαρ γύρω στις ίδιες ώρες τις ίδιες περίπου μέρες μέσα στην εβδομάδα»

«Και που ήξερες ότι δεν θα σπάσουν την ρουτίνα τους σήμερα;» τον ρώτησα, γιατί αυτό ήταν μία τρύπα στο σχέδιο μας.

«Για αυτό βρισκόμαστε εδώ τρία βράδια τώρα. Το σκέφτηκε και ο Μαλχη αυτό και πρότεινε να βρισκόμαστε στο μπαρ όλο το Σαββατοκύριακο και να πιάσουμε τουλάχιστον τους προμηθευτές για να τους ανακρίνεται πίσω στα κεντρικά και να βρούμε τον... “αρχιπρομηθευτή”, ή τους “αρχηπρομηθευτές” αν είναι πολλοί, αλλά δεν νομίζω» είπε ο Ρομπ και έκανε πως καθάριζε ένα ποτήρι μπύρας για να μπει στον ρόλο του.

«Το έχετε πολύ το σχέδιο, ε;» του είπα.

«Ε ντάξει, αφού έχουμε τον Μαλχη ως μυαλό τι περίμενες;» πετάχτηκε ο Φυτ. «Ο αρχηγός είναι ο πιο έξυπνος πράκτορας της Ειδικής Υπηρεσίας. Χωρίς αυτόν δεν ξέρω τι θα κάναμε. Και να σκεφτείς δεν είναι καν μεγάλος ηλικιακά.»

«Καλά δεν είναι και σαν εμάς» είπε ο Ρομπ

«Γιατί πόσο λες εσύ ότι είναι;»

«Λέω ότι είναι τουλάχιστον 40.»

«Τι 40; Εγώ τον έκανα για εικοσάρη. Πως προέκυψε το 40;»

«Εγώ είμαι 25 και τον θυμάμαι από όταν ήμουν 20 και μπήκα στην Ειδική Υπηρεσία, οπότε υπέθεσα πως θα είναι κάποια χρόνια μεγαλύτερος τουλάχιστον, αφού ήταν και τότε αρχηγός» είπε ο Ρομπ και άφησε το ποτήρι που καθάριζε και πήρε το δικό μου ποτήρι. «Δεν σε πειράζει να το αδιάσω, ε; Αφού δεν πίνεις έτσι κι αλλιώς.»

«Ναι άδειασε το, ευχαριστώ.» του απάντησα.

«Εσύ Σπαρκ πόσο τον κάνεις τον Μαλχη;» με ρώτησε ο Φυτ.

«Ε... δεν ξέρω. Θα συμφωνήσω με τον Ρομπ, αφού και το πρόσωπο του Μαλχη, ειδικά το μουστάκι του, προδίδει ότι η ηλικία του είναι αρκετά μεγαλύτερη.»

«Δεν νομίζω ρε φίλε» είπε απογοητευμένος με την απάντηση μου ο Φυτ. «Η συμπεριφορά του πάντως δεν είναι σίγουρα συμπεριφορά σαραντάρη. Λέτε να ξέρει ο Χερμπ πόσο χρονών είναι ο Μαλχη; Πρέπει να τον ρωτήσουμε όταν βρεθεί η ευκαιρία.»

«Να τος» είπε ο Ρομπ διακριτικά. Ο Φυτ και εγώ κοιτάξαμε της είσοδο του μπαρ και είδαμε έναν έφηβο με αφάνα ο οποίος έκατσε σε ένα τραπέζι κοντά στην είσοδο.

«Σίγουρα είναι αυτός;» ρώτησε ο Φυτ.

«Χωρίς αμφιβολία. Τον έχω δει τέσσερις φορές τον τελευταίο μήνα να προμηθεύει ανθρώπους σε αυτό το τετράγωνο.»

«Θες να του πιάσω συζήτηση όπως στους υπόλοιπους ή να τον συλλάβουμε κατευθείαν;» ρώτησα για να μάθω τι πλάνο θα ακολουθούσαμε.

«Φέρτε θα έλεγα την Μπλαιρ να βοηθήσει. Η άντι-μαγεία της θα μας κάνει τη ζωή εύκολη. Αν χρειαστείτε κάτι θα είμαι εδώ και θα παρακολουθώ» είπε ο Ρομπ και πήγε να σερβίρει δύο πελάτες που κάθονταν στην άλλη μεριά του πάσου. Είχα διαβάσει ήδη τις σκέψεις τους, δεν ήταν ύποπτοι, οπότε δεν ασχοληθήκαμε παραπάνω. Τους αφήσαμε απλά στην γωνία τους για να μπορέσει ο Ρομπ να έχει κάποια κάλυψη και να μην μιλάει μόνο σε εμένα.

Είδα την Μπλαιρ να μπαίνει πάλι στο μπαρ, μουρμουρίζοντας και παίζοντας κλασικά με τα δάχτυλα της, και ο Φυτ της έκανε νόημα να έρθει. Αυτή σκόνταψε στην καρέκλα που καθόταν ο έφηβος με την αφάνα. Τον συγχώρεσε βιαστικά και ήρθε προς το μέρος μας. Είδα την έκφραση του εφήβου και ήταν κατακόκκινος και έπαιζε με τα δάχτυλα του προσπαθώντας να την μιμηθεί. Λες να του αρέσει η Μπλαιρ;

«Ήρθα, τι έγινε;»

«Τίποτα. Μόνο που ο τελευταίος μας ύποπτος κάθεται στην καρέκλα στην οποία σκόνταψες» είπε ο Φυτ ειρωνικά.

«Τι; Αυτός; Αποκλείεται. Αφού είναι πολύ μικρός για να είναι έμπορος ναρκωτικών» προβληματίστηκε ο Μπλαιρ.

«Ο Ρομπ λέει πως είναι 100\% αυτός, τι θες να κάνουμε τώρα. Μη μου πεις πως δεν ένιωσες και την αύρα του;»

«Την ένιωσα όταν τον ακούμπησα καταλάθος για να μην σκοντάψω, σίγουρα κλάσης Α, αλλά και πάλι είναι περίεργο...» είπε η Μπλαιρ και φάνηκε σαν να άρχισε να σκέφτεται κάτι. Ποια αύρα; Τι λένε; Μήπως σχετίζεται με την χειραψία μου με τον Μαλχη και τον Κουπ;

«Τι... εννοείται αύρα;» ρώτησα για να μου λυθεί η απορία. Οι δυο τους με κοίταξαν λες και ρώτησα κάτι αυτονόητο.

«Δώσε μου το χέρι σου Σπαρκ» είπε η Μπλαιρ αρπάζοντας το χέρι μου. \textasciitilde Οριακά αύρα κλάσης Β, για αυτό δεν μπορεί να νιώσει την αύρα\textasciitilde  «Είσαι κλάσης Β, για αυτό δεν ξέρεις για ποιο πράγμα μιλάμε. Όσοι άνθρωποι είναι κλάσης Α μπορούν με ένα άγγιγμα του σώματος ασυνείδητα να νιώσουν το πόσο... “δυνατός” ας το πούμε είναι κάποιος με την μαγεία του μέσω της αύρας του. Όπως εσύ ασυνείδητα διάβαζες σκέψεις με ένα άγγιγμα; Ή τουλάχιστον νόμιζες πως μπορούσες να διαβάσεις μόνο έτσι σκέψεις. Μπορεί να σχετίζεται δεν ξέρω, αλλά πάντως, για να μην έχεις ακούσει ποτέ για την αύρα, είσαι κλάσης Β.»

Έσκυψα το κεφάλι μου απογοητευμένος.

«Δεν το λέω για κακό, Σπαρκ. Η ικανότητα σου να διαβάζεις σκέψεις μας έχει χρησιμέψει υπερβολικά μη σου πω σε αυτή την αποστολή. Χάρη σε σένα ξέρουμε ποιους να συλλάβουμε και ποιους όχι σε αυτό το μπαρ, αλλά και όχι μονό. Χθες καθώς πηγαίναμε στο βανάκι δεν βρήκες άλλον έναν προμηθευτή διαβάζοντας τις σκέψεις του;»

«Ναι...» είπα χωρίς να σηκώσω το κεφάλι μου «...απλά απογοητεύτηκα γιατί ήθελα και γω να νιώσω την “αύρα” που λέγατε.»

«Δεν πειράζει Σπαρκ. Άσε τις αύρες και τις σαύρες σε μας. Εσύ κάνε την δουλειά σου τώρα και επιβεβαίωσε μία αν ο αφάνας είναι όντως ύποπτος» μου είπε η Μπλαιρ δείχνοντας τον έφηβο που καθόταν δίπλα στην πόρτα.

Επικεντρώθηκα στο διάβασμα των σκέψεων του υπόπτου μας. \textasciitilde ...πως το έκανε αυτό ρε φίλε με τα δάχτυλα της. Κάτω,αριστερά, δεξιά, πάνω. Τι είναι αυτή η κίνηση; Χρησιμεύει σε κάτι; Μπορεί να είναι κάτι κοριτσίστικο που να μην καταλαβαίνω. Και πάλι όμως ήταν τόσο ωραία... μακάρι να έβρισκα ευκαιρία να της μιλήσω...\textasciitilde

«Όπως υποπτευόμουν» είπα.

«Είναι ύποπτος;» με ρώτησε ο Μπλαιρ.

«Όχι... σε γουστάρει...» είπα και είδα την Μπλαιρ να αηδιάζει με τις λέξεις που ξεστόμησα.

«Όχι, όχι, όχι. Εγώ να ξέρετε δεν τον πλησιάζω, δεν είναι ο τύπος μου. Μόνο χειροπέδες θα του βάλω και τέλος. Αφού είναι και ο τελευταίος για σήμερα, θα φύγουμε όλοι μαζί. Να τον πάτε εσείς μέχρι το...»

«Σε χρειαζόμαστε, όμως.» την διέκοψα. «Χρειαζόμαστε την άντι-μαγεία σου, αφού είναι κλάσης Α. Δεν γίνεται να ρισκάρουμε να ξεσπάσει μάχη εδώ μέσα, θα τραυματιστεί πολύς κόσμος.»

«Έχεις δίκιο... αλλά δεν επιβεβαίωσες ποτέ ότι είναι έμπορος ναρκωτικών. Εγώ δεν τον πλησιάζω αν δεν το επιβεβαιώσεις» είπε η Μπλαιρ και μου γύρισε πλάτη.

Προσπάθησα να ξαναδιαβάσω τις σκέψεις του. Αυτός μιμούνταν ακόμη την Μπλαιρ με τα δάχτυλα του. \textasciitilde ...πώς και δεν έχω δει κανέναν άλλο να κάνει έτσι τα δάχτυλα του στο σχολείο; Λες να είναι όντως κάποιο κοριτσίστικο πράγμα; Η επόμενη μου πελάτης είναι κοπέλα, λες να την ρωτήσω ενώ την προμηθεύω; Μπα, αυτή με συναντάει μόνο για έναν σκοπό, δεν νομίζω να θέλει να απαντήσει κάποια παραπάνω ερώτηση...\textasciitilde . Άκουσα ό,τι ήθελα να ακούσω.

«Είναι έμπορος. Φυτ εσύ το νου σου στην πορτα. Μπλαιρ πάμε» είπα καθώς σηκώθηκα από το σκαμπό μου και η Μπλαιρ με ακολούθησε. «Προτείνω να πας πρώτα εσύ, για να τον ακουμπήσεις ώστε να λειτουργήσει η άντι-μαγεία σου και να του θα του βάλω χειροπέδες.»

«Ίου, δεν τον ακουμπάω. Είναι πολύ μικρός για μένα» είπε η Μπλαιρ αηδιασμένη.

«Δεν είναι θέμα προτιμήσεων τώρα, πρέπει να τον συλλάβουμε»

«Και γιατί δεν πας μόνος σου; Έχω βάλει ήδη στις χειροπέδες αύρα άντι-μαγείας, οπότε αν τον ακουμπήσεις με τις χειροπέδες θα πετύχεις το ίδιο πράγμα» είπε η Μπλαιρ δίνοντας μου τις χειροπέδες.

«Α... τώρα κατάλαβα αυτό που είπες την Πέμπτη... όταν περιέγραφες την μαγεία σου. Άρα βάζεις λίγη από την αύρα σου στις χειροπέδες και μπορούν να χρησιμοποιηθούν για να ουδετεροποιήσουν την μαγεία του άλλου ενώ τις φοράει; Ή τουλάχιστον όσο τις ακουμπάει;»

«Ναι»

«Αυτό δουλεύει και με άλλα αντικείμενα;» συνέχισα να την ρωτάω από περιέργεια.

«Ναι»

«Δηλαδή αν έχω μία ασπίδα, ή μία σφαίρα, ή μία ενδυμασία, ή ένα...»

«Ναι ναι ναι. Πήγαινε τώρα.»

«Εντάξει... απλά ρώτησα» είπα και έσπευσα προς τον έφηβο με την αφάνα.

Προσπάθησα να διαβάσω τις σκέψεις του για να επιβεβαιώσω πάλι πως είναι ο στόχος μας. \textasciitilde ...δεν ξέρω βέβαια αν θα έχει αγόρι. Λες αυτός που έρχεται προς το μέρος μου να είναι το αγόρι της; Ωχ όχι, κρατάει χειροπέδες. Λες να ξέρει ότι κουβαλάω κοκαΐνη; Ή λες απλά να έρχεται να με πλακώσει επειδή μίλησα στην κοπέλα του; Ή μήπως...\textasciitilde  άκουσα ό,τι χρειαζόταν. Είδα τον έφηβο να σηκώνεται διακριτικά και να πηγαίνει προς την είσοδο.

«Που πας; Δεν έχεις παραγγείλει ακόμη κάτι να πιεις.» του φώναξα, αλλά αυτός δεν γύρισε και άρχισε να τρέχει προς την πόρτα. Ο Φυτ προφανώς του έκλεισε την διέξοδο και ο μικρός σκόνταψε μπροστά από τις ρίζες. 

\textasciitilde Τι θέλει τώρα αυτός; Γιατί με ακολουθεί; Λες να ξέρει τι έχω στην δεξιά μου τσέπη;\textasciitilde «Δεν έχω τίποτα πάνω μου, υπόσχομαι» είπε και σήκωσε τα χέρια του.

«Δεν χρειάζεται να αντιστέκεσαι. Έλα απλά μαζί μου» είπα καθώς προσπάθησα να τον ακουμπήσω με τις χειροπέδες. Αυτός όμως πριν προλάβω να τον φτάσω τέντωσε το χέρι του και εκτίναξε μεγάλη ποσότητα φωτιάς προς το μέρος μου, είδα όμως το βέλος του Ρομπ και δεν πανικοβλήθηκα. Η φωτιά άλλαξε τροχιά και πήγε στο ταβάνι του μπαρ, αντί να φτάσει σε εμένα.

«Σου είπα, δεν χρειάζεται να αντιστέκεσαι. Έλα, δώσε μου το χέρι σου»

«Είναι γιατί μίλησα στην κοπέλα σου;»

«Τι; Όχι. Είναι για αυτό...» είπα και έβγαλα την κοκαΐνη από την δεξιά του τσέπη. «Θα μου τα πεις όλα μετά, φόρα τώρα τις χειροπέδες και πάμε στο βανάκι» του είπα καθώς του φόρεσα τις χειροπέδες. Έκανα νεύμα στον Φυτ και την Μπλαιρ να φύγουμε και τους είδα να σηκώνονται και να λένε και στον Ρομπ ότι η δουλειά μας στο Μπαρ τελείωσε.

«Θες να τον πάρω εγώ;» μου είπε ο Φυτ.

«Ό,τι θες. Έτσι κι αλλιώς στο βανάκι εδώ δίπλα πάμε, δεν είναι μακριά»

«Έλα θα τον πάρω εγώ, εσύ έκανες όλη τη δουλειά» μου είπε ο Φυτ και ξεκίνησε να προχωράει προς το βανάκι.

«Εντυπωσιακό Σπαρκ» άκουσα την Μπλαιρ από πίσω μου. «Πώς δεν φοβήθηκες εκεί που σου έριξε φωτιά; Εγώ θα είχα τρέξει από τον φόβο μου»

«Είδα το βέλος του Ρομπ... για αυτό δεν φοβήθηκα»

«Όντως; Εγώ δεν θα το είχα παρατηρήσει καν από το άγχος της στιγμής. Είσαι παρατηρητικός Σπαρκ, μπράβο σου» είπε η Μπλαιρ και μου χαμογέλασε.

«Πάμε παιδιά, τι κάθεστε;» είπε ο Ρομπ που περνούσε δίπλα μας. «Περιμένει ο Χερμπ άδικα, πάμε προς το αμάξι»

«Ναι πάμε στον Χερμπ. Να κάτσω εγώ σήμερα μαζί του στο βανάκι;» ρώτησε η Μπλαιρ γεμάτη χαρά, σαν μικρό παιδί.

«Όχι. Ο Φυτ, είπαμε, πρέπει να είναι στο βανάκι μαζί με τον Χερμπ για να μπορέσει να τον βοηθήσει αν γίνει κάτι και μπορέσει κάποιος από τους υπόπτους μας ξεφύγει. Δεν είναι θέμα προτίμησεις...» της απάντησε ο Ρομπ, συνθλίβοντας όλες τις ελπίδες της Μπλαιρ.

«Δεν έχεις πλάκα...» είπε η Μπλαιρ εκνευρισμένη και πήγε στο αμάξι.

«Πάμε Σπαρκ. Άστη να παραπονιέται. Κάθε μέρα το ίδιο πανηγύρι με αυτή την κοπέλα.» μου είπε ο Ρομπ και άνοιξε το αμάξι για να μπούμε μέσα. «Σπαρκ, πάρε και τον Χερμπ τηλέφωνο και πες του να ξεκινήσει. Θα τον ακολουθήσω από πίσω όπως και τις άλλες φορές.»

«Εντάξει» του απάντησα ενώ έβαζα ζώνη. Σήκωσα το κινητό μου για να τον πάρω, αλλά είδα το βανάκι να βγαίνει από το στενό που βρισκόταν. Ο Χερμπ μας χαμογελούσε με σηκωμένο τον αντίχειρα του.

«Ξεπαρκάρετε και πάμε. Μπορεί προλάβω να πάω στο καζίνο με τον Μαλχη σήμερα» είπε ο Χερμπ κατεβάζοντας το παράθυρο του.

«Αφού ο Μαλχη και ο Σπαρκ θα ανακρίνουν τους υπόπτους σήμερα το βράδυ» του απάντησε ο Ρομπ.

«Τι; Πω... Ρομπ δεν έχεις πλάκα. Μου χάλασες τη διάθεση» είπε ο Χερμπ εκνευρισμένος και έκλεισε το παράθυρο.

Ο Ρομπ έβαλε μπροστά το αυτοκίνητο αδιαφορώντας με τις λέξεις το Χερμπ. Καθώς έβγαινε στον δρόμο το βανάκι, είδα την πίσω πόρτα να ανοίγει. Ο έφηβος με την αφάνα είχε καταφέρει να την ξεκλειδώσει χωρίς να βγάλει τις χειροπέδες του και βγήκαν όλοι από το βανάκι. Οι ελπίδες τους όμως συνθλίφτηκαν όσο γρήγορα όσο εμφανήστικαν, γιατί ο Φυτ χρησιμοποίησε την μαγεία του και τους περικύκλωσε έναν έναν με τις ρίζες του.

«Βλέπεις γιατί χρειάζεται ο Φυτ στο βανάκι;» είπε ο Ρομπ στην Μπλαιρ, αλλά αυτή δεν κοίταξε καν μπροστά για να δει το θέαμα.

Είδα τον Χερμπ να κατεβαίνει από το βανάκι και έβγαλα την ζώνη μου για να πάω να τον βοηθήσω, όμως αυτός μας έκανε νεύμα να κάτσουμε εκεί που βρισκόμασταν. Σε μία στιγμή, από ένας μέσος πενηντάρης, ο Χερμπ ψήλωσε 1,5 μέτρο και οι μυς του έγιναν αρκετές φορές μεγαλύτεροι από έναν μέσο γυμνασμένο άνθρωπο. Ο Χερμπ έχωσε πίσω και τους πέντε υπόπτους στο βανάκι και φάνηκε σαν να φωνάζει κάτι στον Φυτ. Μόλις το φώναξε οι ρίζες του Φυτ περικύκλωσαν ολόκληρο το βανάκι και ο Χερμπ κόντυνε στο πραγματικό του ύψος, με τους μυς του να επιστρέφουν στο κανονικό τους μέγεθος.

«Εντάξει!» μας φώναξε ο Χερμπ καθώς έμπαινε στο βανάκι. «Πάμε!»

«Ξεκινάμε παιδιά, Σπαρκ βάλε πάλι ζώνη» μου είπε ο Ρομπ και χαμήλωσε το χειρόφρενο.

«Εντάξει» του είπα και φόρεσα τη ζώνη μου.

Είμαστε πολύ καλή ομάδα έχω να πω. Σε όποιον τομέα έχει ελλείψεις κάποιος, οι υπόλοιποι μπορούμε να τον καλύψουμε. Εγώ, για παράδειγμα, που δεν μπορώ να εμπλακώ σε μάχη, έχω την Μπλαιρ, τον Ρομπ και τον Φυτ που με βοηθούν έμμεσα με τις δικές τους μαγείες και καταφέρνω να τα βγάλω πέρα. Για τους άλλους να πω την αλήθεια δεν έχω προσέξει και τόσα μειονεκτήματα όσα έχω παρατηρήσει στον εαυτό μου αυτές τις μέρες, αλλά δεν νομίζω να υπάρχει κάτι που ο Χερμπ με την εμπειρία του, ο Ρομπ με την παρατηρητικότητα του, ο Φυτ με την ευελιξία των ικανοτήτων του και εγώ και η Μπλαιρ με τις ιδιαίτερες ικανότητες μας να μην μπορούμε να αντιμετωπίσουμε.

Μου έκανε εντύπωση πάντως αυτό που μου είπε η Μπλαιρ για την αύρα. Πώς και δεν μου το είχε αναφέρει ο Βασι τόσο καιρό; Λες να το θεωρούσε και αυτός αυτονόητο; Ούτε τηλέφωνο δεν με πήρε μετά την αξιολόγηση. Λες να μην πέρασε; Λες να μην θέλει να μου μιλήσει; Δεν νομίζω, αφού αυτός μου έλεγε να τον παίρνω τηλέφωνο όταν μπούμε στην Ειδική Υπηρεσία για να μην χαθούμε. Κάτσε να τον πάρω μία εγώ. Προσπάθησα να τον πάρω αλλά δεν απάντησε.

«Δεν απαντάει... κρίμα...»

«Ποιος Σπαρκ;» με ρώτησε ο Ρομπ.

«Ένας φίλος, δεν τον ξέρετε.»

«Α εντάξει» είπε και συνέχισε να οδηγεί λες και δεν είχε συμβεί τίποτα.

Κοίταξα την Μπλαιρ που καθόταν στην πίσω αριστερή πλευρά του αυτοκινήτου, ενώ εγώ καθόμουν μπροστά και δεξιά. Από τη στιγμή που μπήκε στο αυτοκίνητο κοιτούσε έξω από το παράθυρο και δεν έβγαζε άχνα. Θυμήθηκα πάλι την “αύρα”, όμως κάτι δεν μου κόλλαγε. Πώς γίνεται να μπορούσα να διαβάσω τις σκέψεις του εφήβου ενώ κρατούσα τις χειροπέδες; Δεν θα έπρεπε η μαγεία μου να έχει αχρηστευθεί; Ακόμη, πώς μπόρεσα να διαβάσω τις σκέψεις της Μπλαιρ όταν μου έπιασε το χέρι; Πολύ περίεργο.

«Μπλαιρ, να σε ρωτήσω κάτι;»

«Όχι» μου απάντησε ξερά.

«Ε...ντάξει δεν πειράζει...» είπα και γύρισα πάλι το κεφάλι μου μπροστά. 

Δεν χρειάζεται να γεμίζω το κεφάλι μου με περίεργες σκέψεις τώρα, η δουλειά μου δεν έχει τελειώσει. Όταν γυρίσω στην Ειδική Υπηρεσία πρέπει να ανακρίνω τους 17 υπόπτους... Ενόχους;... Δεν ξέρω πως να τους αποκαλέσω. Τέλος πάντων, πρέπει να τους ανακρίνουμε με τον Μαλχη για να βρούμε τον “αρχηπρομηθευτή”, όπως τον αποκάλεσε πριν ο Ρομπ. Μέχρι τότε ας χαλαρώσω λίγο, αξίζω μία σύντομη ανάπαυση.

\chapter{}

«Ας ξεκινήσουμε λοιπόν Σπαρκ, συγγνώμη που άργησα. Είχα μία δουλειά στο κέντρο της πόλης που χρειαζόταν να τελειώσει σήμερα» είπε ο Μαλχη και άπλωσε το σακάκι του στον καλόγερο που βρισκόταν στην γωνία του δωματίου ανάκρισης.

«Να σε ρωτήσω κάτι Μαλχη πριν αρχίσουμε;»

«Παρακαλώ» είπε και γύρισε να με κοιτάξει. 

Παρατήρησα το περιποιημένο του μουστάκι και σκέφτηκα αυτό που συζητούσαμε με τον Ρομπ και τον Φυτ στο μπαρ. Σίγουρα ο Μαλχη ήταν πάνω από 40 χρονών. Βέβαια, το σώμα του ήταν πολύ γεροδεμένο και δεν μπορούσα να διακρίνω ούτε μία ρυτίδα στο πρόσωπο που. Μάλλον σε αυτό βασίζεται ο Φυτ στην εκτίμηση του ότι είναι εικοσάρης. Εκείνη την στιγμή δεν μπόρεσα να τον ρωτήσω αυτή την ασήμαντη ερώτηση, οπότε προσπάθησα να σκεφτώ κάτι άλλο να ρωτήσω για να μην φανώ περίεργος.

«Εσύ... ως αρχηγός της Ειδικής Υπηρεσίας... τι ακριβώς κάνεις;»

«Λίγο απ’ όλα»

«Δηλαδή;»

«Πάω σε αποστολές όπου χρειάζεται, καταστρώνω σχέδια, οργανώνω τις δουλειές των ομάδων, αν και αυτό τώρα τελευταία το έχω αφήσει περισσότερο στην γραμματέα μας. Τι άλλο; Α, κατασκοπεύω, ανακρίνω... Ναι αυτά. Γιατί, ρωτάς;»

«Τίποτα... απλά η έλλειψη ειδικότητας σου μου κίνησε την περιέργεια... Αλλά για να κάνεις τόσα πολλά πράγματα θα πρέπει να έχεις και μία... “καλή” μαγεία, έτσι δεν είναι;»

«Ω είναι καλή, τουλάχιστον για τα δικά μου γούστα. Έχω μαγεία Όπλων.»

«Μαγεία... Όπλων; Δηλαδή;»

«Μπορώ να δημιουργώ όπλα, τόσο απλό»

«Αυτό δεν μπορεί να το κάνει οποιοσδήποτε;» τον ρώτησα από περιέργεια, γιατί ήθελα να εξηγήσει μου λίγο παραπάνω την μαγεία του.

«Πέτα μου το στυλό από την τσέπη του σακακιού σου. Έχουμε πολλά, έτσι κι αλλιώς, θα σου δώσω άλλο αν χρειάζεται» μου είπε τεντώνοντας το χέρι του έτοιμος να πιάσει ό,τι του ρίξω.

Έχωσα το χέρι μου στην τσέπη μου και του πέταξα το στυλό μου. Στον αέρα, είδα το στυλό μου να αλλάζει πραγματικά μαγικά εμφάνιση και να γίνεται πιστόλι πριν πέσει στα χέρια του Μαλχη.

«Κατάλαβες τώρα;» μου είπε ο Μαλχη πιάνοντας το πιστόλι και γυρνώντας το λες και ήταν καουμπόης. «Καλά δεν χρειαζόταν να το κάνω στον αέρα, ίσως έκανα λίγο επίδειξη, αλλά γενικά μπορώ να μετατρέψω ότι αντικείμενο θέλω σε όπλο. Σπαθί, ασπίδα, πιστόλι, σφαίρες, ό,τι θες, ακόμη και χειροπέδες ξέρω γω που δεν θεωρούνται “τυπικά όπλα”. Όλοι οι πράκτορες της Ειδικής Υπηρεσίας βασίζονται σε μένα για τα πυρομαχικά και τα εργαλεία τους, για αυτό και δεν πρόκειται να με διώξουν ποτέ. Καλά δεν μπορεί και κάποιος να με διώξει, αφού είμαι ο αρχηγός, αλλά κατάλαβες.»

«Ναι...» του απάντησα.

«Κράτα το αυτό» μου είπε και μου πέταξε πίσω το πιστόλι. «Πάρτο σπίτι σου, ξέρω γω, ως εφεδρικό. Ποτέ δεν ξέρεις που θα χρειαστεί»

«Ευ...χαριστώ» είπα και έβαλα το πιστόλι στην αριστερή μεριά της ζώνης μου.

«Ωραία, ας ξεκινήσουμε τώρα. Πάω να φέρω τον πρώτο» σηκώθηκε ο Μαλχη από την καρέκλα του και βγήκε από το δωμάτιο ανάκρισης.

Αυτό το δωμάτιο δεν είχε κάτι το ιδιαίτερο, ήταν απλά ένα μικρό δωμάτιο με ένα τραπέζι στο κέντρο, με δύο καρέκλες, αντικριστά η μία με την άλλη, και έναν καλόγερο για να κρεμάμε τα πράγματα μας στην γωνία του δωματίου. Εγώ δεν είχα βγάλει το σακάκι μου γιατί ήθελα κάτι να κάνω με τα χέρια μου, οπότε τα έβαζα στις τσέπες του σακακιού. Ο Μαλχη επιστρέφοντας στο δωμάτιο είχε μαζί του τον προτελευταίο άντρα που συλλάβαμε την Κυριακή, τον “φίλο” μου που είχα κεράσει μπύρα.

«Είναι επειδή σου είπα ότι είμαι έμπορος φαρμάκων;» με ρώτησε ο άντρας.

«Όχι, αυτό το ήξερα από πριν» του είπα και του χαμογέλασα. «Θα μπω κατευθείαν στο ψητό, γιατί έχουμε άλλους 16 που περιμένουν. Από που προμηθεύεστε;»

«Σιγά μη σου πω» μου απάντησε ο άντρας καθώς ο Μαλχη τον έβαλε να καθίσει στην καρέκλα του ανακρινόμενου.

«Ω θα μας πεις, είτε το θες, είτε όχι...» είπε ο Μαλχη, προσπαθώντας να το παίξει κακός αστυνομικός. «...για αυτό ξεστόμισε τα τώρα και θα σε αφήσουμε να φύγεις»

«Γιατί αν δεν μιλήσω τι θα μου κάνετε;»

«Σπαρκ, ξέρεις τι να κάνεις» μου είπε ο Μαλχη και υπονόησε μάλλον να του διαβάσω τις σκέψεις.

Είδα τον άντρα να αγχώνεται λίγο όταν ο Μαλχη ξεστόμισε αυτές τις λέξεις, ίσως γιατί νόμιζε ότι θα τον πείθαμε να μιλήσει χρησιμοποιώντας βια. Η έκφραση του όμως χαλάρωσε όταν με είδε να κάθομαι ακίνητος στη θέση μου και να τον κοιτάζω επίμονα.

\textasciitilde ...γιατί με κοιτάει έτσι αυτός; Μήπως επειδή δεν ήπια την μπύρα που με κέρασε; Σταμάτα, είναι άβολο...\textasciitilde 

«Πες μας από που προμηθεύεσαι»

\textasciitilde Τι ρωτάει; Νομίζει πως επαναλαμβάνοντας την πρόταση θα με κάνει να ομολογήσω; Που να ήξερε πως η πηγή μας δεν είναι καν στην πόλη\textasciitilde 

«Δεν είναι στην πόλη, ε;»

«Τι; Πως το συμπέρανες αυτό» με ρώτησε ο άντρας αγχωμένος. \textasciitilde ...μη μου πεις πως ξέρουν ήδη και για τον Γιν και απλά μας φέραν όλους εμάς εδώ για να του στήσουν παγίδα;\textasciitilde 

«Σου θυμίζει κάτι το όνομα... Γιν; Ίσως αυτός είναι ο αρχηπρομηθευτής σας;»

«Γιν; ‘Οχι δεν μου λέει κάτι αυτό το όνομα» \textasciitilde ...σκατά. Τον ξέρουν ήδη, άρα όντως είναι παγίδα. Έπρεπε να το είχα σκεφτεί όταν μας είδα και τους 17 στα κελιά. Αν πιάναν μόνο τρεις ή τέσσερις θα έλεγα πως δεν γνωρίζουν κάτι, αλλά το γεγονός ότι μας τσάκωσαν όλους ήταν ύποπτο από την αρχή.\textasciitilde 

«Κάτι μου λέει πως δεν λες την αλήθεια κύριε. Μπορείς να μου εξηγήσεις πώς ξέρεις αυτόν τον Γιν;» του είπα και συνέχισα να τον κοιτάω επίμονα διαβάζοντας τις σκέψεις του.

\textasciitilde Μπορώ να ομολογήσω τώρα, αλλά να χάσω την δουλειά μου ή να παραμείνω πιστός και να...\textasciitilde 

«Μην είσαι τόσο πιστός στους ανθρώπους από πάνω σου. Αν ομολογήσεις τώρα δεν θα σε βάλουμε φυλακή, μόνο ένα μικρό πρόστιμο θα φας» του είπα για να του δώσω μία ψευδαίσθηση ασφάλειας.

«Και τι θα γίνει στους ναρκομανείς που προμηθεύουμε εμείς με την προμήθεια του Γιν; Θα τους αφήσετε να πεθάνουν; Το στερητικό σύνδρομο θα αρχίσει να εξαφανίζει ανθρώπους από δω και απο κει και εσείς μετά θα ψάχνετε να βρείτε αιτία για τους θανάτους τους σαν χαζοί» μου είπε ο άντρας εκνευρισμένος.

«Μην αγχώνεσαι. Έχουμε ήδη αρχίσει να προσφέρουμε σε κάποιους από τους πελάτες αναρωτικά προγράμματα. Το θέμα μας τώρα είναι άλλο...» είπε ο Μαλχη και έσκυψε δίπλα στον αναρινόμενο. «...δουλεύεις για τον Γιν. Ή τουλάχιστον συνεργάζεστε, απ’ ό,τι κατάλαβα. Μήπως ξέρεις και που μπορούμε να τον βρούμε;»

«Σιγά μη σου πω»\textasciitilde ...που να ξέρω κιόλας. Πάντα ο μικρός μας έφερνε την προμήθεια. Δεν ξέρω από που ή πώς έφτανε στα χέρια του, αλλά όταν τον είχα ρωτήσει είχε πει πως ο Γιν είναι εκτός της μεγάλης πόλης. Έχω περιέργεια όμως να δω πως οι κύριοι ανακριτές θα με κάνουν να ξεστομίσω που βρίσκεται ο Γιν χωρίς να το ξέρω καν\textasciitilde 

«Εντάξει, ας μην τον κουράζουμε άλλο» είπα και έκανα νεύμα στον Μαλχη να ηρεμήσει. «Δεν φαίνεται να ξέρει κάτι άλλο. Εξάλλου έχουμε άλλους 16 που περιμένουν, σίγουρα κάποιος από αυτούς θα ξέρει παραπάνω πληροφορίες»

«Πολύ ωραία. Επόμενος!» φώναξε ο Μαλχη ενώ σήκωνε τον ανακρινόμενο για να τον συνοδεύσει έξω από το δωμάτιο, λες και δεν θα πήγαινε έτσι κι αλλιώς να πάρει ο ίδιος τον επόμενο.

«Να σας πω κάτι τελευταίο;» είπε ο άντρας.

«Α τελικά ξέρεις;» είπε ο Μαλχη κρατώντας με το ένα του χέρι τις χειροπέδες του άντρα και με το άλλο τον ώμο του.

«Όχι... αλλά αν ήμουν εσείς, θα πρόσεχα λίγο παραπάνω τον μικρό με την αφάνα...»

«Και γιατί αυτό;» τον ρώτησε ο Μαλχη σπρώχνοντας τον λίγο.

«Τίποτα, απλά φαίνεται να ξέρει περισσότερα από εμάς τους υπόλοιπους. Είναι και από τους λίγους κλάσης Α, απ’ ό,τι ξέρω, οπότε θέλει ιδιαίτερη προσοχή»

«Μην αγχώνεσαι, δεν πρόκειται να μπορέσει να χρησιμοποιήσει την μαγεία του για να ξεφύγει. Έχουμε δικά μας κόλπα για τέτοιου είδους εγκληματίες» είπε ο Μαλχη καθώς συνόδευε τον άντρα έξω από το δωμάτιο, προφανώς υπονοώντας την μαγεία της Μπλαιρ στις χειροπέδες που φορούσαν όλοι οι προμηθευτές.\\\\

Οι επόμενες 15 ανακρίσεις πήγαν όπως η πρώτη. Διάβαζα τις σκέψεις των υπόπτων, αυτοί σκέφτονταν αυτόν τον Γιν, ανέφεραν το στερητικό σύνδρομο, εμείς τους λέγαμε για τα αναρρωτικά προγράμματα και στο τέλος μας έλεγαν να προσέξουμε τον μικρό με την αφάνα. Γιατί όμως; Μήπως μας είχαν στήσει κάποια παγίδα; Όταν 15 ανακρίσεις είναι τόσο όμοιες όμως κάτι πρέπει να σχεδιάζουν. Μήπως είχαν προσυμφωνήσει αν τους συμβεί ποτέ κάτι όπως τώρα να τα ρίξουν όλα στον μικρό; Αλλά οι περισσότεροι λέγαν, όταν διάβαζα τις σκέψεις τους, πως ο μικρός είναι ο μόνος που πιθανότατα ξέρει που και πώς θα βρούμε τον αρχηπρομηθευτή, οπότε μπορεί να μην είναι μπλόφα και να είναι όλα στο μυαλό μου. 

«Σπαρκ! Έλα γρήγορα να δεις κάτι» άκουσα τον Μαλχη να φωνάζει απ’ έξω από το δωμάτιο ανάκρισης. 

Είχε πάει να φέρει τον επόμενο ύποπτο για ανάκριση, αλλά για να με φωνάζει κάτι δεν θα πήγαινε με βάση το σχέδιο. Είχαμε οργανώσει τους υπόπτους σε τέσσερα κελιά. Στα τρία πρώτα, είχαμε πεντάδες από προμηθευτές ναρκωτικών οι οποίοι ήταν κλάσης Β, τους οποίους είχαμε συμφωνήσει να ανακρίνουμε πρώτους. Στο τελευταίο κελί, είχαμε τον έφηβο με την αφάνα και έναν άλλον άντρα, οι οποίοι ήταν και οι δύο κλάσης Α. Λες να απέδρασαν; Να συνεργάστηκαν για να αποφύγουν την ανάκριση;

Όταν έφτασα στο τελευταίο κελί, είδα τον έφηβο με την αφάνα να κάθεται στο παγκάκι του κελιού και τον άλλο άντρα ξαπλωμένο ανάσκελα στο πάτωμα. Από πάνω του ήταν η Ιφι η νοσοκόμα και έβγαζε πάλι αυτό το περίεργο κίτρινο φως από τα χέρια της. Το θέαμα ήταν όμοιο με την εικόνα που είδα πριν κάποιες μέρες στον χώρο αναμονής πριν την αξιολόγηση μου, απλά αντί αυτή τη φορά να ξέρω τον τραυματισμένο, μου ήταν τελείως άγνωστος, οπότε δεν μου προκάλεσε κάποιο συναίσθημα αυτή η εικόνα.

«Ζει;» ρώτησε ο Μαλχη αδιάκριτα.

«Όχι, δεν ξέρω γιατί προσπαθώ ακόμη.» του απάντησε η Ιφι σβήνοντας το φως από τα χέρια της. «Α... γεια σου Σπαρκ...» μου είπε και έσκυψε το κεφάλι.

«Τι εννοείς δεν ζει; Πώς;» είπε εκνευρισμένος ο Μαλχη.

«Δεν ξέρω. Μπορεί να έπαθε κάποιο καρδιακό επεισόδιο από το άγχος του. Πάντως η μαγεία μου δεν έχει επίδραση» του απάντησε η Ιφι με ήπιο τόνο.

«Εσύ το έκανες» είπε ο Βασι και έδειξε τον έφηβο που καθόταν ήσυχα στην γωνία του. Αυτός δεν απάντησε.

\textasciitilde ...έγιναν όλα τόσο ξαφνικά, δεν μπορώ να το εξηγήσω. Ο Γιν φταίει για όλα. Αυτό το δήθεν “κόλπο” που μας έμαθε του στέρησε την ζωή και τώρα εγώ θα κατηγορηθώ χωρίς λόγο.\textasciitilde 

«Σταμάτα να κατηγορείς χωρίς λόγο τον μικρό και πάμε πίσω να τον ανακρίνουμε» είπα στον Μαλχη και άρχισα να πηγαίνω πίσω στο δωμάτιο ανάκρισης.

«Όχι, βαρέθηκα το μουντό δωματιάκι. Ας τον ανακρίνουμε εδώ, δεν έχει κάποια διαφορά» μου είπε ο Μαλχη και άρχισε να περπατάει προς το μέρος του μικρού με την αφάνα. «Ιφι αν μπορείς να μας κάνεις τη χάρη και να φύγεις θα το εκτιμούσα πάρα πολύ»

«Μάλιστα...» είπε η Ιφι και άρχισε να τρέχει για να απομακρυνθεί όσο πιο γρήγορα μπορούσε.

«Λοιπόν. Ας ξεχάσουμε τώρα τον άλλον που έπεσε ξερός στο πάτωμα. Όλοι οι συνάδελφοι σου μιλάνε για έναν Γιν. Μήπως σου θυμίζει εσένα κάτι αυτό το όνομα;» είπε ο Μαλχη και έβαλε το πρόσωπο του μερικά εκατοστά μπροστά από το πρόσωπο του εφήβου.

Αυτός δεν απάντησε.

«Γιν. Δεν ξέρεις τίποτα θες να μου πεις.»

Ο μικρός δεν απάντησε.

«Σπαρκ τι κάνεις, τελείωνε»

«Μάλιστα» είπα και άρχισα να διαβάζω τις σκέψεις του εφήβου.

\textasciitilde ...όσο κάθεσαι τόσο κοντά δεν πρόκειται να σου πω τίποτα παλιο μουστακαλή. Άσε να πάρω μία ανάσα.\textasciitilde 

«Μαλχη κάνε λίγο πίσω» του είπα τραβόντας τον από τον ώμο.

\textasciitilde ...πάλι καλά που ο άλλος έχει αίσθηση του χώρου. Δεν το πιστεύω πάντως πως όλοι τους είπαν για τον Γιν. Πώς γίνεται αυτό; Τι είδους μεθόδους ανάκρισης έχουν αυτοί οι δύο;\textasciitilde 

«Λοιπόν θα μας πεις κάτι; Γιατί όλοι οι συνάδελφοι σου εσένα κατηγορούν. Μίλα λοιπόν, αφού ξέρουμε ότι εσύ έχεις τις πληροφορίες που χρειαζόμαστε» είπε ο Μαλχη και σταύρωσε τα χέρια του.

\textasciitilde Κοίτα τα παλιόπαιδα. Δεν φτάνει που τους μοίραζα την προμήθεια και δεν την κράταγα όλη για τον εαυτό μου, με κατηγορούν κιόλας. Αν μάθουν οι ανακριτές ότι οδηγάω κιόλας χωρίς δίπλωμα θα πάω σίγουρα μέσα...\textasciitilde 

«Αφού δεν μιλάει κανείς, θα πάρω εγώ τον λόγο» είπα και έκανα ένα βήμα προς τον έφηβο. «Μάθαμε πως ο Γιν δεν βρίσκεται στην πόλη. Μήπως ξέρεις που βρίσκεται;»

«Που θες να ξέρω. Εγώ απλά τον βλέπω μία φορά τον μήνα για να πάρω την προμήθεια και να την μοιράσω στους υπόλοιπους» \textasciitilde ...δεν λέω και ψέματα. Μία φορά τον μήνα τον βλέπω, απλά όχι στην πόλη, οπότε δεν μπορεί να με κατηγορήσει τώρα.\textasciitilde 

«Και που ακριβώς τον βλέπεις μία φορά τον μήνα;»

«Στην... πόλη»

«Δεν μου ακούστηκε και πολύ πειστική η απάντηση του Σπαρκ» πετάχτηκε ο Μαλχη.

«Επειδή δεν είναι αλήθεια. Μπορείς να μου πεις τώρα που ακριβώς βρίσκεται ο Γιν και πως προμηθεύεσαι από αυτόν» είπα στον έφηβο, σκύβοντας για να μπορέσω να τον κοιτάξω στα μάτια.

\textasciitilde ...και να μάθαινες που βρίσκεται ο Γιν θα σου έπαιρνε 3 ώρες ταξίδι, πολύ κουραστικό κάθε φορά το ταξίδι για το χωριό του...\textasciitilde 

«Α ναι θυμήθηκα, ένας συνάδελφος σου μας είπε πως ο Γιν είναι σε ένα χωριό περίπου 3 ώρες μακριά από την πόλη» του είπα, προφανώς ψέματα, για τον πείσω να πει αυτά που σκεφτόταν.

«Ποιος το είπε αυτό; Αφού μόνο εγώ πάω μέχρι εκεί για να...»

«Άρα είπες ψέματα» του είπα χαμογελώντας. «Για πες μας τώρα πως ακριβώς θα πάμε σε αυτό το χωριό.»

\textasciitilde ...τα έκανα μαντάρα. Τουλάχιστον δεν τους είπα και τι μαγεία έχει ο Γιν...\textasciitilde 

«Θα μας πεις μικρέ μη βρίσω το σπανάκι σου;» είπε ο Μαλχη προφανώς εκνευρισμένος. «Δεν έχουμε και άπειρο χρόνο. Έχουν περάσει μεσάνυχτα και ακόμη δουλεύουμε.»

\textasciitilde ...εντάξει μη με πλακώσεις κιόλας κύριε σκλυρέ ανακριτή...\textasciitilde «Θα πάρετε την εθνική οδό που πάει βόρεια μέχρι το Α34, όπου θα βγείτε. Εκεί γύρω είναι ένα εγκαταλελειμμένο χωριό που φαίνεται και από τον αυτοκινητόδρομο. Εκεί είναι ο Γιν. Εντάξει; Να φύγω τώρα;»

«Όχι. Τουλάχιστον όχι μέχρι να μας πεις τι μαγεία έχει ο Γιν, αφού μόνο εσύ πρέπει να ξέρεις» είπε ο Μαλχη και ήρθε δίπλα μου χωρίς να σκύψει. Εγώ σηκώθηκα για να φτάσω στο επίπεδο του.

\textasciitilde χεχε που να ήξεραν ότι ο Γιν έχει δύο μαγείες τώρα και όχι μία. Θα τους πω την μία όμως για να τους αιφνιδιάσει...\textasciitilde 

«Δύο μαγείες;» μου βγήκε αυθόρμητα

«Τι; Πώς το ή...ξερες;»\textasciitilde ...κάτι πάει λάθος...\textasciitilde 

«Δύο μαγείες;» είπε και ο Μαλχη, όχι τόσο προβληματισμένος, αλλά εκνευρισμένος.

«Ναι ρε, δύο μαγείες. Μαγεία Φύσης και... μαγεία Γης. Χρησιμοποιεί την μαγεία του για να φτιάξει την προμήθεια μας και μετά μας την πουλάει φθηνά. Για αυτό είμαστε τόσο μεγάλη μπιζνα. Αυτό θέλατε να μάθετε; Να φύγω τώρα;»

«Πάμε Σπαρκ» είπε ο Μαλχη καθώς έβγαινε από το κελί και παίζοντας με το μουστάκι του λες και σκεφτόταν κάτι. Τον ακολούθησα και βγήκαμε από το κελί.

«Και μένα θα με αφήσετε εδώ;» ρώτησε εκνευρισμένος ο έφηβος με την αφάνα.

«Για τώρα, ναι. Είσαι και κλάσης Α, οπότε δεν γίνεται να σε αφήσουμε να τριγυρίζεις χωρίς να έχει κλείσει αυτή η υπόθεση...» είπε ο Μαλχη και έκλεισε την πόρτα του κελιού. «Ειδοποίησε και τους υπόλοιπους» συνέχισε «Αύριο 12:00 θα πάμε οι έξι μας στο χωριό του αυτού του Γιν»

\chapter{}

«Αδελφέ, κατέβασε τα πόδια σου. Δεν βλέπω τον καθρέφτη»

«Σιγά μωρέ. Τι τον χρειάζεσαι τώρα τον δεξί καθρέφτη έτσι κι αλλιώς.»

«Για να κάνω προσπεράσεις;»

«Αφού από για να προσπεράσεις πρέπει να πας στην αριστερή λωρίδα. Που κολλάει ο δεξιός καθρέφτης;»

«Πώς θα μπω μετά την προσπέραση πάλι στην μεσαία λωρίδα; Πραγματικά δεν έχεις οδηγήσει ποτέ αμάξι;»

«Α ναι... ουπς...»

«Επίσης χτυπάει το τηλέφωνο σου αρκετή ώρα τώρα. Σήκωσε το γιατί πρέπει να είναι κάτι σημαντικό για να επιμένει τόσο όποιος σε παίρνει.»

«Ωχ όντως, δεν το είχα πάρει πρέφα. Παρακαλώ! Α, εσύ είσαι κορίτσι μου; Τι κάνεις, πώς είσαι; Καιρό έχουμε να τα πούμε... Ναι ρε πλάκα κάνω, τι θες;... Θα σε βάλουν στην ομάδα μου σε δύο εβδομάδες; Υπέροχα νέα, τώρα δεν χρειάζεται να ανησυχήσεις ποτέ ξανά τώρα που θα είμαστε μαζί και στη δουλειά. Όταν συναντηθούμε θα σε συγχαρώ για την πρόοδο σου. Συνέχισε την προσπάθεια που έχεις κάνει μέχρι τώρα και όλα θα... Ναι ναι και για τον άλλο έμαθα, αλλά πολύ περίεργος μου ακούγεται. Δεν νομίζω να υπάρχει κάποιος που να έχει τέτοιες ικανότητες, αλλά τι να πω. Αν υπάρχει όντως θα μας κάνει την ζωή πιο εύκολη... Εντάξει ναι καταλαβαίνω... Τι;... Καλά δεν πειράζει θα μου τα πεις με λεπτομέρειες όταν βρεθούμε από κοντά. Αντίο τώρα... Έλα αντίο»

«Κορίτσι μου; Δεν είσαι χωρισμένος;»

«Υπάρχουν πολλά πράγματα που δεν ξέρεις για μένα αδελφέ»

«Όπως, ποια σε πήρε τηλέφωνο τώρα ας πούμε.»

«Ο Ρομπ θα την θυμάται σίγουρα, γιατί την είχε σαν μικρή αδελφή»

«Πω όχι αυτή. Δεν την αντέχω μπαμπά, μου έσπαγε τα νεύρα.»

«Λίγα τα λόγια σου για αυτή την κοπέλα Ρομπ. Δεν ξέρεις τι έχει περάσει.»

«Τι έχει περάσει; Αν επιτρέπεται...»

«Θες να μάθεις αδελφέ;»

«Ναι ρε Χερμπ, δεν γίνεται να μου κινείς την περιέργεια ενώ οδηγώ και να με αφήνεις στον αέρα. Και κατέβασε επιτέλους τα πόδια σου»

«Ντάξει μη με φας κιόλας. Η κοπέλα αυτή είναι κόρη της γυναίκας μου... αλλά όχι δική μου. Η γυναίκα μου, όπως ξέρεις, δεν ήταν πιστή σε μένα. Την είχα πιάσει με δύο άλλους άντρες όπως ξέρεις στο σπίτι μου και αυτός ήταν και ο λόγος που χωρίσαμε.»

«Γιατί μου λες πράγματα που ήδη ξέρω;»

«Για τους μικρούς πίσω, μήπως ακολουθεί κάποιος την ιστορία. Τέλος πάντων, μερικές μέρες μετά από εκείνο το συμβάν έμαθα ότι είναι έγκυος. Το πρόβλημα ήταν πως αυτή δεν ήξερε ποιανού είναι το παιδί. Η ίδια δεν είχε καμία πρόθεση να το κρατήσει, το ίδιο βέβαια και οι δύο άλλοι άντρες. Από εγωισμό και μόνο πήρα πρωτοβουλία και μεγάλωσα το παιδί σαν δικό μου, σαν κόρη μου»

«Και πόσο καιρό γίνεται αυτό;»

«Πόσο χρονών είναι τα παιδιά σου; Νομίζω κάπου ενδιάμεσα από τον πρώτο και τον δεύτερο σου γιο έγινε το συμβάν, οπότε γύρω στα είκοσι χρόνια.»

«Και μου το λες τώρα;»

«Δεν με έχεις ρωτήσει και ποτέ...»

«Πώς να σε ρωτήσω; Αφού δεν το ήξερα καν»

«Ναι αλλά δεν έχεις έρθει ούτε μία φορά σπίτι μου τα τελευταία... είκοσι χρόνια.»

«Τι λες μωρέ, σίγουρα θα... θα...»

«Δίκιο έχει θείε. Η μόνη ανάμνηση που έχω με σένα σπίτι μας είναι όταν ήμουν τριών ή τεσσάρων και είχε μόλις γεννήσει η θεία το πρώτο σας παιδί»

«Δίκιο έχετε, εντάξει. Αλλά και πάλι, θα μπορούσες να μου το είχες πει. Πώς κατάφερες όμως να τα βγάλεις πέρα με την δουλειά σου και δύο παιδιά;»

«Όπως και εσύ»

«Ναι αλλά εγώ είχα την σύντροφο μου να με βοηθάει στις δουλειές του σπιτιού, αυτό εννοώ ρε Χερμπ.»

«Α ναι. Απλά δεν κοιμόμουν. Δηλαδή κοιμόμουν... απλά πολύ λίγο. ‘Οταν ανακάλυψα τον καφέ ειδικά κοιμόμουν ελάχιστα. Το πρωί προφανώς δούλευα. Το μεσημέρι... δούλευα. Όλη μέρα βασικά δούλευα. Το ωράριο μας αλλάζει ανάλογα με την μέρα αλλάζει, οπότε τις ώρες που ήμουν σπίτι έκανα τις κατάλληλες δουλειές έτσι ώστε να μην χρειαστεί τα παιδιά μου να  κάνουν κάτι»

«Ναι αλλά όταν έγινα 15 μπαμπα άρχισες να μου λες να κάνω εγώ τις δουλειές...»

«Προφανώς ρε Ρομπ. Προφανώς κάποτε θα χρειαζόμουν και την βοήθεια σας γιατί θα κατέρρεα. Η μικρή πάντως ήταν λίγο μπελάς, σε αντίθεση με τον Ρομπ. Όταν ήταν πολύ μικρή δεν βοηθούσε καθόλου στο σπίτι, ήθελε όλο βόλτες με τις φίλες της, ακριβά ρούχα, κοσμήματα, και ας μη μιλήσω για τα αγόρια της...»

«Δεν με ενδιαφέρει η προσωπική ζωή της μικρής σου τώρα. Θα σου κάνω μία ερώτηση μόνο. Της έχεις πει την αλήθεια;»

«Προφανώς. Αυτό θα σου έλεγα τώρα αλλα με διέκοψες... Πρέπει να ήταν γύρω στα 14 ή 15 όταν της είπα την αλήθεια. Από εκείνη τη στιγμή, ήταν λες και όλη της η προσωπικότητα. Ήταν σαν να άρχισε να με εκτιμάει πολύ περισσότερο. Ένιωθε μάλλον πως χρειαζόταν να μου ανταποδώσει αυτό που έκανα. Γυρνούσα κάθε μέρα από την δουλειά και είχε φτιάξει ήδη φαγητό, το σπίτι δεν ήταν ποτέ βρώμικο, τα ρούχα μας ήταν πάντα καθαρά και σιδερωμένα. Για κάποιο λόγο όμως παράτησε και το αγόρι της εκείνη τη περίοδο. Ποτέ δεν κατάλαβα γιατί. Ήταν, πάντως, λες και είχα μία προσωπική βοηθό, αλλά αυτή η βοηθός ήταν η “κόρη” μου, η οποία εκτός από όλα αυτά πήγαινε και διάβαζε παράλληλα για το σχολείο. Είναι πραγματικά ένας άγγελος...»

«Σιγά τον άγγελο...»

«Εντάξει ρε Ρομπ. Το μόνο αρνητικό της είναι η συμπεριφορά. Μπορεί να σου σπάει κάποιες φορές τα νεύρα, αλλά πρέπει να καταλάβεις πόσο σε έχει βοηθήσει αυτή η κοπέλα»

«Σιγά την βοήθεια»

«Να σου πω Ρομπ, θα σε πετάξω από το παράθυρο αν δεν συμμορφωθείς»

«Θέλω να σε δω να προσπαθείς...»

«Ήρεμα οι δυο σας. Τρέμει το αμάξι και δεν μπορώ να οδηγήσω»

«Σιγά ρε αδερφέ. Μία στροφή έκανα για να προειδοποιήσω τον Ρομπ»

«Δεν με νοιάζει»

«Καλά, ούτε εμένα με νοιάζει...»

«Κάτσε. Θείε Χερμπ... αυτή η κοπέλα μένει ακόμη μαζί σας;»

«Γιατί ρωτάς μικρέ;»

«Επειδή είπες έχεις καιρό να την δεις και μετά της είπες πως θα είστε μαζί “και” στην δουλειά, μαγική λέξη το “και”»

«Πάντα αυτό της λέω. Αφού με παίρνει κάθε τρεις ώρες τηλέφωνο, τι να κάνω. Καλά έχει τους λόγους της μάλλον. Αγχώνεται πάρα πολύ για μένα λόγω της δουλειάς μου επειδή με βλέπει κάθε μέρα να πέφτω στον καναπέ κουρασμένος με το που γυρίζω. Τώρα όμως που θα είναι μαζί μου ελπίζω να καταλάβει πως η δουλειά μας δεν είναι και τίποτα. Έτσι δεν είναι ρε Ρομπ;»

«Να, ξέρω γω. Είμαι σε τελείως άλλο τομέα από σένα μπαμπά, οπότε δεν μπορώ να συγκρίνω...»

«Καλά εντάξει, μη με φας κιόλας»

«Και αυτή η κοπέλα θείε, αν επιτρέπεται, πως λέγεται;»

«Την ονόμασα Μπλαιρ»

«Μπλαιρ;»

«Tι; Δεν σου αρέσει;»

«Τίποτα, απλά...»

«Απλά τι;»

«Τίποτα... θα συνεχίσω να διαβάζω...»

\chapter{}

«Σπαρκ, τι προτείνεις να κάνουμε όταν εντοπίσουμε αυτόν τον Γιν;» μου είπε ο Μαλχη που καθόταν στο πίσω κάθισμα και σημείωνε κάτι στο μπλοκάκι του.

«Τι να κάνουμε; Δεν... ξέρω...»

«Έλα ρε Σπαρκ. Πρέπει να μπορείς να φτιάχνεις και μόνος σου σχέδια. Δεν θα έχετε πάντα εμένα να σας οργανώνω σαν ομάδα»

«Και γιατί εγώ; Γιατί όχι ο Ρομπ για παράδειγμα;» είπα και έδειξα τον Ρομπ που οδηγούσε.

«Επειδή εσύ είσαι ο πιο έξυπνος στην ομάδα...»

«Πώς το κατέληξες σε αυτό το συμπέρασμα; Με ξέρεις μόνο τέσσερις μέρες...»

«Από την υπόλοιπη ομάδα σας, μόνο ο Ρομπ μπορεί ίσως να πλησιάσει τον τρόπο σκέψης σου» είπε ο Μαλχη και έκλεισε το μπλοκάκι του.

«Και σου ξαναλέω, με ξέρεις μόνο τέσσερις μέρες. Από που βγαίνουν αυτά τα συμπεράσματα; Που ξέρεις πως η Μπλαιρ δεν είναι κάποια κρυφή ιδιοφυία;»

«Δεν νομίζω, θα το είχα καταλάβει. Πες μου τώρα Σπαρκ, τι σχέδιο έχεις σκεφτεί;» είπε και σταύρωσε τα χέρια του ο Μαλχη περιμένοντας την απάντηση μου.

«Δεν... έχω σκεφτεί...»

«Σκέψου τώρα;»

«Πώς;»

«Τι “πώς”; Φτιάξε ένα σχέδιο. Πες μου για παράδειγμα ποιοι θα πάνε να να βρουν τον Γιν, ποιοι θα κάτσουν πίσω, ποιοι θα περιπολούν την περιοχή, τέτοια πράγματα. Οι λεπτομέρειες μετά συνήθως είναι αυτοσχεδιασμός»

Γύρισα και κοίταξα τον αυτοκινητόδρομο για να πάρω έμπνευση. Η δουλειά που μου είχε βάλει ο Μαλχη δεν ήταν και όσο εύκολη όσο το κάνει να ακούγεται, καθώς χρειαζόταν να σκεφτώ κάθε συνδυασμό ζευγαριών για κάθε πόστο, οι οποίοι είναι έξι επί πέντε, επί... πάρα πολλοί τέλος πάντων. Δεν έχω χρόνο να κάνω όλους τους συνδυασμούς στο μυαλό μου. Αρκεί απλά να βρω ποιοι ταιριάζουν μεταξύ τους και θα αυτοσχεδιάσουμε όπως είπε ο Μαλχη.

Ως ζευγάρια, σίγουρα πιστεύω εγώ με τον Ρομπ θα ήμασταν καλός συνδυασμός, αλλά αν ξεσπάσει κάποια μάχη δεν θα βρεθούμε σε κίνδυνο πολύ εύκολα εξαιτίας μου. Από τους άλλους τέσσερις, πιστεύω άλλο ένα καλό δίδυμο για να ψάξει τον Γιν θα ήταν η Μπλαιρ με τον Χερμπ. Βέβαια, πιστεύω πως στην ομάδα αναζήτησης του Γιν θα πρέπει να ανήκω και εγώ, για να μπορέσω να του διαβάσω τις σκέψεις όταν τον συναντήσουμε, μη μας έχει στήσει κάποιου είδους παγίδα. Αλλά αυτό χαλάει την ιδέα ότι θα είμαστε σε ζευγάρια. Δεν πειράζει όμως, τώρα που το σκέφτομαι, καθώς δεν χρειάζεται να μείνουν δύο άτομα πίσω στα αυτοκίνητα.

Ωραία, οπότε θα πάω εγώ με την Μπλαιρ και τον Χερμπ να ψάξουμε τον Γιν. Τώρα, από τους υπόλοιπους θα χρειαστεί σίγουρα ο Φυτ να είναι κάπου τριγύρω, σε περίπτωση που συμβεί κάτι απρόοπτο. Ο Ρομπ επίσης θα ήταν μία καλή ιδέα να περιπολεί την υπόλοιπη περιοχή για τυχόν άλλες παγίδες, αφού ως κατάσκοπος είναι μάλλον ο πιο παρατηρητικός. Οπότε, αυτό αφήνει τον Μαλχη να κάτσει πίσω στα αυτοκίνητα. Γιατί όχι, καλό μου φαίνεται το σχέδιο.

«Λοιπόν...»

«Επιτέλους...» με διέκοψε ο Μαλχη. «...για πες»

«Θα χωριστούμε σε τρεις άνισες ομάδες. Εγώ με την Μπλαιρ και τον Χερμπ θα αναζητήσουμε τον Γιν στο εγκαταλελειμμένο χωριό. Ο Ρομπ και ο Φυτ θα περιπολούν τριγύρω και θα είναι σαν δευτερεύον γραμμή άμυνας σε περίπτωση που συμβεί κάτι, ή σε περίπτωση που πέσουμε παγίδα»

«Χμ»

«Τι έγινε; Δεν σου αρέσει η ιδέα μου;»

«Όχι, απλά... εγώ τι θα κάνω;» με ρώτησε προβληματισμένος ο Μαλχη.

«Εσύ... θα κάτσεις πίσω να προσέχεις τα αυτοκίνητα. Δεν πιστεύω πως η μαγεία σου θα μας είναι πολύ χρήσιμη αν ξεσπάσει μάχη» του απάντησα για να του εξηγήσω τον συλλογισμό μου.

«Έτσι πιστεύεις; Εντάξει. Μη μου κλαίγεσαι όμως μετά αν ξεσπάσει μάχη και δεν μπορείτε να αποκλιμακώσετε την κατάσταση. Εξάλλου ο στόχος μας έχει δύο είδη μαγείας, αν το ξέχασες. Μπορεί να είναι πιο σκληρός απ’ ό,τι νομίζεις...» είπε ο Μαλχη και άνοιξε πάλι το μπλοκάκι του.

«Άρα θα ακολουθήσουμε το σχέδιο μου;» τον ρώτησα επειδή δεν κατάλαβα.

«Ναι ναι, κάνετε ό,τι θέλετε» μου απάντησε με ειρωνικό ύφος.

Λες να έχει δίκιο ο Μαλχη; Μπορεί να υποτιμάμε τον στόχο μας. Αλλά θα έχω μαζί μου δύο ειδικούς στις μάχες, οπότε τι μπορεί να πάει στραβά; Σε περίπτωση που ο Γιν προσπαθήσει να χρησιμοποιήσει τη μαγεία του για να ξεφύγει ή να μας τραυματίσει, η Μπλαιρ μπορεί να του την αχρηστεύσει και ο Χερμπ μπορεί να τον ελέγξει με τη δύναμη του και να του βάλει χειροπέδες. Αν πάλι μας ξεφύγει, ο Φυτ ή ο Ρομπ θα τον αναλάβουν, αφού θα περιπολούν στην γύρω περιοχή. Τι μπορεί να πάει στραβά;

«Α34» είπε ο Ρομπ όταν είδε την πινακίδα της εξόδου. «Φτάσαμε σχεδόν. Βλέπει κανείς το χωριό;»

«Να μάλλον εκεί» είπε ο Μαλχη μας έδειξε επτά σπιτάκια δύο από τα οποία μοιάζαν με θερμοκήπια, ενώ τα υπόλοιπα ήταν γκρεμισμένα. «Ας σταματήσουμε εκεί για να ξεκινήσουμε την αναζήτηση»

«Μάλιστα» είπε ο Ρομπ και μπήκε στον δρόμο που πήγαινε προς αυτά τα σπιτάκια.

Κοίταξα πίσω από περιέργεια να δω που βρίσκονταν οι υπόλοιποι, γιατί αγχώθηκα μήπως δεν μας είδαν που αλλάξαμε δρόμο. Μόλις είδα, όμως, το αμάξι τους αρκετά μέτρα πίσω να αντιγράφει την πορεία που έπαιρνε ο Ρομπ, μου έφυγε το άγχος. Ο Χερμπ είναι πραγματικά πολύ καλός οδηγός. Τις τρεις προηγούμενες μέρες οδηγούσε το βανάκι με τους υπόπτους που συλλαμβάναμε στο μπαρ. Σήμερα οδηγεί το δεύτερο μας αυτοκίνητο, καθώς είμαστε έξι αυτή τη φορά και όχι πέντε, για να μπορούμε να χωρέσουμε όλοι σε ένα αμάξι. Πραγματικά τι θα κάναμε χωρίς αυτόν. Ο Ρομπ πάλι δεν πάει πίσω. Μας έχει βοηθήσει και αυτός στο θέμα της μετακίνησης. Λες η Μπλαιρ οδηγεί; Σίγουρα οδηγεί, αφού την είχα δει να κρατάει κλειδιά αυτοκινήτου έξω από το καφενείο που είχα πάει με τον Βασι πριν την αξιολόγηση. Καλά, σιγά μην αφήσει ο Χερμπ την Μπλαιρ να οδηγήσει, θα τον πιάσει ο εγωισμός του. Η Μπλαιρ όμως δεν φαίνεται να παραπονιέται, τουλάχιστον όσο είναι με τον Χερμπ. Ποια να είναι η σχέση τους άραγε; Γιατί έτρεξε να τον αγκαλιάσει εκείνη την μέρα μετά την αξιολόγηση; Θα τους ρωτήσω ίσως κάποια άλλη στιγμή.

«Να παρκάρω εδώ;» ρώτησε ο Ρομπ ανάβοντας τα αλάρμ.

«Ναι πήγαινε δίπλα στο πρώτο σπίτι και άστο εκεί» είπε ο Μαλχη και άνοιξε το παράθυρο του. Βγήκε μισός έξω από το παράθυρο και έκανε νόημα στον Χερμπ να σταματήσει δίπλα στο πρώτο σπίτι.

Όταν κατεβήκαμε από τα αυτοκίνητα, εξήγησα και στους υπόλοιπους το σχέδιο μας.

«Θα είμαι με τον Χερμπ επιτέλους, τι ωραία» είπε η Μπλαιρ σαν μικρό παιδί.

«Άρα εγώ θα κάθομαι απλά σε μία απόσταση και θα περιμένω σινιάλο;» είπε ο Φυτ ο οποίος δεν είχε καταλάβει το πλάνο μας.

«Όχι... θα έλεγα να περιπλανιέσαι και να ψάχνεις για παγίδες μαζί με τον Ρομπ. Έτσι κι αλλιώς όλα τα σπίτια δίπλα είναι, αν σε χρειαστούμε θα μας ακούσεις» του απάντησα εξηγώντας λίγο παραπάνω τον συλλογισμό μου.

«Εντάξει κατάλαβα. Δεν νομίζω να υπάρχουν και πολλές παγίδες όμως. Μόνο σε εκείνα τα θερμοκήπια πιθανότατα θα υπάρχει κάποιος. Τα υπόλοιπα κτίρια φαίνεται να είναι γκρεμισμένα και εγκαταλελειμμένα. Δεν πιστεύω πως αξίζει να ψάξουμε»

«Ποτέ δεν ξέρεις» πετάχτηκε ο Μαλχη. «Κάνε αυτό που σου λέει ο Σπαρκ τώρα και θα βλέπουμε»

«Καλά εντάξει...» είπε ο Φυτ και χρησιμοποίησε τις ρίζες του σαν σκαλοπάτια για να ανέβει πάνω στο πρώτο από τα κτίρια. «Έλα Ρομπ, ανέβα» συνέχισε ο Φυτ και ο Ρομπ πήδηξε και χρησιμοποίησε την μαγεία τροχιάς του για να φτάσει στην ταράτσα του κτιρίου. Μα καλά, η μαγεία τροχιάς δεν επηρεάζεται από την βαρύτητα;

«Εγώ θα έλεγα να ξεκινούσαμε την αναζήτηση μας από τα θερμοκήπια» πρότεινε ο Χερμπ.

«Ου ναι, καλή ιδέα» απάντησε η Μπλαιρ και με κοίταξε. «Τι λες Σπαρκ;»

«Πηγαίνετε επιτέλους. Δεν θέλω να κάθομαι για πολύ ώρα εδώ» πετάχτηκε ο Μαλχη διώχνοντας μας με το χέρι του.

«Ωραία πάμε» είπα και προχώρησα προς τα θερμοκήπια πρώτος. Σύντομα με ακολούθησαν και η Μπλαιρ με τον Χερμπ.

Όταν μπήκαμε στο πρώτο θερμοκήπιο στον δρόμο μας, βλέπαμε πράσινα φυτά παντού. Η μυρωδιά θύμιζε την μυρωδιά των σακούλων που είχαν κατασχεθεί από τους εμπόρους ναρκωτικών που συλλάβαμε. Η Μπλαιρ έσκυψε να ακουμπήσει ένα φυτό, με το που το ακούμπησε όμως, το φυτό μαράθηκε και τα φύλλα έπεσαν στο πάτωμα.

«Αυτά τα φυτά είναι αποτέλεσμα μαγείας» είπε η Μπλαιρ.

«Πώς το ξέρεις;» την ρώτησα.

«Αφού κατέρρευσε με το που το ακούμπησα. Ελπίζω να μη ξέχασες ότι έχω αντι-μαγεία. Αν τα φυτά ήταν κανονικά, αυτό το φυτό δεν θα πέθαινε» μου είπε η Μπλαιρ και μου έδειξε τα φύλλα που είχαν πέσει στο πάτωμα.

«Άρα λες ο στόχος μας να είναι κοντά;»

«Δεν είπα αυτό. Η μαγεία μπορεί να λειτουργεί και σε μεγάλες αποστάσεις. Πώς νομίζεις ότι λειτουργούν οι χειροπέδες που τους δίνω την αύρα μου για να μην μπορούν οι κρατούμενοι να χρησιμοποιήσουν την μαγεία τους;»

«Ναι... δεν το σκέφτηκα...» είπα και έσκυψα το κεφάλι.

«Δεν σε κατηγορώ. Απλά...»

«Συγγνώμη φίλε, ποιοι είστε εσείς; Και γιατί καταστρέφετε τα φυτά μου;» άκουσα μία φωνή από πίσω μου. Γύρισα ξαφνιασμένος και είδα έναν κοντό νέο Ασιάτη, με δεμένα τα χέρια του πίσω από την πλάτη. Φορούσε ενδυμασία κηπουρού και ένα ψάθινο καπέλο. 

«Είμαστε...»

«Είμαστε νέοι στην περιοχή» με διέκοψε ο Χερμπ. «Ψάχναμε μέρος να μείνουμε γιατί το προηγούμενο μας χωριό κάηκε και δεν έχουμε πλέον μέρος να μείνουμε» είπε προφανώς ψέματα ο Χερμπ. «Είδαμε πως η πόρτα του θερμοκηπίου σας ήταν ανοιχτή και είπαμε να εισέλθουμε μήπως πετυχαίναμε κάποιον που να μας εξυπηρετήσει»

«Να σας εξυπηρετήσω πώς; Ένας απλός κηπουρός είμαι» του απάντησε ο Ασιάτης.

«Μήπως θα μπορούσατε να μας δώσετε ένα μέρος να μείνουμε για κάποιες μέρες, καθώς ψάχνουμε καινούργιο σπίτι ηλεκτρονικά;» τον ρώτησε ο Χερμπ.

«Όχι» του απάντησε ξερά ο Ασιάτης.

«Γιατί;»

«Δεν πείστηκα με την ιστορία σου. Κανένα χωριό δεν κάηκε πρόσφατα εδώ γύρω, οπότε μου λες ψέματα»

«Σπαρκ κάνε κάτι, μη κάθεσαι και κοιτάς σαν χαζός...» μου ψιθύρισε η Μπλαιρ.

«Αν θέλετε βέβαια, μπορείτε να μείνετε σε κάποιο από τα εγκαταλελειμμένα σπίτια. Εγώ πάντως δεν πρόκειται να σας φιλοξενήσω...» είπε ο Ασιάτης και μας γύρισε πλάτη. «Αν έχετε την καλοσύνη τώρα φύγεται από το θερμοκήπιο μου. Δεν θέλω να πεθάνουν άλλα φυτά» \textasciitilde ...γιατί θα πέσει το κέρδος μου.\textasciitilde 

«Μπορώ να σας κάνω μία ερώτηση;» είπα διακριτικά στον Ασιάτη.

«Όχι» \textasciitilde Τι θέλει τώρα αυτός; Δεν τους είπα να φύγουν;\textasciitilde 

«Θα την κάνω έτσι κι αλλιώς. Τα φυτά αυτά, γιατί έχουν τέτοια μυρωδιά;» του είπα και τον κοίταξα επίμονα για να τον αγχώσω.

\textasciitilde Γιατί με κοιτάει έτσι αυτός; Δεν έχει ξαναμυρίσει ναρκωτικά στη ζωή του; Από την άλλη φαίνεται μικρός, μπορεί να μην έχει μυρίσει.\textasciitilde 

«Μήπως είναι ναρκωτικά;» του είπα. Αυτός γύρισε και με κοίταξε.

\textasciitilde Άρα έχει μυρίσει. Οπότε γιατί με ρωτάει; Και γιατί δεν παραξενεύεται ο πατέρας του; Λες να μην είναι ο πατέρας του;\textasciitilde  «Ναι μαριχουάνα είναι. Αν θέλετε σας πουλάω λίγο» είπε ο Ασιάτης.

«Ου ναι, θα ήθελα λίγο» πετάχτηκε η Μπλαιρ και τον πλησίασε.

«Εσύ; Δεν το περίμενα να πω την αλήθεια. Ορίστε...» είπε ο Ασιάτης γελώντας και της έδωσε μία διάφανη σακούλα που είχε στην τσέπη του.

«Μόνο τόσο; Λίγο μου φαίνεται» είπε η Μπλαιρ όσο κοιτούσε την σακούλα.

«Κανονική ποσότητα είναι αυτή...» \textasciitilde ...σίγουρα έχει ξανακάνει μαριχουάνα αυτή; Δεν φαίνεται να γνωρίζει και τις ποσότητες στις οποίες πωλείται. Και να μη μιλήσω για τον τρόπο που κρατάει την σακούλα\textasciitilde 

«Πώς σας λένε παρεμπιπτόντως;» είπε η Μπλαιρ συνεχίζοντας να κοιτάει την σακούλα.

\textasciitilde ...να τους πω το κανονικό μου όνομα; Μου φαίνονται πολύ ύποπτοι...\textasciitilde 

«Είχαμε ακούσει για έναν Γιν που μένει εδώ γύρω. Μήπως τον ξέρετε κύριε κηπουρέ;» τον ρώτησε η Μπλαιρ πριν προλάβει ο Ασιάτης να απαντήσει την προηγούμενη της ερώτηση. Αυτός φάνηκε να αγχώνεται λίγο.

\textasciitilde Πώς ξέρουν για μένα; Λες να είναι πράκτορες της Ειδικής Υπηρεσίας; Τώρα που τους παρατηρώ λίγο παραπάνω, γιατί έχουν και οι τρεις όμοια σακάκια και όπλα στις ζώνες τους; Είμαι χαζός, έπρεπε να το είχα σκεφτεί από την πρώτη στιγμή\textasciitilde 

«Άρα τον ξέρετε, για να αγχώνεστε» είπε η Μπλαιρ και πέταξε κάτω την σακούλα. «Σπαρκ, για πες, τι κάνουμε;»

«Αυτός είναι μάλλον ο Γιν» της απάντησα.

\textasciitilde Τι λέει αυτός; Πώς ξέρει το όνομα μου; Λες να είναι παγίδα; Πρέπει να φύγω διακριτικά από εδώ\textasciitilde  «Συγγνώμη αλλά δεν έχω χρόνο για...»

«Ούτε εμείς “φίλε”. Οπότε για έλα μαζί μας...» είπε η Μπλαιρ και πήγε να τον ακουμπήσει έτσι ώστε να αχρηστεύσει την μαγεία του. Ο Γιν όμως την έριξε κάτω πριν προλάβει να τον πλησιάσει, τυλίγοντας ρίζες στα πόδια της. Ο Χερμπ έτρεξε να του βάλει χειροπέδες, σπρώχνοντας με στην προσπάθεια προς τα πίσω γιατί ήμουν στην μέση, αλλά αυτό που έγινε μετά δεν θα μπορούσε κανείς να το προβλέψει. Πριν ο Χερμπ φτάσει κοντά στον Ασιάτη, πριν προλάβει καν να ενεργοποιήσει την μαγεία του, μία τεράστια ρίζα μαγείας Φύσης του είχε διαπεράσει το στήθος. Αίματα έτρεχαν στο πάτωμα σαν καταρράκτης. Ο Χερμπ δεν μπορούσε να κουνηθεί. Ένιωσα δύο ρίζες να τυλίγουν και τα δικά μου πόδια. Προσπάθησα να τα κουνήσω, αλλά δεν μπορούσα, οι ρίζες με είχαν ακινητοποιήσει τελείως κάτω από τον αστράγαλο. Ο Γιν άρχισε να τρέχει προς την έξοδο του θερμοκηπίου.  Έβγαλα το όπλο μου και προσπάθησα να τον πυροβολήσω στα δικά του πόδια, για να του τα αχρηστεύσω, αλλά αστόχησα όλες μου τις σφαίρες.

Ευτυχώς για μένα, είδα το βέλος του Ρομπ να εμφανίζεται στα πόδια του Γιν, οπότε ανακουφίστηκα καθώς ο Ρομπ και ο Φυτ μας είχαν ακούσει μάλλον και έσπευσαν να βοηθήσουν. Η ανακούφιση μου αυξήθηκε όταν είδα το βέλος του Ρομπ να αλλάζει κατεύθυνση, τον Γιν να “σκοντάφτει”, τις ρίζες του Φυτ να περικυκλώνουν τον κοντό Ασιάτη και τον Μαλχη να του βάζει χειροπέδες. Τα ποδια μου ήταν πλέον ελεύθερα από τις ρίζες του Γιν. Όλα τα φυτά στο θερμοκήπιο ξαφνικά ξεράθηκαν. Σκέφτηκα ότι μπορεί να έχει να κάνει με την άντι-μαγεία της Μπλαιρ, αφού τώρα ο Γιν φορούσε χειροπέδες με την αύρα της.

Η ανακούφιση μου όμως δεν κράτησε για πολύ. Γύρισα προς το μέρος της Μπλαιρ και το θέαμα που είδα ήταν σοκαριστικό. Ένας άντρας με μία τεράστια τρύπα στο σώμα του με την διάμετρο της ρίζας που τον είχε χτυπήσει στο στήθος, να είναι ξαπλωμένος μπρούμυτα πάνω σε μία κόκκινη λίμνη από αίμα στην οποία επέπλεαν ξεραμένα φύλλα μαριχουάνας. Από δίπλα του βρισκόταν η Μπλαιρ γονατιστή, η οποία χτυπούσε τις γροθιές της και το κεφάλι στην κόκκινη λίμνη, λερώνοντας τον εαυτό της, αλλά και εμένα που την πλησίασα, με το αίμα του συναδέλφου μας. 

Όταν έφτασα δίπλα της και έσκυψα δίπλα της αυτή σταμάτησε να ξεσπάει στο πάτωμα και με κοίταξε. Τα μάτια της είχαν κοκκινήσει. Η θλίψη της θα μπορούσε να συγκριθεί με την δική μου όταν έχασα τον πατέρα μου. Την είδα να σηκώνεται με αίματα να τρέχουν από τα μαλλιά της, τις γροθιές της, τα γόνατα της, όλο της το σώμα.

«Εσύ...» μου είπε πριν προλάβω να σηκωθώ. «Εσύ... φταις για όλα...» ξεφώνησε με σοβαρό τόνο και με κλώτσησε. Εγώ έπεσα με την πλάτη στα ξεραμένα φυτά.

«Ξέρεις τι έχει κάνει αυτός ο άνθρωπος για μένα; Όχι βέβαια, που να ξέρεις» είπε ο Μπλαιρ πιάνοντας με από το κολάρο του σακακιού μου.

«Ηρέμησε Μπλαιρ, δεν έκανα κάτι εγώ. Ο Γιν...»

«Προφανώς και δεν μπορείς να καταλάβεις τον πόνο μου αυτή τη στιγμή. Αυτός ο άνθρωπος ήταν σαν πατέρας για μένα και τον έχασα μέσα σε κλάσματα δευτερολέπτων, αλλά που να ξέρεις εσύ πως είναι να χάνεις τον πατέρα σου...» μου είπε η Μπλαιρ και με πέταξε πάλι κάτω. Προσπάθησα να διαβάσω τις σκέψεις της, αλλά δεν σκεφτόταν κάτι εκείνη την ώρα. Το μόνο που έβγαζε προς τα έξω ήταν απογοήτευση. Δεν καταλαβαίνω γιατί έριξε το φταίξιμο σε μένα, εγώ δεν έκανα κάτι άμεσα. Μήπως επειδή εγώ είχα καταστρώσει το σχέδιο; Και πάλι όμως, δεν μπορούσα να προβλέψω ότι θα γινόταν κάτι τέτοιο. Είδα με την άκρη του ματιού μου τον Μαλχη και τον Ρομπ να πλησιάζουν τον Χερμπ.

«Κρίμα» είπε και γύρισε προς το μέρος μου. «Έλα Σπαρκ, σήκω, πάμε στα αυτοκίνητα. Μπλαιρ και συ το ίδιο».

Η Μπλαιρ δεν του απάντησε, κοιτούσε το πτώμα με απογοήτευση. Ο Ρομπ σήκωσε τον Χερμπ και τον πέταξε στον αέρα, για να μπορέσει να του αλλάξει την τροχιά και να τον αιωρήσει μέχρι το αυτοκίνητο, αντί να τον μουβαλάει. Είδα ένα δάκρυ να κυλάει στο μάγουλο του Ρομπ, αλλά η έκφραση του είχε παραμείνει σοβαρή. Η Μπλαιρ έμεινε όρθια στην κόκκινη λίμνη και κοιτούσε τα ξεραμένα φύλλα που επέπλεαν στο αίμα του Χερμπ.

«Σου είπα έλα» φώναξε ο Μάλχη και την τράβηξε προς το μέρος του.

«Άσε με, μπορώ να περπατήσω και μόνη μου...» είπε και άρχισε να περπατάει με γοργό βήμα προς τα αυτοκίνητα, παίζοντας με τα δάχτυλα της και μουρμουρώντας, όπως ακριβώς έκανε και εκείνη τη μέρα στον διάδρομο, μετά την αξιολόγηση μας.

«Έλα σήκω και συ επιτέλους. Πάμε πίσω γρήγορα να δούμε αν προλάβουν να κάνουν κάτι οι νοσοκόμες μας, αν και αμφιβάλω...» μου είπε ο Μαλχη και μου έδωσε το χέρι του για να με σηκώσει. Εγώ τον έπιασα και τραβήχτηκα για να σηκωθώ. «Βλέπεις γιατί σου είπα μη δένεσαι με τους συναδέλφους σου;»

«Ναι...» του απάντησα και έσκυψα το κεφάλι.

\chapter{}

«Δεν μπορώ να κάνω κάτι. Έχει περάσει πολύ ώρα από τον τραυματισμό του. Συγγνώμη, αλλά ο κύριος Χερμπ δεν θα μπορέσει να είναι πλέον μαζί μας» είπε η Ιφι και έκλεισε τα μάτια του Χερμπ που ήταν ξαπλωμένος στο κρεβάτι του ιατρείου. Οι υπόλοιποι τέσσερις ήμασταν σκόρπιοι μέσα και έξω από το ιατρείο, αλλά η πόρτα του δωματίου ήταν ανοιχτή, οπότε όλοι μας ακούσαμε τα λόγια της Ιφι που επιβεβαίωναν ότι ο Χερμπ ήταν νεκρός. Η Μπλαιρ δεν μπορούσε να με κοιτάξει καν στα μάτια. Δεν ξέρω γιατί έριχνε το φταίξιμο σε μένα, λες και σκότωσα εγώ τον Χερμπ.

«Είστε ελεύθεροι να φύγετε αν θέλετε κύριοι... και κυρίες» είπε η Ιφι ενώ μάζευε τα πράγματα της.

«Μπλαιρ, πραγματικά λυπάμαι...» προσπάθησα να της πω, αλλά αυτή με έδιωξε και έφυγε από το ιατρείο.

«Σπαρκ... άστην. Θέλει τον χρόνο της» μου είπε η Ιφι και έκλεισε την τσάντα της. «Αν μπορείτε μην αναφέρεται την αποστολή ή τον Χερμπ όταν είστε μαζί της για τις επόμενες μέρες»

«Εντάξει» είπα και βγήκα από το ιατρείο. 

Ο Ρομπ και ο Φυτ κάθονταν σε έναν καναπέ. Ο Ρομπ είχε γείρει τελείως μπροστά, είχε το αριστερό του χέρι να καλύπτει το στόμα του και με το δεξί του χέρι σημείωνε στο κινητό του κάτι. Ο Φυτ καθόταν χαλαρός στον καναπέ, στερεώνοντας το κεφάλι του στο δεξί του χέρι και φαινόταν λες και κάτι σκέφτεται. Πραγματικά όμως, ο Ρομπ τι γράφει συνέχεια στο κινητό του;

«Ρομπ»

«Ναι;» μου απάντησε χωρίς να αλλάξει πόζα.

«Τι κάνεις συνέχεια στο κινητό σου, αν επιτρέπεται;»

«Σημειώνω»

«Ναι αυτό το έχω καταλάβει. Εννοώ τι ακριβώς σημειώνεις;»

«Τα πάντα»

«Τι εννοείς;»

«Πληροφορίες για όλους τους ανθρώπους που συναναστρέφομαι, όλους τους οποίους κατασκοπεύω, όλους τους οποίο βλέπω στον δρόμο και τους θεωρώ περιέργους. Τα πάντα»

«Γιατί;»

«Πότε δεν ξέρεις που θα χρησιμεύσει» μου απάντησε.

«Δεν σε νοιάζει... που ο Χερμπ...»

«Προφανώς και με νοιάζει, αφού είναι ο μπάρμπας μου, αλλά...»

«Ο Χερμπ είναι ο πατέρας σου;» φώναξα αυθόρμητα. 

Ο Ρομπ και ο Φυτ να με κοιτάξανε περίεργα.

«Συγγνώμη...» είπα και έσκυψα το κεφάλι μου.

«Δεν πειράζει» με καθησύχασε ο Ρομπ. «Και γω θα ξαφνιαζόμουν» μου είπε και έκλεισε το κινητό του. «Λοιπόν. Πάμε στο γραφείο;»

«Δεν γίνεται να πετάς έτσι τυχαία ότι ο Χερμπ ήταν ο πατέρας σου, ειδικά αυτή τη στιγμή, και να μην το εξηγείς παραπάνω» του είπα εκνευρισμένος. «Και η Μπλαιρ; Ποια η σχέση της με τον Χερμπ; Γιατί έχει πάθει τέτοια κρίση;»

«Αυτό είναι συζήτηση για άλλη ώρα. Τώρα πάμε στο γραφείο, μας περιμένει ο Μαλχη. Σήκω Φυτ» είπε ο Ρομπ και σηκώθηκαν και οι δύο ταυτόχρονα από τον καναπέ.

Όσο προχωράγαμε στον διάδρομο κανείς μας δεν μιλούσε. Ήμασταν τρεις μαυροντυμένοι άντρες που είχαν χάσει τον τέταρτο άντρα της ομάδας τους, οπότε δεν είχαμε και πολλά να πούμε. Όχι ότι κάθε ομάδα είχε τέσσερις άντρες και μία κοπέλα. Από ότι ακούγεται, στην Ειδική Υπηρεσία υπάρχουν μικτές ομάδες με διάφορες αναλογίες φύλων, αλλά και ομάδες με μόνο άντρες οι γυναίκες. Όμως, κάθε ομάδα είχε έναν ανακριτή, που από ότι έχω καταλάβει είναι περισσότερο το “μυαλό” της ομάδας και λιγότερο ανακριτής, έναν κατάσκοπο ο οποίος μαζεύει πληροφορίες για τις υποθέσεις που μελετά η ομάδα, και τρεις ειδικούς στις μάχες, για τυχόν αποστολές. Αυτό το σχήμα, λένε, δίνει στις ομάδες ευελιξία στην οργάνωση σχεδίων. Όμως εμείς τώρα είχαμε έλλειψη στους ειδικούς στις μάχες, μετά την τελευταία μας αποστολή. Λες όντως να φταίω εγώ; Λες να έχει δίκιο η Μπλαιρ; Μήπως έπρεπε να σκεφτώ πιο καλά το σχέδιο; Μήπως έπρεπε να σκεφτώ πως ο θάνατος κάποιο μας δεν είναι τόσο απίθανος; 

Οι τύψεις μου είχαν αυξηθεί αρκετά, αλλά μόλις είδα την νέα κορνίζα δίπλα από αυτή του πατέρα μου, έξω από το γραφείο μας, έμεινα άφωνος. “Χερμπ Γκαμπλερ, ειδικός στις μάχες”.

«Σου είπα θα μάθεις σύντομα...» μου είπε ο Ρομπ και με χτύπησε φιλικά στον ώμο καθώς περνούσε από πίσω μου για να μπει στο δωμάτιο.

Όλες, δηλαδή, οι κορνίζες σε αυτόν τον διάδρομο ήταν... νεκροί πράκτορες; Πώς γίνεται αυτό; Ο διάδρομος είναι τεράστιος, τουλάχιστον 8 λεπτά περπάτημα. Όλοι αυτοί οι πράκτορες θέλει να μου πει πως θυσίασαν τη ζωή τους για την Ειδική Υπηρεσία; Λες δηλαδή αυτός ο διάδρομος να είναι το... νεκροταφείο της Ειδικής Υπηρεσίας;

«Σπαρκ κάτι έχουμε πει» εμφανίστηκε ο Μαλχη ξαφνικά από πίσω μου και με τίναξε. «Καλά δεν ήθελα να σε τρομάξω, αλλά περιμένω αρκετή ώρα να έρθετε, για να σας πω τις εξελίξεις με τον Γιν»

«Με τον Γιν; Τον ανέκρινες μόνος σου;»

«Αυτό είναι το πρόβλημα, πάμε μέσα για να το πω σε όλους» μου είπε και μπήκαμε μαζί μέσα στο γραφείο.

Μέσα, στον έναν καναπέ κάθονταν ο Ρομπ και ο Φυτ, με σχεδόν τις ίδιες πόζες που είχαν έξω από το Ιατρείο, μόνο που αυτή την φορά ο Ρομπ είχε και τα δύο χέρια να καλύπτουν το στόμα του. Δίπλα στην πόρτα, γυρισμένη πλάτη, καθόταν η Μπλαιρ, διπλωμένη πάνω στον καναπέ χωρίς τα παπούτσια της, με τα χέρια της γύρω από τα πόδια της, λυγισμένη δεξιά έτσι ώστε να ακουμπάνε τα πλευρά και το κεφάλι της στην πλάτη του καναπέ. Αυτή τη φορά δεν μουρμούριζε τίποτα, ούτε κουνούσε τα χέρια της με τον χαρακτηριστικό της τρόπο.

«Κάθισε στο γραφείο Σπαρκ» μου είπε ο Μαλχη και στάθηκε στο κέντρο του δωματίου για να μπορεί να μας κοιτάζει όλους ταυτόχρονα. «Λοιπόν. Η προηγούμενη αποστολή ήταν μία, πως να το πω τώρα... αποτυχία. Προφανώς ξέρετε ότι χάσαμε τον Χερμπ, έναν από τους πιο έμπειρους πράκτορες μας, αλλά υπάρχει και άλλο ένα γεγονός που συνέβει όταν γυρίσαμε από την αποστολή μας. Σπαρκ, εσένα ίσως σου θυμίσει κάτι αυτή η εικόνα...»

Ο Μαλχη ήρθε πιο κοντά και μου έδειξε στο κινητό του το κελί στο οποίο είχε τοποθετήσει τον Γιν για περαιτέρω ανάκριση. Ο Γιν όμως δεν καθόταν στο παγκάκι ή στο κρεβάτι του κελιού, ήταν ξαπλωμένος στο πάτωμα και μία νοσοκόμα, όχι η Ιφι αυτή την φορά, προσπαθούσε να τον θεραπεύσει με παρόμοιο τρόπο με την Ιφι, βγάζοντας κίτρινο φως από τα χέρια της. Λες να είναι κάποιου είδους μαγείας αυτό το κίτρινο φως τελικά; Προφανώς και είναι, δεν ξέρω γιατί δεν το σκέφτηκα την πρώτη φορά. Αυτό που με προβληματίζει όμως είναι το γεγονός ότι μέσα σε δύο μέρες είχε γίνει το ίδιο περιστατικό σε δύο διαφορετικά άτομα, στον προμηθευτή κλάσης Α, που ήταν λίγο μεγαλύτερος ηλικιακά από εμένα, και στον Γιν, που φαινόταν να ήταν και αυτός περίπου 20 χρονών. Λες να έχει να κάνει με την ηλικία; Μήπως τα κελιά έχουν να κάνουν κάτι με αυτά τα συμβάντα; Μάλλον όχι, θα είχαν πάθει κάτι και οι υπόλοιποι 16 ανακρινόμενοι, εκτός αν επηρεάζει μόνο ανθρώπους κλάσης Α...

«Τι σκέφτεσαι Σπαρκ;» με ρώτησε ο Μαλχη.

«Τίποτα...» του απάντησα και συνέχισα τις σκέψεις μου.

«Θα μας δείξεις και εμάς ή θα μας αφήσεις στον αέρα;» είπε ο Φυτ εκνευρισμένος.

«Έρχομαι τώρα...» του απάντησε ο Μαλχη και πήγε προς τον καναπέ τους. Οι δυο τους φάνηκαν προβληματισμένοι.

«Τι ακριβώς βλέπουμε;» ρώτησε ο Φυτ.

«Ο Γιν, ευτυχώς ή δυστυχώς για εμάς, είναι νεκρός» είπε ο Μαλχη. 

Παρατήρησα την Μπλαιρ που αναστέναξε όταν άκουσε την πρόταση του Μαλχη. Λογικά θυμήθηκε τον Χερμπ.

«Θες να δεις και εσύ Μπλαιρ;» την ρώτησε ο Μαλχη.

Αυτή δεν του απάντησε.

«Καλά, εσύ χάνεις» της είπε και έβαλε το κινητό του στην τσέπη του. «Καταλαβαίνετε, λοιπόν, πως η υπόθεση αυτή δεν μπορεί κλείσει, καθώς δεν είχαμε την δυνατότητα να ανακρίνουμε τον κύριο ύποπτο. Αλλά η υπόθεση δεν μπορεί να μείνει και ανοιχτή, καθώς, θα το πω λαϊκά, δεν έχουμε ιδέα από που να την ξαναπιάσουμε για να συνεχιστή η έρευνα. Άρα η υπόθεση αυτή μένει μισάνοιχτη στους φακέλους της Ειδικής Υπηρεσίας μέχρι να ανακαλυφθεί κάτι καινούργιο. Τώρα...»

«Μισάνοιχτη;...» αναστέναξε η Μπλαιρ χωρίς να αλλάξει την στάση που βρισκόταν.

«Τι πρόβλημα υπάρχει;» την ρώτησε ο Μαλχη.

«Μπορεί εσύ να μην βλέπεις το πρόβλημα, αλλά αυτό δεν σημαίνει ότι δεν υπάρχει ένα τεράστιο πρόβλημα...» του απάντησε η Μπλαιρ, μένοντας πάλι στην ίδια στάση, αφήνοντας μερικά δάκρια να κυλήσουν στα μάγουλα της.

«Δεν σε καταλαβαίνω. Μπορείς να με βοηθήσεις να καταλάβω και να το λύσουμε μαζί;» της είπε ο Μαλχη και προσπάθησε να την πλησιάσει.

«Μην τολμήσεις και με αγγίξεις...» σηκώθηκε τον απείλησε η Μπλαιρ. «...και μην κάνεις τον ηλίθιο. “Αποτυχία”; “Μισάνοιχτη”; Πώς μπορείς να λες τέτοιες λέξεις ενώ έχουμε χάσει ένα από τα πιο σημαντικά μέρη της ομάδας μας;»

«Πιο σημαντικά; Ο Χερμπ το μόνο που πρόσφερε στην ομάδα σας ήταν μέγεθος και εμπειρία. Τίποτα παραπάνω» της είπε ο Μαλχη. «Εκτός αν συμπεριλαμβάνεις τις βόλτες μου με τον Χερμπ στο καζίνο. Τότε ναι, ο Χερμπ ήταν πολύ σημαντικός για μένα, γιατί ήταν ο μόνος που ήταν όσο πορομένος όσο εγώ στον τζόγο και...»

«Δεν θα μπορούσε να με νοιάζει λιγότερο ο τζόγος» σηκώθηκε η Μπλαιρ και πλησίασε τον Μαλχη. «Δεν έχεις ιδέα τι έχει κάνει αυτός ο άνθρωπος. Για σένα, ο Χερμπ μπορεί να ήταν απλά ένα πιόνι σου, αλλά για μένα ήταν κάτι πολύ παραπάνω. Ό,τι και να κάνεις από εδώ και πέρα δεν μπορείς ούτε να τον αντικαταστήσεις, ούτε να γεμίσεις το κενό που έχω αυτή τη στιγμή στην καρδιά μου. Ο μόνος λόγος που μπήκα στην Ειδική Υπηρεσία ήταν ο Χερμπ. Τώρα δεν έχω πλέον λόγω να βρίσκομαι εδώ, οπότε αντίο» είπε η Μπλαιρ και έφυγε από το γραφείο.

«Δεν θα της πεις κάτι Ρομπ;» είπα, αφού ο Ρομπ βρισκόταν στην ίδια ακριβώς κατάσταση με την Μπλαιρ, οπότε υπέθεσα πως μπορεί να την βοηθήσει.

«Δεν χρειάζεται, θέλει απλά τον χρόνο της. Κάνει πολλές τέτοιες συναισθηματικές εκρήξεις και πάντα κάποιος βρίσκεται να την ηρεμήσει. Συνήθως αυτός ο κάποιος είναι ο Χερμπ ή κάποια άλλη κοπέλα που μπορεί να την καταλάβει» μου είπε ο Ρομπ και σηκώθηκε.

«Μη φύγεις ακόμη. Θέλω να σας γνωρίσω και το καινούργιο μέλος της ομάδας σας» είπε ο Μαλχη και μας έδειξε την πόρτα του γραφείου. 

Από την πόρτα μπήκε ένα γνώριμο πρόσωπο. Ήταν ο Κουπ. Άρα εκείνη τη μέρα πέρασε την αξιολόγηση, για αυτό ήταν χαρούμενος. Γιατί, όμως, είναι στην ομάδα μας; Δεν έχει ήδη ομάδα;

«Χαίρεται. Λέγομαι Κουπ, ειδικός στις μάχες, είμαι προφανώς κλάσης Α και έχω μαγεία Ελαστικότητας. Μπορώ να ρωτήσω εγώ κάτι πριν μου κάνετε ερωτήσεις;»

«Πες το» του είπε ο Μαλχη.

«Αυτή η κοπέλα που έφυγε εκνευρισμένη από το γραφείο τι πρόβλημα έχει;»

«Αν γνωρίζαμε θα στο λέγαμε, αλλά δυστυχώς δεν έχουμε ιδέα» του απάντησε ο Μαλχη προσπαθώντας να μας αντιπροσωπεύσει όλους, αλλά προφανώς ήταν ο μόνος που δεν είχε καταλάβει τι γινόταν στο μυαλό της Μπλαιρ.

«Καλά... Ωχ Σπαρκ, δεν σε είδα. Ώστε είσαι και εσύ σε αυτή την ομάδα; Ενθουσιάστηκα τώρα» είπε ο Κουπ ο οποίος δεν είχε καμία ιδέα του τι είχε συμβεί προηγουμένος.

«Και εμείς...» τον ειρωνεύτηκε ο Φυτ και σηκώθηκε για να τον καλωσορίσει στην ομάδα.

Όσο οι τέσσερις τους συζητούσαν μεταξύ τους για τις μαγείες και τις ειδικότητες τους, εγώ καθόμουν στο γραφείο και σκεφτόμουν. Προφανώς, το μόνο πράγμα που μπορούσα να σκεφτώ ήταν η Μπλαιρ. Η κατάσταση της ήταν τόσο όμοια, αλλά ταυτόχρονα τόσο διαφορετική από την δική μου κατάσταση όταν είχε πεθάνει ο δικός μου πατέρας. Τόσο όμοια γιατί προφανώς και οι δύο χάσαμε την πατρική μας φιγούρα στα ξαφνικά, χωρίς προειδοποίηση, όπως για παράδειγμα να είχαν κάποια αρρώστια, στην οποία περίπτωση θα ήμασταν και οι δύο ψυχολογικά περισσότερο προετοιμασμένοι.

Από την άλλη, εγώ έμαθα τα νέα από τρίτους, ενώ η Μπλαιρ ήταν μπροστά τη στιγμή του θανάτου, ανίκανη να βοηθήσει τον Χερμπ καθώς ήταν παγιδευμένη από τη μαγεία Φύσης του Γιν. Λες να έχει τύψεις; Να νιώθει ότι μπορούσε να κάνει κάτι; Και τότε γιατί ρίχνει το φταίξιμο σε μένα; Μήπως έπρεπε εγώ να είχα κάνει κάτι; Τι να είχα κάνει όμως. Το περιστατικό αυτό έγινε τόσο γρήγορα που ούτε ο Ρομπ δεν θα μπορούσε να σταματήσει τις ρίζες του Γιν με την μαγεία Τροχιάς του. Αν μπορούσα όμως να κάνω κάτι και απλά έχω την ψευδαίσθηση ότι ήμουν ανίκανος;

Η Μπλαιρ επέστρεψε στο γραφείο με την Ιφι να την συνοδεύει.

«Ώστε επέστρεψες...» είπε ο Μαλχη.

Η Μπλαιρ τον αγνόησε και κάθισε στον καναπέ με παρόμοια στάση όπως προηγουμένως. Η Ιφι κάθισε δίπλα της και άρχισαν να συζητούν χαμηλόφωνα. Δεν μπορούσα να διακρίνω τι λέγανε, αλλά μάλλον η Ιφι προσπαθούσε να την καθησυχάσει.

«Δεν μας μιλάς, ε;» την ειρωνεύτηκε ο Μαλχη.

«Μαλχη σταμάτα» του είπε ο Φυτ. «Είπαμε θέλει τον χρόνο της»

«Όσο δεν μιλάει όμως με εκνευρίζει, μη βρίσω το σπανάκι της» του απάντησε ο Μαλχη. «Τέλος πάντων, πάω στο καζίνο. Θέλει κανείς να έρθει;»

Κανείς δεν δέχτηκε.

«Ποιον κοροϊδεύω, μόνο ο Χερμπ θα ερχόταν. Αντίο παίδες, τα λέμε αύριο» είπε ο Μαλχη και έφυγε από το γραφείο.

«Επιτέλους έφυγε. Δεν μπορώ να τον καταλάβω αυτόν τον άνθρωπο. Είναι πολύ ψυχρός. Δεν γίνεται να μην μπορεί να καταλάβει πως είμαι συναισθηματικά φορτισμένη; “Αποτυχία” ήταν όντως η αποστολή, αλλά δεν μπορείς να μιλάς με λιγότερο...» άρχισε να κλαίει η Μπλαιρ.

Η Ιφι την αγκάλιασε χωρίς να πει κάτι και με κοίταξε με την άκρη του ματιού της. Της είχε πει κάτι η Μπλαιρ για μένα; Δεν μπορεί, αφού εγώ δεν έκανα τίποτα. Στο μυαλό της Μπλαιρ, όμως, μπορεί να είχα κάνει. Δεν μπορώ να γνωρίζω πως σκέφτεται ο καθένας, αν και... αυτή είναι κυριολεκτικά η μαγεία μου, ή τουλάχιστον το κόλπο που κάνω με την μαγεία μου.

«Σπαρκ, εμείς πάμε για έναν καφέ. Θες να έρθεις;» μου πρότεινε ο Κουπ.

«Όχι, θέλω να κάτσω λίγο να σκεφτώ» του απάντησα

«Δεν πειράζει. Τα λέμε αύριο» μου είπε ο Κουπ και οι τρεις τους μας αποχαιρέτησαν καθώς έβγαιναν από το δωμάτιο.

Επέλεξα να μείνω στο γραφείο για έναν λόγο, να προσπαθήσω να διαβάσω τις σκέψεις της Μπλαιρ ώστε να καταλάβω σε τι οφείλεται η συναισθηματική της έκρηξη. Εκείνη τη στιγμή μου είχε πει πως ο Χερμπ ήταν “σαν” πατέρας της, όχι ακριβώς πατέρας. Τι μπορεί να σημαίνει αυτό; Ξεκίνησα να διαβάζω της σκέψεις της για να μου λυθεί η απορία.

\textasciitilde ...είμαι άχρηστη, ούτε έναν απλό κοντοστουπι δεν μπορώ να αφοπλίσω. Και που με οδήγησε αυτό; Δεν γίνεται αυτός ο άνθρωπος να έχει τόσα πολλά για μένα και εγώ να μη μπορώ να τον βοηθήσω στην πιο κρίσιμη στιγμή της ζωής του...\textasciitilde 

«Μπλαιρ;» της είπα.

«Ναι»

«Τι... είναι... ήταν ο Χερμπ για σένα;»

Δεν μου απάντησε.

\textasciitilde ...ο Χερμπ για μένα; Δεν μπορείς να καταλάβεις εσύ. Δεν έχεις την παραμικρή ιδέα. Προφανώς και τώρα προσπαθείς να διαβάσεις τις σκέψεις μου για να καταλάβεις, αλλά σε ξέρω καλά, τέτοια κόλπα δεν δουλεύουν σε εμένα...\textasciitilde 

«Εντάξει συγνώμη Μπλαιρ, σταματάω» της είπα και σηκώθηκα. «Θέλετε να σας αφήσω μόνες σας να συζητήσετε;» τους πρότεινα καθώς πήγαινα προς την πόρτα.

«Όχι Σπαρκ δεν πειράζει...» ξεκίνησε να λέει η Ιφι.

«Ναι, αν μπορείς» την διέκοψε η Μπλαιρ.

«Καλά...» είπε η Ιφι και έσκυψε το κεφάλι της.

«Συγγνώμη για ότι έγινε σήμερα Μπλαιρ. Ελπίζω να μπορέσω να επανορθώσω κάπως...»

«Απλά φύγε...» με έδιωξε ξεκάθαρα η Μπλαιρ. Η Ιφι φάνηκε απογοητευμένη.

Δεν της απάντησα κάτι, βγήκα απλά από το δωμάτιο, έκλεισα την πόρτα και άρχισα να πηγαίνω σπίτι μου.

\chapter{}

«Καλημέρα Σπαρκ» μου είπε η γραμματέας.

«Καλημέρα. Μήπως έχουμε κάτι καινούργιο σήμερα;»

«Ο Μαλχη δε με έχει ενημερώσει για κάποια καινούργια υπόθεση για την ομάδα σας, οπότε μπορείς να πας στο γραφείο σου και να κάτσεις προς το παρόν»

«Έχουν περάσει τρεις μέρες όμως. Πως γίνεται να μην έχει εμφανιστεί κάτι καινούργιο;»

«Υπάρχουν τέτοιες περιόδους για κάθε ομάδα στην Ειδική Υπηρεσία. Μερικές φορές πρέπει να είσαι ευγνώμων που έχεις την ευκαιρία να χαλαρώσεις παρά να παραπονιέσαι...» μου είπε η γραμματέας και συνέχισε να γράφει στον υπολογιστή της.

«Καλημέρα...» άκουσα από πίσω μου.

Γύρισα ξαφνιασμένος και είδα την Ιφι να στέκεται πίσω μου με το βλέμμα της στραμμένο στο πάτωμα.

«Άργησες σήμερα...» μου είπε χωρίς να σηκώσει το βλέμμα της.

«Καλά 10 λεπτά αργότερα ήρθα και δεν είναι ότι έχουμε τίποτα σημαντικό να κάνουμε αυτές τις μέρες...»

«Τίποτα σημαντικό, ε...»

«Γιατί, σου είπε ο Μαλχη καμία νέα υπόθεση;»

«Όχι δεν εννοώ αυτό...»

«Θα μου πεις, εσύ είσαι νοσοκόμα οπότε δεν συμβάλεις στις υποθέσεις άμεσα, άρα ο Μαλχη δεν έχει λόγο να σε...»

«Θες να περπατήσουμε μαζί μέχρι το γραφείο σου, Σπαρκ; Έτσι κι αλλιώς, όπως λες κι εσύ, δεν έχουμε τίποτα σημαντικό να κάνουμε» είπε η Ιφι και με κοίταξε στα μάτια, παίζοντας με τα χέρια της πίσω από την πλάτη της και κουνώντας τα πόδια της αμήχανα.

«Και αν τραυματιστεί κανείς; Δεν πρέπει να είσαι στην βάρδια σου;»

«Έχουμε πάρα πολλές νοσοκόμες με μαγεία Θεραπείας στην Ειδική Υπηρεσία, οπότε δεν νομίζω να υπάρχει έλλειψη ανθρώπινου δυναμικού...»

«Καλά, αφού το λες εσύ, πάμε» της είπα και αρχίσαμε να περπατάμε στον διάδρομο προς το γραφείο.

Οι κορνίζες τις τελευταίες μέρες μου φαίνονταν πολύ αδιάφορες. Της είχα συνηθίσει; Όχι δεν ήταν αυτό, αφού κάθε μέρα βλέπω όλο και κάποιο καινούργιο πρόσωπο που δεν είχα παρατηρήσει τις προηγούμενες φορές. Λες να έχασαν τη σημασία τους για μένα όταν έμαθα τον λόγω ύπαρξης τους μετά την πρώτη μου αποστολή; Γιατί όμως; Οι κορνίζες απ’ ό,τι έχω καταλάβει υπάρχουν για να τιμούνται όσοι έχουν χάσει την ζωή τους υπηρετώντας την Ειδική Υπηρεσία. Γιατί μου φαίνεται πλέον αδιάφορο αυτό; Λες τελικά να είμαι εγώ ο ψυχρός και όχι ο Μαλχη; Δεν νομίζω, ο Μαλχη είναι όντως πιο ασυναίσθητος από εμένα, αλλά ακολουθώντας τις συμβουλές του θα γίνω και εγώ σαν αυτόν.

«Ιφι» της απευθύνθηκα.

«Παρακαλώ» τινάχτηκε γιατί και αυτή φαινόταν σαν να σκέφτεται κάτι όσο περπατάγαμε.

«Αν είμαι διακριτικός, τι έγινε τελικά με την Μπλαιρ;»

«Α... είναι...»

«Ρωτάω γιατί τις τελευταίες μέρες απλά έρχεται στο γραφείο και απλά κάθεται διπλωμένη στους καναπέδες χωρίς να μου μιλάει. Θέλω να μάθω αν μπορώ να κάνω κάτι για να την βοηθήσω»

«Δεν... μπορείς να βοηθήσεις. Θέλει απλά τον χρόνο της»

«Ναι αλλά έχουν περάσει 3 μέρες. Πόσο χρόνο ακόμη;»

«Όταν έρθει η στιγμή θα το καταλάβεις, ας... την αγνοούμε για τώρα και ας συγκεντρωθούμε στις δουλειές μας...»

Η τελευταία αυτή πρόταση της Ιφι δεν μου φάνηκε λογική. Δεν μπορώ να αγνοήσω μια τόσο προφανέστατα προβληματική κατάσταση, και προφανώς δεν μπορώ να αφήσω αυτή τη κατάσταση να συνεχίσει, γιατί θα δημιουργούσε στη συνέχεια άλλα προβλήματα.

Καθώς περνάγαμε έξω από τον χώρο αναμονής, είδα τον Φυτ, τον Μαλχη και τον Κουπ να συζητάνε.

«Ωχ ο Σπαρκ και η Ιφι. Καλημέρα παιδιά» μας είπε ο Κουπ.

«Καλημέρα ομάδα» είπα και τους χαμογέλασα.

«Σας έχω νέα υπόθεση σήμερα Σπαρκ» μου είπε ο Μαλχη.

«Υπέροχα νέα. Αυτό σημαίνει πως θα αρχίσουμε δουλειά από σήμερα;»

«Όχι. Μέχρι στιγμής είμαστε στο στάδιο της έρευνας. Ο Ρομπ κατασκοπεύει τους υπόπτους μας»

«Και αυτό τι σημαίνει;» τον ρώτησα από περιέργεια

«Σημαίνει ότι εσείς οι υπόλοιποι δεν χρειάζεται να κάνετε κάτι μέχρι να επιστρέψει και έχουμε τις πληροφορίες που χρειαζόμαστε» μου απάντησε.

«Όμως, την προηγούμενη φορά...»

«Την προηγούμενη φορά είχαμε κάνει την έρευνα μας πριν εσύ και η Μπλαιρ γίνετε μέλη στην ομάδα»

«Θα μας πεις τουλάχιστον τι αφορά η υπόθεση;»

«Αυτό βεβαίως. Θέλουμε πρώτα όμως να κάνουμε λίγη προπόνηση στη δίπλα αίθουσα τώρα που είναι πρωί και θα έρθουμε σε λιγάκι»

«Εντάξει, θα σας αφήσουμε να κάνετε την δουλειά σας»

«Ιφι να είσαι και εσύ στο γραφείο, μήπως χρειαστεί να θεραπεύσεις κάποιον από εμάς, αν και ελπίζω να μην χρειαστεί»

«Μάλιστα...»

«Ευχαριστώ πολύ Ιφι» είπε ο Μαλχη και της έκλεισε το μάτι. «Ας πηγαίνουμε τώρα»

«Και εμείς» είπα και άρχισα να περπατάω προς το γραφείο. Η Ιφι με ακολούθησε συντόμως.

Στην υπόλοιπη διαδρομή μας προς το γραφείο δεν μιλήσαμε καθόλου. Δεν ξέρω γιατί μου είπε η Ιφι να πάμε μαζί στο γραφείο. Δεν είμαστε καν “φίλοι”, την ξέρω μόνο μερικές μέρες, όσο και τους υπόλοιπους στην ομάδα μου, αλλά με την Ιφι ειδικά δεν έχουμε συζητήσει ποτέ. Λες να... δεν νομίζω. Μπορεί απλά να θέλει να δει και αυτή την Μπλαιρ. Φαίνεται να είναι γύρω στην ίδια ηλικία με εμένα και την Μπλαιρ, οπότε μπορεί να γνωρίζονται από πιο παλιά.

Όσο περπατούσαμε, μου ήρθε στο μυαλό η ανορθόδοξη σχεδίαση του διαδρόμου που περπατούσαμε. Δεν βγάζει νόημα που είναι τόσο μεγάλος και έχει μόνο τέσσερις πόρτες, οι τρεις εκ των οποίων είναι στο τέλος του διαδρόμου και οδηγούν σε γραφεία ομάδων. Γιατί να μην έχουν τοποθετηθεί στην αρχή του διαδρόμου και η αίθουσα εξάσκησης να έχει μετακινηθεί λίγο πιο εκεί. Αν και η πόρτα της αίθουσας εξάσκησης θα μπορούσε να είναι και αυτή προς την αρχή του διαδρόμου. Είναι λες και αυτός που σχεδίασε τον διάδρομο τον έκανε επίτηδες τόσο μεγάλο για να μας κάνει να υπεραναλύουμε τις σκέψεις στο μυαλό μας πριν πάμε στα γραφεία των ομάδων μας. Επίσης, γιατί τα γραφεία στο τέλος του διαδρόμου είναι μόνο τρία; Δεν νομίζω να υπάρχουν μόνο τρεις ομάδες στην Ειδική Υπηρεσία. Λες να υπάρχουν κι άλλοι τέτοιοι διάδρομοι ή, ακόμη καλύτερα, τα γραφεία των άλλων ομάδων βρίσκονται κάπου κρυφά που μπορείς να μπεις μόνο με μαγεία;

Φτάνοντας στο γραφείο, το βλέμμα μου στράφηκε προς την κορνίζα του πατέρα μου. “Δεν ξέρω ποιος είμαι”, “Δεν ξέρω που είμαι”. Ακόμη δεν έχω καταλάβει αυτές τις τελευταίες λέξεις του πατέρα μου. Λες να τις έχω παρερμηνεύσει; Θυμάμαι να λέει και κάποιες άλλες λέξεις ενδιάμεσα που δεν μπορώ να τις ανακαλέσω.

«Τι σε προβληματίζει Σπαρκ;» με ρώτησε η Ιφι.

«Τίποτα»

«Ωχ, αυτός ο ανακριτής έχει ίδιο επίθετο με εσένα. Μήπως έχετε κάποια συγγένεια;» είπε η Ιφι και μου έδειξε την κορνίζα του πατέρα μου. «Τι σύμπτωση που είναι και έξω από το...»

«Είναι ο πατέρας μου. Δυστυχώς τον έχασα σε μικρή ηλικία και δεν το έχω ξεπεράσει ακόμη»

«Άρα... γιατί δεν μιλάς στην Μπλαιρ;»

«Τι;»

«Γιατί δεν της εξηγείς πως ξέρεις πως νιώθει;»

«Τι εννοείς;»

«Αφού και εσύ έχεις χάσει τον πατέρα σου, άρα καταλαβαίνεις τα τυχόν συναισθήματα που μπορεί να έχει. Άρα θα μπορούσες να την προσεγγίσεις και να της εξηγήσεις πως θα το ξεπεράσει και πως δεν θα κάθεται όλη μέρα και να κλαίει την μοίρα της»

«Οι καταστάσεις μας  όμως είναι διαφορετικές»

«Τι εννοείς;»

«Δεν... μπορώ να το εξηγήσω, αλλά στο μυαλό μου...»

«Ναι;»

«Άστο. Πάμε να κάτσουμε στο γραφείο και να περιμένουμε τους υπόλοιπους. Αν μου έρθει κάτι να πω στην Μπλαιρ θα ανοίξω συζήτηση» της είπα και άνοιξα την πόρτα του γραφείου.

Μέσα βρισκόταν ήδη η Μπλαιρ και είχε πάρει την χαρακτηριστική της διπλωμένη πόζα στον καναπέ. Εγώ κάθισα στο γραφείο και άνοιξα το κινητό μου για να δω αν μου έχει στείλει κάποιος μήνυμα. Η Ιφι κάθισε δίπλα στην Μπλαιρ, ακριβώς όπως την προηγούμενη φορά. Στο κινητό μου είδα ένα μήνυμα από τον Βασι που έγραφε “Συγγνώμη που δεν σου έχω απαντήσει τις τελευταίες μέρες, ήμουν σε μία αποστολή εκτός πόλης και έπρεπε να ήμουν όλη μέρα σε εγρηγορση. Ελπίζω τώρα που γύρισα να μπορέσουμε να μιλήσουμε καθόλου”. Ευτυχώς ο Βασι ήταν καλά. Μετά τον θάνατο του Χερμπ είχα αρχίσει να αγχώνομαι για αυτόν, αλλά και για την μητέρα μου που είναι σε αποστολή εδώ και πάνω από μία εβδομάδα. Είχε φύγει από το σπίτι κάποιες μέρες πριν την αξιολόγηση μου και δεν έχουμε μιλήσει από τότε. Ελπίζω να είναι και αυτή καλά, γιατί δεν θέλω να χάσω το μόνο μέρος της οικογένειας μου που έχει μείνει.

«Σπαρκ!» με φώναξε η Ιφι. «Η Μπλαιρ θέλει να σου πει κάτι» 

«Τι έγινε» αναρωτήθηκα.

«Συγγνώμη για τον τρόπο που σου έχω φερθεί τις τελευταίες μέρες» μου είπε η Μπλαιρ χωρίς να με κοιτάξει ούτε να αλλάξει την στάση της

«Από που ήρθε αυτό;» παραξενεύτηκα.

«Δεν δέχεσαι τη συγχώρεση μου;»

«Όχι όχι, προφανώς και τη δέχομαι, απλά είναι πολύ ξαφνική, λες και είμαστε σε ταινία και ο σεναριογράφος ήθελε να προχωρήσει κάπως η πλοκή, αλλά δεν μπορούσε να βρει πως»

«Ήξερα ότι είσαι χαζός» μου είπε η Μπλαιρ και άρχισε να κλαίει.

«Συγγνώμη» είπα κλείνοντας το κινητό μου και πλησιάζοντας την. «Δεν ήθελα να... δεν κατάλαβα βασικά τι σε πείραξε. Μία πλάκα προσπάθησα να κάνω»

«Η Ιφι μου είπε... πως έχεις χάσει και εσύ... και... ήθελα να...» είπε η Μπλαιρ κλαίγοντας με το κεφάλι της ανάμεσα στα πόδια της για να μην μπορούμε να διακρίνουμε το πρόσωπο της.

«Δεν πειράζει. Σου είπα δέχομαι την συγχώρεση σου. Μπορείς να δεχτείς και εσύ τη δική μου συγχώρεση; Δεν χρειάζεται να τσακωνόμαστε πλέον» είπα και την αγκάλιασα. Αυτή δεν έκανε κάποια κίνηση. Κοίταξα την Ιφι και φάνηκε λίγο προβληματισμένη.

Όταν χαλάρωσα την αγκαλιά μου, η Μπλαιρ με κοίταξε επιτέλους μετά από τόσο καιρό στα μάτια, με ένα χαμόγελο στο πρόσωπο της και δάκρυα να κυλούν στα μάγουλα της.

«Έχεις δίκιο. Δεν φταις εξάλλου σε τίποτα. Εγώ είμαι η άχρηστη της υπόθεσης...»

«Όχι βέβαια. Μία δασκάλα μου είχε πει παλιά πως όλοι οι άνθρωποι είναι χρήσιμοι με τον δικό τους τρόπο»

«Και αυτό πως υποτίθεται θα με βοηθήσει να είμαι χρήσιμη;»

«Αυτό σου λέω, είσαι είδη “χρήσιμη”. Χωρίς εσένα και τη μαγεία σου δεν θα μπορούσαμε να είχαμε πιάσει όλους αυτούς τους προμηθευτές στην πρώτη μας αποστολή. Χωρίς την εξωστρέφεια σου, που την έχεις χάσει τις τελευταίες μέρες, δεν θα μπορούσαμε να...»

«Τι σου έκανε η εξωστρέφεια μου τώρα;»

«Τίποτα, για καλό το...»

«Πάλι τσακώνεστε;» με διέκοψε ο Μαλχη.

«Όχι, βασικά... Πώς τελειώσατε τόσο γρήγορα;» ρώτησα τον Μαλχη.

«Γρήγορα; Έχει περάσει μισή ώρα από όταν σας αφήσαμε»

«Ναι, αυτό εννοώ. Μόνο τόση λίγη ώρα προπόνηση;»

«Είσαι κλάσης Β εσύ θα μου πεις, που να καταλάβεις» μου είπε ο Μαλχη και άνοιξε το κινητό του για να δει κάτι. Ο Φυτ και ο Κουπ είχαν μπει ήδη στο δωμάτιο και είχαν καθίσει στον άδειο καναπέ, χωρίς να βγάλουν άχνα.

«Λοιπόν...» είπε ο Μαλχη και έκλεισε το κινητό του. «Να μπω κατευθείαν στο ψητό;»

«Γιατί τι έγινε;» τον ρώτησε η Μπλαιρ «Επίσης που στο καλό είναι ο Ρομπ; Έχω να τον δω αρκετές μέρες»

«Α ώστε μας μιλάς τελικά...» την ειρωνεύτηκε ο Μαλχη.

«Μαλχη δεν είναι ώρα για τέτοια» του είπε η Ιφι για να τον ηρεμήσει.

«Άσε μας και εσύ, ό,τι θέλω θα κάνω, μη βρίσω το σπανάκι σου. Λοιπόν θα τα πω συνοπτικά, γιατί έτσι κι αλλιώς δεν έχουμε και πολλές πληροφορίες. Την Δευτέρα, όταν γυρνούσα το βράδυ από το καζίνο, είδα δύο κουκουλοφόρους άντρες που φαίνονταν πολύ ύποπτοι...»

«Τι εννοείς “φαίνονταν” ύποπτοι;» τον διέκοψε η Μπλαιρ.

«Θα με αφήσεις να τελειώσω; Είχα μια διαίσθηση, ρε παιδί μου, πως να το πω για να το καταλάβεις; Τους ακολούθησα, λοιπόν, μέχρι που φτάσαμε σε μία πολυκατοικία που δεν φαινόταν να έχει κάτι παράξενο. Υπάρχει μεγάλη πιθανότητα αυτό να ήταν απλά το σπίτι τους. Ανέθεσα στον Ρομπ εκείνο το βράδυ να συνεχίσει να τους παρακολουθήσει στενά την επόμενη μέρα και αν του κινήσουν υποψίες, να συνεχίσει να τους κατασκοπεύει, μέχρι να καταλήξει σε ένα συμπέρασμα. Από ό,τι βλέπετε, ο Ρομπ δεν είναι μαζί μας σήμερα, οπότε οι κουκουλοφόροι ήταν όντως ύποπτοι για κάτι. Αυτό που θέλω από όλους σας είναι αν είστε σε εγρήγορση, γιατί σε περίπτωση που επιστρέψει ο Ρομπ, μπορεί να χρειαστεί να φύγετε κατευθείαν για αποστολή. Ερωτήσεις;»

«Ακόμη δεν καταλαβαίνω πώς ή γιατί τους θεωρείς ύποπτους;» είπε η Μπλαιρ. «Σου έκαναν κάτι αυτοί οι άνθρωποι; Προσπάθησαν να σε κλέψουν; Να σε τραυματίσουν μήπως;»

«Όχι, απλά μία διαίσθηση» της απάντησε ο Μαλχη.

«Μαλχη όντως δεν βγάζει νόημα αυτό, μπορείς να μας εξηγήσεις λίγο παραπάνω τι εννοείς όταν λες “μία διαίσθηση”;» του είπε ο Φυτ προβληματισμένος.

«Δεν εννοώ κάτι παραπάνω, μη βρίσω το σπανάκι σας. Απλά ο τρόπος που ήταν ντυμένοι οι τύποι ήταν λες και ανήκουν σε κάποια αίρεση ή κάτι»

«Και τι το παράνομο έχουν οι αιρέσεις;» του είπε ο Μπλαιρ.

«Ποτέ δεν ξέρεις τι περίεργο μπορεί να κρύβεται πίσω από μία θρησκεία ή από μία αίρεση, οπότε για να είμαστε σίγουροι θα εξετάσουμε αυτούς τους άντρες και θα δούμε αν είναι “καλοί” ή “κακοί”. Άλλη ερώτηση;»

Δεν του απάντησε κανείς.

«Ωραία λοιπόν. Αν με χρειαστείτε κάτι θα είμαι στο γραφείο μου»

\chapter{}

«Καλημέρα Σπαρκ» μου είπε η γραμματέας.

«Καλημέρα. Μήπως έχουμε κάτι καινούργιο σήμερα;»

«Ο Μαλχη δε με έχει ενημερώσει για κάποια καινούργια υπόθεση για την ομάδα σας, οπότε μπορείς να πας στο γραφείο σου και να κάτσεις προς το παρόν»

«Έχουν περάσει πέντε μέρες όμως. Πως γίνεται να μην έχει εμφανιστεί κάτι καινούργιο;»

«Μη με κάνεις να επαναλαμβάνω τον εαυτό μου Σπαρκ. Αν είσαι τόσο περίεργος γιατί δεν πας να ρωτήσεις τον Μαλχη εσύ; Το γραφείο είναι εδώ δίπλα ξέρεις» μου είπε η γραμματέας και μου έδειξε μία πόρτα στα δεξιά της που δεν είχα παρατηρήσει ποτέ.

«Αυτή η πόρτα ήταν από πάντα εδώ;» τη ρώτησα από περιέργεια.

«Μη ρωτάς χαζές ερωτήσεις, θα νευριάσεις τον Μαλχη»

«Έχεις ένα δίκιο... αλλά πώς ξέρεις ότι ο Μαλχη είναι σίγουρα στο γραφείο του; Προχθές τον πέτυχα με τον Φυτ και τον Κουπ έξω από την αίθουσα εξάσκησης»

«Είναι σίγουρα μέσα στο γραφείο του. Η Ιφι μπήκε πριν λίγη ώρα λέγοντας ότι είχε κάτι να του ανακοινώσει και άκουσα την φωνή του να της λέει να περάσει μέσα. Δεν την είδα να βγαίνει, όμως, οπότε μπες διακριτικά» μου είπε η γραμματέας και συνέχισε να γράφει στον υπολογιστή της.

Μπαίνοντας στο γραφείο του Μαλχη, μία σφαίρα πέρασε ξυστά μπροστά από το πρόσωπο μου.

«Βλέπεις τι γίνεται αν με νευριάζεις. Μπορεί να τραυματίσω άνθρωπο. Μπες μέσα και κλείσε την πόρτα!» φώναξε ο Μαλχη και μπήκα γρήγορα μέσα από τον φόβο μου. «Α, εσύ είσαι Σπαρκ. Τι θέλεις; Μη μου πεις πως θες να αλλάξεις κι εσύ ειδικότητα όπως η μαντάμ από εδώ;»

«Να αλλάξω ειδικότητα; Γιατί αν το θέλω αυτό;»

«Που θες να ξέρω, ρώτα την φιλενάδα σου από δω» μου είπε και έδειξε με το όπλο του την Ιφι. «Θέλει, λέει, να γίνει κατάσκοπος και να ενταχθεί στην ομάδα σου»

«Τι;» φώναξα γιατί αιφνιδιάστηκα.

«Και εγώ αυτή την αντίδραση είχα, απλά μου έφυγε και μία σφαίρα από τον εκνευρισμό μου. Οι ειδικότητες σας δίνονται με βάση έναν μόνο χαρακτηριστικό σας, τις μαγείες σας. Δεν μπορούμε να έχουμε “κατάσκοπους” με μαγεία Θεραπείας, όπως ταυτόχρονα δεν μπορούμε να έχουμε νοσοκόμες με μαγεία... ξέρω γω, Φωτιάς ή κάτι...»

«Αυτό όμως είναι πολύ άδικο για αυτούς που έχουν άλλες ικανότητες εκτός από την μαγεία τους» τον διεκοψε η Ιφι. «Γιατί ο ρόλος μας να κρίνεται μόνο με βάση ένα κριτήριο; Γιατί να μη μπορεί ένας παρατηρητικός άνθρωπος, που έχει μία άσχετη μαγεία, να γίνει κατάσκοπος;»

«Γιατί η μαγεία σου δεν είναι άσχετη, αυτό είναι το πρόβλημα. Δεν γίνεται να έχεις μαγεία Θεραπείας και να μην είσαι στο ιατρικό προσωπικό» της απάντησε ο Μαλχη.

«Μα, το ιατρικό προσωπικό περιλαμβάνει πολύ περισσότερα άτομα από όσα χρειάζονται στην πραγματικότητα. Αν δεν υπήρχε έλλειψη προσωπικού δεν υπήρχε περίπτωση να σου έκανα τέτοια πρόταση»

«Και γιατί κατάσκοπος; Πιστεύεις ότι θα σου ταιριάζει περισσότερο;»

«Ναι, Ιφι, εξήγησε το λίγο, γιατί ούτε εγώ μπορώ να σε καταλάβω» είπα.

«Είναι για να... ναι γιατί πιστεύω ότι θα μου ταιριάζει. Είμαι πολύ παρατηρητική και διακριτική στη ζωή μου, οπότε...»

«Οπότε;» την διέκοψε ο Μαλχη. «Η μαγεία σου δεν μπορεί να σε βοηθήσει στην κατασκοπία. Πώς σκοπεύεις να γεφυρώσεις το χάσμα μεταξύ των άλλων κατασκόπων και του εαυτού σου;»

«Απλά δώσε μου μία ευκαιρία. Και για να ξέρεις, ούτε του Ρομπ η μαγεία βοηθάει στην κατασκοπία, παρόλα αυτά είναι κατάσκοπος»

«Άσε μας που η μαγεία του Ρομπ δεν τον βοηθάει στην κατασκοπεία»

«Πειράζει να σας διακόψω για λίγο;» τους χώρισα διακριτικά.

«Τι θες;»

«Ήρθα για να δω αν έχουμε κανένα νέο από τον Ρομπ. Πέντε μέρες είναι πολύς χρόνος, δεν νομίζεις;»

«Ισχύει...» μου είπε ο Μαλχη και άνοιξε το κινητό του. «Έπρεπε να είχε στείλει τουλάχιστον κανένα μήνυμα, καμία φωτογραφία, κάτι... Θα τον πάρω ένα τηλέφωνο» είπε και σήκωσε το τηλέφωνο στο αυτί του.

Το τηλέφωνο χτύπησε αρκετές φορές, αλλά καμία απάντηση από την άλλη γραμμή. Ο Μαλχη προσπάθησε να πάρει και δεύτερη φορά, αλλά μάταια. Αντί να ξαναπάρει τηλέφωνο, άρχισε κάτι να ψάχνει στο κινητό του.

«Μήπως έχεις καμία ιδέα για την τοποθεσία του;» πρότεινα.

«Τώρα αυτό θα κοίταζα... τι; Τελευταία ανανέωση δύο μέρες πριν; Κάτι πάει λάθος» είπε ο Μαλχη προβληματισμένος.

«Τι έγινε;»

«Δεν έχουμε χρόνο, μπορεί να είναι ήδη αργά, αλλά χρειάζεται να πάμε να τσεκάρουμε» σηκώθηκε ο Μαλχη και άνοιξε βιαστικά την πόρτα του γραφείου του. Απ’ έξω από το γραφείο βρισκόταν ο Κουπ και μίλαγε με την γραμματέα.

«Α, καλημέρα Μαλχη, μήπως θες...»

«Δεν έχουμε χρόνο Κουπ. Εσύ και ο Σπαρκ πρέπει να έρθετε μαζί μου επειγόντως» είπε ο Μαλχη και έσπευσε προς την είσοδο τρέχοντας. Ο Κουπ και εγώ τον ακολουθήσαμε.

Ο Μαλχη μας είπε πως το ρανταρ στο κινητό του έδειχνε ότι ο Ρομπ βρισκόταν σε μία ταράτσα κοντά στο σπίτι των δύο κουκουλοφόρων. Το πρόβλημα ήταν ότι η τοποθεσία του δεν είχε ανανεωθεί εδώ και δύο μέρες, που μπορούσε να σημαίνει δύο πράγματα, απ’ ότι κατάλαβα: είτε ότι ο Ρομπ είχε στηθεί σε μία ταράτσα δύο μέρες και κατασκόπευε από εκεί, οπότε το ρανταρ δεν χρειαζόταν να ανανεώσει την τοποθεσία του, αφού δεν κουνιόταν, είτε ότι ήταν... νεκρός. Αυτά τα δύο σενάρια, σε συνδυασμό ότι δεν απάντησε στις κλήσεις του Μαλχη μας οδήγησαν στο συμπέρασμα ότι ο Ρομπ πιθανότατα είχε τραγικό τέλος και τρέχαμε να επιβεβαιώσουμε αν οι υποθέσεις μας ήταν σωστές. 

Όταν φτάσαμε στο κτίριο που έδειχνε το ρανταρ του Μαλχη, ο Κουπ μας ανέβασε εκεί με την μαγεία Ελαστικότητας του στην ταράτσα. Δεν αργήσαμε πολύ να βγάλουμε το τελικό συμπέρασμα μας. Είδαμε τον Ρομπ ξαπλωμένο μπρούμυτα στην άκρη της ταράτσας. Το κινητό του ήταν στα χέρια του, όπως πάντα, αλλά με δύο τρύπες λες και το είχε δαγκώσει φίδι. Στην γάμπα του είχε πάλι ένα σημάδι από δόντια ερπετού.

«Δηλητηρίαση...» είπε ο Μαλχη και σήκωσε το τηλέφωνο του Ρομπ. «Δεν ανοίγει. Διέλυσαν και την μνήμη του κινητού, μη βρίσω το σπανάκι τους. Είναι καλοί Σπαρκ»

Ο Μαλχη έβαλε στην τσέπη του το κινητό του Ρομπ και έσκυψε να πιάσει ένα πετραδάκι. Ο Κουπ και εγώ παραξενευτήκαμε γιατί νομίζαμε ότι θα σήκωνε το πτώμα του Ρομπ για να τον επιστρέψουμε πίσω στην Ειδική Υπηρεσία. Έπειτα ο Μαλχη πέταξε το πετραδάκι από την ταράτσα και το κοίταξε όσο έπεφτε στο πάτωμα, όπου ευτυχώς δεν χτύπησε κανέναν περαστικό ή κάποιο αυτοκίνητο.

«Πάμε παιδιά. Κουπ, κατέβασε μας έναν-έναν και πάρε και τον Ρομπ» διέταξε ο Μαλχη.

«Μάλιστα» του απάντησε ο Κουπ χαμογελαστός και ακολούθησε την εντολή του Μαλχη, κατεβάζοντας μας με την μαγεία του στο  πεζοδρόμιο.\\\\

Όταν επιστρέψαμε στην Ειδική Υπηρεσία είδαμε την Ιφι και την Μπλαιρ να συζητάνε και να γελάνε. Τα χαμόγελα τους αυτά, όμως, δεν θα διαρκούσαν για πολύ ακόμη, καθώς όταν διασχίσαμε την περιστρεφόμενη πόρτα και μας παρατήρησαν, έτρεξαν και οι δυο τους προς το μέρος μας.

«Τι έγινε;» ρώτησε η Ιφι.

«Δεν βλέπεις;» την ερωτεύτηκε ο Μαλχη.

«Να προσπαθήσω τουλάχιστον...»

«Τι να προσπαθήσεις; Ο άνθρωπος είναι εδώ και δύο μέρες νεκρός, ούτε με θαύμα δεν σώνεται» την διέκοψε ο Μαλχη.

Ο Κουπ άφησε κάτω τον Ρομπ και η Μπλαιρ έπεσε στα γόνατα. Ακούμπησε το κρύο πρόσωπο του Ρομπ και ξέσπασε σε κλάματα. Όταν είδα αυτή την εικόνα μου ήρθε στον μυαλό ότι δεν έτυχε, και δεν πρόκειται προφανώς, να ρωτήσω τον Ρομπ τι σχέση έχει με την Μπλαιρ. Η αντίδραση της Μπλαιρ μου επιβεβαιώνει πως ο Ρομπ ήταν κοντινός της άνθρωπος, αν όχι συγγενής της. Συνέχεια όμως τους θυμάμαι να τσακώνονται. Αυτές τις μέρες που είχαμε περάσει όλοι μαζί, σαν ομάδα, από την ημέρα της αξιολόγησης μου μέχρι και το ταξίδι στο “χωριό” του Γιν, όποτε μιλούσαν οι δυο τους, υπήρχε μία ένταση μεταξύ τους, λες και ήταν αδέρφια... Λες; Αφού ο Χερμπ δεν είναι πατέρας της Μπλαιρ. Αν ήταν πατέρας της, η Μπλαιρ δεν θα του αναφερόταν ονομαστικά. Βασικά... ο Ρομπ πως αναφερόταν στον Χερμπ; Πως γίνεται να μην έχω παρατηρήσει ένα τόσο βασικό πράγμα.

«Τώρα τι κάνουμε αρχηγέ;» ρώτησε ο Κουπ.

«Θα βρούμε έναν αντικαταστάτη και θα αρχίσουμε να ερευνούμε πιο προσεκτικά εκείνους τους άντρες» του απάντησε ο Μαλχη. 

«Μήπως θα μπορούσα να... είμαι εγώ η αντικαταστάτρια;» πρότεινε η Ιφι

«Ποια πιστεύεις ότι είσαι... για να λες τέτοιες ανοησίες...» ακούστηκε να λέει η Μπλαιρ ενώ ήταν σκυμμένη στο πάτωμα πάνω από το σώμα του Ρομπ. «Μία απλή νοσοκόμα είσαι, δεν μπορείς να συγκριθείς σε κανέναν τομέα με τον Ρομπ. Μπορεί ο Ρομπ να ένας ψυχρός ηλίθιος σαν τον Μαλχη, αλλά ήταν έξυπνος. Πώς πιστεύεις ότι μπορείς να γεφυρώσεις το κενό μεταξύ σου και του Ρομπ, ε;»

«Άντε πάλι» της απάντησε η Ιφι. «Πραγματικά, κανείς δεν πιστεύει σε μένα στην Ειδική Υπηρεσία; Και να σου πω κυρά μου, πολύ μεγάλο στόμα έχεις. Δεν με έχεις δει να κάνω τίποτα εκτός από το να σας θεραπεύω με την μαγεία μου, οπότε δεν ξέρω με τι κριτήριο με κρίνετε όλοι σας και...»

«Αρκετά!» φώναξε ο Μαλχη, εκνευρισμένος προφανώς από τον τσακωμό τους. «Προς το παρόν, χρειάζεται να γίνει μία παύση σε αυτή την υπόθεση για προφανείς λόγους. Θα σας βρω κατάσκοπο όσο το δυνατόν γρηγορότερα μην ανησυχείτε. Μέχρι τότε, κάντε ό,τι κάνατε τις υπόλοιπες μέρες, που δεν είχατε αναλάβει σαν ομάδα κάποια υπόθεση ακόμη»

«Να ενημερώσουμε τον Φυτ;» ρώτησε ο Κουπ.

«Κάνε ό,τι θες, λες και με ενδιαφέρει» του απάντησε ο Μαλχη και πήγε στο γραφείο του, αφήνοντας μας με τον Ρομπ στην είσοδο. Εκείνη τη στιγμή μας πλησίασαν δύο νοσοκόμες με ένα φορείο.

«Θα τον αναλάβουμε εμείς αυτόν» είπε μία από τις δύο νοσοκόμες και φόρτωσαν τον Ρομπ για να το πάρουν από την είσοδο, όπου μάλλον ενοχλούσε.

«Γιατί σε μένα... γιατί...» άκουσα την Μπλαιρ να λέει χωρίς να έχει σηκωθεί από το πάτωμα.

«Ξέρεις...» πήγε να της μιλήσει ο Κουπ.

«Τι έγινε Κουπ;» τον διέκοψε η Μπλαιρ.

«Δεν χρειάζεται...»

«Δεν χρειάζεται τι, Κουπ;»

«...τίποτα... άστο» είπε ο Κουπ και έφυγε από δίπλα μας. Τον ακολούθησε και η Ιφι. Εγώ γονάτισα δίπλα στην Μπλαιρ και προσπάθησα να την καθησυχάσω.

«Κοίτα, δεν ξέρω πως νιώθεις αυτή τη στιγμή, αλλά θα ήθελες μήπως να...»

Η Μπλαιρ σηκώθηκε και με αγκάλιασε πριν προλάβω να τελειώσω την πρόταση μου. Ένιωθα τα δάκρυα της να πέφτουν στο πρόσωπο και στους ώμους μου. Μετά από λίγο την αγκάλιασα και εγώ.

\textasciitilde ...δεν πρέπει να χάσω και εσένα Σπαρκ. Σε παρακαλώ, ζήσε. Είσαι ο μόνος άνθρωπος που έχω...\textasciitilde 

Οι σκέψεις της Μπλαιρ εκείνη τη στιγμή ήταν... για μένα; Πώς γίνεται αυτό; Την ξέρω μόνο μερικές εβδομάδες. Λες και αυτή να... αποκλείεται...

\chapter{}

«Καλημέρα Σπαρκ» μου είπε η γραμματέας.

Δεν της απάντησα. Πήγα κατευθείαν στο γραφείο του Μαλχη εκνευρισμένος. Είχαν περάσει δύο μέρες από τότε που χάσαμε τον Ρομπ και δεν είχαμε ακόμη κάποιο νέο για τον αντικαταστάτη του, ούτε όμως για την υπόθεση. Όταν άνοιξα την πόρτα αυτή τη φορά δεν χρειάστηκε να αποφύγω σφαίρα του Μαλχη, τον είδα όμως προβληματισμένος στο γραφείο του να επεξεργάζεται μερικά χαρτιά. Στο πλάι του ήταν η Ιφι, η οποία χαμογελούσε μέχρι τα αυτιά με τις εκφράσεις που έκανε ο Μαλχη.

«Να περάσω;» ρώτησα διακριτικά.

«Έλα εδώ Σπαρκ» είπε ο Μαλχη και μου έγνεψε με το χέρι του να πλησιάσω.

Στα δύο χαρτιά που κράταγε ο αρχηγός υπήρχαν εικόνες από πέντε διαφορετικούς άντρες, τρεις εκ των οποίων είχαν αναγραμμένα τα ονόματα τους κάτω από τις φωτογραφίες. Στο τρίτο χαρτί που ήταν ακουμπισμένο στο γραφείο υπήρχε μία λίστα από δραστηριότητες που ανέγραφε στην κορυφή “Ρουτίνα”.

«Δεν καταλαβαίνω...» είπα στον Μαλχη.

«Θα σου εξηγήσω εγώ» με διέκοψε η Ιφι. «Αυτοί είναι οι πέντε ύποπτοι άντρες στόχοι μας» μου είπε και έδειξε τα δύο χαρτιά που κρατούσε ο Μαλχη. «Στην αριστερή σελίδα είναι οι τρεις άντρες που μένουν στην πόλη, οι δύο εκ των οποίων είναι οι άντρες που συνάντησε ο Μαλχη καθώς επέστρεφε από το καζίνο πριν μία εβδομάδα. Τα ονόματα τους, μπορεί να σου είναι αδιάφορα, αλλά είναι Στριτ Ρατ και Γιατς Κεβιτς. Οι δύο αυτοί άντρες ξεκινούν από την πολυκατοικία τους νωρίς το πρωί και περπατάνε 40 λεπτά μέσα στην πόλη, ντυμένοι ως κουκουλοφόροι, όπως μας είχε περιγράψει ο Μαλχη, μέχρις ότου να φτάσουν στο σπίτι ενός τρίτου κουκουλοφόρου, ο οποίος λέγεται Ματεο. Δεν μπόρεσα να διακρίνω το επίθετο του τρίτου άντρα, καθώς και τις δύο μέρες που τους παρακολούθησα, με το που τον συναντούσαν, έμπαιναν στο αυτοκίνητο του και πήγαιναν σε ένα μικρό μοναχικό σπιτάκι, μισή ώρα οδικώς μακριά από την πόλη. Σε εκείνο το σπίτι, συναντιόντουσαν με τους άλλους δύο άντρες. Αυτών δεν μπόρεσα να βρω τα ονόματα τους επειδή έξω από το σπίτι εκείνο δεν αντάλλαζαν καμία λέξη. Το μόνο χαρακτηριστικό που παρατήρησα ήταν ότι οι τρεις άντρες σήκωσαν και τις δύο μέρες τα μανίκια στο δεξί τους χέρι και κούνησαν τα βραχιόλια τους, τα οποία πρόλαβα να βγάλω φωτογραφία εδώ» είπε η Ιφι και μου έδειξε στο κινητό της το χέρι ενός άντρα, ο οποίος φορούσε υπερβολικά πολλά πολύχρωμα βραχιόλια. «Όλη την υπόλοιπη μέρα, το σπίτι έμενε σχετικά ήσυχο, τουλάχιστον τόσο ήσυχο που δεν μπορούσα να διακρίνω τι λέγανε οι άντρες μέσα στο σπίτι. Όταν νύχτωνε, όμως, τα πράγματα δεν ήταν και τόσο ήσυχα. Χωρίς κάποια ένδειξη, ξαφνικά μέσα από το σπίτι πέντε διαφορετικά χρωματισμένες ακτίνες φωτός εκτινάσσονται προς τον ουρανό. Να, κοίτα...» μου έδειξε πάλι το κινητό της και ήταν πραγματικά ένα θέαμα. Ακτίνες κίτρινου, πράσινου, μπλε, κόκκινου και γκρι φωτός να φτάνουν από τη γη μέχρι τα σύννεφα. Το σπίτι, παρ’ όλ’ αυτά, δεν φαινόταν να είχε υποστεί κάποια ζημιά από αυτές τις ακτίνες φωτός. «Αυτό το φαινόμενο διαρκούσε περίπου μία ώρα» συνέχισε η Ιφι. «Όταν έσβηναν τα φώτα, οι τρεις άντρες από την πόλη επέστρεφαν στο σπίτι του Ματεο με το αυτοκίνητο τους και στη συνέχεια ο Στριτ και ο Γιατς περπατούσαν μέχρι το δικό τους σπίτι. Αυτά τα ολίγα έγιναν αυτές τις δύο μέρες. Ερωτήσεις;» μου είπε η Ιφι και μου χαμογέλασε.

«Έχω εγώ μία ερώτηση» της είπε ο Μαλχη. «Πως ξέρεις πως η “ρουτίνα” τους αυτή είναι καθημερινή;»

«Πιστεύω πως είναι, μία διαίσθηση όπως λες και εσύ» τον ειρωνεύτηκε η Ιφι.

«Αλλά και να είναι η πραγματική ρουτίνα τους αυτή, αυτό δεν σε κάνει κατάσκοπο επειδή κατάφερες να βρεις μερικές πληροφορίες για κάποιους. Μπορεί και ο Ρομπ να είχε βρει τις ίδιες πληροφορίες αλλά να μην πρόλαβε να μας τις μεταφέρει» της συνέχισε ο Μαλχη.

«Ναι αλλά ο Ρομπ έπαιξε με την τύχη του. Έκατσε πολλές μέρες στο πόστο του και τον ανακάλυψαν» του είπε η Ιφι και πήγε στον καναπέ για να κάτσει.

«Ευτυχώς που δεν είναι η Μπλαιρ εδώ, θα έτρωγες κράξιμο...» της είπε ο Μαλχη και σηκώθηκε.

«Δεν με νοιάζει τι θα μου έλεγε, εγώ την δουλειά μου έκανα. Τώρα θέλω να δω πως θα χειριστεί η υπόλοιπη ομάδα την υπόθεση»

«Σιγά που έκανες τη δουλειά σου. Σαχλαμάρες έκανες, μη βρίσω το σπανάκι σου...»

«Μπορεί κάποιος να μου εξηγήσει ήρεμα και πιο αργά τι συμβαίνει;» τους διέκοψα γιατί είχα χάσει τον ρυθμό της συζήτησης.

«Τι δεν καταλαβαίνεις ακριβώς;» με ρώτησε ρητορικά ο Μαλχη. «Η Ιφι κατασκόπευσε τους κουκουλοφόρους, χωρίς την έγκριση μου να σημειώσω, και τώρα μας δείχνει τις πληροφορίες που σύλλεξε»

«Ναι αλλά... γιατί είπε η υπόλοιπη ομάδα;»

«Γιατί νομίζει ότι είναι μέλος της ομάδας σου»

«Δεν μπορώ να γίνω;» πετάχτηκε η Ιφι. «Δεν σας έχω δείξει πως είμαι ικανή κατάσκοπος;»

«Όχι»

«Μαλχη, είναι όντως εντυπωσιακές οι πληροφορίες που μάζεψε η Ιφι σε τόσο μικρό χρονικό διάστημα. Δεν γίνεται να την απορρίπτεις τόσο απλά. Γιατί να μην της δώσεις μία ευκαιρία ως μόνιμο μέλος της ομάδας μου με... “υβριδικές” ικανότητες;» είπα και κοίταξα την Ιφι, οποία είχε κοκκινίσει από το σχόλιο μου.

«Όχι» μου απάντησε ξερά ο Μαλχη.

«Θες να βάλουμε ένα στοίχημα τότε;» του πρότεινα, γιατί ο τζόγος θα ήταν ο μόνος τρόπος που ο Μαλχη θα δεχόταν την πρόταση μου. «Θα αναλάβουμε εγώ, η Ιφι και η ομάδα μου αυτή την υπόθεση. Αν καταφέρουμε να την λύσουμε, χωρίς ταυτόχρονα να έχουμε απώλειες, η Ιφι θα ενταχθεί στην ομάδα μου. Αν δεν τα καταφέρουμε, η Ιφι δεν θα σε ξαναενοχλήσει για αυτό το θέμα. Τι λες;»

«Εντάξει» μου απάντησε ο Μαλχη σαν γνήσιος τζογαδόρος, μετά από λίγη σκέψη βέβαια. «Το αφήνω πάνω σας. Αν θέλετε ξεκινήστε και σήμερα, δεν θα σας κρατήσω καθόλου πίσω. Απλά να σου θυμίζω, στην προηγούμενη αποστολή σου δυσκολευτήκατε να ακινητοποιήσετε ένα άτομο κλάσης Α. Θέλω να δω πραγματικά πως θα χειριστείτε 5 διαφορετικά άτομα»

«Δεν χρειάζεται να αγχώνεσαι για αυτό. Έχουμε την υπόλοιπη ομάδα για τέτοιου είδους δουλειές» του απάντησα χαμογελώντας.

«Καλά, ό,τι νομίζεις. Φύγετε τώρα» μας έδιωξε ο Μαλχη από το γραφείο του.

«Ιφι πάμε στο δικό μου γραφείο για να συζητήσουμε το σχέδιο;» είπα καθώς άνοιγα την πόρτα. Αυτή μου έγνεψε θετικά.

«Ωχ Σπαρκ; Πάλι στο γραφείο του Μαλχη;»

«Δεν έχουμε χρόνο Κουπ... δηλαδή έχουμε, δεν τραυματίστηκε κανείς. Θέλω, όμως, όλη την ομάδα στο γραφείο όσο το δυνατόν συντομότερα»

«Μάλιστα φίλε!» μου χαμογέλασε ο Κουπ.

\chapter{}

«Ξέρετε που θα ήταν ωραία για να σταματήσουμε;»

«Που;»

«Έχει σε κανένα μισάωρο ένα χωριό με ένα μίνι μάρκετ, οπότε μπορούμε να πάρουμε και τίποτα μικρό να φάμε, αν χρειαστεί. Βέβαια πρέπει να περάσουμε πρώτα αυτό το βουνό...»

«Γιατί το είπες τόσο απογοητευμένος;»

«Τίποτα, απλά έχει δύσκολες και επίφοβες στροφές. Αν χάσει κανείς την προσοχή του εδώ, μπορεί να πέσει από κανέναν γκρεμό, για αυτό ήθελα να οδηγήσω εγώ σήμερα. Τα παιδιά δεν έχουν την απαραίτητη εμπειρία οδήγησης για τέτοιες διαδρομές»

«Γιατί, εγώ δεν μπορούσα να οδηγήσω ρε Ριντ;»

«Όχι, ούτε εσύ δεν μπορούσες να οδηγήσεις αδελφέ, επειδή είσαι αδέξιος. Μία λάθος κίνηση σου και θα μας έστελνες στην μετά θάνατον ζωή»

«Καλά καλά, εγώ δεν σε πρόσβαλα»

«Δεν ήθελα να... συγγνώμη αν σε πρόσβαλα»

«Δεν πειράζει ρε αδερφέ, πλάκα κάνουμε. Ακόμη να μάθεις το χιούμορ μου»

«Είναι λίγο δύσκολο θείε να καταλάβουμε πότε κάνεις πλάκα και πότε αστειεύεσαι, οπότε...»

«Οπότε τι;»

«Δεν ξέρω, απλά ίσως να είσαι πιο... εκφραστικός»

«Τι εννοείς όταν λες εκφραστικός ρε Βασι;»

«Εννοώ να αλλάζεις κάπως τον τόνο της φωνής σου, ξέρω γω, ή τις κινήσεις των χεριών σου, ώστε να καταλαβαίνει ο άλλος πότε κάνεις πλάκα»

«Δίκιο έχει μπαμπα. Και εγώ στο έχω πει αρκετές φορές ότι δεν καταλαβαίνω πότε κάνεις πλάκα, αλλά με αγνοείς...»

«Πότε σε αγνόησα εγώ βρε Ρομπ;»

«Θες να αρχίσω να απαριθμώ τις ημερομηνίες;»

«Ναι»

«Που να τις θυμάμαι τώρα βρε μπαμπά...»

«Βλέπεις»

«Τι; Το ότι δεν θυμάμαι τις ημερομηνίες δε σημαίνει ότι έχεις δίκιο»

«Ρομπ, Χερμπ, σταματήστε. Δεν χρειάζεται κάθε συζήτηση μας να οδηγείται σε τσακωμό»

«Καλά...»

«Κοίτα ποιος μιλάει. Ριντ, εσύ δεν έχεις αρχίσει σήμερα τις περισσότερες διαφωνείες μας;»

«Όχι»

«Τι “όχι”. Αφού... κι άλλο τηλέφωνο;»

«Δεν είναι το δικό σου αδερφέ, αφού έχει άλλον ήχο κλήσης. Μάλλον είναι του Βασι, γιατί μου ακούγεται γνώριμο»

«Ναι, δικό μου είναι, μισό λεπτό. Παρακαλώ... Ναι... Α, εσύ είσαι. Συγγνώμη, περνάμε από ένα βουνό τώρα, δεν πιάνουμε καλό σήμα...»

«Ποιος είναι;»

«Ένας φίλος μου και του αδερφού μου. Τι σε νοιάζει θειε;»

«Τίποτα, απλά είμαι περίεργος...»

«...ναι είναι και ο αδελφός μου εδώ... να σε βάλω ανοιχτή ακρόαση; Θα σε ακούσουν όλοι... εντάξει αφού δεν σε πειράζει»

«Με έβαλες;»

«Ναι, για πες μας τώρα τι θες Κουπ. Σχετίζεται με αυτό που ήθελε να μας πει η Μπλαιρ όταν γυρνούσαμε;»

«Όχι...»

«Τι έγινε τότε;»

«Οι γονείς μου... πως να το πω τώρα... δεν είμαι καλός με τα λόγια... τους πάτησε φορτηγό»

«Τι;»

«Τι;»

«Ναι... δεν... ξέρω πως να... έχασα και την αδελφή μου πριν τρεις μήνες και... τώρα...»

«Έλα ρε Κουπ, μην κλαις. Αφού έχεις εμάς. Όταν γυρίσουμε από το ταξίδι μας θα είμαστε δίπλα σου»

«Μα έχω μείνει... μόνος. Δεν έχω πλέον οικογένεια...»

«Αυτό σου λέω βρε χαζούλη, θα γίνουμε εμείς η οικογένεια σου. Κάνε λίγη υπομονή και θα είμαστε δίπλα σου»

«Δεν ξέρω αν μπορώ να κάνω υπομονή πλέον...»

«Μη τολμήσεις και κάνεις αυτό που σκέφτομαι. Απλά μην, δεν αξίζει. Αν χρειαστείς συντροφιά αυτές τις μέρες, η μητέρα μας είναι στο σπίτι μας...»

«Δεν είναι ρε Βασι, μη του λες ψέματα του ανθρώπου τέτοια ώρα»

«Αχ, ναι... συγγνώμη Κουπ. Ξέχασα ότι η μαμά είναι σε επαγγελματικό ταξίδι πάλι. Αυτό όμως δεν αναιρεί τίποτα. Αν χρειαστείς κάποιον να μιλήσεις, εμείς θα είμαστε πάντα εκεί για σένα, έτσι δεν είναι αδερφέ;»

«Ναι... αυτό που είπε ο Βασι»

«Σας ευχαριστώ παιδιά. Δεν σας πήρα όμως μόνο για αυτό. Ήθελα επίσης να σας πω επίσης να προσέχετε στο δρόμο. Δεν θέλω να χάσω και τους πιο κοντινούς μου φίλους από την παρέα...»

«Πες του ότι εμείς οδηγούμε, δεν περπατάμε, οπότε δεν γίνεται να μας πατήσει φορτηγό»

«Αμάν ρε Χερμπ»

«Τι; Άλλαξα τον τόνο της φωνής μου τώρα. Δεν ήταν καλή η στιγμή, ε;»

«Πάλι καλά που το κατάλαβες και μόνος σου...»

«Δεν πρόκειται να πάθουμε κάτι Κουπ, μην αγχώνεσαι. Οδηγάει ο πατέρας οπότε...»

«Κοίτα δω ενοχλητικό πλάσμα... Έλα εδώ να σε σκοτώσω βρε μυγάκι»

«Κατσίκα, Ριντ... Ριντ; ΡΙΝΤ ΤΙ ΚΑΝΕΙΣ...»

\chapter{}

«Είσαι έτοιμος Κουπ;»

«Ναι. Εσύ το έχεις το κουμπί κάτω από τα ρούχα σου;»

«Ναι βεβαίως. Να τους χτυπήσω;»

«Χτύπα» μου είπε ο Κουπ και χτύπησα την πόρτα του σπιτιού.

Καθώς άνοιξε η πόρτα, από μέσα βγήκαν ένα σωρό ποντίκια, τα οποία έτρεξαν προς το δάσος. Ο άντρας που μας άνοιξε μας κοίταξε προβληματισμένος. Εμείς κάναμε την χαρακτηριστική κίνηση που περιέγραψε η Ιφι, σηκώσαμε τα δεξιά μας μανίκια και κουνίσαμε τα βραχίολια μας. Ο άντρας μας άφησε να μπούμε μέσα στο σπίτι, όμως το πρόσωπο του είχε ακόμη υποψίες για εμάς.

Το σπίτι δεν είχε και πολλά πράγματα μέσα. Όταν μπήκαμε μπορούσαμε να διακρίνουμε το σαλόνι, στο οποίο βρίσκονταν δύο άντρες, ο ένας στον καναπέ και ο άλλος χάιδευε ένα ποντίκι στο πάτωμα. Διαγώνια δεξιά ήταν η κουζίνα του σπιτιού στην οποία βρισκόταν ένας τρίτος άντρας και μαγείρευε. Τέλος, υπήρχε ένα στρογγυλό τραπέζι στην αριστερή πλευρά του σπιτιού με πέντε καρέκλες, στο οποίο καθόταν ο τελευταίος άντρας.

\textasciitilde ...από που ήρθαν αυτοί τώρα; Θα δω αν τους ξέρει κανένας...\textasciitilde \  προσπάθησα να διαβάσω τις σκέψεις του άντρα που μας έβαλε μέσα.

«Ποιοι είναι αυτοί ρε Γιατς; Γιατί τους άφησες να μπουν;» είπε ένας άντρας που καθόταν στον καναπέ.

«Έχουν και αυτοί βραχιόλια, οπότε υπέθεσα πως ξέρουν και αυτοί τη γυναίκα, έτσι δεν είναι;» μας απευθύνθηκε ο άντρας που μας άνοιξε την πόρτα

«Εμείς... τη γυναίκα...» δεν ήξερα τι να πω.

«Δεν γνωρίζουμε για ποια γυναίκα μιλάτε φίλτατοι» είπε ο Κουπ. «Εμείς είμαστε αγρότες από ένα διπλανό χωριό. Ήρθαμε σε εσάς γιατί η ατυχία μας είναι τρομερή τον τελευταίο χρόνο. Βλέπετε, το αγρόκτημα μας έχει καεί δύο φορές τους τελευταίους μήνες. Την πρώτη φορά χάσαμε όλη μας τη σοδειά και την δεύτερη όλα μας τα ζώα»

«Άρα δεν αναμειγνύονται στα γεγονότα της πόλης...» είπε ο άντρας που χάιδευε το ποντίκι.

«Βεβαίως και όχι. Η μόνη σχέση μας με την πόλη είναι οι συναλλαγές που κάνουμε για να μπορέσουμε να βγάλουμε το ψωμί μας, αν με πιάνετε» του απάντησε ο Κουπ.

«Και γιατί ήρθατε εδώ; Πώς μας βρήκατε; Και ποιος σας έδωσε αυτά τα βραχιόλια;»

«Ε...» κόλλησα πάλι.

«Βλέπετε, ο φίλος μου και εγώ, αφού είδαμε πως η τύχη μας ήταν τόσο άτυχοι, είπαμε να στραφούμε σε μεταφυσικές οντότητες, στην θρησκεία δηλαδή» με έσωσε πάλι ο Κουπ. «Σας είδαμε από μακριά ένα πρωί που μαζευτήκατε σε αυτό το σπίτι φορώντας χιτώνες και κουκούλες και είπαμε μήπως μπορείτε να μας βοηθήσετε»

«Ενδιαφέρον η Ιστορία σας» είπε και σηκώθηκε ο άντρας που καθόταν στον καναπέ. «Βέβαια, δεν ξέρω αν μπορούμε να σας βοηθήσουμε. Βλέπετε, είμαστε όντως μία αίρεση, αλλά δεν προσκυνάμε κάποιον θεό, ή τουλάχιστον δεν πιστεύουμε ότι είναι θεά, αλλά οι μαρτυρίες μας μπορεί να πείσουν κάποιον ότι είναι...»

«Τι εννοείται με αυτό, καλέ μου κύριε;» τον ρώτησε ο Κουπ από περιέργεια.

«Βλέπετε...» κοίταξε τους υπόλοιπους άντρες για να πάρει την έγκριση τους να συνεχίσει. Δεν τον σταμάτησε κανείς, οπότε ήρθε πιο κοντά μας και άρχισε να μας εξηγεί. «Η γυναίκα, που ανέφερε πιο πριν ο φίλος μου που σας άνοιξε, είναι το... “είδωλο” μας, δεν ξέρω πως να το περιγράψω. Ίσως καλύτερη ονομασία θα ήταν: η “μητέρα” μας. Πριν περίπου είκοσι χρόνια, όταν ήμασταν όλοι μας ακόμη μωρά στο μαιευτήριο, το νοσοκομείο που μας φιλοξενούσε έπιασε φωτιά. Η μητέρα μας έσωσε και μας μεγάλωσε σε αυτό εδώ το σπίτι, μαζί με όσα άλλα παιδιά είχε σώσει, μέχρι την ηλικία τον οκτώ. Μία μέρα όμως, εξαφανίστηκε χωρίς καμία προειδοποίηση, και άφησε εμάς τα 10, ίσως και 15, παιδιά πίσω στο “σπίτι” της. Μας είχε μάθει βέβαια πολλά πράγματα, για την ζωή αλλά και για την μαγεία, οπότε μπορούσαμε να ζήσουμε μόνοι μας. Μερικά παιδιά, βέβαια, γιατί βλέπετε δεν βγάζαμε και πολλά λεφτά, έγιναν εγκληματίες, δολοφόνοι, για να μπορέσουμε να βγάλουμε και εμείς το ψωμί μας, όπως είπατε και σεις. Αυτοί, όπως βλέπετε, δεν είναι πλέον μαζί μας, μερικοί έχουν πεθάνει και άλλοι είναι στη φυλακή»

«Και η “αίρεση” σας τι κάνει ακριβώς;» ρώτησε συμπληρωματικά ο Κουπ, επειδή είδε πως ο άντρας μονολογούσε χωρίς να φτάνει σε κάποιο συμπέρασμα.

«Τίποτα το ιδιαίτερο... το μόνο “παράξενο”, θα έλεγε κανείς, είναι το βράδυ που...» κοίταξε πάλι τους γύρω του για να πάρει την έγκριση τους. Δεν τον σταμάτησε κανείς. «...που προσπαθούμε να την καλέσουμε πίσω...»

«Τι εννοείς με αυτό;» προβληματίστηκε ο Κουπ. «Μήπως αυτά τα περίεργα φώτα το βράδυ έχουν να κάνουν κάτι με αυτό;»

«Αυτά τα φώτα είναι ακριβώς αυτό. Μπορεί να ξέρετε ήδη για την αύρα της μαγείας που όλοι έχουμε μέσα μας αλλά μόνο οι άνθρωποι κλάσης Α μπορούν να την αισθανθούν και να την χρησιμοποιήσουν»

«Ναι έχουμε ακούσει. Βέβαια, σαν κλάσης Β, εγώ και ο φίλος μου μπορεί να μην καταλάβουμε και τόσο τη συνέχεια, γιατί δεν έχουμε νιώσει ποτέ αύρα μαγείας» αστειεύτηκε ο Κουπ.

«Δεν πειράζει, θα σας το πω λίγο πιο απλοϊκά τότε, αφού το έφερε η κουβέντα. Αν συγκεντρώσεις την αύρα της μαγείας σου στην άκρη του σώματος σου, μπορείς να δεις το χρώμα της. Να, κάπως έτσι...» μας είπε ο άντρας και ένα κόκκινο φως άρχισε να λάμπει γύρω του.

Μου θύμισε την μέρα που η μητέρα μου μου είχε πει πως οι κεραυνοί μου ήταν ασταθείς και όταν κοιτάχτηκα στον καθρέφτη, είχα ένα κίτρινο στρώμα κεραυνών γύρω μου. Αυτός όμως δεν είχε κάποια μαγεία γύρω του, απλά ένα κόκκινο φως, το οποίο μας ενοχλούσε λίγο τα μάτια.

«Αν, τώρα, εξάγω από τα χέρια μου την αύρα...» μας είπε και άρχισε να λάμπει παραπάνω το χέρι του «...μπορώ να δημιουργήσω μια πηγή φωτός η οποία μπορεί να φτάσει τα 300.000 λούμεν, για να σας το πω με φυσικούς όρους, μήπως το καταλάβετε καλύτερα, το οποίο μπορεί να φτάσει μέχρι τα σύννεφα» είπε και έδειξε το ταβάνι. Όταν τελείωσε την πρόταση του, όλη του η αύρα που είχε γύρω του και μας τύφλωνε εξαφανίστηκε. «Χρησιμοποιούμε όλοι μαζί, λοιπόν, αυτά τα φώτα με την ελπίδα ότι θα μας δει η μητέρα και θα γυρίσει πίσω. Ερωτήσεις;»

«Ναι» του είπε ο Κουπ. «Γιατί είναι τόσο σημαντική αυτή η μητέρα για εσάς;»

«Μα καλά χαζός είσαι;» νευρίασε ο άντρας.

«Όχι, άκουσα την ιστορία σου, απλά... 12 χρόνια προσπαθείτε να την καλέσετε πίσω με αυτό το κολπάκι και δεν έχει επιστρέψει. Δεν πιστεύεται ότι είναι λίγο υπερβολή; Τι το ιδιαίτερο έχει αυτή η γυναίκα;» ρώτησε ο Κουπ.

«Τι το ιδιαίτερο;» κοίταξε τους υπόλοιπους άντρες για να πάρει έγκριση να συνεχίσει. «Έχει... επτά μαγείες, ίσως και περισσότερες»

«Τι;» φωνάξαμε εγώ και ο Κουπ επειδή αιφνιδιαστήκαμε. «Πώς γίνεται αυτό;» συμπλήρωσε ο Κουπ.

«Δεν θα σας πω όλα μας τα μυστικά σε μία μέρα. Αν θέλετε ελάτε και αύριο για να μάθετε. Πάντως για αυτό σας είπα ότι θα νομίζεται ότι είναι θεά, αλλά στην πραγματικότητα δεν είναι. Υπάρχει...»

«Σταμάτα Ματεο, πολύ μίλησες» τον διέκοψε ο άντρας που χάιδευε το ποντίκι στο σαλόνι.

«Και να σας ρωτήσω» πετάχτηκε ο Κουπ. «Εσείς δεν είστε κάποιοι εγκληματίες, όπως τα... “αδέλφια” σας;»

Οι άντρες τον κοίταξαν.

«Όχι βέβαια...» \textasciitilde ...αν εξαιρέσεις όλες τις ληστείες που έχω κάνει με τον Γιατς και τον φόνο που έκανα επί πληρωμή για έναν πλούσιο...\textasciitilde \  έλεγαν οι σκέψεις του άντρα που μας εξηγούσε τόση ώρα την ιστορία του.

«Πολύ ωραία» είπε ο Κουπ και το σπίτι ησύχασε.

Προσπάθησα  να διαβάσω και τις σκέψεις τον άλλων ανδρών, ξεκινώντας από αυτόν που μαγείρευε και αυτόν που καθόταν στο τραπέζι.

\textasciitilde ...μπραβο Ματεο. Η αμηχανία σου θα μας δώσει...\textasciitilde 

\textasciitilde ...εγκληματίες είμαστε βέβαια, απλά σιγά μην σας το πούμε κιόλας. Σας ξέρουμε μισή ώρα και θες να μάθεις και τα προσωπικά μας...\textasciitilde 

Άρα όντως κάτι κρύβουν αυτοί οι άντρες. Ο Κουπ είχε χαλαρώσει και θα μου έκανε νεύμα να φύγουμε, αλλά του έκανα νεύμα να κάτσουμε λίγο ακόμη. Το σχέδιο μας ήταν απλό: εγώ και ο Κουπ θα πιάναμε συζήτηση με τους άντρες για να συμπεράνουμε αν είναι ένοχοι κάποιου εγκλήματος ή όχι. Αν χρειαζόμασταν ενισχύσεις, θα είχα ένα κουμπί πάνω μου το οποίο θα έδινε σινιάλο στους υπόλοιπους που θα βρίσκονταν απ’ έξω. Πήγα ασυναίσθητα να πιάσω το κουμπί που βρισκόταν στη δεξιά μου τσέπη, αλλά δεν βρισκόταν εκεί.

«Μήπως ψάχνεις αυτό;» με ρώτησε ο άντρας που καθόταν στο πάτωμα και χάιδευε το ποντίκι, δείχνοντας μου ένα φίδι που είχε στο στόμα του το κουμπί μου. «Δεν μοιάζει με κινητό... Τι είναι αυτό για πες μας;» \textasciitilde ...να δούμε αν είναι πράκτορες τελικά αυτοί...\textasciitilde 

«Αυτό είναι...» δεν ήξερα τι να πω.

Ο άντρας με το ποντίκι έγνεψε θετικά στον άντρα που μας άνοιξε την πόρτα, ο οποίος βρισκόταν από πίσω μου. Προσπάθησα να διαβάσω τις σκέψεις του \textasciitilde ...και τώρα η μαχαιριά...\textasciitilde 

«Σπαρκ όχι!» με έσπρωξε ο Κουπ και έπεσα πάνω στον άντρα με το ποντίκι. «Πάτα το» μου φώναξε, εννοώντας το κουμπί που έτρωγε το φίδι, αλλά ήταν ήδη αργά. Είχε χαλάσει ο μηχανισμός του κουμπιού και όσες φορές και να το πατούσα δεν έκανε τον ήχο που έπρεπε.

«ΣΠΑΡΚ ΤΙ ΚΑΝΕΙΣ;» μου είπε ο Κουπ και τέντωσε το χέρι του με την μαγεία ελαστικότητας του για να με πιάσει και να με φέρει κοντά του. «Πάμε έξω» είπε και έσπασε την πόρτα. Οι άντρες τον ακολούθησαν, ρίχνοντας του βολές με τις μαγείες τους, τις οποίες ο Κουπ απέφυγε.

«Σπαρκ, Κουπ; Γιατί δεν πατήσατε...» πήγε να πει ο Φυτ.

«Δεν έχουμε χρόνο» είπε ο Κουπ και με άφησε πίσω από το αυτοκίνητο. «Ελάτε γρήγορα, έχουμε μάχη»

«Ο Σπαρκ είναι καλά;» ρώτησε η Μπλαιρ.

«Άστον με την Ιφι για τώρα. Θα δούμε τι θα κάνουμε με αυτόν σε λίγο» είπε ο Κουπ και έβγαλε ένα ζευγάρι χειροπέδες που είχε κρυμμένο κάτω από τα ρούχα του και πήγε προς τους άντρες, αποφεύγοντας πάλι τις βολές μαγείων που του έριχναν. Η Μπλαιρ και ο Φυτ τον ακολούθησαν.

«Σπαρκ... όλα καλά;» με ρώτησε η Ιφι.

Δεν της απάντησα. Είχα μείνει άφωνος. Πώς κατάφερε να μου κλέψει το κουμπί; Γιατί δεν το κατάλαβα; Γιατί ήμουν... άχρηστος; Μέσα στο σπίτι δεν κατάφερα να αρθρώσω ούτε λέξη από το άγχος μου. Όλη τη δουλειά την έκανε ο Κουπ. Το μόνο που έκανα ήταν να διαβάσω τις σκέψεις των ανδρών και να μας βάλω σε ακόμη πιο δύσκολη κατάσταση από ότι θα έπρεπε να βρισκόμαστε. Ίσως... δεν είμαι κατάλληλος για ανακριτής... Ο Κουπ καλύτερη δουλειά έκανε από εμένα στην εξαγωγή πληροφοριών από τους άντρες, χωρίς να είναι η ειδικότητα του.

«Σπαρκ, μίλα μου» με έπιασε από τα μπράτσα η Ιφι κλαίγοντας. \textasciitilde ...πες μου τι έγινε μέσα στο σπίτι, Σπαρκ...\textasciitilde 

Δεν της απάντησα.

«Αχ!»

«Φυτ!» άκουσα την Μπλαιρ να λέει και την είδα να πηγαίνει μπροστά του με μία ασπίδα για να τον προστατέψει.

«Είναι πάρα πολλοί» άκουσα την φωνή του Κουπ από μακριά.

«Σπαρκ... κάνε κάτι... σε παρακαλώ»\textasciitilde ...πρέπει να πάω να θεραπεύσω τον Φυτ...\textasciitilde \  είπε η Ιφι και με παράτησε πίσω από το αυτοκίνητο.

Εκτός από άχρηστος στην ανάκριση, είμαι και άχρηστος στην μάχη. Το μόνο που μπορούσα να κάνω ήταν να πυροβολήσω τους άντρες με το πιστόλι μου...

«Επιτέλους σηκώθηκες...» μου είπε η Μπλαιρ ενώ καθόταν μπροστά από την Ιφι και τον Φυτ με την μεγάλη της ασπίδα στην οποία είχε μάλλον δώσει λίγη από την αύρα της, καθώς οι βολές των ανδρών που την χτύπαγαν γίνονταν σκόνη. «Θα κάνεις κάτι ή...»

Εκείνη τη στιγμή έβγαλα το πιστόλι μου και άρχισα να πυροβολάω τους άνδρες, αυτοί όμως απέφευγαν τις σφαίρες μου με ευκολία. Ο Κουπ εκείνη τη στιγμή έπιασε έναν από τους τέσσερις άντρες και του έβαλε χειροπέδες. Πάλι άχρηστος είμαι. Δεν μπορώ ούτε με το πιστόλι μου να τους πετύχω.

«Σπαρκ, συνέχισε να ρίχνεις!» μου φώναξε ο Κουπ, ο οποίος έκανε μανούβρες γύρω από τις σφαίρες μου και τις μαγείες των ανδρών για να μπορέσει να τους πλησιάσει.

Να συνεχίσω να ρίχνω, ε; Μα γιατί; Τι θα κερδίσουμε με αυτόν τον τρόπο;

«Μα καλά δεν μπορείς να διαβάσεις τις κινήσεις τους με την μαγεία σου;» με ρώτησε η Μπλαιρ.

Να διαβάσω τις κινήσεις τους; Μέσω των σκέψεων τους; Πώς δεν το είχα σκεφτεί ποτέ αυτό; Είμαι όντως άχρηστος... Ας το δοκιμάσω όμως.

\textasciitilde ...έρχεται για μένα τώρα αυτός ο αγρότης. Θα κάνω μερικά βήματα πίσω και θα του ρίξω φωτιά στα μάτια...\textasciitilde 

Πυροβόλησα λίγο πιο πίσω από εκείνον τον άντρα και τον πέτυχα στο αριστερό του πόδι, όταν έκανε το πίσω βήμα που έλεγαν οι σκέψεις του. Ο Κουπ τύλιξε το χέρι του γύρω από τον άντρα και του έβαλε χειροπέδες.

«Μπράβο Σπαρκ. Άλλοι τρεις» μου είπε ο Κουπ και πήδηξε σαν ελατήριο προς τον άντρα με το ποντίκι στο χέρι του.

\textasciitilde ...έλα πάνω μου και θα το μετανιώσεις, έχω δύο δηλητηριώδης φίδια που θα σε δαγκώσουν στα πόδια τη στιγμή που έρθεις κοντά μου...\textasciitilde 

Πήρα το ρίσκο και πυροβόλησα μερικά μέτρα κάτω από τα πόδια του Κουπ και, ευτυχώς για μένα, δύο φίδια εμφανίστηκαν μέσα από το χώμα και πήγαν προς τον Κουπ, χωρίς να τον φτάσουν βέβαια γιατί οι σφαίρες μου τα πέτυχαν. Ο Κουπ κλώτσησε με την φόρα που είχε από το άλμα του τον άντρα με το ποντίκι και τον έβαλε κάτω.

«Άλλοι δύο!» μου φώναξε ο Κουπ ενώ του έβαζε χειροπέδες. «Αν μπορείτε φέρτε μου κι άλλες χειροπέδες...»

«Δεν χρειάζεται» άκουσα τον Φυτ από πίσω μου και είδα τις ρίζες του να τυλίγουν τους άλλους δύο άντρες, ακινητοποιώντας τους τελείως. «Ρίξε τους και δύο σφαίρες στα πόδια ρε Σπαρκ για να είμαστε σίγουροι πριν τους βάλουμε χειροπέδες»

«Μάλιστα...» του απάντησα και πυροβόλησα τους δύο τυλιγμένους άντρες στα πόδια. «Είσαι...»

«Ναι καλά είμαι. Ευτυχώς που ήρθε και η Ιφι μαζί, γιατί μπορεί θα μπορούσαμε να είχαμε μεγαλύτερο θέμα στα χέρια μας...»

«Ναι αλλά είσαι ζωντανός, οπότε μην λες βλακείες» τον διέκοψε η Μπλαιρ και του έδωσε δύο ζευγάρια χειροπέδες. «Άντε τελειώνετε για να γυρίσουμε πίσω τώρα που είναι ακόμη μεσημέρι»

«Γιατί έχεις κάποιο ραντεβού;» την ρώτησε ο Φυτ παιχνιδιάρικα καθώς πήγαινε προς τον Κουπ και τους άντρες.

«Και να είχα, δεν... τέλος πάντων. Πάω να ξεκινήσω το αμάξι. Ιφι εσύ πήγαινε φέρε το βανάκι που είναι πιο πίσω στο δάσος για να τους βάλουμε μέσα»

«Εντάξει. Σπαρκ, μήπως θες...»

«Ο Σπαρκ θα πάει να βοηθήσει τον Κουπ και τον Φυτ να συνοδέψουν τους άντρες στο βανάκι»

«Καλά...» είπε η Ιφι νευριασμένη και έφυγε. Εγώ κοίταξα την Μπλαιρ.

«Τι με κοιτάς, πήγαινε βοήθα» μου είπε και άνοιξε την πόρτα του αυτοκινήτου.

Όταν πλησίασα τον Κουπ και τον Φυτ, υπήρχε ένα πρόβλημα. Φαινόταν λες και όλοι οι άντρες είχαν λιποθυμήσει. Έβλεπα τον Κουπ να έχει πιάσει έναν άντρα από το κολάρο και να του μιλάει, αλλά ο άντρας ήταν αναίσθητος.

«Τι έγινε ρε φίλε» είπε ο Κουπ. «Ξαφνικά τους είδα να κάνουν κάτι περίεργες κινήσεις με τα χέρια τους και στη συνέχεια έπεσαν όλοι τέζα στο πάτωμα λέγοντας “Μητέρα, ελπίζω να μας ακούσεις” ή κάτι τέτοιο»

«Νιώθω ότι κάπου το έχω ξαναδεί αυτό...» είπε ο Φυτ «...απλά δεν ξέρω αν ο τύπος είχε κάνει κάποια χειρονομία, γιατί είδα μόνο το αποτέλεσμα. Κάτσε να δούμε αν έχουν παλμό» είπε και έπιασε τον καρπό ενός άντρα που ήταν στο πάτωμα.

«Τι έγινε;» ρώτησα προβληματισμένος.

«Όσα ξέρεις εσύ ξέρουμε και εμείς» μου είπε ο Κουπ. «Τους είδαμε να κάνουν κάτι κινήσεις με τα χέρια τους και να πέφτουν...»

«Ναι αυτό το άκουσα. Εννοώ, είναι... ζωντανοί;»

«Όχι» μου είπε ο Φυτ. «Τουλάχιστον αυτοί οι δύο που έλεγξα, αλλά μάλλον ούτε οι άλλοι είναι. Τουλάχιστον αυτή τη φορά δεν τραυματίστηκε σοβαρά κανείς από εμάς...»

«Ναι...» είπα και κοίταξα την Μπλαιρ. «Μήπως να...»

«Ναι θα τους πάρω εγώ δύο-δύο» είπε ο Κουπ και τύλιξε τα χέρια του γύρω από δύο άντρες για να τους κουβαλήσει μέχρι το βανάκι. «Αν μπορεί κάποιος να φέρει και τον πέμπτο θα με ευχαριστούσε πολύ»

«Θα τον πάρω εγώ Σπαρκ, μην αγχώνεσαι» μου είπε ο Φυτ και σήκωσε έναν άντρα στους ώμους του. «Όπα! Πάμε στο βανάκι τώρα»

«Σίγουρα δεν...»

«Βοήθησες αρκετά, δεν χρειάζεται» μου είπε ο Φυτ και μου χαμογέλασε.

Βοήθησα αρκετά; Δεν ξέρω αν το λέει για καλό ή για κακό. Εξάλλου εγώ προκάλεσα όλον αυτό τον πανικό εδώ έξω, αλλά και μέσα στο σπίτι των ανδρών, λόγω της αφέλειας μου. Ίσως όντως να έκανα χειρότερη την κατάσταση με την παρουσία μου. Χωρίς εμένα, ίσως να πήγαινε πιο ομαλά η αποστολή...

«Γιατί είσαι ακόμη εδώ Σπαρκ;» μου είπε ο Κουπ που επέστρεψε για να πάρει τους δύο εναπομείναντες άντρες.

«Κάτι σκεφτόμουν...»

«Σκέφτεσαι και στο αμάξι, πάμε» μου είπε και σήκωσε τους άντρες.

«Πάμε...»

\chapter{}

«...και, μόλις τους βάλαμε χειροπέδες και χαλαρώσαμε, αυτοί κάνανε κάποιες περίεργες κινήσεις με τα χέρια τους και είπαν “Μητέρα, μακάρι να μας ακούσεις” ή κάτι τέτοιο» εξήγησε ο Κουπ στον Μαλχη.

«Χμ» προβληματίστηκε ο Μαλχη.

«Δεν είναι όντως περίεργο;» συμπλήρωσε ο Φυτ. «Οι εγκληματίες που έχουμε πιάσει στις τελευταίες δύο αποστολές να πέφτουν ξεροί χωρίς κάποια εξήγηση»

«Μπορεί να υπάρχει εξήγηση...» του απάντησε ο Μαλχη και άρχισε να ψάχνει στα συρτάρια του γραφείου του.

«Τι ψάχνεις ρε Μαλχη;» τον ρώτησε η Ιφι.

«Είστε μικροί εσείς, δεν πρόκειται να το θυμάστε...» είπε και έβγαλε έναν φάκελο από το πάνω πάνω συρτάρι του.

«Τι είναι αυτό;» ρώτησε πάλι η Ιφι.

«Αυτό είναι ο φάκελος της υπόθεσης Εμπρισμού Νοσοκομείων. Πριν είκοσι χρόνια, τρία νοσοκομεία στην πόλη έπιασαν φωτιά την ίδια μέρα. Πότε σας είπαν οι κουκουλοφόροι πως τους έσωσε η “μητέρα” τους από το μαιευτήριο;»

«Πριν είκοσι χρόνια...» ψιθύρισα.

«Και πόσα παιδιά σας είπαν πως ήταν;»

«Αυτό... δεν παρατήρησα... πάντως λιγότερα από είκοσι σίγουρα» του απάντησα.

«Χμ...» σκέφτηκε ο Μαλχη. «Οι αριθμοί και οι τοποθεσίες δεν ταιριάζουν»

«Τι εννοείς δεν ταιριάζουν;» τον ρώτησε η Ιφι. «Αφού οι χρονολογίες είναι ίδιες, γιατί σε προβληματίζει τόσο; Μη μου πεις πως την ίδια χρονιά κάηκαν και άλλα νοσοκομεία;»

«Όχι δεν είναι αυτό, απλά...» σκέφτηκε για λίγο ο Μαλχη. «Θα μπορούσε όντως να είναι αυτό που λες»

«Δηλαδή;»

«Δεν εννοώ ότι κάηκαν κι άλλα νοσοκομεία, αυτό θα το καταλαβαίναμε. Εννοώ πως τα νοσοκομεία δεν κάηκαν την ίδια μέρα»

«Και πάλι δεν σε καταλαβαίνω» αγανάκτησε η Ιφι.

«Οι συνολικοί αριθμοί που έχω είναι πως περίπου 40 παιδιά χάθηκαν στις πυρκαγιές. Όμως τα νοσοκομεία που έπιασαν φωτιά ήταν σε διαφορετικές μέρες μέσα στη χρονιά. Μήπως ρωτήσατε τους κουκουλοφόρους αν θυμούνται να αυξάνεται το πλήθος τους σε εκείνο το σπίτι που τους μεγάλωσε η “μητέρα” τους;

«Δε... ρωτήσαμε» του απάντησα απογοητευμένος.

«Δεν πειράζει. Αν ισχύει όμως αυτό που είπα, σημαίνει ότι η γυναίκα αυτή είχε πολλά τέτοια σπίτια στα οποία μεγάλωνε μικρά παιδιά»

«Δεν καταλαβαίνω που θες να το πας Μαλχη...» του είπα και κοίταξα γύρω μου τους υπόλοιπους από την ομάδα μου οι οποίοι ήταν εξίσου προβληματισμένοι.

«Λοιπόν, πρέπει να σας εξηγήσω όλη την ιστορία, γιατί οι περισσότεροι έχετε χαθεί. Είπαμε, πριν 20 χρόνια έπιασαν φωτιά τρία νοσοκομεία και χάθηκαν 40 μωρά. Μέχρι εδώ έχετε καταλάβει όλοι;»

Του γνέψαμε όλοι θετικά.

«Ωραία. Οι έρευνες μας πριν 20 χρόνια δεν μπόρεσαν να καταλήξουν σε κάποιο συμπέρασμα. Θεωρήσαμε πως τα μωρά είχαν καεί στην φωτιά ή κάτι τέτοιο, που κατά την γνώμη μου ήταν χαζό, αλλά τι να κάνουμε, αφού δεν είχαμε άλλα στοιχεία...» σταμάτησε για λίγο ο Μαλχη και κοίταξε για λίγο τον φάκελο που είχε ανοιχτό, για να θυμηθεί μάλλον την συνέχεια της υπόθεσης. «Πριν... 12 χρόνια, το ίδιο περιστατικό ξανασυνέβει. Υπήρξαν πυρκαγιές στα ίδια τρία νοσοκομεία. Πάλι χάθηκαν γύρω στα 30 με 40 μωρά. Αυτή τη φορά, όμως, ήμασταν προετοιμασμένοι. Είχαμε μόνιμους κατασκόπους να περιπολούν αυτά τα τρία νοσοκομεία, αλλά και τα υπόλοιπα νοσοκομεία στην πόλη και στα διπλανά χωριά. Σε ένα από αυτά τα νοσοκομία, ο κατάσκοπος μας είχε καταφέρει να εντοπίσει την γυναίκα που... “έσωσε”; “απήγαγε”; Δεν ξέρουμε ακόμη. Πάντως την ακολούθησε μέχρι την εγκαταλελειμμένη εκκλησία της πόλης, αυτή που είναι κοντά στο καζίνο αν ξέρετε...»

Τον κοιτάξαμε όλοι προβληματισμένοι, γιατί προφανώς κανείς μας δεν είχε πλησιάσει αυτή τη μεριά της πόλης.

«Τέλος πάντων, δεν πειράζει... Δυστυχώς αυτός ο κατάσκοπος δεν είναι πλέον μαζί μας, αλλά ευτυχώς για εμάς πρόλαβε να μας στείλει τις πληροφορίες που χρειαζόμασταν. Το ίδιο βράδυ κιόλας, είχα στείλει μία ομάδα στην εκκλησία εκείνη για να ερευνήσουν την υπόθεση. Αυτή η ομάδα, Σπαρκ, ήταν η ομάδα του πατέρα σου»

Δεν του είπα κάτι, περίμενα να δω που ήθελε να καταλήξει.

«Εκείνη τη μέρα, ο πατέρας σου και ο Χερμπ ήταν οι μόνοι που βγήκαν ζωντανοί από την εκκλησία. Τις επόμενες μέρες, στείλαμε άλλες δύο ομάδες σε ξεχωριστές μέρες. Από αυτές τις ομάδες, όμως, κανείς δεν επέζησε να μας πει την ιστορία. Ο πατέρας σου Σπαρκ ήθελε να ξαναπάει μέσα στην εκκλησία μόνος του, γιατί νόμιζε πως είχε λύσει την υπόθεση. Εκείνη την ημέρα δεν φαινόταν να ήταν ψυχολογικά καλά. Συμπεριφερόταν λες και κάποιος κοντινός του φίλος τον είχε προδώσει ή κάτι. Ο Χερμπ προσπαθούσε να τον φέρει στα λογικά του, όμως ο πατέρας σου ήταν αποφασισμένος. Τα υπόλοιπα νομίζω τα ξέρεις...»

«Ναι... και πάλι όμως. Γιατί σταματήσατε να ερευνάτε αυτή την υπόθεση μετά από όλα αυτά;» ρώτησα από περιέργεια.

«Κοίτα, δεν ήθελα να ρισκάρω να χάσω άλλους πράκτορες έτσι, οπότε βάλαμε μία παύση σε αυτή την υπόθεση και την αφήσαμε μισάνοιχτη»

«Μισάνοιχτη...» αναστέναξε η Μπλαιρ.

«Μα αυτό δεν βγάζει νόημα. Γιατί να μην συνεχίσεται την έρευνα; Αυτή δεν είναι η δουλειά της Ειδικής Υπηρεσίας; Να εγγυόμαστε την ασφάλεια των πολιτών;» τον ρώτησα εκνευρισμένος.

«Τώρα με πιάνεις...» ψιθύρισε η Μπλαιρ στον εαυτό της.

«Αν δεν σκόπευα να συνεχίσω την έρευνα δεν θα σας το ανέφερα...» μου απάντησε ο Μαλχη.

«Τώρα είναι πολύ αργά» του φώναξα καθώς τον πλησίαζα. «Εγώ σε ρωτάω γιατί 12 χρόνια τώρα δεν κάνεις κάτι» χτύπησα το χέρι μου στο γραφείο του.

«Γιατί ήταν επικίνδυνο. Τέλος τώρα, ώρα να επικεντρωθούμε στο παρόν»

«Όχι Μαλχη, δεν μπορούμε να επικεντρωθούμε στο παρόν, αφού λες χαζομάρες» ξανακοπάνησα το χέρι μου στο γραφείο του και αυτός με κοίταξε χωρίς φόβο.

«Το έβγαλες από μέσα σου;»

«Όχι, γιατί αν η υπόθεση “Εμπρησμού”, ή όπως την είπες, ήταν υπόθεση που έπρεπε να μείνει μισάνοιχτη, και αν όντως αυτή η γυναίκα έχει 7 μαγείες, γιατί δεν βγήκε ούτε πριν 20, αλλά ούτε πριν 12 χρόνια δημόσια να ανακοινώσει “Εγώ έσωσα τα παιδιά από τη φωτιά”. Προφανώς μιλάμε για απαγωγή. Οι έρευνες έπρεπε να συνεχίσουν για να μάθουμε τον σκοπό της απαγωγής και...»

«Δεν με άφησες να τελειώσω» με διέκοψε ο Μαλχη. «Οι έρευνες μας δεν μπόρεσαν να βρουν ακριβώς τον λόγο της απαγωγής, οπότε σταματήσαμε την ενεργή έρευνα. Βέβαια, υπάρχουν ακόμη κατάσκοποι στα νοσοκομεία της πόλης μήπως συμβεί πάλι το ίδιο συμβάν, αλλά χλωμό το βλέπω...»

«Που θες να καταλήξεις;» τον ρώτησα γιατί η υπομονή μου τελείωνε.

«Κουπ, Φυτ. Μήπως παρατηρήσατε τις χειρονομίες που εκάναν οι άντρες πριν πέσουν ξεροί;»

«Ναι ήταν κάπως έτσι» είπε ο Κουπ και έβαλε την αριστερή του παλάμι πάνω από την δεξιά. Στη συνέχεια, έφερε και τις δύο παλάμες του παράλληλα και κάθετα με το έδαφος. «Και τέλος το αντίθετο με το πρώτο...»

«Σταμάτα!» του φώναξε ο Μαλχη. «Θες να ρισκάρεις να πέσεις και εσύ ξερός;»

«Ουπς ναι, δεν το σκέφτηκα» γέλασε ο Κουπ. «Άρα λες ότι οι χειρονομίες να είναι ο λόγος που έπεσαν όλοι τους ξεροί;»

«Ναι. Αν παρατηρήσουμε και τους υπόλοιπους που συναντήσαμε αυτές τις μέρες και κατέληξαν ξαπλωμένοι στο πάτωμα όταν βρέθηκαν σε δύσκολη κατάσταση, τι κοινό είχαν όλοι τους;» μας ρώτησε ο Μαλχη.

«Ήταν όλοι...»

«Ήταν όλοι 20 χρονών» με δικέκοψε. «Καλά εκτός από τον προμηθευτή ναρκωτικών στο κελί μαζί τον έφηβο με την αφάνα. Αυτός μάλλον έμαθε την χειρονομία από αυτόν τον Γιν»

«Ναι αλλά γιατί να δώσει κάποιος στα καλά καθούμενα τη ζωή του;» ρώτησε η Ιφι.

«Δεν μπορείς να μαντέψεις;» είπε ο Μαλχη και έβγαλε από το χαμηλότερο συρτάρι του γραφείου του ένα βιβλίο.

«Οδηγός της Μαγείας... Μαλχη Σορντ...» διάβασα τον τίτλο και τον συγγραφέα του βιβλίου ανοιχτόφωνα. «Έχεις γράψει βιβλίο για την μαγεία;»

«Όχι εγώ, ο παππούς μου, ο οποίος έφτιαξε και την Ειδική Υπηρεσία. Βέβαια, αυτό το βιβλίο δεν δημοσιεύτηκε ποτέ, οπότε αυτό είναι το μοναδικό αντίγραφο που υπάρχει» μου απάντησε ο Μαλχη και άνοιξε το βιβλίο. «Προς το τέλος του βιβλίου, υπάρχει ένα απαγορευμένο... “ξόρκη”; Δεν ξέρω πως να το περιγράψω, πάντως είναι ακριβώς αυτές οι κινήσεις που έκανε ο Κουπ. Αριστερή παλάμη πάνω από την δεξιά, στη συνέχεια και οι δύο παλάμες κατακόρυφα, λες και κάνεις προσευχή, και τέλος η δεξιά παλάμη πάνω από την αριστερή. Να, κοίτα...» μου έδειξε την εικόνα που περιέγραφε τις κινήσεις. «Από κάτω γράφει: “Κάθε χρήστης μαγείας μπορεί να χρησιμοποιήσει αυτές τις κινήσεις για να δώσει την μαγεία του, μαζί με όλη του την αύρα, σε οποιοδήποτε άλλο χρήστη μαγείας, με τίμημα την ζωή του και δεδομένου ότι ο άλλος χρήστης μαγείας δεν έχει ήδη την μαγεία που του κληρονομείται, στην οποία περίπτωση η αύρα και η μαγεία απλά χάνεται. Αυτές οι χειρονομίες είναι πολύ επικίνδυνες και η χρήση τους πρέπει να αποφεύγεται”. Τι έχεις να πεις;» μου απευθύνθηκε ο Μαλχη.

«Πώς;»

«Τι πώς;»

«Πώς το ανακάλυψαν αυτό; Και με τόση λεπτομέρεια...»

«Δεν ξέρω, και δεν με απασχολεί, αλλά θέλω να σας πω την θεωρία μου. Η γυναίκα αυτή με τις 7 μαγείες πιθανότατα συλλέγει μαγείες. Ίσως απήγαγε τα μωρά από τα νοσοκομεία για να τα εκβιάσει μετά και να τους πάρει την μαγεία»

«Δεν βγάζει νόημα αυτό...» πετάχτηκε η Μπλαιρ. «...αφού οι τελευταίοι άντρες είχαν κάνει ολόκληρη θρησκεία για την γυναίκα»

«Έχεις δίκιο, είπα χαζομάρα» παραδέχτηκε ο Μαλχη.

«Μήπως...» πρότεινα «...δεν εκβιάζει τα παιδιά, αλλά τα μεγαλώνει με την ελπίδα ότι θα την αγαπήσουν σαν “μητέρα” και τους μαθαίνει αυτές τις κινήσεις με τα χέρια έτσι ώστε όταν βρεθούν σε δύσκολη κατάσταση από μόνοι τους να της δώσουν τη μαγεία;»

«Πώς... σου ήρθε αυτός ο συλλογισμός;» με ρώτησε ο Μαλχη. «Αν και δεν είναι κακή σκέψη. Θα μπορούσε όντως να γίνεται κάτι τέτοιο. Τα στοιχεία μας στις τελευταίες σας αποστολές μας οδηγούν σε ένα τέτοιο σενάριο τώρα που το σκέφτομαι. Ωραία λοιπόν, Τι θέλετε να κάνουμε;»

«Γιατί να μην πάω στην εκκλησία για να ελέγξω τον χώρο;» πρότεινε η Ιφι.

«Δεν με άκουγες όταν εξηγούσα την υπόθεση;» νευρίασε ο Μαλχη. «Είναι επικίνδυνα εκεί. Μόνο ο Χερμπ κατάφερε να επιζήσει στη τελική και... δεν μας είπε και ποτέ τι έγινε εκεί μέσα πλάκα πλάκα. Γιατί όμως;»

«Άφησε με να πάω τότε για να μάθουμε»

«Όχι!» χτύπησε το χέρι του στο γραφείο.

«Γιατί;»

«Μη βρίσω το σπανάκι σου, είπα όχι»

«Γιατί;»

«Ναι Μαλχη, γιατί;» πετάχτηκε ο Φυτ. «Αφού η Ιφι μας έδειξε στην προηγούμενη υπόθεση πως μπορεί να συγκριθεί με έναν κορυφαίο κατάσκοπο της Ειδικής Υπηρεσίας, τον Ρομπ, οπότε γιατί να μην το δοκιμάσουμε»

«Καλά, ξέρω όσες φορές και να πω όχι σε αυτή τη γυναίκα, αυτή θα κάνει το δικό της. Κάντε ό,τι θέλετε. Πάρτε τα χαρτιά και αφήστε με ήσυχο. Εγώ πάω στο καζίνο τώρα που είναι ακόμη απόγευμα. Αν με χρειαστείται κάτι, πάρτε με τηλέφωνο» είπε ο Μαλχη καθώς σηκωνόταν και μου έδωσε τα χαρτιά. «Πάμε τώρα όλοι έξω από το γραφείο μου»

«Μάλιστα» χαμογέλασε ο Κουπ.

«Εγώ πάω στο γραφείο να εξετάσω την υπόθεση. Θέλει να έρθει κανείς μαζί μου;» πρότεινα.

«Εγώ» μου χαμογέλασε η Ιφι.

«Εμείς λέγαμε να ξεκουραστούμε για σήμερα» μου είπε ο Φυτ, μιλώντας και για τον Κουπ και την Μπλαιρ, οι οποίοι έγνεψαν θετικά. «Μετά τη σημερινή μάχη έχουμε κουραστεί...»

«Λέει ο τύπος που ήταν ξαπλωμένος τη μισή μάχη...» σχολίασε η Μπλαιρ.

«Μπορείτε απλά να φύγετε από το γραφείο μου;» εκνευρίστηκε ο Μαλχη.

«Καλά καλά... μη μας βρίσεις και το σπανάκι κιόλας...» σχολίασε η Μπλαιρ και σηκώθηκε.

Βγήκαμε όλοι από το γραφείο του Μαλχη και η Ιφι και εγώ αρχίσαμε να περπατάμε τον μακρύ διάδρομο προς το γραφείο μας. Πυρκαγιές σε νοσοκομεία, ε; Απαγωγές παιδιών, μεταφορά μαγείας από έναν άνθρωπο σε άλλον... τι άλλο; Ενδιαφέρον υπόθεση πάντως. Δεν ξέρω πως ο πατέρας μου δεν μπόρεσε να την λύσει, ίσως είναι πιο δύσκολη από ότι φαίνεται. Και πάλι όμως, 7 μαγείες; Ο άντρας νομίζω μας είχε πει ότι τώρα μπορεί να έχει και περισσότερες. Άρα ήξεραν το αποτέλεσμα των χειρονομιών αυτών; Θυμήθηκα τις χειρονομίες που κάνει συχνά η Μπλαιρ όταν μουρμουρίζει. Αυτές όμως δεν έχουν καμία σχέση με αυτές που μου έδειξε ο Μαλχη στο βιβλίο του. Λες να είναι κάποιο άλλο... “ξόρκη”; Δεν νομίζω, αφού η Μπλαιρ το κάνει συνέχεια, χωρίς να συμβαίνει κάτι.

Είμαι περίεργος, όμως, να μάθω τον τρόπο που ανακαλύφθηκε το αποτέλεσμα αυτών των κινήσεων, και με τόση λεπτομέρεια... Ο παππούς του Μαλχη πρέπει να ήταν περίεργος άνθρωπος. Αλλά και πάλι, ποιο το νόημα να έχεις πάνω από μία μαγεία; Ο Γιν, που θυμάμαι ο έφηβος μας είχε πει ότι έχει δύο μαγείες, χρησιμοποίησε μόνο την μία μαγεία του εκείνη τη μέρα. Μήπως, όσες περισσότερες μαγείες έχεις, κερδίζεις κάποιο πλεονέκτημα; Δεν μπορώ να καταλάβω. Ελπίζω η απορία μου να λυθεί όταν λυθεί και η υπόθεση...

«Φτάσαμε Σπαρκ. Συγκεντρώσου τώρα» μου είπε η Ιφι καθώς μπαίναμε στο γραφείο. «Θέλω να μου πεις που είναι η εκκλησία και πως μοιάζει η γυναίκα»

«Που θες να ξέρω;» της απάντησα.

«Αφού εσύ έχεις τα χαρτιά...»

«Α ναι, συγγνώμη. Η εκκλησία είναι... κοίτα εδώ» της έδειξα τον χάρτη που ήταν κολλημένος στην γωνία ενός εγγράφου. «Το καζίνο πρέπει να είναι αυτό το μωβ σηματάκι που έχει στην κάτω αριστερή γωνία ο χάρτης και η Ειδική Υπηρεσία πρέπει να είναι...»

«Άστο θα βάλω στο κινητό μου οδηγίες για το καζίνο και θα βρω τον δρόμο μου από εκεί» με διέκοψε η Ιφι. «Τώρα βρες πως μοιάζει η γυναίκα με τις πολλές μαγείες»

«Να, μάλλον αυτή είναι» της έδειξα την μόνη φωτογραφία που περιείχε άνθρωπο. Στην φωτογραφία αυτή εικονιζόταν μία μοναχή με ολόμαυρη ένδυση. Το πρόσωπο της φαινόταν αρκετά νέο, βέβαια η εικόνα ήταν αρκετά παλιά, οπότε υπάρχει πιθανότητα η γυναίκα να έχει περισσότερα σημάδια ηλικίας αφού έχουν περάσει 12 χρόνια.

«Ωραία φεύγω»

«Τι; Τόσο γρήγορα;» της είπα και κάθισα στην καρέκλα του γραφείου μου.

«Ναι»

«Και γιατί ήρθες ως εδώ;»

«Για παρέα...» μου χαμογέλασε και έκλεισε την πόρτα.

\chapter{}

Τι κάνει τόση ώρα; Γιατί δεν έχει επιστρέψει; Έχουν περάσει έξι ώρες από τότε που έφυγε η Ιφι για να πάει στην εκκλησία και έχουν αρχίσει να μου δημιουργούνται τύψεις. Γιατί την άφησα να πάει σε ένα τόσο επικίνδυνο μέρος μόνη της; Δεν έχω κοιτάξει καν τους φακέλους της υπόθεσης εδώ και έξι ώρες, δεν μπορώ να σταματήσω να σκέφτομαι την Ιφι. Λες να... μου αρέσει; Μα γιατί; Πώς μου ήρθε τόσο ξαφνικά αυτό το συναίσθημα; Κοίτα, εμφανισιακά δεν είναι και άσχημη, τα κυματιστά της μαλλιά είναι πανέμορφα και σε συνδυασμό με τα πράσινα μάτια της... Τι σκέφτομαι, άντε μη βρίσω το σπανάκι μου. Γιατί το λέει συνέχεια ο Μαλχη αυτό; Από που προκύπτει καν, γιατί να θες να βρίσεις το σπανάκι κάποιου; 

«Ακόμη εδώ είσαι Σπαρκ;» άκουσα τη φωνή της, αλλά όταν σήκωσα το κεφάλι δεν είδα κανέναν μπροστά μου, οπότε ξανακοίταξα τα χαρτιά που είχα απλωμένα στο γραφείο μου. Ακούω φαντάσματα μάλλον, αλλά γιατί το φάντασμα να έχει τη φωνή της Ιφι; Γιατί έχω κολλήσει τόσο με αυτήν την κοπέλα;

«Σπαρκ;» άκουσα πάλι τη φωνή της. «Εδώ δίπλα σου είμαι...»

Έστρεψα το βλέμμα μου αριστερά και την είδα να κάθεται όρθια και να περιμένει να την παρατηρήσω.

«Ωχ... συγγνώμη, είχα... αφαιρεθεί» της είπα προσπαθώντας να δικαιολογήσω τον εαυτό μου.

«Έχει πάει 11:53 η ώρα. Γιατί δεν έχεις γυρίσει σπίτι σου ακόμη;» με ρώτησε.

«Γιατί...» ήθελα να πω γιατί σε περίμενα, αλλά δεν μπόρεσα να τελειώσω την πρόταση μου.

«Τέλος πάντων. Κοίτα τι βρήκα...» μου έδειξε μία φωτογραφία της μοναχής γυναίκας που ψάχναμε, μέσα στην εγκαταλελειμμένη εκκλησία. Η γυναίκα καθόταν στα γόνατα της στην φωτογραφία και φαινόταν σαν να προσευχόταν. Γύρω της υπήρχαν πολλά διαφορετικά φώτα, σαν αυτά που παρήγαγαν οι κουκουλοφόροι, αλλά είχαν πολύ μικρότερη ένταση και ήταν πολλά περισσότερα σε πλήθος. Μπορούσα να μετρήσω 16 διαφορετικές ακτίνες φωτός γύρω από τη γυναίκα, όλες με διαφορετικά χρώματα. «Τι λες;»

«Πώς; Δεν ήταν υποτίθεται επικίνδυνα;»

«Δεν μπήκα μέσα στην εκκλησία χαζούλη. Από την οροφή την τράβηξα αυτή τη φωτογραφία. Είχε μία τρύπα και...»

«Ευτυχώς που γύρισες...»

«Γιατί;»

«Γιατί... ας πάρουμε τον Μαλχη να του πούμε ότι βρήκαμε τη γυναίκα που ψάχναμε... κάτσε λίγο. Μόνο προσευχή έκανε αυτή η γυναίκα; Δεν έχεις κάποια άλλη φωτογραφία;»

«Όχι, δεν μπόρεσα. Όταν έφτασα στην εκκλησία, σκαρφάλωσα στην ταράτσα με τη βοήθεια του Φυτ, τον οποίο είχα σήρει μαζί μου για να με βοηθήσει και μετά...»

«Άρα δεν ήσουν μόνη σου;» την ρώτησα προβληματισμένος.

«Όχι βρε χαζούλη, θα ήταν χαζό να πήγαινα μόνη μου σε ένα τόσο κακοφημισμένο από τον Μαλχη μέρος»

«Εντάξει, πάλι καλά... και τι έγινε μετά;»

«Μετά... βρήκα ένα άνοιγμα στην οροφή και παρακολούθησα την μοναχή από εκεί. Αυτή έκανε επί τρεις ώρες προσευχή και έλεγε κάτι για...»

«Συγγνώμη που σε διακόπτω, αλλά να πάρω μία τηλέφωνο τον Μαλχη για να έρθει να του εξηγήσουμε την κατάσταση;» την ρώτησα.

«Πάρ’ τον. Τι να σου πω...» εκνευρίστηκε η Ιφι.

«...Ναι έλα Μαλχη, ο Σπαρκ είμαι... Ναι επέστρεψε η Ιφι... Ναι αν μπορείς έλα στο γραφείο μας, θα σε περιμένουμε... Εντάξει αντίο»

«Έρχεται;»

«Ναι είπε στον δρόμο είναι. Συνέχισε»

«Έκανε τρεις ώρες που λες η μοναχή προσευχή και έλεγε κάτι για τα παιδιά της, ότι τα “προηγούμενα” παιδιά έχουν στραφεί στο “κακό” και ότι ελπίζει “αυτά” τα παιδιά να στραφούν στο “καλό” ή κάτι τέτοιο, δεν θυμάμαι ακριβώς. Να σου πω την αλήθεια μετά τη μία ώρα σταμάτησα να προσέχω τα λόγια της γιατί βαρέθηκα. Τρεις ώρες είναι λίγο πολλές δεν νομίζεις;»

«Ναι... και μετά τις τρεις ώρες, τι έγινε;»

«Μετά άνοιξε μία πόρτα στην αριστερή μερία της εκκλησίας, η οποία από πίσω είχε κάποιες σκάλες που οδηγούσαν μάλλον σε υπόγειο. Περίμενα μερικές ώρες για να δω αν θα επιστρέψει στο κύριο χολ της εκκλησίας, αλλά μάλλον πήγε για ύπνο στο υπόγειο, οπότε είπα να επιστρέψω με τον Φυτ στην Ειδική Υπηρεσία, μήπως και ήσουν εσύ ή ο Μαλχη ακόμη στα γραφεία σας»

«Και ο Φυτ;»

«Μιλάει με τη γραμματέα. Γιατί ρωτάς;»

«Τίποτα, κάτι σκεφτόμουν... Τώρα τι κάνουμε;»

«Αφού ξέρουμε που βρίσκεται η μοναχή, θα μπορούσαμε να την προσεγγίσουμε» μου απάντησε η Ιφι και κάθισε στον αριστερό καναπέ, τινάζοντας τα μαλλιά της.

«Να την προσεγγίσουμε;»

«Σίγουρα είσαι καλά σήμερα Σπαρκ; Δεν βλέπω να στροφάρει πολύ το μυαλό σου....»

«Όχι δεν προβληματίστηκα με την λέξη που είπες, απλά... πιστεύεις ότι είναι καλή ιδέα να μπούμε σε αυτή την εγκαταλελειμμένη εκκλησία, από την οποία να σημειώσω πως η πλειοψηφία των πρακτόρων που μπήκε, δεν βγήκε; Δεν νομίζεις ότι είναι λίγο...»

«Για αυτό δεν θα πάμε μόνο οι πέντε μας. Θα προτείνουμε στον Μαλχη να καλέσει τουλάχιστον άλλες δύο ομάδες για ενισχύσεις. Όποιος και να είναι ο αντίπαλος μας δύσκολα θα μπορέσει να αντιμετωπίσει 15 πράκτορες ταυτόχρονα. Τι λες;»

«Και όσοι είναι κλάσης Β;»

«Τι λες μωρέ. Αφού στην προηγούμενη αποστολή ακόμη και εσύ που είσαι κλάσης Β κατάφερες να βοηθήσεις με τον δικό σου μοναδικό τρόπο. Δεν χρειάζεται να σκέφτεσαι τέτοια πράγματα τώρα»

«Ήρθα» είπε ο Μαλχη καθώς μπήκε στο γραφείο μαζί με τον Φυτ. «Έφερα και τον φίλο σας επειδή τον πέτυχα στην είσοδο. Λοιπόν, για πείτε μου τι έγινε...»

Η Ιφι του εξήγησε όσα είχε πει και σε εμένα για την γυναίκα, από την μακροσκελή προσευχή της μέχρι την είσοδο της στο υπόγειο. Ο Μαλχη την άκουγε με προσοχή.

«...πιστεύω πως αν είναι να την προσεγγίσουμε αύριο, θα χρειαστούμε τρεις ομάδες αντί για μία»

«Να την προσεγγίσουμε; Στην εκκλησία;» προβληματίστηκε ο Μαλχη.

«Ναι, γιατί όχι; Θα είμαστε πιο σίγουροι έτσι» του απάντησε η Ιφι. «Έχεις κάποια άλλη ιδέα;»

«Δίκιο έχει η Ιφι, γιατί αν ο αντίπαλος μας δεν είναι τελικά ένας, αλλά πολλοί, πρέπει να μπορέσουμε να έχουμε ενισχύσεις» πρόσθεσα.

«Τι σε κάνει να νομίζεις ότι υπάρχουν πολλά άτομα στην εγκαταλελειμμένη εκκλησία;» με ρώτησε ο Μαλχη.

«Η προσευχή της γυναίκας. Αφού αναφέρεται σε “προηγούμενα” παιδιά και σε “αυτά” τα παιδιά, υπάρχει πιθανότητα τα δεύτερα να είναι κάπου κρυμμένη στην εκκλησία» του απάντησα.

«Χμ» σκέφτηκε για λίγο ο Μαλχη. «Γιατί να μην το κάνουμε όπως τις υπόλοιπες φορές; Να πάει μία ομάδα, τεσσάρων πχ ατόμων, πρώτη στην εκκλησία και να σαρώσουν τον τόπο και αν χρειαστούν ενισχύσεις, θα έχουμε ένα κουμπί ή κάτι για να στείλει σινιάλο στους υπόλοιπους απ’ έξω από την εκκλησία να εισέλθουν. Τι λέτε; Έτσι ταυτόχρονα δεν ρισκάρουμε κιόλας τον πιθανό θάνατο όλων μας από ξαφνική ενέδρα»

«Ναι...» του απάντησα. «Ποιοι θα πάνε όμως πρώτοι στην εκκλησία;»

«Εμείς οι τέσσερις φυσικά. Ο Φυτ είναι εξαιρετικός στις μάχες, οπότε αν χρειαστούμε προστασία για λίγη ώρα μέχρι να έρθουν οι ενισχύσεις θα είμαστε καλυμμένοι. Εγώ πάλι μπορώ να βοηθήσω στις μάχες, καθώς έχω εμπειρία. Η Ιφι πρέπει να είναι αναγκαστικά εκεί, γιατί είναι η μόνη που έχει καθαρή εικόνα για το τι ψάχνουμε. Τέλος, εσύ Σπαρκ πρέπει να είσαι εκεί για να διαβάσεις τις σκέψεις της γυναίκας όταν την προσεγγίσουμε»

«Καλό μου ακούγεται» είπε ο Φυτ.

«Και εμένα» συμπλήρωσε και η Ιφι.

«Άρα θα ακολουθήσουμε αυτό το πλάνο» είπε ο Μαλχη. «Οι υπόλοιποι προφανώς θα είναι απ’ έξω για να μας βοηθήσουν, αλλά μπορείτε να μου κάνετε μία χάρη;» μας είπε και τον κοιτάξαμε όλοι. «Τι λέτε να χρησιμοποιήσουμε το όπλο της γυναίκας αυτής εναντίον της;»

«Τι εννοείς;» ρώτησε η Ιφι προβληματισμένη.

«Εννοώ πως, αν βρεθεί κάποιος από εμάς σε δύσκολη κατάσταση, να δοκιμάσει να κάνει αυτές τις κινήσεις που συζητούσαμε το απόγευμα με τα χέρια του και να δώσει τη μαγεία του σε όποιον νομίζει εκείνη τη στιγμή πως βρίσκεται στην καλύτερη κατάσταση. Τι λέτε;»

Κοίταξα την Ιφι και τον Φυτ. Ήταν και οι δυο τους το ίδιο προβληματισμένοι με εμένα.

«Περίεργη πρόταση, αλλά εντάξει Μαλχη, αν χρειαστεί θα το κάνουμε» του είπε η Ιφι.

«Άσε που δεν θα χρειαστεί να, αφού θα έχετε εμένα εκεί να σας προστατεύω» πρόσθεσε χαριτολογώντας ο Φυτ.

«Άρα αυτό είναι ναι από όλους; Σπαρκ;» με ρώτησε ο Μαλχη.

«Ναι...»

«Θα το πω και στους άλλους όταν έρθει η ώρα. Πάμε σπίτια μας για την ώρα και σε δύο μέρες, που θα έχω οργανώσει τις ομάδες, θα πάμε όλοι μαζί στην εκκλησία και θα “προσεγγίσουμε”, όπως λέει και η Ιφι, τη γυναίκα. Προς το παρόν δεν χρειάζομαι τίποτα από εσάς αυτές τις δύο μέρες, μπορείτε να χαλαρώσετε»

«Αχ τι ωραία» είπε η Ιφι.

«Μπορούμε να ερχόμαστε για προπόνηση, έτσι;» ρώτησε ο Φυτ.

«Δεν θα σου πω τι να κάνεις. Κάνε ότι θες αυτές τις δύο μέρες. Απλά μεθαύριο σας θέλω όλους στις 12 το πρωί στην εκκλησία. Εντάξει;»

«Μάλιστα κύριε» είπε ο Φυτ μιμούμενος τη χαρούμενη φωνή του Κουπ όταν λέει τις ίδιες λέξεις.

«Ωραία. Πάμε σπίτια μας τώρα, έχει πάει 12:15 και είμαστε ακόμα εδώ...»

\chapter{}

«Πότε πέρασαν δύο μέρες;» είπε ο Φυτ που καθόταν στη θέση του συνοδηγού.

«Ναι όντως... δεν θυμάμαι καν τι έκανα αυτές τις δύο μέρες...» σχολίασα.

«Μάλλον καθόσον σπίτι όπως όλοι μας, αφού ο Μαλχη μας είχε δώσει το ελεύθερο να κάνουμε ό,τι θέλουμε» μου είπε ο Φυτ.

«Ναι μάλλον αυτό έκανα... Ιφι, πότε φτάνουμε; Έχει πάει 12:15, θα μας κράξει ο Μαλχη που αργήσαμε»

«Δεν φταίω εγώ που είχε τόση κίνηση» μου απάντησε εκνευρισμένη η Ιφι. «Να εδώ είμαστε»

«Πω, τώρα όμως θα πρέπει να βρεις και να παρκάρεις...» σχολίασε ο Φυτ.

«Ελάτε από δω!» μας φώναξε ο Μαλχη και μας έκανε νόημα να σταματήσουμε. «Αφήστε το όπου να ναι το αμάξι! Έχουμε κλείσει έτσι κι αλλιώς τους δρόμους τριγύρω!»

«Ε για αυτό είχε τόση κίνηση...» ψιθύρισε στον εαυτό της η Ιφι. «Εντάξει σταματάω!» του φώναξε πίσω.

Όταν κατεβήκαμε από το αμάξι είδαμε τους υπόλοιπους 12 πράκτορες να μας περιμένουν μαζί με τον Μαλχη. Μαζί με εμάς, ήμασταν συνολικά 16 πράκτορες της Ειδικής Υπηρεσίας, όλοι ντυμένοι με τα κλασικά μας μαύρα σακάκια και τα όπλα μας στη μέση μας. Στα τριγύρω στενά δεν υπήρχε κανένα ίχνος άλλου ανθρώπου. Λογικά ο Μαλχη είχε βάλει κάποιους ακόμη πράκτορες να περιπολούν την περιοχή για να μην έχουμε εμπλοκές από τρίτους. Εκτός αν βέβαια η περιοχή γύρω από την εκκλησία έχει κλειστεί με κάποιου είδους μαγείας. Τι μαγεία θα μπορούσε να είναι αυτή; Μαγεία Απώθησης; Μαγεία Ασπίδας; Μπα, πολύ συγκεκριμένες ακούγονται αυτές οι μαγείες για να είναι πραγματικές. Μπορεί απλά να έχουν βάλει απαγορευτικές πινακίδες και ταινίες και να το σκέφτομαι πολύ.

«Είστε έτοιμη» μας είπε ο Μαλχη όταν τον πλησιάσαμε.

«Βεβαίως» του είπε ο Φυτ.

«Έχει εξηγήσει στους υπόλοιπους το σχέδιο;» ρώτησα τον Μαλχη για σιγουριά.

«Ναι»

«Όλο το σχέδιο;»

«Για χαζό με έχεις;» μου απάντησε ρητορικά.

«Καλά... Ας ξεκινήσουμε«» είπα και το βλέμμα μου στράφηκε στην Μπλαιρ και στον Κουπ, οι οποίοι ήταν σχετικά μακριά από εμάς εκείνη τη στιγμή και φαίνονταν και οι δυο τους αγχωμένοι. Σκέφτηκα να διαβάσω τις σκέψεις τους από περιέργεια.

\textasciitilde ...γιατί να πάνε μόνοι τους; Είναι πολύ ριψοκίνδυνο. Δεν πέθαναν σχεδόν τρεις ομάδες μέσα σε αυτή την εκκλησία; Πραγματικά ποιος σκέφτηκε αυτό το σχέδιο;...\textasciitilde \  σκεφτόταν η Μπλαιρ.

\textasciitilde ...ελπίζω να πάνε όλα καλά. Ελπίζω με τον Μαλχη και τον Φυτ στην ομάδα τους να είναι καλυμμένοι...\textasciitilde \  σκεφτόταν ο Κουπ.

Ήθελα να τους πλησιάσω για να τους καθησυχάσω, αλλά με πρόλαβε ο Μαλχη.

«Τι κοιτάς Σπαρκ;» με ρώτησε.

«Τίποτα...»

«Μη μου πεις ότι διαβάζεις τις σκέψεις τον συναδέλφων σου πάλι; Σου είπα χάνεις χρόνο όταν το κάνεις αυτό, ειδικά στις αποστολές» εκνευρίστηκε ο Μαλχη.

«Καλά...»

«Πάμε μέσα τώρα. Η Ιφι και ο Φυτ έχουν μπει ήδη στην εκκλησία» μου είπε και έκανε ένα περίεργο νεύμα με το χέρι του στους υπόλοιπους πράκτορες που βρίσκονταν εκεί.

«Τι ήταν αυτό;» τον ρώτησα από περιέργεια.

«Μην σε νοιάζει, οι υπόλοιποι με κατάλαβαν» μου απάντησε και πέρασε την σπασμένη πόρτα της εκκλησίας. Τον ακολούθησα.

Μέσα η εκκλησία ήταν μία φυσιολογική ορθόδοξη εκκλησία. Μερικές καρέκλες ήταν σπασμένες, ο μισός γυναικωνίτης είχε γκρεμιστεί, οι εικόνες μπροστά από το ιερό είχανε σημάδια, αλλά κατά τα άλλα σου θύμιζε μία τυπική εκκλησία. Κοίταξα πάνω για να δω που ήταν αυτή η τρύπα στην οροφή που έλεγε η Ιφι και είδα τον τρούλο της εκκλησίας να λείπει, αλλά στο πάτωμα δεν υπήρχε κάποιος μικρός κρατήρας ή σημάδι γενικά. Γιατί λείπει ο τρούλος; Μήπως κάποιος τον έβγαλε από πάνω από την εκκλησία επίτηδες; Και ο γυναικωνίτης γιατί είναι γκρεμισμένος; Πολλές απορίες. Πήγα πίσω από το ιερό για να ελέγξω αν κρυβόταν κανείς. Το μόνο που είδα, όμως, ήταν γκρεμισμένες καρέκλες και αγιογραφίες κατεστραμμένες από τα φυσικά φαινόμενα.

«Σπαρκ υπάρχει κάτι εκεί πίσω;» μου φώναξε ο Μαλχη.

«Δεν βλέπω κάτι» του απάντησα με ίδια ένταση

«Ωραία, έλα να πάμε στο υπόγειο τότε»

Ήμασταν γρήγοροι, δεν μπορώ να πω. Ελέγξαμε το υπόγειο της εκκλησίας σε λιγότερο από 5 λεπτά και δεν χρειάστηκε να χρησιμοποιήσει κανείς τη μαγεία του. Μόλις επανενώθηκα με τους υπόλοιπους, ο Φυτ και η Ιφι φαίνονταν όλοι τους αγχωμένοι. Μου πέρασε από το μυαλό να διαβάσω τις σκέψεις τους, αλλά απέρριψα αυτή τη σκέψη γιατί θα εκνευριζόταν πάλι ο Μαλχη.

Ο Φυτ άνοιξε την πόρτα που οδηγούσε στο υπόγειο. Πίσω από την πόρτα ήταν προφανώς οι σκάλες που μας είχε αναφέρει η Ιφι.

«Ποιος θέλετε να πάει πρώτος;» ρώτησε ο Μαλχη.

«Να πάω εγώ;» πρότεινε ο Φυτ.

«Ναι αυτό θα έλεγα και γω» είπε ο Μαλχη. «Ίσως πήγαινε εσύ μπροστά και πάρε και αυτή την ασπίδα» συνέχισε και του πέταξε ένα πετραδάκι από το πάτωμα που στον αέρα το μετέτρεψε σε ασπίδα με την μαγεία όπλων του. «Από πίσω σου θα έλεγα να είναι η Ιφι για να μπορεί να σε θεραπεύσει με την μαγεία της. Ο Σπαρκ και εγώ θα είμαστε πίσω σας»

«Εντάξει» του απάντησε ο Φυτ και άρχισε να κατεβαίνει τις σκάλες. Η Ιφι τον ακολούθησε.

«Μετά από εσάς...» έκανε ένα πλαϊνό βήμα και μου παραχώρησε την σειρά του ο Μαλχη.

«Ευχαριστώ» του απάντησα και άρχισα να κατεβαίνω τις σκάλες.

Η διαδρομή προς το υπόγειο δεν είχε και τίποτα το ιδιαίτερο. Το μόνο παράξενο ήταν ότι οι σκάλες δεν είχαν κανένα σημάδι ζημιάς. Αντίθετα ήταν καθαρές και γυαλιστερές σε επίπεδο που μπορούσες να δεις μία θολή αντανάκλαση του εαυτού σου. Όσο κατεβαίναμε ύλπιζα να διακρίνουμε κάποιον να μιλάει, αλλά υπήρχε τόση ησυχία που μπορούσε κανείς να ακούσει τον παλμό της καρδιάς του. Όταν τελείωσαν οι σκάλες, υπήρχε ακόμη μία ξύλινη πόρτα που μας χώριζε μάλλον από τον χώρο που θα βρισκόταν η μοναχή με τις πολλές μαγείες.

«Να την ανοίξω;» ρώτησε ο Φυτ τον Μαλχη.

«Με δύναμη...» σχολίασε ο Μαλχη.

Ο Φυτ μας έκανε νόημα να κάνουμε όλοι άκρη, πήρε λίγη φόρα και έσπασε με τον ώμο του την πόρτα. Οι υπόλοιποι ήμασταν σε ετοιμότητα να σηκώσουμε τα όπλα μας αν χρειαζόταν, αλλά πίσω από την πόρτα δεν βρισκόταν ούτε ψυχή. Ήταν απλά ένα άδειο δωμάτιο, χωρίς επίπλωση, χωρίς πρίζες, χωρίς τίποτα. Το μόνο που υπήρχε στο δωμάτιο, εκτός από τα συντρίμμια της πόρτας, ήταν ένα κερί στο κέντρο του δωματίου. Ο Φυτ και η Ιφι το πλησίασαν και, ευτυχώς για αυτούς, δεν ήταν κάποια παγίδα. Το περιεργάστηκαν λίγο παραπάνω, αφού ήταν το μόνο στοιχείο που είχαμε σε αυτό το δωμάτιο.

Περίεργο πάντως που δεν έχουμε ακούσει κανέναν. Είμαστε σε εγκαταλελειμμένη εκκλησία ναι, αλλά η Ιφι είχε υποψίες πως η μοναχή με τις 7 μαγείες θα βρισκόταν εδώ. Και τα παιδιά που ανέφερε στην προσευχή της που βρίσκονται; Λες να βρίσκονται και αυτά σε κάποιο σπίτι εκτός της πόλης;... Λες να κρύβονται χρησιμοποιώντας μαγεία; Δεν ξέρω τι μαγεία, αλλά θα μπορούσε κάποιος με μία μαγεία όπως της γραμματέως της Ειδικής Υπηρεσίας να προσπαθεί να μας ξεγελάσει. Πάρα πολλά σενάρια θα μπορούσαν να συμβαίνουν αυτή τη στιγμή.

Το βλέμμα μου στράφηκε προς την Ιφι και τον Φυτ, οι οποίοι ήταν σκυμμένοι μπροστά στο κερί, γυρισμένοι πλάτη προς εμένα και τον Μαλχη. Προχώρησα λίγο προς το μέρος τους γιατί ήθελα να διαβάσω τις σκέψεις τους, τώρα που ήμασταν πιο χαλαροί.

\textasciitilde ...γιατί κοιτάμε τόσο προσεκτικά αυτό το κερί; Χαζός είναι ο Φυτ; Αν μας κάνουν ενέδρα τώρα δεν θα μπορούμε να αμυνθούμε...\textasciitilde \  σκεφτόταν η Ιφι.

Έσκυψα δίπλα στον φυτ και άρχισα και εγώ να περιεργάζομαι το κερί ενώ ταυτόχρονα άρχισα να διαβάζω τις σκέψεις του.

\textasciitilde ...προστασία, προστασία, προστασία, προστασία...\textasciitilde \  σκεφτόταν ο Φυτ.

Τι γλυκούλης. Ακόμη και στην πιο ήρεμη στιγμή της αποστολής, που είχαμε εξετάσει ήδη όλο το μέρος που ήρθαμε να εξετάσουμε, είχε στο μυαλό του την ασφάλεια μας. Προσπάθησα διακριτικά να διαβάσω τις σκέψεις του Μαλχη, χωρίς να τον κοιτάξω, για να μην του κινήσω υποψίες.

\textasciitilde ...και τώρα αρχίζει το σόου...\textasciitilde 

«Ποιο;» σηκώθηκα γρήγορα και τον κοίταξα. Είχε ήδη βγάλει το πιστόλι του και μας σημάδευε. Πριν προλάβουν οι άλλοι δύο να γυρίσουν και να τον κοιτάξουν, ο Μαλχη είχε στείλει ήδη δύο σφαίρες στα πνευμόνια τους. Εγώ ευτυχώς πρόλαβα και απέφυγα την τρίτη σφαίρα που προοριζόταν για μένα.

«Προφανώς και διάβασες τις σκέψεις μου» μου είπε εκνευρισμένος, με την κάννη του πιστολιού του να με αντικρίζει. «Φυτ, Ιφι, γρήγορα. Ο Σπαρκ είναι προδότης. Δώστε μου τις μαγείες σας για να μπορέσω να τον ακινητοποιήσω»

«Μά... λιστα» είπε ο Φυτ, φτύνοντας αίμα στην προσπάθεια του να μιλήσει.

Είχα μείνει άφωνος από το σχόλιο του Μαλχη. Έστρεψα το βλέμμα μου στους δύο τραυματισμένους. 

«Όχι παιδιά, μη. Μη τον ακούτε» είπα και έτρεξα να τους αποτρέψω από το να κάνουν τις κινήσεις. Ο Μαλχη χωρίς δεύτερη σκέψη με πυροβόλησε στο αριστερό μου πόδι και έπεσα ανάμεσα τους, δίπλα στο κερί. Όταν σήκωσα το κεφάλι μου τους είδα να πέφτουν και οι δύο ξεροί στο πάτωμα. Το βλέμμα μου στράφηκε στο άψυχο πρόσωπο της Ιφι, το οποίο είχε πάρει μία γκρίζα απόχρωση και τα μάτια της είχαν χάσει το πρασινωπό τους χρώμα. Προσπάθησα να σηκωθώ και να πατήσω το δεξί μου πόδι, αλλά ο Μαλχη χρησιμοποίησε τη μαγεία Φύσης του Φυτ για να με τυλίξει με ρίζες και να με κρατήσει ξαπλωμένο στο πάτωμα.

«Κανείς δεν θα σε πιστέψει» μου ψιθύρισε ενώ μου θεράπευε ταυτόχρονα το πόδι.

«Τι... εννοείς;»

«Τόλμα να σηκωθείς και να μη με ακούσεις και θα δεις τι θα γίνει...» με απείλησε και πάτησε το κουμπί του σινιάλου.

Ένιωσα τις ρίζες του να χαλαρώνουν και να μη με τυλίγουν τόσο σφιχτά. Εκείνη τη στιγμή σηκώθηκα και τον έσπρωξα μακριά μου, βγάζοντας ταυτόχρονα το πιστόλι από τη ζώνη μου.

«Γιατί προσπαθείς καν» μου είπε και πυροβόλησε τον εαυτό του στο μπούτι.

Έμεινα άφωνος. Γιατί να το κάνει αυτό; Γιατί με θεράπευσε και με άφησε να σηκωθώ; Και γιατί πυροβόλησε τον εαυτό του; Θα μπορούσε άνετα να με είχε σκοτώσει αν πραγματικά το ήθελε.

Τον σημάδεψα με το πιστόλι μου. Εκείνη τη στιγμή, όμως, είδα την Μπλαιρ και τον Κουπ να κατεβαίνουν.

«Πάλι καλά που ήρθατε παιδιά. Ο Μαλχη...»

«Μη τον ακούτε, Μπλαιρ, Κουπ. Ο Σπαρκ είναι προδότης. Κοιτάξτε τι μας έχει κάνει. Σκότωσε πρώτα την Ιφι και τον Φυτ και τώρα θέλει να σκοτώσει εμένα»

«Όχι, μην τον ακούτε, δεν λέει...» κοίταξα γύρω μου και η σκηνή του εγκλήματος ήταν γεμάτη στοιχεία εναντίον μου. Κάποιος που δεν βρισκόταν εκεί την στιγμή του γεγονότος προφανώς θα σκεφτόταν πως ο μόνος άνθρωπος που δεν είναι τραυματισμένος προκάλεσε όλον αυτό τον σαματά.

Κοίταξα την Μπλαιρ και τον Κουπ, οι οποίοι με σημάδευαν με τα όπλα τους και έρχονταν προς το μέρος μου. Η έκφραση του Κουπ, από χαρούμενη που είναι πάντα, είχε γίνει τώρα πιο σοβαρή από ποτέ. Η Μπλαιρ φαινόταν λες και προσπαθούσε να κρατήσει πίσω δάκρυα. Όταν με πλησίασαν, εγώ χαμήλωσα το όπλο μου, για να μην κινήσω παραπάνω υποψίες.

«Παιδιά, είναι όλα μία παρεξήγηση, μπορώ να...»

Με ακούμπησε η Μπλαιρ, μάλλον για να αχρηστέψει τη μαγεία μου, αλλά είμαι κλάσης Β, οπότε γιατί να το κάνει αυτό;

\textasciitilde ...η αύρα του έχει γίνει κλάσης Α τώρα, από το οριακό κλάσης Β. Πώς γίνεται αυτό; Λες να εκβίασε τον Φυτ και την Ιφι να του δώσουν τις μαγείες τους; Κάτι περίεργο έχει γίνει εδώ...\textasciitilde 

«Όχι Μπλαιρ, δεν είναι αυτό που σκέφτεσαι»

«Πώς ξέρεις τι σκέφτομαι; Αφού σε...»

«Δεν έχω χρόνο για τέτοιες εξηγήσεις, απλά βοήθή με τώρα να σκοτώσουμε τον Μαλχη»

«Γιατί θες να σκοτώσεις τον Μαλχη; Γιατί να σε βοηθήσουμε; Έχουμε μάτια, δεν είμαστε χαζοί. Βάλε τις χειροπέδες τώρα και πάμε πίσω»

«Όχι Μπλαιρ» την διέκοψε ο Μαλχη ξαπλωμένος, πιάνοντας το πόδι του. «Ο Σπαρκ δεν θα μπει σε κελί στην Ειδική Υπηρεσία. Θα τον στείλουμε κατευθείαν στις φυλακές εκτός της πόλης»

«Μάλιστα αρχηγέ» είπε ο Κουπ και τύλιξε το χέρι του με την μαγεία του γύρω μου, για να μην προσπαθήσω να αποδράσω. Στη συνέχεια με σήκωσε και άρχισε να με κουβαλάει σαν σακούλα του σούπερ μάρκετ. Είχα ήδη χειροπέδες στα χέρια μου, γιατί χρειαζόταν και αυτό; Ο Μαλχη σηκώθηκε και άρχισε να ανεβαίνει τις σκάλες κουτσαίνοντας. Κοίταξα πίσω για να δω τα πτώματα δίπλα στο κερί.

«Η Ιφι και ο Φυτ;» τους ρώτησα.

«Γιατί σε ενδιαφέρουν τόσο;» σχολίασε η Μπλαιρ. «Μη μιλάς τώρα γιατί...» άρχισε να κλαίει μόνη της. Ο Κουπ την αγκάλιασε με το άλλο του χέρι.

«Γιατί δεν με πιστεύετε;» τους ρώτησα.

Δεν μου απάντησε κανείς. Δεν ξέρανε τι να πούνε προφανώς, αφού είχαν μείνει άφωνοι από το σκηνικό που είδαν.

‘Οταν φτάσαμε πίσω στο κύριο χολ της εκκλησίας, είδα τους υπόλοιπους πράκτορες να με κοιτάζουν όλοι περίεργα. Ο Μαλχη τους είπε ακόμη μία φορά το ψέμα που είχε σκαρφιστεί και αυτοί τον πίστεψαν χωρίς δεύτερη σκέψη και έφυγαν όλοι για τα αμάξια τους απ’ έξω από την εκκλησία.

«Μα καλά ανέβηκες όλες τις σκάλες κουτσαίνοντας;» είπα ειρωνικά στον Μαλχη.

Αυτός με κοίταξε δήθεν νευριασμένος και κάθισε στο πάτωμα.

«Δεν θα σχολιάσεις τίποτα ρε ψεύτη;» του φώναξα. «Γιατί το κάνεις αυτό; Τι θα κερδίσεις;»

«Ήρεμα Σπαρκ» μου είπε ο Κουπ. «Τελείωσαν όλα τώρα. Μη μιλάς άλλο»

«Μα...»

«Δεν με ενδιαφέρει. Να τα πεις στους φίλους σου στην φυλακή»

Κοίταξα τον Μαλχη, ο οποίος είχε ξαπλώσει στο πάτωμα και περίμενε να έρθει κάποια νοσοκόμα να τον θεραπεύσει.

«Γιατί δεν χρησιμοποίεις τη μαγεία της Ιφι για να θεραπεύσεις τον εαυτό σου;» τον ρώτησα.

«Α τώρα εκτός από προδότης είσαι και σχιζοφρενής; Κουπ σε παρακαλώ πάρ’ τον από εδώ, δεν θέλω να τον ακούω άλλο»

«Μάλιστα» του απάντησε με το σοβαρό του ύφος ο Κουπ και με κουβάλησε μέχρι έξω.

«Κουπ, πρέπει να με ακούσεις, δεν έκανα κάτι εγώ, ο Μαλχη μας είχε στήσει παγίδα και...»

«Πες τα στους φίλους σου τους κρατούμενους αυτά» με διέκοψε. «Τώρα ησυχία»

Γιατί δεν με πιστεύει κανείς;

\chapter{}

«Δεν είναι δίκαιο αυτό που κάνετε. Είμαι αθώος»

«Καλά καλά, φόρα αυτό τώρα...» μου είπε ο φρουρός που με είχε συνοδεύσει ως το κελί μου και μου έδωσε ένα βραχιόλι.

«Τι είναι αυτό;»

«Βραχιόλι που περιέχει αύρα αντιμαγείας, για να μη φοράς χειροπέδες όλη μέρα» μου εξήγησε.

«Καλά...» του είπα και μου φόρεσε το βραχιόλι. «Βγάλε μου τις χειροπέδες τώρα»

«Μη με διατάζεις...» μου είπε και ξεκλείδωσε το ενοχλητικό μέταλλο που είχα στα χέρια εδώ και μία ώρα.

«Δεν μπορώ να καταλάβω όμως, γιατί να μην έχω δικαίωμα να ανακριθώ ή τουλάχιστον να περάσω από δίκη»

«Μετά από αυτό που έκανες; Πολλά ζητάς. Και τα 100 χρόνια φυλάκισης που σου δώσανε λίγα είναι...» μου είπε ο φρουρός και με έσπρωξε με μίσος μέσα στο κελί μου, ρίχνοντας με στο πάτωμα.

«Και που ξέρεις ότι είμαι όντως ένοχος;»

«Θες να μου πεις ότι ο Μαλχη λέει ψέματα;»

«Ναι»

«Χμ, για να δω ποιον θα πιστέψω τώρα, τον αρχηγό της Ειδικής Υπηρεσίας... ή έναν τυχαίο πράκτορα; Νομίζω καταλαβαίνεις ότι η ερώτηση είναι ρητορική. Εις το επανιδείν» με αποχαιρέτησε και έκλεισε την μεταλλική πόρτα.

Το κελί μου δεν είχε και πολλά πράγματα, μόνο ένα κρεβάτι, μία τουαλέτα, μία βρύση και ένα παράθυρο στον απέναντι από την πόρτα τοίχο. Κατά τα άλλα όμως είχε άπλετο χώρο να περπατήσει κανείς, παρόλο που ήταν κελί για ένα άτομο. Προσπάθησα να κοιτάξω έξω από το παράθυρο, αλλά αυτό ήταν πολύ ψιλά για μένα. Το μόνο που μπορούσα να δω ήταν τα γκρίζα σύννεφα στον ουρανό που έκρυβαν τον ήλιο.

Κάθισα στο πάτωμα, διπλωμένος όπως η Μπλαιρ εκείνες τις μέρες, με τα πόδια μου λυγισμένα και τα χέρια μου γύρω από τα γόνατα μου, ακουμπόντας το κεφάλι μου στον τοίχο πίσω μου. Είχε δίκιο ο Μαλχη, δεν με πίστεψε κανείς. Και γιατί να με πιστέψουν; Ένας τυχαίος πράκτορας είμαι, όπως είπε και ο φρουρός.

Γιατί όμως ο Μαλχη μας έστησε παγίδα; Γιατί είναι προδότης; Τι έχει να κερδίσει από αυτό; Δεν μπορώ να καταλάβω. Επίσης, πώς μπορεί να συνδέεται ο Μαλχη με αυτή τη γυναίκα που ψάχναμε; Μήπως είναι σύνεργος της και για αυτό η υπόθεση “Εμπρησμού Νοσοκομείων”, ή όπως την είπε, είναι ένα ψέμα; Δεν μπορεί, αφού η Ιφι είδε με τα μάτια της τη μοναχή, η οποία ήταν η ίδια με την φωτογραφία πριν 12 χρόνια, και τα νούμερα των παιδιών και των ετών είχαν μία συνοχή, οπότε πώς... και γιατί... δεν μπορώ να καταλάβω.

Γιατί μπήκα καν στην Ειδική Υπηρεσία; Δεν είχα κάποιο ιδιαίτερο κίνητρο. Το μόνο ίσως που θα μπορούσε κανείς να μετρήσει σαν κίνητρο θα ήταν η εύρεση της αιτίας του θανάτου του πατέρα μου, και έπαιξε και λίγο ρόλο η πίεση που μου είχε ασκήσει ο Βασι όλα τα χρόνια που τον ήξερα και βέβαια η ανακάλυψη της ικανότητας μου να διαβάζω σκέψεις με την μαγεία μου, αλλά δεν είχα κάποιο αποσκοπό, όπως το να γίνω ο καλύτερος ανακριτής ή κάτι. Δεν μπορούσα με την καμία να ξεπεράσω τον πατέρα μου στη δουλειά του... λες... ο Μαλχη να... 12 χρόνια πριν... Οι χρονολογίες ταιριάζουν, οι εμπειρίες μας συμπίπτουν, απλά εγώ κατάφερα να επιζήσω την ενέδρα του Μαλχη. Αλλά αν ήταν όντως ο Μαλχη που τους έκανε ενέδρα όταν ο πατέρας μου πήγε με την ομάδα του εκείνη τη μέρα στην εκκλησία, πώς τον άφησε να επιζήσει την πρώτη φορά χωρίς να τον βάλει φυλακή; Και ο Χερμπ; Πώς και δεν είχε αναφέρει, όπως είπε και ο Μαλχη, τι είχε συμβεί στην ομάδα τους εκείνη τη μέρα; Δεν βγάζει κανένα νόημα. Νιώθω λες και είμαι σε εφιάλτη ή κάποιου είδους κακογραμμένο βιβλίο με ανεξήγητα κενά στην ιστορία του.

Και η Ιφι και ο Φυτ; Γιατί δεν με πίστεψαν; Αφού εγώ ήμουν δίπλα τους και ο Μαλχη από πίσω μας. Θα έβγαζε περισσότερο νόημα να ακούσουν εμένα παρά τον Μαλχη, αφού αν τους είχα τραυματίσει εγώ θα με είχαν δει. Καλά, όχι ότι προσπάθησα και πολύ να τους σταματήσω, αφού με είχε αφήσει άφωνο το ψέμα του Μαλχη εκείνη τη στιγμή. Και πάλι όμως, δεν θα το είχαν παρατηρήσει αν εγώ ήμουν ο ένοχος σε αυτό το έγκλημα; Αφού... δεν το επαναλαμβάνω πάλι, γίνομαι σπαστικός. Καλά και να το ξαναπώ, δεν είναι κανείς εδώ να ενοχλήσω. Όλοι οι “φίλοι” μου νομίζουν ότι είμαι εγκληματίας. Νομίζουν ότι τους έχω προδώσει, ότι σκότωσα εσκεμμένα δύο από τους λίγους φίλους που είχα χωρίς κάποιον φανερό λόγο, ενώ στην πραγματικότητα εγώ είμαι ο προδομένος. Αυτοί με έχουν προδώσει επειδή κανείς δεν κάθισε να με ακούσει, να ακούσει τι έχω να πω εγώ για τα γεγονότα που συνέβησαν στο υπόγειο. Γιατί δεν δικαιούμαι τουλάχιστον μία δίκη, να μπορέσω να αποδείξω την αθωότητα μου μέσω αυτής; Τι λέω ρε φίλε, όλες οι προσπάθειες μου στη τελική θα ήταν μάταιες, αφού αν ο Μαλχη σχετίζεται όντως με αυτή τη γυναίκα αυτό σημαίνει ότι καλύπτει τον εαυτό του εδώ και τουλάχιστον 20 χρόνια. Ποιος ξέρει τι άλλα κόλπα βρύβει στο μανίκι του.

Εκείνη τη στιγμή, είδα το κεφάλι του φρουρού από πριν να περνάει απ’ έξω από την πόρτα του κελιού μου.

«Συγγνώμη!» του φώναξα.

«Τι θες;»

«Μήπως θα μπορούσα να έχω... μία κιμωλία; Ή και παραπάνω αν μπορείτε...» του ζήτησα.

«Γιατί, θες να αποδράσεις;»

«Με κιμωλίες; Σε παρακαλώ, ακούς τον εαυτό σου;»

«Καλά καλά, πάω να ρωτήσω αν έχουμε κιμωλίες, αν και περίεργο πρώτο αίτημα για φυλακισμένος...» μου είπε ο φρουρός και έφυγε από την πόρτα.

Να αποδράσω με κιμωλίες; Τι χαζή απάντηση ήταν αυτή. Μπορεί, βέβαια, να το είπε αντανακλαστικά, αφού φρουρός φυλακής είναι, όλο και κάποιος θα έχει προσπαθήσει να αποδράσει με αντικείμενα που φαίνονται αθώα με την πρώτη ματιά. Κιμωλίες όμως; Αν προσπαθήσω να γκρεμίσω τον τοίχο με τις κιμωλίες θα καταλήξω να σπάω το χέρι μου αντί τον τοίχο.

«Ορίστε» άκουσα τη φωνή του φρουρού και ένα κουτί με κιμωλίες πετάχτηκε μέσα από το παραθυράκι της μεταλλικής πόρτας του κελιού μου.

«Τόσο γρήγορα;» τον ρώτησα.

Δεν υπήρχε απάντηση.

«Καλά... ευχαριστώ» είπα και άνοιξα το κουτί. Μέσα του υπήρχαν δέκα άσπρες κιμωλίες, οι μισές ήταν σπασμένες, οι άλλες μισές ήταν χρησιμοποιημένες και το κουτί ήταν γεμάτο σκόνη. «Ποιος χρησιμοποιεί κιμωλίες στη φυλακή; Υπάρχουν και άλλοι περίεργοι σαν εμένα; Τέλος πάντων, δεν έχω χρόνο... δηλαδή έχω, αλλά θέλω να το τελειώσω όσο πιο γρήγορα μπορώ για να μπορώ να συμπεράνω κάποια πράγματα.»

Άρχισα να γράφω στον δεξί τοίχο του κελιού μου, ο οποίος ήταν ο πιο κενός, καθώς βρισκόταν πάνω από το κρεβάτι μου. Ήθελα να σημειώσω όλα όσα γνωρίζω για τον Μαλχη, και ίσως και τους υπόλοιπους, έτσι ώστε να καταλάβω επιτέλους γιατί θα μπορούσε ο Μαλχη να κάνει κάτι τέτοιο. Ξεκίνησα με ένα κυκλάκι στο κέντρο, στο οποίο έγραψα μέσα το όνομα του Μαλχη. Στη συνέχεια αφιέρωσα ακόμη ένα κυκλάκη στη γυναίκα με τις επτά μαγείες και ένωσα αυτά τα δύο κυκλάκια με ένα ερωτηματικό, γράφοντας από δίπλα “Παιδιά και καμένα νοσοκομεία”. Μετά, έγραψα τα ονόματα όλων όσον ήξερα από την Ειδική Υπηρεσία, ακόμη και των γονιών μου, γύρω από τα δύο πρώτα κυκλάκια και άρχισα να συνδέω τα ονόματα όπου υπήρχε λογική σύνδεση. Έκανα ένα βήμα πίσω για να πάρω μία ανάσα και να παρατηρήσω τι είχα γράψει λίγο καλύτερα.

«Σκατά...» ψιθύρισα στον εαυτό μου.

Δεν είχα κατορθώσει τίποτα. Το μόνο που είχα κάνει ήταν να σχηματίσω ένα δεκάγωνο με ονόματα γύρω από δύο κύκλους.

«Θα προσπαθήσω πάλι αύριο...» είπα και ξάπλωσα στο κρεβάτι.

Όπως ακούμπησα την πλάτη μου στο σκληρό στρώμα του κρεβατιού μου, μου ήρθαν στο μυαλό τα τελευταία λόγια του πατέρα μου. “Δεν ξέρω ποιος είμαι... Δεν ξέρω που είμαι...”. Λες να έχει να κάνει με τον Μαλχη; Αν σχετίζονται με τον Μαλχη όμως αυτά τα λόγια, γιατί να μιλάει σε πρώτο πρόσωπο ο πατέρας μου; Δεν μπορώ να καταλάβω. Γύρισα στο πλάι και έγραψα “Δεν ξέρω ποιος είμαι” δίπλα από το κρεβάτι μου, τραβώντας μια γραμμή προς το όνομα του πατέρα μου, το οποίο ευτυχώς το είχα γράψει χαμηλά και δεν χρειάστηκε να σηκωθώ.

«Πας για ύπνο από τώρα;» άκουσα την φωνή του φρουρού απ’ έξω από την πόρτα και είδα τα μάτια του μέσα από το παραθυράκι.

«Όχι... κάτι σκέφτομαι»

«Ελπίζω να μην σκέφτεσαι να αποδράσεις...» μου απάντησε και έφυγε από την μεταλλική πόρτα.

Γιατί είναι συνέχεια έξω από το κελί μου αυτός ο φρουρός; Άλλα κελιά δεν έχει να περιπολήσει; Ο Μαλχη είπε ότι αυτές οι φυλακές εκτός της πόλης είναι δεσμευμένες για τους πιο επικίνδυνους εγκληματίες, οπότε γιατί δεν περιπολεί και τα υπόλοιπα κελιά. Εκτός αν έχουν πολλούς φρουρούς σε κάθε όροφο, τότε θα έβγαζε νόημα.

Είμαι ολομόναχος όμως πλέον, χωρίς κάποιον να μιλήσω. Όλοι μου οι φίλοι έχουν πιστέψει τον Μαλχη, ο οποίος πιθανότατα έχει πείσει και την οικογένεια μου... δηλαδή την μητέρα μου, δεν έχω άλλη οικογένεια στην πόλη. Ο Κουπ δεν είχε πει όταν μας γνώρισε ότι δεν είχε και αυτός οικογένεια, ή είναι η ιδέα μου; Γιατί θυμάμαι κάποιον να μου λέει ότι ο Κουπ έχει χάσει τους γονείς του και την αδελφή του και μετά την φωνή του Βασι να μονολογεί, αλλά δεν μπορώ να θυμηθώ το μέρος και την ημερομηνία; Μπορεί θυμάμαι και λάθος, δεν ξέρω. Πάντως για την Μπλαιρ ξέρω σίγουρα πως έχει χάσει τον... “πατέρα” της; Την πατρική εικόνα της, όπως αναφερόταν στον Χερμπ, αλλά και τον Ρομπ, ο οποίος ως γιος του Χερμπ λογικά θα ήταν “αδελφός” της. Αυτές οι σχέσεις μου όμως έχουν χαλάσει τώρα. Πώς θα μπορέσω να ξαναντικρίσω στα μάτια αυτούς τους δύο μετά από όσα έγιναν, και δεν είναι καν δικό μου λάθος. Βασικά... όχι λάθος... δεν μου έρχεται η λέξη τώρα. 

Τίποτα δεν μπορώ να σκεφτώ σήμερα, έχω μείνει άφωνος από την εξέλιξη των πραγμάτων. Από αύριο θα αρχίσω να απαριθμώ όλα τα πιθανά σενάρια μεταξύ όλων των ανθρώπων που έχω στον τοίχο μου, μέχρι να καταλήξω σε συμπεράσματα. Μπορεί να μου πάρει μέρες, βδομάδες, μήνες, αλλά αν δε μπορέσω να καταλάβω το “γιατί”, δεν θα μπορέσω και να ησυχάσω.

«22059!» άκουσα απ’ έξω από την πόρτα μία γνώριμη φωνή. Η πόρτα στη συνέχεια άνοιξε. «Κατέβα στην τραπεζαρία για βραδινό»

«Τόσο νωρίς; Είναι ακόμη... τι ώρα είναι;»

«Έξι, με το ρολόι. Τελείωνε τώρα, έχουμε κι άλλους να κατεβάσουμε» μου είπε και προχώρησε στο δίπλα κελί «22060!...»

Εγώ ακολούθησα την διαταγή του και πήγα στην τραπεζαρία. Στην αρχή νόμιζα πως χάθηκα, γιατί νόμιζα πως απλά έπρεπε να κατέβω τις σκάλες του ορόφου μου και θα βρισκόμουν ως δια μαγείας στην τραπεζαρία, αλλά ευτυχώς για μένα βρήκα τον δρόμο μου ακολουθόντας το πλήθος, ή τουλάχιστον έτσι νόμιζα. Μου φάνηκε περίεργο που η τραπεζαρία ήταν τόσο μακριά από το κελί μου, προσπάθησα όμως να απορρίψω αυτή τη σκέψη λέγοντας στον εαυτό μου ότι υπάρχουν υπερβολικά πολλοί κρατούμενοι στην φυλακή. Στη συνέχεια, όμως, κατάλαβα ότι οι άντρες αυτοί δεν πήγαιναν στην τραπεζαρία, αλλά κάπου κρυφά. 

Δυστυχώς, όταν το κατάλαβα, ήταν ήδη πολύ αργά για μένα. Είχαμε φτάσει ήδη σε ένα κελί που είχε μέσα τρία κρεβάτια και άπλετο χώρο για να φιλοξενήσει κανείς έως και είκοσι άτομα. Οι άντρες που ακολουθούσα μέχρι εκείνη τη στιγμή μπήκαν μέσα στο κελί από την ανοιχτή του πόρτα. Εγώ στάθηκα για λίγο απ’ έξω από την πόρτα προσπαθώντας να αναλύσω την κατάσταση στην οποία είχα μπλεχτεί.

«Ο κοντοστούπης απ’ έξω;» άκουσα μία βαριά φωνή μέσα από το μεγάλο κελί. «Θα μπεις μέσα ή θα κάτσεις και θα μας κοιτάς;»

«Εγώ... δεν...»

«Πλάκα κάνω, δεν έχεις επιλογή» είδα έναν ψηλό άντρα να με πλησιάζει και να με τραβάει μέσα στο κελί, κλείνοντας ταυτόχρονα και την πόρτα. «Για πες μας, λοιπόν, είσαι καινούργιος εσύ;»

«Φρέσκος...» του απάντησα σκύβοντας το κεφάλι. Δεν μπορούσα καν να τον αντικρίσω από τον φόβο μου. «Σήμερα το μεσημέρι...»

«Σήμερα το μεσημέρι; Τι έκανες ρε μεγάλε και μπήκες φυλακή μέσα σε λίγες ώρες;»

«Δεν... έκανα κάτι...»

«Ναι όλοι αυτό λένε φιλαράκο. Πες μας τουλάχιστον γιατί σε κατηγόρησαν»

«Να... ήμουν σε ένα υπόγειο... και... με κατηγορούν ότι σκότωσα δύο πράκτορες της Ειδικής Υπηρεσίας και τραυμάτισα τον αρχηγό της...»

«Είσαι καλός εσύ» είπε ειρωνικά και γελάσανε όλοι οι υπόλοιποι άντρες σε αυτό το κελί. «Πρέπει να έχεις κάποια καλή μαγεία τότε»

«Μαγεία;»

«Ναι, τι;»

«Τίποτα, απλά...»

«Απλά τι; Μη μου πεις ότι είσαι κλάσης Β, γιατί όταν σε τράβηξα πριν ένιωσα την αύρα σου και ήσουν κλάσης Α. Πες μας τώρα τι μαγεία έχεις» μου είπε ο ψηλός άντρας και άρχισαν να με κοιτάνε όλοι επίμονα.

«Έχω... μαγεία Κεραυνών»

«Τι; Όντως;»

«Ναι...»

Άρχισαν να γελάνε όλοι μαζί μου. Δεν κατάλαβα τον λόγο. Είπα κάτι αστείο;

«Άρα δεν έχεις καλή μαγεία» μου είπε ο ψηλός άντρας.

«Γιατί;»

«Πλάκα μου κάνεις; Οι κεραυνοί στις μάχες τι μπορούν να καταφέρουν; Όσους ήξερα που είχαν μαγεία Κεραυνών μόνο το κινητό τους μπορούσαν να φορτίζουν με τη μαγεία τους και στις μάχες όλο έχαναν. Εκτός αν μπορείς να μου αποδείξεις ότι δεν ισχύει αυτό» μου είπε και ένιωσα την γροθιά του να μου συνθλίβει το πηγούνι.

Το χτύπημα του με έριξε πάνω στην πόρτα, κάνοντας έναν δυνατό θόρυβο. Έβαλα το χέρι μου στο μάγουλο μου για να δω αν τα δόντια μου ήταν στη θέση τους, που ευτυχώς ήταν, αλλά ήμουν σίγουρος πως άμα με χτυπούσε μερικές ακόμη φορές θα χρειαζόταν να πω αντίο στο στόμα μου για πάντα.

«Έλα ρε κοντούλη, θα κάνεις κάτι;»

«Ό...χι» προσπάθησα να του απαντήσω, αλλά τα δόντια μου πονούσαν τόσο από το πρώτο του χτύπημα που δεν μπορούσα να κάνω κάτι.

Ο άντρας με πλησίασε, αφού είδε πως δεν ήμουν πρόθυμος να παλέψω, και μου έχωσε μία γονατιά ανάμεσα στα μπούτια μου, ρίχνοντας με μπρούμυτα στο πάτωμα.

«Τι γίνεται εδώ;» άκουσα τη φωνή του φρουρού που περνούσε συχνά έξω από το κελί μου. «Ωχ, τι κάνεις εσύ εδώ;» είπε και άνοιξε την πόρτα. «Δεν θα πάτε να φάτε εσείς; Θα κάτσετε στο κελί σας να συζητήσετε πάλι;»

«Ναι» του απάντησε ο ψηλός άντρας.

«Καλά, κάντε ό,τι θέλετε, δεν είμαι εγώ υπέυθυνος αυτού του τομέα. 22059, πάμε στο κελί σου» μου είπε ο φρουρός και με σήκωσε.

Προχωρήσαμε πίσω από τον δρόμο που ακολούθησα τους άντρες για να φτάσω στο μοιραίο κελί τους. Ο φρουρός δεν είχε την πρόθεση να βγάλει άχνα στη διαδρομή μας.

«Ευχ...αριστώ...» προσπάθησα να σπάσω τον πάγο.

«Την δουλειά μου έκανα, μη με ευχαριστείς» μου είπε και άνοιξε το τηλέφωνο του. «...Έλα, ο Πιτ είμαι. Μπορείς να στείλεις μία νοσοκόμα στον 22059... Εντάξει, γεια» είπε και έκλεισε το τηλέφωνο του.

«Πιτ, ε;»

«Ναι, γιατί;»

«Λέγο...μαι Σπα...κ»

«Εντάξει και τι θες να κάνω;»

«Δεν... θες να γνωρι...στούμε;»

«Όχι. Και μη μιλάς πολύ μέχρι να σε θεραπεύσει η νοσοκόμα. Μπορεί να σου φύγει το σαγόνι»

Δεν του απάντησα. Φαινόταν ότι δεν είχε διάθεση για γνωριμίες εκείνη τη στιγμή.

Όταν φτάσαμε στο κελί μου, μία νοσοκόμα με θεράπευσε με μαγεία Θεραπείας, όπως είχε και η Ιφι... 

«Είστε εντάξει κύριε. Εγώ φεύγω τώρα...» είπε και μας άφησε μόνους η νοσοκόμα, η οποία με θεράπευσε πολύ πιο γρήγορα από την ταχύτητα που θυμάμαι να μας θεραπεύει η Ιφι.

«Φεύγω και γω... Σπακ. Αν θες τίποτα πες το τώρα»

«Ναι θέλω κάτι»

«Κι άλλες κιμωλίες;»

«Όχι όχι... μήπως θα μπορούσατε να μου φέρνετε τα γεύματα μου εδώ;»

«Στο κελί σου;»

«Ναι. Βλέπεις, δεν θέλω να ξαναμπλέξω με συμμορίες...»

«Τι νομίζεις ότι είμαστε; Ξενοδοχείο; Να έβρισκες τον δρόμο για την τραπεζαρία, δεν είναι και τόσο δύσκολο. Υπάρχουν και ταμπέλες ξέρεις...» μου είπε και μου έδειξε τις σκάλες, όπου υπήρχε μία ταμπέλα η οποία έλεγε “Τραπεζαρία” και είχε ένα βελάκι διαγώνια προς τα κάτω. «Αν θες να φας πήγαινε. Αν θες να κάτσεις στο κελί σου πάλι, κάτσε, δεν θα σου πω τι να κάνεις, αρκεί να μην αρχίσεις καμία συμμορία και συ και αρχίσεις να πλακώνεις άλλους...» μου είπε ο Πιτ και έφυγε.

Με βλέπει ως άτομο που μπορώ να αρχίσω συμμορία; Ένας μέσος εικοσάρης είμαι, χωρίς μυς ή κάποια ιδιαίτερη μαγεία. Αλλά γιατί μου λένε όλοι ότι είμαι κλάσης Α; Πρώτα οι σκέψεις της Μπλαιρ και τώρα ο ψηλός άντρας που με πλάκωσε. Μήπως... έχω εξελιχθεί; Δεν έχω ακούσει ποτέ κάποιον να μεγαλώνει και να γίνεται πιο “δυνατός” με τη μαγεία του. Είχα στο μυαλό μου ότι είτε γεννιέσαι “αδύναμος”, κλάσης Β, είτε “δυνατός”, κλάσης Α. Πώς γίνεται αυτό; Είναι λες και είμαι όντως πρωταγωνιστής σε κάποιο βιβλίο, το οποίο βέβαια έχει πολλά κενά. Γιατί να συμβούν όλα αυτά χωρίς εξήγηση; Δεν μπορώ να καταλάβω.

«Το φαγητό σου» άκουσα τον φρουρό από πίσω μου.

«Είσαι γρήγορος...» σχολίασα και πήρα το δίσκο από τα χέρια του «...αλλά γιατί;»

«Επειδή το ζήτησες. Δεν θα σου φέρνω κάθε μέρα όμως. Φάε για τώρα και θα δούμε...» είπε και εξαφανίστηκε όσο γρήγορα εμφανίστηκε. Εγώ άρπαξα την ευκαιρία για λίγη ησυχία και κάθισα στο πάτωμα του κελιού μου για να άρχισα να τρώω.

\chapter{}

«22059» άκουσα το φρουρό να χτυπάει την πόρτα.

Δεν του απάντησα. Ήξερα ότι αν ήθελε να μου αφήσει το φαγητό μου θα άνοιγε την πόρτα του κελιού μου και θα το άφηνε στην άκρη του δωματίου.

«Απάντα μου» είπε χτυπώντας πάλι την πόρτα, εγώ όμως δεν του απάντησα. «Είσαι ελεύθερος να φύγεις. Κάποιοι πράκτορες της Ειδικής Υπηρεσίας πλήρωσαν για να σε ελευθερώσουμε»

Τι κακόγουστο αστείο ήταν αυτό; Μετά από τόσο καιρό κάποιος με θυμήθηκε; Μη γελάσω. Ούτε η ίδια μου η μητέρα δεν απαντούσε στα τηλεφωνήματα που έκανα μία φορά την εβδομάδα και τώρα θέλει να με πείσει ότι κάποιος νοιάζεται για μένα.

«Ποιοι πράκτορες;» τον ρώτησα χωρίς να σηκωθώ από την γωνία στην οποία καθόμουν διπλωμένος

«Έλα να μάθεις» είπε και άνοιξε την πόρτα.

Εγώ δεν σηκώθηκα, έστρεψα απλά το βλέμμα μου στην πόρτα και τότε ήταν που είδα δύο γνώριμα πρόσωπα, τα οποία δεν είχαν κανένα σημάδι ηλικίας. Ήταν ο Βασι και ο Κουπ, δύο άνθρωποι που έχω να τους δω στο ίδιο δωμάτιο από πριν την αξιολόγηση μου στην Ειδική Υπηρεσία.

«Είστε σίγουροι ότι αυτός είναι ο Σπαρκ; Δεν του μοιάζει και πολύ...» σχολίασε ο Βασι στον φρουρό.

«Δεν ξέρω, τα χαρτιά του αυτό λένε» είπε ο φρουρός και τους έδειξε το χαρτί που κρατούσε.

«Γιατί είναι τόσο κοκαλιάρης;... Και το μούσι;» ρώτησε ο Βασι.

«Δεν τον έχω δει ποτέ να ξυρίζεται αυτά τα χρόνια, μπορεί να μην ξέρει πως. Εδώ και 10 χρόνια δεν έχει βγει από το κελί του και αναγκαζόμουν κάθε τρίτη μέρα να του φέρνω φαγητό για να μην πεθάνει της πείνας. Τον ακούω συνέχεια να γράφει με την κιμωλία του στους τοίχους, και όποτε δεν γράφει, απλά κάθεται στην γωνία όπως κάθεται τώρα...»

«Άρα αυτό εξηγεί το γεγονός ότι όλοι οι τοίχοι είναι άσπροι...» σχολίασε ο Βασι και κοίταξε γύρω του. «Έχει σημάδια από κιμωλίες και στο πάτωμα. Ερωτηματικά, νεκροκεφαλές... Τι είναι όλα αυτά ρε Σπαρκ;»

«Γιατί;» τον ρώτησα χωρίς να σηκωθώ.

«Τι πράγμα;» μου απάντησε προβληματισμένος.

«Γιατί είστε εδώ;»

«Για να σε βγάλουμε από εδώ. Μαζέψαμε επιτέλους τα λεφτά και...»

«Και σας πήρε 10 χρόνια;» σχολίασα και σηκώθηκα.

«Ναι βλέπεις δεν ήταν λίγα τα λεφτά για...»

«Ναι αλλά γιατί; Γιατί με βοηθάτε; Αφού...»

«Σπαρκ, κοίτα με...» μου είπε ο Βασι και με έπιασε από τους ώμους. «Ξέρω ότι είσαι αθώος. Δεν ξέρω γιατί, αλλά έχω μία διαίσθηση ότι κάποιος σε έχει ενοχοποιήσει. Το συζήτησα με τον Κουπ και την Μπλαιρ και είπαμε να μαζέψουμε οι τρεις μας λεφτά για να σε ελευθερώσουμε, ώστε να μας πεις τι πραγματικά έγινε»

«Και γιατί δεν σήκωσες ούτε ένα τηλέφωνο που σε πήρα;» τον ρώτησα εκνευρισμένος.

«Βλέπεις... δεν θα μας άφηναν να μιλήσουμε αρκετή ώρα στο τηλέφωνο, οπότε ύλπιζα να συζητούσαμε όταν σε βγάζαμε από...»

«Ευχαριστώ Βασι, αλλά αν είναι ακόμη ένα σχέδιο του Μαλχη αυτό, εγώ δεν έχω όρεξη για συνομιλίες»

«Του Μαλχη; Που κολλάει ο Μαλχη;» με κοίταξε προβληματισμένος ο Βασι.

«Φρουρέ, σε παρακαλώ βγάλε μου το βραχιόλι για να πάω σπίτι μου»

«Μη με διατάζεις...» είπε και ξεκλείδωσε το βραχιόλι.

«Δεν θα μας πεις τουλάχιστον ένα ευχαριστώ;» με ρώτησε ο Βασι.

«Ναι έχει δίκιο, Σπαρκ» είπε και ο Κουπ.

«Δεν έχω καμία όρεξη σας είπα. Πρέπει να τελειώσω μία μπίζνα που έχω αφήσει ανοιχτή»

«Τι μπίζνα;» με ρώτησε ο Κουπ προβληματισμένος.

Δεν του απάντησα. Προχώρησα προς την έξοδο της φυλακής. Στον δρόμο είδα τον ψηλό άντρα που με είχε τραυματίσει την πρώτη μέρα.

«Ωχ, κοιτάξτε ποιος βγήκε από το κελί του. Φύτρωσε και μουσάκι. Ήρθες για δεύτερο γύρο;» με ρώτησε. Εγώ δεν του απάντησα και περπάτησα μερικά μέτρα μακριά του στον διάδρομο.

«Αρχηγέ κοίτα, δεν έχει βραχιόλι» είπε ένας άλλος άντρας πίσω από τον ψηλό άντρα και όλοι έκαναν ένα βήμα πίσω καθώς τους προσπέρασα. Με φοβήθηκαν επειδή μπορούσα να χρησιμοποιήσω τώρα ελεύθερα την μαγεία μου, ενώ αυτοί δεν μπορούσαν; Προσπάθησα να διαβάσω τις σκέψεις του ψηλού άντρα που αποκάλεσαν αρχηγό για να μάθω ως προς τι ο φόβος.

\textasciitilde ...ώστε βγήκε αυτός ε; Πρέπει να έχει πολύ καλούς φίλους...\textasciitilde 

Πολύ καλούς φίλους; Μακάρι. Αν ήταν τόσο καλοί οι φίλοι μου δεν θα τους έπαιρνε 10 χρόνια να με βγάλουν από τη φυλακή, ή ακόμη καλύτερα δε θα με είχαν προδώσει τη στιγμή του εγκλήματος. Θα είχαν κάτσει να ακούσουν τουλάχιστον τι έχω να πω, την δική μου μεριά της ιστορίας, θα είχαν παλέψει για να γίνει δίκη με εμένα και τον Μαλχη, όμως τίποτα από αυτά δεν συνέβει. Και το καλύτερο απ’ όλα; Ήρθαν μετά από 10 χρόνια να μου πουν ότι “νομίζουν” ότι είμαι αθώος. Ή πώς το είπε; “Δεν ξέρω γιατί, αλλά έχω μία διαίσθηση ότι κάποιος σε έχει ενοχοποιήσει και το συζήτησα με τους υπόλοιπους, μπλα μπλα...”, χαζομάρες δηλαδή. Κάποιο κόλπο του Μαλχη θα είναι πάλι για να με εκβιάσει και μένα να του δώσω την μαγεία μου. Αν και δεν εκβίασε τον Φυτ και την Ιφι για να του δώσουν τη μαγεία τους. Περισσότερο τους παραπλάνησε παρά τους εκβίασε είναι η αλήθεια, αλλά ποιος νοιάζεται. Ένας είναι ο στόχος μου τώρα, να σκοτώσω τον Μαλχη. Δεν με νοιάζει πώς, δεν με νοιάζει αν θα ξαναμπώ άλλα 100 χρόνια φυλακή, η δικαιοσύνη πρέπει να νικήσει.

Καθώς βγήκα από το τείχος της φυλακής και έφτασα στο εξωτερικό πάρκινγκ, είδα ένα μοναχικό αμάξι. Στην θέση του οδηγού καθόταν η Μπλαιρ, που ούτε αυτή είχε αλλάξει εμφανισιακά εδώ και 10 χρόνια. Είχε την χαρακτηριστική κοτσίδα της και το σακάκι της Ειδικής Υπηρεσίας και φαινόταν λες και περίμενε κάποιον, λογικά τον Βασι, τον Κουπ και... εμένα; Το βλέμμα της στράφηκε πάνω μου. Με το που με κοίταξε, απέφυγα την οπτική επαφή, καθώς ήθελα να επιστρέψω σπίτι μόνος, χωρίς να με περισπάσει κανείς από τον στόχο μου. Δεν άκουσα την Μπλαιρ να μου λέει τίποτα καθώς περπατούσα, αλλά ήμουν σίγουρος πως με παρατηρούσε καθώς προχωρούσα προς τον δρόμο μόνος μου, φορώντας μία σκισμένη φανέλα και ένα κοντό πορτοκαλί παντελόνι, όχι από τη σχεδίαση του, αλλά επειδή το είχα κόψει για να μπορώ να έχω κάτι να σβήνω τις κιμωλίες όταν έκανα λάθος συλλογισμό. Μπορεί βέβαια να μην με έχει αναγνωρίσει καν και να είναι προβληματισμένη που ένας μουσάτος φυλακισμένος με γδαρμένα ρούχα πηγαίνει μόνος του προς την πόλη.

Όλα αυτά τώρα είναι ασήμαντοι συλλογισμοί. Δέκα χρόνια έχω ένα πράγμα στο μυαλό μου, τον θάνατο του Μαλχη, την απόδοση της δικαιοσύνης. Με τους συλλογισμούς μου δεν κατάφερα να φτάσω πολύ μακριά, παρόλο που είχα άπλετο χρόνο και απαρίθμησα όντως κάθε πιθανό σενάριο που θα μπορούσε να συμβαίνει. Το πιο ρεαλιστικό σενάριο είναι αυτό που σκεφτόμουν από την αρχή, ο Μαλχη να συνδέεται με αυτή τη γυναίκα, είτε οικογενειακός, είτε ως φίλοι, και να της έδειξε τις χειρονομίες μεταφοράς μαγείας από έναν άνθρωπο σε άλλο και αυτή η γυναίκα, με την βοήθεια του Μαλχη, απήγαγε παιδιά και τα φρόντισε μέχρις ότου αυτά να της δώσουν την μαγεία τους, με τον έναν τρόπο ή τον άλλο.

Δεν καταλαβαίνω όμως ο Μαλχη τι έχει να κερδίσει από όλο αυτό; Πήρε τις μαγείες της Ιφι και του Φυτ, ναι, αλλά τι θα μπορούσε να κερδίσει καλύπτοντας την γυναίκα με τις πολλές μαγείες; Χρήματα; Δόξα; Ποιο μπορεί να είναι το κίνητρο του για να συνεχίζει να την βοηθάει; Βέβαια δεν ξέρω τι έχει γίνει μέσα σε αυτά τα δέκα χρόνια, αλλά δεν νομίζω να έχουν σταματήσει το σχέδιο τους. Πρέπει να πάω σπίτι γρήγορα για να αλλάξω ρούχα και να σκεφτώ τι θα κάνω. Κρίνοντας από τη θέση του ηλίου είναι μεσημέρι, γύρω στις 4:00 θα έλεγα, και θέλω άλλα 50 λεπτά περίπου για να φτάσω στην πόλη με τα πόδια. Πρέπει να φτάσω σπίτι μου γρήγορα για να τον προλάβω.

\chapter{}

«Για να δούμε...» είπα και σήκωσα το πατάκι του σπιτιού μου. Θυμάμαι από μικρός ότι κρύβαμε το κλειδί κάτω από το πατάκι, αλλά δεν θυμάμαι για κάποιο λόγο να ξεκλειδώνω την πόρτα. Μπορεί να έχει περάσει πολύς καιρός και να έχουν σβηστεί οι βασικές μου αναμνήσεις, ποιος ξέρει. Πάντως, το γεγονός ότι το κλειδί ήταν έξω από το σπίτι σήμαινε πως το σπίτι ήταν άδειο και η μητέρα μου δεν είχε γυρίσει από την τρέχουσα αποστολή της. Καλύτερα για μένα, δεν έχω όρεξη να εξηγώ στην μητέρα μου τις προθέσεις μου αυτή τη στιγμή.

Όταν άνοιξα την πόρτα, ένα αεράκι νοσταλγίας με διαπέρασε. Βρισκόμουν μετά από 10 χρόνια στο σπίτι το οποίο μεγάλωσα. Ακούμπησα τα παπούτσια μου στη γωνία όπου τα ακουμπούσα πάντα από συνήθεια. Προχώρησα λίγο πιο μέσα και είδα το κυκλικό τραπέζι στο οποίο καθόμουν κάθε μέρα όταν ήμουν μικρός και διάβαζα. Ακόμη δεν έχω καταλάβει που χρησιμεύει η Ιστορία, άχρηστο μάθημα. Κοντά στο τραπέζι ήταν και η κουζίνα, η οποία δεν είχε αλλάξει καθόλου από την τελευταία φορά που τη θυμόμουν. Δεν είχε και τίποτα το ιδιαίτερο, μία τυπική κουζίνα με έναν ηλεκτρικό φούρνο, έναν νεροχύτη και τέσσερα μάτια ήταν, αλλά μπορούσα να ακούσω τους ήχους που θα έκαναν οι κατσαρόλες και τα μαχαίρια αν οι γονείς μου ήταν εδώ και μαγείρευαν βραδινό.

Γιατί μου έρχονται τώρα τέτοιες σκέψεις; Αποσπώμαι από τον στόχο μου, μη χτυπήσω το σπανάκι μου. Έτσι δεν θα έλεγε και ο Μαλχη; Έχουν περάσει και τόσα χρόνια, δεν θυμάμαι και τι ακριβώς έλεγε, αλλά ήταν σίγουρα κάτι με σπανάκι.

Συνέχισα το τουρ του σπιτιού μου περπατώντας προς το δωμάτιο μου. Πριν φτάσω καν όμως στο δωμάτιο μου, είδα τον μοιραίο καθρέφτη, τον καθρέφτη που άλλαξε την πορεία της ζωής μου, τον καθρέφτη μπροστά στον οποίο έμαθα ότι μπορώ να διαβάζω σκέψεις. Σταμάτησα για λίγο να κοιτάξω τον εαυτό μου. Πώς είχα γίνει έτσι; Γιατί είμαι τόσο λεπτός; Πότε φύτρωσε αυτό το μούσι; Έχω κοντίνει και λίγο, μπορεί επειδή καθόμουν συνέχεια διπλωμένος στο πάτωμα και έκλαιγα την μοίρα μου. Και το μαλί μου γιατί είναι τόσο μακρύ; Θα μου πεις, έχουν περάσει 10 χρόνια χωρίς να έχω φροντίσει τον εαυτό μου οπότε... πάλι αποσπώμαι.

«Ξύπνα Σπαρκ!» είπα και χαστούκισα τον εαυτό μου για να συγκεντρωθώ.

Μπαίνοντας στο δωμάτιο μου, άνοιξα το συρτάρι του κομοδίνου μου και βρήκα αυτό που έψαχνα, το εφεδρικό πιστόλι  που μου είχε δώσει πριν τις πρώτες μου ανακρίσεις ο Μαλχη. Δεν θα ήταν ειρωνικό να πεθάνει από την ίδια του την μαγεία; Όταν σήκωσα το όπλο, βρήκα από κάτω κάτι το οποίο με προβλημάτισε, το κινητό μου. Γιατί βρισκόταν εκεί; Πώς έφτασε στο κομοδίνο μου; Δεν μπορώ να ανακαλέσω την ανάμνηση ότι μου το παίρνουν. Μπορεί να έχει περάσει πολύς καιρός. Το άνοιξα για να δω αν είχε μπαταρία και... ήταν γεμάτο; Όχι πλήρως, αλλά στο 93\% δεν το λες και άδειο. Το άνοιξα για να δω αν είχα κανένα εισερχόμενο μήνυμα, όμως δεν υπήρχε ούτε μία ειδοποίηση. Είδα τουλάχιστον την ώρα και ήταν 5:27. Απογοητευμένος, έβαλα το κινητό μου στην τσέπη και το πιστόλι κάτω από τα ρούχα μου βγαίνοντας από το δωμάτιο. Δεν άλλαξα ρούχα βέβαια, αλλά πάλι αποσπάστηκα. Πρέπει να τελειώσω όσο το δυνατόν γρηγορότερα με τον Μαλχη για να μπορέσω να ηρεμήσω.

Έβαλα πάλι βιαστικά τα παπούτσια μου, τα οποία έβγαλα χωρίς λόγο, αφού κάθησα 2 λεπτά στο σπίτι, και βγήκα από το σπίτι, κλείδωσα την πόρτα και έφυγα τρέχοντας για τα κεντρικά της Ειδικής Υπηρεσίας. Ήξερα ότι αν ήταν κάπου ο Μαλχη αυτή τη στιγμή θα ήταν στην Ειδική Υπηρεσία. Τώρα αν θα ήταν στην αίθουσα εξάσκησης, στο γραφείο του ή στο γραφείο κάποιας ομάδας δεν με ενδιέφερε και τόσο. Το σχέδιο μου ήταν απλό: θα έμπαινα, θα έκανα τη δουλειά μου, και θα έφευγα με οποιονδήποτε τρόπο, δηλαδή είτε να με συλλάβουν εκείνη τη στιγμή, είτε να “αποδράσω”.

Καθώς έτρεχα, ένιωσα μερικές σταγόνες να πέφτουν στην ήδη λερωμένη μου φανέλα. Κοίταξα στον ουρανό και ο ήλιος είχε καλυφθεί τελείως από ένα στρώμα γκρίζων σύννεφων.

«Ατυχία...» ψιθύρισα στον εαυτό μου όσο έτρεχα. Δεν πειράζει όμως. Αφού μπόρεσα να αντέξω δέκα χρόνια μόνος με τις σκέψεις μου μέσα σε ένα γκρίζο κελί, γιατί να μην μπορώ να αντέξω λίγες σταγόνες μέχρι να φτάσω στην Ειδική Υπηρεσία; Η ένταση της βροχής όλο και αυξανόταν καθώς πλησίαζα τα κεντρικά, πράγμα το οποίο θα μου έκανε δύσκολη την έξοδο αν χρειαζόταν να ξεφύγω. Δεν πειράζει όμως, και η φυλακή δεν ήταν άσχημη. Εξάλλου, δεν θα μπορεί κάποιος να με προδώσει στην φυλακή.

Φτάνοντας στα κεντρικά της Ειδικής Υπηρεσίας, έσπρωξα γρήγορα την περιστρεφόμενη πόρτα και έτρεξα προς την γραμματέα. Αυτή φάνηκε προβληματισμένη βλέποντας έναν άντρα με γδαρμένα ρούχα να την πλησιάζει, αλλά διατήρησε την ψυχραιμία της και δεν σηκώθηκε από τη θέση της.

«Καλωσήρθατε. Πείτε μου το όνομα σας και πώς μπορώ να σας βοηθήσω» μου είπε η γραμματέας.

«Που είναι ο Μαλχη;»

«Ποιος;» μου απάντησε η γραμματέας.

«Μη με κοροϊδεύεις, ξέρω ότι πιθανότατα είναι πίσω από αυτή την πόρτα» της απάντησα και έδειξα την πόρτα του γραφείου του.

«Όχι κύριε, δεν καταλάβατε. Δεν υπάρχει κάποιος Μαλχη στην Ειδική Υπηρεσία» μου απάντησε αυτή με τη σειρά της.

«Είπα μη με κοροϊδεύεις» φώναξα και χτύπησα το χέρι μου στο γραφείο της. Όλοι οι τριγύρω πράκτορες έστρεψαν το βλέμμα τους πανω μου. «Ξέρω ότι ο Μαλχη είναι ο αρχηγός της Ειδικής Υπηρεσίας, αλλά ξέρω και πράγματα που δεν ξέρεις εσύ»

«Πώς ξέρεις ότι... τέλος πάντων. Δώσε μου το όνομα σου και πώς μπορώ να σε βοηθήσω» είπε ψύχραιμα η γραμματέας.

«Σου είπα, πες μου απλά που είναι ο Μαλχη»

«Γιατί; Δεν μπορώ να δίνω τέτοιες πληροφορίες σε ξένους»

«Δεν είμαι ξένος» είπα εκνευρισμένος.

«Ε πες μου πώς σε λένε τότε για να μπορέσω να σε βοηθήσω»

«Σπαρκ Θοτς»

«Σπαρκ... Θοτς...» είπε καθώς έγραφε το όνομα μου στον υπολογιστή της. «Συγγνώμη κύριε, αλλά δεν μπορώ να σας βοηθήσω»

«Τι εννοείς δεν μπορείς να με βοηθήσεις;» την έπιασα από το κολάρο του πουκάμισου που φόραγε. «Πες μου που είναι ο Μαλχη μην...»

«Κάποιος; Βοήθεια!» φώναξε η γραμματέας και δύο πράκτορες με έπιασαν από τα χέρια και άρχισαν να με σέρνουν προς την είσοδο.

«Αφήστε με κάτω!» φώναξα απελπισμένος. Αυτοί δεν μου απάντησαν.

\textasciitilde ...αύρα κλάσης Α. Αν ήθελε να μας επιτεθεί θα το είχε κάνει ήδη. Γιατί δεν επιτίθεται;...\textasciitilde \  σκεφτόταν ο ένας πράκτορας.

Προφανώς και δεν μπορούσα να επιτεθώ με την μαγεία μου, αφού δεν είμαι ειδικός στις μάχες, ένας απλός ανακριτής είμαι... ήμουν... δεν έχει σημασία. Το σχέδιο μου απέτυχε γιατί δεν έμεινα ήρεμος. Γιατί δεν μπόρεσα να κρατήσω και γω την ψυχραιμία μου όπως η γραμματέας; Αν και...

«Ωπ, καλησπέρα Ιφι» είπε ένας από τους δύο πράκτορες που με έσερναν.

Ιφι; Πιθανότατα κάποια συνονόματη θα ήταν, δεν υπάρχει περίπτωση να... Έστρεψα το βλέμμα μου προς την κοπέλα που μίλησε ο πράκτορας και είδα τα κυματιστά μελαχρινά μαλλιά της, τα πράσινα μάτια της, το άψυχο πρόσωπο που θυμόμουν σε εκείνο το υπόγειο να χαμογελάει για ακόμη μία φορά. Πώς; Τι συμβαίνει; Είμαι ζωντανός; Πέθανα από την πείνα στην φυλακή; Που... είμαι...

«Ιφι...» της απευθύνθηκα.

Αυτή δεν μου απάντησε.

Σηκώθηκα γρήγορα κλώτσησα τον έναν πράκτορα ανάμεσα από τα μπούτια του.

«Τι κάνεις;» μου είπε ο άλλος πράκτορας, τον οποίο αφόπλισα ρίχνοντας του μία μπουνιά στο στομάχι με το αριστερό μου χέρι και μία ακόμη στο σαγόνι με το δεξί μου χέρι.

«Ιφι... μου έλειψες. Δεν ξέρω πώς, ή πότε, ή... γιατί, αλλά ευτυχώς είσαι ζωντανή» είπα καθώς την αγκάλιασα σφιχτά. Αυτή δεν έδειχνε ότι θέλει να με αγκαλιάσει πίσω. Έκανα ένα βήμα πίσω και την έπιασα από τους ώμους.

\textasciitilde ...ποιος είναι αυτός τώρα; Τι θέλει από μένα;...\textasciitilde 

«Εγώ είμαι Ιφι» απάντησα στις σκέψεις της.

«Συγγνώμη, αλλά... ποιος είσαι;»

«Ποιος... είμαι;... Ο Σπαρκ. Δεν με θυμάσαι; Εγώ και εσύ ήμασταν στην ίδια ομάδα πριν αρκετά χρόνια και σε μία αποστολή... πέθανες»

\textasciitilde ...τι βλακείες λέει αυτός;\textasciitilde «Θα το θυμόμουν αν ήμασταν στην ίδια ομάδα, και σίγουρα θα το θυμόμουν αν είχα πεθάνει...» με ειρωνεύτηκε και έδιωξε τα χέρια μου από τους ώμους της. «Συγγνώμη τώρα, αλλά...»

«Ιφι άκουσε με» την έπιασα από το χέρι.

«Άσε με» φώναξε και τράβηξε το χέρι της.

«Τι γίνεται εδώ; Ιφι είσαι καλά;» άκουσα μία γνώριμη φωνή και ένιωσα ρίζες να τυλίγουν τα πόδια μου. Έσκυψα και προσπάθησα να τις βγάλω από τα πόδια μου, αλλά η προσπάθεια μου ήταν μάταιη. 

«Ναι» είπε και τον πλησίασε. «Απλά αυτός ο άντρας...»

«Σε ενοχλεί, ε; Δεν πειράζει θα τον καθαρίσω εγώ»

Κοίταξα πάνω και είδα άλλο ένα πρόσωπο που με προβλημάτισε. Δεν ήταν ο Μαλχη, αλλά ο Φυτ. Πώς;...

«Φυτ!» του είπα προσπαθώντας να ακουστώ χαρούμενος.

«Σε ξέρω;»

«Ναι είμαι ο Σπαρκ... δεν με θυμάσαι;»

«Όχι» είπε και χαλάρωσε τη λαβή των ριζών του από τα πόδια μου.

«Ευχαριστώ που με άφησες ελεύθερο» είπα καθώς τον πλησίαζα. «Βλέπεις...»

Ο Φυτ έβγαλε ρίζες από τα χέρια του και τις τύλιξε γύρω από τη μέση μου, πιάνοντας και τα χέρια μου.

«Γιατί δεν με ακούς; Άσε με να σου μιλήσω λίγο»

«Όχι» είπε με πέταξε έξω από την πόρτα με την μαγεία του. Που την βρήκε τόση δύναμη;

Έξω από το κτίριο έβρεχε πολύ πιο έντονα από πριν. Στην προσπάθεια μου να σηκωθώ, γλίστρησα δύο φορές στο πεζοδρόμιο που ήταν γεμάτο νερά. Όταν σηκώθηκα, είδα τον Φυτ μπροστά μου, έξω από την πόρτα και αυτός, να με κοιτάζει επίμονα.

«Ευτυχώς που ήρθες. Τώρα θα με ακούσεις;»

«Όχι...» είπε και με κλώτσησε σε έναν λάκκο που βρισκόταν πίσω μου. «...και μην ξανάρθεις» είπε και επέστρεψε μέσα στο κτίριο.

Σηκώθηκα από το λάκκο και κάθησα για λίγη ώρα όρθιος έξω από τα κεντρικά της Ειδικής Υπηρεσίας σκεπτόμενος. Πώς γίνεται αυτό; Να ζει η Ιφι και ο Φυτ; Πώς γίνεται να μην με θυμάται κανείς στην Ειδική Υπηρεσία, ενώ ο Βασι και ο Κουπ με θυμούνται; Λες να έχει να κάνει με την εμφάνιση μου; Μα όχι, δεν μπορεί να μην μπορούν να ανακαλέσουν ούτε το όνομα μου, ενώ ο Βασι και ο Κουπ το θυμόντουσαν. Τι... μου συμβαίνει;

Ξεκίνησα να περπατάω προς το σπίτι μου, απογοητευμένος που η αποστολή μου είχε αποτύχει. Κοίταξα τον εαυτό μου στη τζαμαρία ενός μαγαζιού. Γύρω μου υπήρχαν μικρές σπίθες κεραυνών, όπως εκείνη τη μέρα που πέθανε ο πατέρας μου. Αυτό, θυμάμαι ότι είχε πει η μητέρα μου συμβαίνει όταν έναν άνθρωπος είναι ψυχολογικά ασταθής, ή κάτι ανάλογο. Λες να σχετίζεται με την αύρα; Γιατί δεν μου το είχε εξηγήσει πότε η μητέρα μου έτσι; Με ότι και να σχετίζεται, είναι επικίνδυνο ενώ βρέχει να έχω ένα ηλεκτρικό πεδίο γύρω μου, πότε δεν ξέρεις τι μπορεί να συμβεί, αν και δεν έχω ακούσει κανέναν κεραυνό ακόμη.

Τη στιγμή που το σκέφτηκα αυτό, ακούστηκε ένας κεραυνός και είδα ένα άσπρο φως από τον ουρανό.

«Τέλεια...» ψιθύρισα στον εαυτό μου. «Τι άλλο μπορεί να μου συμβεί;»

Στο βάθος είδα ένα γνώριμο αυτοκίνητο να πλησιάζει. Από μακριά φαινόταν λες και μέσα υπήρχαν τρεις φιγούρες. Το αυτοκίνητο αυτό σταμάτησε σε ένα φανάρι. Που το είχα ξαναδεί αυτό το αυτοκίνητο και γιατί η προσοχή μου στράφηκε σε αυτό; Καθώς πλησίασα το σταματημένο αυτό αυτοκίνητο, κατάλαβα γιατί μου φάνηκε γνώριμο. Στην θέση του οδηγού βρισκόταν η Μπλαιρ και στις πίσω θέσεις κάθονταν ο Βασι και ο Κουπ. Τι έκαναν τέτοια ώρα στον δρόμο; Ήταν στις φυλακές τόση ώρα; Μπορεί, αφού δεν τους είχα δει να με προσπερνάνε όσο περπατούσα σπίτι μου από τις φυλακές.

Το φανάρι τους άναψε πράσινο, το αυτοκίνητο προχωρούσε προς το μέρος μου. Προσπάθησα να τους κάνω νόημα για να σταματήσουν, αλλά η Μπλαιρ όχι μόνο δε σταμάτησε, αλλά προσπέρασε έναν λάκκο με ταχύτητα και πετάχτηκε προς τα πάνω μου ένα κύμα από λάσπη. Κοίταξα το αυτοκίνητο για να δω αν κάποιος κοιτάξει πίσω, αλλά κανείς δεν με κοίταζε. Κανείς δεν πρόσεξε καν τον πεζό που είχαν μόλις βρέξει, έστω κι από περιέργεια.

Απογοητευμένος τώρα με όλους τους ανθρώπους που θεωρούσα “φίλους”, συνέχισα να περπατώ προς το σπίτι μου. Όλα ήταν μάταια πλέον. Κανείς δεν είχε την υπομονή να με ακούσει, ήμουν μούσκεμα και είχα άλλα πέντε λεπτά περπάτημα μέσα στην βροχή για να φτάσω σπίτι μου. Τι άλλο θα μου συμβεί...

«Αααα!»

\chapter{}

Η βροχή είχε σταματήσει όταν άνοιξα πάλι τα μάτια μου. Κοίταξα γύρω μου και είδα τρεις άγνωστους να με περιτριγυρίζουν και να με κοιτούν επίμονα. Όταν σηκώθηκα και ξανάρχισα να περπατάω ξανά προς το σπίτι μου, όλοι τους έκαναν μερικά βήματα προς τα πίσω και άρχισαν να συζητούν προβληματισμένοι, δεν μπορούσε όμως να με ενδιαφέρει λιγότερο τι συζητούσαν. Δεν είχα απομακρυνθεί και πολύ πριν ένας από τους αγνώστους αυτούς με πλησίασε και με σταμάτησε, ακουμπώντας το χέρι του στον ώμο μου.

«Είστε καλά κύριε;»

«Καλύτερα από ποτέ...» τον ερωτεύτηκα και συνέχισα να περπατώ.

«Πώς γίνεται αυτό; Αφού μόλις σας χτύπησε κεραυνός» είπε κόβοντας μου τον δρόμο και δείχνοντας τα σύννεφα στον ουρανό.

«Κεραυνός;»

«Ναι. Σας είδαμε από μακριά και καλέσαμε ένα ασθενοφόρο για να σας...»

«Χίλια συγγνώμη αλλά δεν έχω όρεξη για φάρσες. Αν με είχε χτυπήσει κεραυνός δεν θα μπορούσα να περπατάω αυτή τη στιγμή, οπότε κάτι πρέπει να μου κρύβεται...» του απάντησα και συνέχισα να προχωράω.

«Μα σας λέω την αλήθεια, γιατί δεν με ακούτε;»

Δεν του απάντησα. Ούτε καν προσπάθησα να διαβάσω τις σκέψεις του για να συμπεράνω αν μου έλεγε αλήθεια ή όχι. Ήθελα απλά να πάω στο σπίτι μου και να σκεφτώ την επόμενη μου κίνηση. Αυτός, βλέποντας μάλλον ότι η υγεία μου δεν ήταν σε άμεσο κίνδυνο, δεν ξαναπροσπάθησε να με σταματήσει και επέστρεψε στους άλλους δύο. Κεραυνός; Πολύ καλό, πραγματικά πώς τους ήρθε. Κοίταξα τα χέρια μου και στη συνέχεια τα πόδια μου προσεκτικά. Δεν είχα τραυματιστεί πουθενά, ούτε εξωτερικά, ούτε εσωτερικά, οπότε σίγουρα λένε ψέματα.

Και οι σπίθες κεραυνών που είχα πριν; Ξανακοίταξα τα χέρια μου. Είχαν εξαφανιστεί, οπότε δεν υπάρχει περίπτωση να με πέτυχε κεραυνός, αφού είμαστε σε πόλη, περπατάμε ανάμεσα σε ουρανοξύστες, άρα πάλι έχω δίκιο, λένε ψέματα. Τι έχουν να κερδίσουν όμως από αυτό; Θυμήθηκα τα γεγονότα που συνέβησαν στο υπόγειο εκείνη τη μέρα. Τι είχε να κερδίσει ο Μαλχη από αυτό; Ακόμη δεν μπορώ να καταλάβω. Καλύτερα να σταματήσω να υπεραναλύω τις σκέψεις μου μέχρι να φτάσω σπίτι.\\\\

Ξεκλειδώνοντας την πόρτα και βγάζοντας τα παπούτσια μου αυτή τη φορά μου ήρθε μία αίσθηση απογοήτευσης. 

«Από όταν ήμουν μικρός, ήμουν μία απογοήτευση. Πώς γίνεται να είναι και οι δύο γονείς μου κλάσης Α και εγώ να γεννιέμαι κλάσης Β; Πώς γίνεται να μου λένε πως έχω αύρα κλάσης Α, ενώ δεν μπορώ ούτε έναν κεραυνό να παράξω; Πώς... απλά πώς; Είμαι απλά άχρηστος.

Και η ειδικότητα μου ως ανακριτής; Στις περισσότερες αποστολές όχι μόνο δεν μπόρεσα να ανακρίνω κανέναν, αφού όλο και κάποιος άλλος αναλάμβανε την δουλειά του ανακριτή επειδή εγώ δεν μπορούσα να σταθώ στο ύψος των περιστάσεων. Άσε που ήμουν και ο λόγος που η ομάδα είχε προβλήματα, ένα βάρος που οι υπόλοιποι έπρεπε να κουβαλήσουν στις αποστολές επειδή η ικανότητα μου να ανταποκριθώ στις μάχες ή γενικότερα στις αλληλεπιδράσεις μου με τους ενόχους ήταν άθλια.»

Περπάτησα προς το σαλόνι, ακούμπησα το όπλο μου στο τραπεζάκι και κάθησα στον καναπέ.

«Και να μη μιλήσω για τα σχέδια μου...» συνέχισα να μιλάω στον εαυτό μου. «Όλα τους, ακόμη και το σημερινό, ήταν της μορφής: πηγαίνει μία ομάδα, διαβάζω τις σκέψεις του ύποπτου και καταλήγουμε σε συμπέρασμα. Όμως ποτέ δεν πήγαινε έτσι. Όλο κάπως έβρισκα έναν τρόπο να τα κάνω μαντάρα»

Έσκυψα προς τα εμπρός και ακούμπησα τις παλάμες μου απογοητευμένος στο πρόσωπο μου. Ένα δάκρυ μου ξέφυγε και κύλησε στο αριστερό μου μάγουλο.

«Όσοι νόμιζα πως με βοηθούσαν... ήταν απλά διπρόσωποι ηθοποιοί. Όσοι άνθρωποι νόμιζα ότι ήταν φίλοι μου, δεν με αναγνωρίζουν καν πλέον. Ακόμη και οι τρεις “φίλοι” μου, που πλήρωσαν για να βγω από την φυλακή, δεν μπορούσαν καν να με κοιτάξουν στα μάτια, ή να σταματήσουν τουλάχιστον το αυτοκίνητο τους ενώ με βλέπουν ότι περπατάω μόνος μου μέσα στη βροχή και να με συνοδεύσουν. Σαν να μην ξέρουν καν... ποιος είμαι...

Και τα δέκα χρόνια στη φυλακή; Τι ηλιθιότητες καθόμουν και έγραφα στους τοίχους, τι συνδέσεις και χαζομάρες, τι υποθέσεις και κουραφέξαλα. Δέκα χρόνια χωρίς νόημα, γεμάτα άχρηστες σκέψεις και συνειρμούς που δεν οδηγούσαν πουθενά. Και αυτός ο Πιτ ήταν πολύ ύπουλος. Με τάιζε μία φορά στις τρεις μέρες. Ποιο το νόημα; Καλύτερα να είχα πεθάνει της πείνας, από το να ζω οριακά και να υποφέρω...»

Έβγαλα το κινητό μου από την τσέπη μου και κάλεσα το μόνο άλλο άτομο που θα μπορούσε να με βοηθήσει, την μητέρα μου. Η προσπάθεια μου όμως ήταν μάταιη. Εκνευρισμένος, σηκώθηκα, πήρα φόρα και πέταξα το κινητό μου στον τοίχο, το οποίο εξερράγη σε δεκάδες μικρά κομμάτια κατά την πρόσκρουση. Κάθησα πάλι πίσω στον καναπέ, έπιασα το πιστόλι μου και άρχιζα να το επεξεργάζομαι.

«Ακόμη και η μητέρα μου είναι... προδότρα. Μάλλον ο Μαλχη θα της είχε πει το γνωστό του ψέμα, ότι σκότωσα δύο πράκτορες και τον τραυμάτισα στην αποστολή μας στην εκκλησία. Αλλά η ίδια η μητέρα μου να γυρνάει πλάτη στο δικό της παιδί και να μην έχει την διάθεση να με ακούσει, ούτε όταν ήμουν στη φυλακή, αλλά ούτε και τώρα, είναι λίγο... παράξενο. Αλλά, αφού η μητέρα μου δουλεύει στην Ειδική Υπηρεσία, δεν είδε πως ο Φυτ και η Ιφι είναι ζωντανοί;

Άλλο και τούτο. Πώς γίνεται να είναι ζωντανοί; Αφού τους είδα ζωντανά να κάνουν τις χειρονομίες που μας είχε δείξει ο Μαλχη και να χάνουν τον παλμό τους. Πώς γίνεται επίσης να μην έχουν χάσει τις μαγείες τους. Ο Φυτ, για να με διώξει από την Ειδική Υπηρεσία χρησιμποιήσε τις ρίζες του για να με ακινητοποιήσει και στη συνέχεια... 

Πραγματικά δεν ξέρω που είμαι πλέον. Δεν ξέρω ποιος είμαι πλέον. Δεν ξέρω ποιο είναι το νόημα πλέον. Κανείς δεν με χρειάζεται, δεν έχω φίλους, δεν έχω οικογένεια, δεν έχω τίποτα. Δεν υπάρχει κάτι πλέον που μπορώ να κάνω για να προσφέρω στην κοινωνία. Αν η ύπαρξη μου είναι όντως ένα πλην για όλους τους υπόλοιπους, γιατί να μην σταματήσω να υπάρχω;»

Κοίταξα την κάννη του πιστολιού μου, η οποία με τη σειρά της κοίταζε επίμονα.

«Και στο κάτω κάτω, ποιος θα λυπηθεί; Ο Μαλχη; Η μητέρα μου; Η Ιφι, ο Φυτ ή κάποιος από τους άλλους τρεις που δεν με αναγνωρίζουν καν; Ας γελάσω...» είπα και σήκωσα το πιστόλι στο ύψος του κεφαλιού μου. «Πατέρα... νομίζω πως κατάλαβα τις λέξεις σου επιτέλους...» είπα και τράβηξα την σκανδάλη...\newpage\newpage

...όμως το πιστόλι δεν ήταν γεμάτο. Κατέβασα το χέρι μου και κοίταξα το όπλο. «Ακόμη και να αυτοκτονήσω δεν μπορώ. Ακόμη σε αυτό το σενάριο ήμουν... άχρηστος» είπα και σηκώθηκα από τον καναπέ. Πήγα στο δωμάτιο μου με το πιστόλι στο χέρι. Αν υπήρχαν κάπου σφαίρες, θα ήταν εκεί που είχα βρει το όπλο μου, στο κομοδίνο. Ανοίγοντας το, βρήκα έναν γεμιστήρα ο οποίος ήταν μισοάδειος. Πώς δεν τον είχα παρατηρήσει την πρώτη φορά; Πώς δεν είχα παρατηρήσει ότι το όπλο ήταν άδειο; 

«Πάλι αποσπώμαι. Μία σφαίρα χρειάζομαι...» είπα και γέμισα το πιστόλι μου. Το βλέμμα μου στράφηκε στο παράθυρο μου. Τα σύννεφα είχαν εξαφανιστεί και ο ουρανός είχε μαυρίσει. Κάθησα στο κρεβάτι μου για μία τελευταία φορά και κοίταξα το φεγγάρι, μόνο του στον ουρανό, χωρίς αστέρια ή σύννεφα να το περιτριγυρίζουν πλέον. 

«Τι σύμπτωση είναι αυτή...» είπα και κοίταξα το όπλο μου. «...τι κάνω ρε φίλε;» είπα και έπιασα το κεφάλι μου. Δεν μπορούσα να τελειώσω τον εαυτό μου, μετά από την πρώτη μου αποτυχία. Αυτή την φορά ήξερα πως αν το δάχτυλο μου τραβούσε τον μοχλό του όπλου μου, όλα θα είχαν τελειώσει, αλλά αυτό δεν με ευχαριστούσε. Εκείνη τη στιγμή δεν ήξερα τι συναίσθημα να νιώσω. Απογοήτευση; Θλίψη; Θυμό; Κανένα ονοματισμένο συναίσθημα δεν ταίριαζε στην κατάσταση που βρισκόμουν

«Θα σου δώσω άλλη μία ευκαιρία Σπαρκ» είπα στον εαυτό μου και σηκώθηκα από το κρεβάτι. «Θα πας στο καζίνο. Υπάρχει υψηλή πιθανότητα να πετύχεις τον Μαλχη εκεί, καθώς κάθε δεύτερη μέρα παράταγε την δουλειά του για να τζογάρει. Αν δεν μπορέσεις ούτε σήμερα, ούτε μέσα στις επόμενες μέρες να τον βρεις και να τελειώσεις την δουλειά, τότε έχεις το ελεύθερο να προχωρήσεις με αυτό που πήγες να κάνεις. Εντάξει;»

«Εντάξει» απάντησα στον εαυτό μου «...απλά περίμενε να δω πόσα λεφτά έχω κρυμμένα...» είπα και έψαξα λίγο παραπάνω στο κομοδίνο. «300; Τώρα αυτά είναι πολλά ή λίγα στην σημερινή οικονομία... ποιος νοιάζεται, είμαι έτοιμας Σπαρκ»

«Ωραία, πάμε» είπα και έσπευσα προς το καζίνο.

\chapter{}

«Πάνω που νόμιζα ότι... θα ήμουν μαζί του όλη την υπόλοιπη μου ζωή, τώρα...»

«Ηρέμησε Μπλαιρ, δεν είναι ώρα για...»

«Δεν είναι ώρα για τι πράγμα, Ιφι; Δεν έχεις την παραμικρή ιδέα τι έχει κάνει ο Χερμπ για μένα, και ο Ρομπ δεν πάει και πολύ πίσω. Το ελάχιστο που μπορώ να κάνω είναι να είμαι εδώ αυτή τη στιγμή»

«Και που είσαι εδώ και κάθεσαι τι προσφέρεις; Αφού δεν βοηθάς κανέναν, απλά κλαις και παραπονιέσαι. Έχω άδικο;»

«Όχι... δεν έχεις. Είμαι ανίκανη. Αυτό ήθελες να ακούσεις;»

«Δεν εννοούσα αυτό...»

«Και τι εννοούσες; Πλάκα πλάκα δεν σε έχω ακούσει να κλαις ούτε μία φορά μετά το δυστύχημα. Γιατί; Μήπως τελικά δεν ήσασταν τόσο κοντά; Μήπως έχεις και άλλους άντρες στη ζωή σου;»

«Ωχ, ακόμη εδώ είστε; Γιατί είστε στο πάτωμα; Τέλος πάντων, ξέχασα το κινητό μου εδώ και ήρθα να το πάρω. Που είναι τώρα... Α, να το. Παρεμπιπτόντως, θα πάω με τον Φυτ για καφέ, μήπως θα θέλατε να έρθετε;»

«Με τον Φυτ, ε... χμ...»

«Όχι ευχαριστώ Κουπ. Έχω άλλα πράγματα στο μυαλό μου τώρα. Αν θες πάρε την Ιφι και... κάτσε, γιατί ψιθύρισες στον εαυτό σου όταν είπε ο Κουπ ότι θα πάει με τον Φυτ για καφέ;»

«Εγώ... δεν...»

«Αν είναι αυτό που σκέφτομαι, είσαι ύπουλη. Πολύ ύπουλη. Μακάρι να ήταν ο Σπαρκ εδώ...»

«Δεν είναι αυτό που νομίζεις...»

«Για να χρειάζεται να το ξεκαθαρίσεις, νομίζω πως είναι ακριβώς αυτό που νομίζω»

«Κυρίες, μην τσακώνεστε...»

«Πώς να μην τσακωνόμαστε μετά από αυτό που έμαθα μόλις»

«Γιατί, τι... έμαθες;»

«Α καλά, άστο Κουπ. Δεν κατάλαβες τι παίχτηκε τώρα με την Ιφι»

«Να σου πω την αλήθεια... όχι. Δεν σας καταλαβαίνω εσάς τις γυναίκες»

«Τέλος πάντων. Μη με φουντώσεται άλλο, θέλω λίγο χρόνο να το επεξεργαστώ»

«Να επεξεργαστείς τι;»

«Είσαι σοβαρός Κουπ; Θες να μου πεις ότι δεν έχεις την παραμικρή ιδέα τι θα μπορούσα να σκέφτομαι αυτή τη στιγμή;»

«Κοίτα... έχω, αλλά πρέπει να χαλαρώσεις»

«Να χαλαρώσω; Έχεις μήπως την παραμικρή ιδέα...»

«Ναι βασικά. Έχω την παραμικρή ιδέα πως είναι να πεθαίνει η οικογένεια σου. Η αδελφή μου έφυγε πριν κάποιους μήνες από μία σπάνια αρρώστια που δεν μπορούσαν να αντιμετωπίσουν οι γιατροί. Οι γονείς μου; Πέθαναν και αυτοί πρόσφατα σε αυτοκινητιστικό δυστύχημα. Στη συνέχεια, έγινε και αυτό σήμερα και... δεν ξέρω πραγματικά πως να το πάρω. Έχω μείνει μόνος σε αυτή τη ζωή. Οι μόνοι άνθρωποι που έχω είναι οι φίλοι μου, εσείς»

«Ω, Κουπ... συγγνώμη. Δεν ήξερα»

«Δεν πειράζει Μπλαιρ. Ήθελα να καταλήξω στο ότι δεν μπορούμε να αλλάξουμε το παρελθόν. Μπορούμε μόνο να ανατρέξουμε τις αναμνήσεις που κάναμε και να κρίνουμε τον εαυτό μας για τα λάθη μας. Για αυτό, αντί να σκέφτεσαι “τι θα μπορούσα να είχα κάνει αλλιώς”, ζήσε στο παρόν, εγώ αυτό άρχισα να κάνω. Προσπαθώ να κάνω τις καλύτερες επιλογές τώρα έτσι, ώστε ο μελλοντικός μου εαυτός να μην με κρίνει, αλλά να ευχαριστιέται τις αναμνήσεις του. Προσπαθώ να περάσω παραπάνω ποιοτικό χρόνο με τους ανθρώπους που μου έχουν μείνει, καθώς δεν ξέρω πότε η ατυχία μου θα τους βγάλει από τη ζωή μου, οπότε αν βγει φύγει τελικά κάποιος από τη ζωή μου, τουλάχιστον θα έχω τις ωραίες αναμνήσεις μου μαζί του, χωρίς να χρειάζεται να μετανιώνω για τίποτα άλλο»

«Κουπ...»

«Έλα ρε Μπλαιρ, μην κλαις... γιατί με αγκαλιάζεις ξαφνικά;»

«Δεν ήξερα... ότι είσαι φιλόσοφος»

«Φιλόσοφος; Απλός φίλος σου είμαι, και καθόλου σοφός»

«Ευχαριστώ για τα λόγια σου»

«Τίποτα... να σε αγκαλιάσω και εγώ ή θα μου νευριάσεις;»

«Τι χαζομάρες λες;»

«Μία... πλάκα έκανα... Θέλετε να πάμε τώρα; Εδώ είναι πολύ μουντό το κλίμα, όχι και ότι καλύτερο για να ξεχαστεί κανείς»

«Πάμε, αλλά... Ιφι»

«Ναι;»

«Μην τολμήσεις και κάνεις αυτό που νομίζω»

«Γιατί;»

«Επειδή δεν έχει περάσει καν μία εβδομάδα»

«Γιατί, αν είχε περάσει 7 μέρες τότε θα...»

«Όχι, ούτε τότε θα ήταν αποδεκτό. Αν είχαν περάσει χρόνια, τότε θα μπορούσες να το σκεφτείς. Δεν έχει περάσει ούτε καν μία ολόκληρη μέρα από το...»

«...γιατί με κοιτάς έτσι επίμονα;»

«Τίποτα, απλά θέλω να μου πεις βρε Ιφι... τι συναισθήματα έχεις πραγματικά για τον Σπαρκ;»

«Δεν είναι προφανές;»

«Ω, καθόλου. Αν ήταν προφανές δεν θα σε ρωτούσα»

«Έχεις δίκιο...»

«Ακόμη δεν καταλαβαίνω τι γίνεται»

«Μη μιλάς Κουπ. Άσε με να την αναλάβω εγώ, γιατί είναι πολύ σοβαρό αυτό το ζήτημα»

«Γιατί είναι τόσο σοβαρο;»

«Ακούς τον εαυτό σου βρε Ιφι; Πώς γίνεται να μην σου φαίνεται καθόλου περίεργη η σκέψη που έκανες με τον Φυτ; Ούτε λιγάκι;»

«Ε, λιγάκι...»

«Τότε, πες μου όντως τι συναισθήματα έχεις για τον Σπαρκ, γιατί απ ό,τι φαίνεται τόσο καιρό βλέπουμε έναν ψεύτικο εαυτό σου»

«Εγώ να φεύγω...»

«Όχι Κουπ, δεν θα φύγεις. Βασικά, όχι, πάμε όλοι μαζί να συζητήσουμε αυτό το ζήτημα. Κουπ, ξέρεις που είναι ο Μαλχη;»

«Δεν είπε ότι θα πήγαινε στο καζίνο μόνος του;»

«Άντε πάλι. Τέλος πάντων, πάμε χωρίς τον Μαλχη»

«Τον Μαλχη... χμ...»

«Ούτε που να το σκέφτεσαι Ιφι...»

\chapter{}

«Καλησπέρα σας, είστε καινούργιος στο καζίνο;»

Δεν της απάντησα, ήμουν συγκεντρωμένος στην αναζήτηση του στόχο μου.

«Κύριε;»

Σάρωσα με το βλέμμα μου όλα τα τραπέζια που μπορούσα να δω από την ρεσεψιόν, αλλά ο Μαλχη δεν ήταν πουθενά. Λες να ήρθα νωρίς τελικά; Ή μήπως ήταν μία από τις μέρες που ο Μαλχη δεν πήγαινε στο καζίνο;

«Ψάχνετε κάποιον, κύριε;»

«Σε μένα μιλάτε;» γύρισα και μίλησα στην κοπέλα που καθόταν πίσω από τον πάγκο της ρεσεψιόν.

«Προφανώς, αφού είστε ο μόνος που στέκεστε στην είσοδο...» με ειρωνεύτηκε. «Τέλος πάντων, που θα θέλατε να καθίσετε;»

«Δεν μπορώ να εισέλθω και μετά να διαλέξω σε ποιο τραπέζι να κάτσω;»

«Ναι, δεν εννοώ αυτό. Ρωτάω αν θα θέλατε να μείνετε στον δημόσιο όροφο του καζίνο, όπου έχουμε δέκα τραπέζια με παιχνίδια για να διαλέξετε, ή μήπως θα θέλατε να σας συνοδεύσουν σε κάποιο από τα ιδιωτικά μας δωμάτια πόκερ ή μπλάκτζακ» μου είπε η ρεσεψιονίστ. «...αν και από την εμφάνιση σας, νομίζω ο δημόσιος όροφος θα σας ταιριάζει περισσότερο...» ψιθύρισε στον εαυτό της.

«Θα μείνω εδώ» της απάντησα, καθώς ο όροφος που βρισκόμασταν είχε τον περισσότερο κόσμο, καθώς και ορατότητα στην μόνη είσοδο και έξοδο του καζίνο, οπότε αν έμπαινε ή έβγαινε ο Μαλχη, θα τον μπορούσα να τον δω.

«Οπότε θα είστε στον δημόσιο όροφο. Πολύ ωραία. Απλά για να ξέρετε...»

«Μπορώ να σας ρωτήσω κάτι;» την διέκοψα.

«Βεβαίως»

«Μήπως ξέρετε κάποιον... Μαλχη»

«Δεν μπορώ να απαντήσω σε τέτοιες προσωπικές ερωτήσεις»

«Δεν πειράζει» της είπα και προσπάθησα να διαβάσω τις σκέψεις της.

\textasciitilde ...ποιον ψάχνει αυτός τώρα; Μήπως είναι ληστής και αυτός ο Μαλχη έχει κερδίσει κάποιο μεγάλο χρηματικό ποσό τελευταία; Θα το αναφέρω στην ασφάλεια μόλις φύγει από την είσοδο...\textasciitilde «Τι σας έλεγα λοιπόν. Α, ναι. Απλά για να ξέρετε, τα περισσότερα τραπέζια έχουν ελάχιστο ποσό συμμετοχής όταν πρωτοκάθεστε που αντιστοιχεί σε μία πράσινη μάρκα»

«Πράσινη μάρκα;»

«Αχ, πόσο αφελής, δεν σας ρώτησα με τι χρηματικό ποσό είστε πρόθυμος να ξεκινήσετε, ώστε να σας δώσουμε τις αντίστοιχες μάρκες.»

«Πάνω μου έχω 300...» είπα και έβγαλα τα τρία κατοστάρικα που είχα στην τσέπη μου. «...οπότε δώστε μου τις αντίστοιχες μάρκες»

«Θα τα στοιχιματίσεται όλα;» ρώτησε με σοβαρό ύφος η κοπέλα.

«Ναι... γιατί όχι»

«Καλά...» είπε και άρχισε να ετοιμάζει τις μάρκες μου. «Απλά είμαι υποχρεωμένη να σας ενημερώσω πως το καζίνο μακροχρόνια πάντα κερδίζει, οπότε πρέπει να παίζεται υπεύθυνα, καθώς το παιχνίδι μπορεί να γίνει εθιστικό»

Δεν της απάντησα.

«Ορίστε οι μάρκες σας. Είκοσι κόκκινες και οκτώ πράσινες. Οι κόκκινες αντιστοιχούν σε 5 και οι πράσινες σε 25 χρηματικές μονάδες»

«Άρα, για να κάτσω σε τραπέζι... πρέπει να πληρώσω 25;»

«Ακριβώς» μου απάντησε χαμογελώντας η κοπέλα στην ρεσεψιόν. «Καλή διασκέδαση»

«Ευχαριστώ...» είπα και προχώρησα προς το πρώτο τραπέζι που βρήκα, στο οποίο θα είχα ακόμη ορατότητα όλου του καζίνο. Στο τραπέζι αυτό έπαιζαν μπλάκτζακ δύο ακόμη άντρες, οι οποίοι είχαν καθίσει στις δύο απέναντι καρέκλες του τραπεζιού και δεν μίλαγαν μεταξύ τους. Κάθησα στην καρέκλα δίπλα στον δεξιό παίκτη, από την οποία μπορούσα να παρατηρώ και την είσοδο του καζίνο, και άφησα την τέταρτη καρέκλα άδεια.

«Καλώς Ήρθατε κύριε. Περιμένετε λίγο να ολοκληρώσουμε αυτό το χέρι και θα μπορείτε να μπείτε και εσείς στο παιχνίδι» μου είπε ο κρουπιέρης και μοίρασε ένα φύλλο στον άντρα δεξιά μου, ο οποίος χτύπησε το τραπέζι με τα δάκτυλα του, το οποίο απ ό,τι θυμάμαι είναι σινιάλο για αυτή ακριβώς την ενέργεια. 

Δεν μπορούσε όμως να με νοιάζει λιγότερο τι παιχνίδι θα έπαιζα και αν θα κέρδιζα ή όχι. Το μόνο που είχα στο μυαλό μου εκείνη τη στιγμή ήταν να βρω τον Μαλχη. Είχα σαρώσει όμως τόσες φορές τα τραπέζια στον όροφο που βρισκόμουν, αλλά δεν τον είχα διακρίνει. Λες όντως να μην ήρθε σήμερα; Ή μήπως βρίσκεται σε κάποιο ιδιωτικό δωμάτιο; Ποιος ξέρει, πρέπει πάντως να είμαι συνεχώς σε εγρήγορση, καθώς μπορεί να μπει ή να βγει οποιαδήποτε στιγμή από τις πόρτες του καζίνο. Είχα την εντύπωση όμως ότι το καζίνο ήταν ένας τελείως κλειστός χώρος όπου δεν έμπαινε καθόλου φως από τον έξω κόσμο και δεν υπήρχαν ρολόγια μέσα στα δωμάτια, για να ξεχνιούνται οι παίκτες και να χάνουν την αίσθηση του χρόνου. Μπορεί να ήταν θεωρία συνωμοσίας, ποιος ξέρει.

«Εικοσιμία. Χάσατε και οι δύο σας συγγνώμη»

«Κοίτα την τύχη του. Δεν πειράζει, θα πάρω πίσω τα λεφτά μου στον επόμενο γύρω» είπε ο ξέμπαρκος στο τραπέζι μας άντρας. «100»

«Και γω το ίδιο» είπε ο άντρας δίπλα μου «50»

Για λίγη ώρα δεν μίλησε κανείς στο τραπέζι. Εγώ σάρωνα επίμονα τα υπόλοιπα τραπέζια και την είσοδο, για να μην χάσω τον Μαλχη αν εμφανιστεί.

«Θα ποντάρεις;» μου απευθύνθηκε ο κρουπιέρης.

«Αχ ναι...» του απάντησα «...δεν σας πλήρωσα όμως. Ορίστε τα 25...»

«Δεν χρειάζεται. Είναι τραπέζι μπλάκτζακ. Εδώ υπάρχει ελάχιστη κατανάλωση των τεσσάρων πράσινων μάρκων, ή είκοσι κόκκινων, που αντιστοιχούν στο ίδιο ποσό, 100 χρηματικές μονάδες»

«Α... τότε ποντάρω τέσσερις πράσινες»

«Δεν εννοούσα σε ένα χέρι»

«Δεν πειράζει, μοίρασε τα φύλα»

«Πολύ καλά...»

Καθώς μοίραζε τα φύλα ο κρουπιέρης, παρατήρησα λίγο πιο προσεκτικά το πρόσωπο του. Μου ήταν γνώριμο, αλλά δεν μπορούσα να θυμηθώ από που.

«Μπορώ να σε ρωτήσω κάτι;» του απευθύνθηκα και προσπάθησα να διαβάσω τις σκέψεις του.

\textasciitilde ...τι θέλει τώρα αυτός;...\textasciitilde «Παρακαλώ, ρωτήστε ότι θέλετε»

«Πώς σε λένε;»

«Ρομπ» \textasciitilde ...γιατί θέλει να μάθει πώς με λένε;...\textasciitilde 

«Ρομπ;» τον κοίταξα λίγο καλύτερα. «Μήπως ξέρετε κάποιον... Χερμπ;»

«Χερμπ ε; Όχι δεν νομίζω» είπε και έδωσε ένα ακόμη φύλλο στον άντρα δεξιά μου.

«Άρα απλά ήταν το μυαλό μου...»

«Τι εννοείται;»

«Μήπως... ήσουν πράκτορας στην Ειδική Υπηρεσία παλαιότερα;»

«Πράκτορας; Μακάρι. Η μαγεία μου δεν είναι για τέτοια πράγματα»

«Τι μαγεία έχεις;»

«Μαγεία φωτιάς, αλλά είμαι κλάσης Β, οπότε είμαι λίγο κατώτερος»

«Ώστε είναι όντως απλά μία σύμπτωση...» ψιθύρισα στον εαυτό μου.

«Θα παίξεις;» με ρώτησε ο άντρας που καθόταν δίπλα μου.

«Α, ναι... 13, ε; Μη μου δώσεις φύλλο»

«Τι κάνεις;» με ρώτησε πάλι ο άντρας. «Ξέρεις να παίζεις;»

«Να σου πω την αλήθεια... όχι»

«Τι; Πώς... καλά δεν πειράζει»

«Εγώ δεν θέλω κάρτα» είπε ο τρίτος άντρας και ο κρουπιέρης γύρισε τις κάρτες του και μας έδειξε ότι είχε δύο άσους.

«Μας σκότωσες μουσάτε. Τώρα αν βγάλει ένα εννιάρι χάσαμε όλοι» είπε εκνευρισμένος ο άντρας δίπλα μου

Ο κρουπιέρης γύρισε την επόμενη κάρτα. Ήταν ρήγας και ο κρουπιέρης έχασε, καθώς είχε συμπληρώσει 22.

«Τυχερός ήσουν» μου είπε πάλι ο άντρας δίπλα μου. «Μην το ξανακάνεις αυτό όμως...»

«Γιατί; Αφού αν το έκανα, θα έχανα, και πιθανότατα θα χάνατε και εσείς, αν η επόμενη κάρτα ήταν καλή για τον κρουπιέρη» του απάντησα.

«Δεν θέλω δικαιολογίες τώρα, απλά παίξε» μου είπε και γύρισε το σώμα του ευθεία.

Είχα αποσπαστεί από τον στόχο μου. Πόνταρα 25 από τα 100 που είχα κερδίσει και έστρεψα κατευθείαν το βλέμμα μου στα υπόλοιπα τραπέζια στο καζίνο. Γιατί δεν είναι ο Μαλχη εδώ; Μήπως απλά είμαι άτυχος;\\\\

«Εικοσιμία, χάσατε και οι τρεις σας» είπε ο κρουπιέρης και μάζεψε τις μάρκες που είχαμε μπροστά μας.

«Είσαι άτυχος τελικά μουσάτε. Έχεις χάσει 14 γύρους στη σειρά. Πάλι καλά που στοιχιματίζεις 5 και 10 την φορά, αλλιώς θα είχες φύγει από το καζίνο τώρα»

Δεν του απάντησα, ήμουν συγκεντρωμένος στον στόχο μου.

«Δεν έχεις να πεις κάτι, ε; Τέλος πάντων, εγώ φεύγω παιδιά. Έχει πάει 8:30 και δουλεύω αύριο, οπότε θα σας καληνυχτίσω»

«Καλό βράδυ» είπε ο κρουπιέρης και άρχισε να μας μοιράζει φύλλα.

«Θα φύγω και γω μετά από αυτό το χέρι» είπε ο άντρας που καθόταν στην μακρινή καρέκλα.

«Εντάξει. Εσείς;» μου απευθύνθηκε ο κρουπιέρης.

«Ίσως... αλλάξω τραπέζι. Δεν θέλω να είμαι μόνος μου» του απάντησα.

«Σας καταλαβαίνω, θέλετε να έχετε παρέα. Ωραία λοιπόν. Αφού έχετε 20 και 19, υποθέτω δεν θέλετε άλλες κάρτες...»

«Ναι» του απαντήσαμε ταυτόχρονα.

«Ωραία» είπε και γύρισε έναν βαλέ δίπλα στον άσο που βλέπαμε ήδη. «Συγγνώμη, εικοσιμία» είπε και μας πήρε τις τελευταίες μάρκες.

«Καλή συνέχεια κύριοι» μας αποχαιρέτησε χωρίς πολλα πολλά ο τρίτος άντρας.

«Καλό βράδυ» του είπε ο κρουπιέρης «...και σε εσάς, καλή τύχη στο επόμενο σας παιχνίδι»

«Σε ευχαριστώ...» είπα και σηκώθηκα, παίρνοντας τις δώδεκα πράσινες μάρκες που μου είχαν μείνει. Ξόδεψα μισή ώρα χωρίς να καταφέρω τίποτα. Ούτε τον Μαλχη βρήκα, ούτε λεφτά κέρδισα στην τελική. Νιώθω λες και είμαι κακός κατάσκοπος. Πώς το έκαναν η Ιφι και ο Ρομπ; Μπορούσαν μέσα σε λίγες ώρες να εξάγουν τόσες πολλές πληροφορίες για τα παρακολουθούμενα άτομα. Εγώ βέβαια δεν παρακολουθώ κάποιον. Ίσως μία καλύτερη λέξη θα ήταν το “αναμένω” κάποιον. 

Κάθισα σε ένα τραπέζι πόκερ, στο οποίο κάθονταν άλλοι τέσσερις παίκτες, τρεις άντρες και μία καλόγρια, ντυμένη με ολόμαυρο ράσο. Τι κάνει μία καλόγρια στο καζίνο; Και γιατί παίζει πόκερ; Πάλι αποσπώμαι. Κάθισα στο τραπέζι και έδωσα στον κρουπιέρη μία πράσινη μάρκα.

«Σας ευχαριστώ» μου είπε. «Ας ξεκινήσουμε το επόμενο χέρι» ανακοίνωσε και άρχισε να μοιράζει από δύο κάρτες στον καθένα μας.

«Αυτός έχασε πάνω από 10 συνεχόμενους γύρους στο μπλάκτζακ πριν λίγο...» κρυφάκουσα έναν άντρα να ψιθυρίζει στον διπλανό του.

«Τι; Πώς γίνεται να είσαι τόσο άτυχος;»

«Δεν ξέρω. Λες να μπορέσουμε να τον μπλοφάρουμε και να του κλέψουμε καμιά πράσινη μάρκα;»

«Φαίνεται λίγο αφηρημένος, μπορεί να δουλέψει»

«Καλύτερα μπλόφαρε εσύ που έχεις και περισσότερα λεφτά από μένα»

«Εντάξει, δεν έχω και καλό χέρι, οπότε ταιριάζει η περίσταση»

«Το ξέρετε ότι ακούγεστε;» τους είπε η καλόγρια που καθόταν δίπλα τους. «Κάντε ότι θέλετε, απλά να ξέρετε ότι σας ακούει, οπότε δεν θα δουλέψει αν μπλοφάρεις»

«Και που ξέρεις ρε καλόγρια ότι δεν μπλοφάρω δεν τα λόγια μου;»

«Δεν το ξέρω»

«Ε τότε φόρα καλύτερα το μαντήλι σου και μη μιλάς»

Η καλόγρια κράτησε την ψυχραιμία της και δεν μίλησε ξανά στον άντρα. Πόνταρε ήρεμα τρεις πράσινες μάρκες.

«Ε τότε εγώ ποντάρω 9 πράσινες»

«Εγώ πάω πάσο» είπε ο δεύτερος άντρας, που με είχε κρίνει που έχασα 14 γύρους στο μπλάκτζακ, με ένα χαμόγελο στα χείλη του και πέταξε τις κάρτες του.

«Τα βάζω όλα μέσα» είπα και έσπρωξα τις 11 μάρκες που μου είχαν μείνει. Άρχισα να κοιτάζω πάλι τριγύρω για να βρω τον Μαλχη, αλλά τίποτα. Έκανα λάθος που άλλαξα τραπέζι, γιατί τώρα βρίσκομαι σχεδόν στο κέντρο του ορόφου και χρειάζεται να γυρνάω όλο μου το σώμα για να μπορέσω να παρατηρώ τα τραπέζια πίσω μου. Τουλάχιστον μπορούσα να παρατηρώ ακόμη την είσοδο από την θέση μου.

«Εγώ πάω πάσο» είπε ο άντρας δίπλα μου και πέταξε κι αυτός τις κάρτες του στον κρουπιέρη.

«Σε ταιριάζω» είπε η καλόγρια και κοίταξε τον άντρα δίπλα της, του οποίου ήταν η σειρά να παίξει. Ο άντρας αυτός είχε ιδρώσει και κοίταζε τις κάρτες του επίμονα.

\textasciitilde ...δεν μπορεί να μπλοφάρει και αυτός. Δεν έχει δει καν τις κάρτες του και τα βάζει όλα μέσα; Κι ας μη μιλήσω για το γεγονός ότι ο κρουπιέρης δεν έχει εμφανήσει ούτε μία κάρτα κάτω. Δεν μπορώ όμως να τα ποντάρω όλα όταν έχω μόνο 2 καρό και 3 μπαστούνι. Άστο καλύτερα...\textasciitilde «πάω και γω πάσο» είπε και πέταξε της κάρτες του ο τελευταίος άντρας.

«Άρα όντως δεν είχες τίποτα...» του είπε η καλόγρια.

«Παίξε τώρα και μη μιλάς» είπε εκνευρισμένος ο άντρας.

«Καλά...» απογοητεύτηκε η καλόγρια γύρισε τις κάρτες της. «Βαλές και ντάμα καρό. Για να δούμε τι έχεις εσύ νέε;» είπε και γύρισα της κάρτες μου χωρίς να τις κοιτάξω καν. Ήμουν αφοσιωμένος στον στόχο μου.

«Άσους;» φώναξε ο άντρας δίπλα της. «Η τύχη του πρωτάρη...» με ειρωνεύτηκε. Δεν του απάντησα.

Ο κρουπιέρης γύρισε τις τρεις πρώτες κάρτες.

«Άσος, 10 και 9 καρό, έχω φλας» χάρηκε η καλόγρια. «Δύσκολα θα κερδίσεις μικρέ»

Δεν της απάντησα.

«Επόμενη κάρτα...» ανακοίνωσε ο κρουπιέρης και έβγαλε ένα εξάρι.

«Οι πιθανότητες σου λιγοστεύουν μικρέ» γέλασε η καλόγρια, η οποία είχε σχεδόν κερδίσει.

Δεν της απάντησα. Είχα γυρίσει και κοίταζα τα πίσω τραπέζια προσεκτικά, μήπως μου είχε ξεφύγει κάποιος. Ο κρουπιέρης παράλληλα έδειξε στους υπόλοιπους την τελευταία κάρτα.

«Άσος. Τι ατυχία είναι αυτή, μη βρίσω το σπανάκι μου» άκουσα την καλόγρια να εκνευρίζεται.

Έστρεψα αμέσως το βλέμμα μου πάνω της.

«Γιατί με κοιτάς έτσι; Νίκησες, πάρε τις μάρκες σου και άσε με» είπε η καλόγρια και άρχισε να παίζει με τις υπόλοιπες 6 μάρκες που της είχαν μείνει.

«Μπορώ να σας κάνω μία ερώτηση;» είπα καθώς ο κρουπιέρης μας μοίραζε το επόμενο χέρι.

«Παρακαλώ» μου απάντησε η καλόγρια.

«Τι εννοείτε όταν λέτε “μη βρίσω το σπανάκι μου”;»

«Α, αυτό; Είναι μία έκφραση που λέω για να μην μου ξεφεύγουν λάθος λέξεις όταν είμαι στο καζίνο»

«Πάω πάσο κρουπιέρη» είπα και πέταξα τις κάρτες μου όταν ήταν σειρά μου να παίξω. «Και... γιατί είσαι στο καζίνο;»

«Έχω μία οικογένεια να φροντίσω»

«Οικογένεια;»

«Ναι. Υιοθετημένη μεν, αλλά δεν παύει να είναι οικογένεια»

«Έχετε και σύζυγο που σας βοηθάει;»

«Όχι, μόνη μου μεγαλώνω τα παιδιά»

«Μισό... δεν είναι λίγο επικίνδυνο να έχετε οικογένεια μόνη σας και να είστε στο καζίνο;»

«Συνήθως κερδίζω, απλά σήμερα ήρθες εσύ και μου χάλασες το σερί...» είπε εκνευρισμένη η καλόγρια.

«Α... συγγνώμη» της απάντησα και της πέταξα δύο πράσινες μάρκες. «Ελπίζω αυτό να βοηθήσει»

«Ωχ, σε ευχαριστώ πολύ παιδί μου. Μακάρι ο Θεός να είναι μαζί σου σε ό,τι περνάς»

«Σας ευχαριστώ» της απάντησα.

Η καλόγρια αυτή σίγουρα συνδέεται κάπως με τον Μαλχη. Εξάλλου, ο Μαλχη είναι το μόνο άτομο που έχω να ακούσει να λέει αυτή την έκφραση με το σπανάκι. Το σχέδιο αλλάζει: μόλις φύγει η καλόγρια από το καζίνο θα την ακολουθήσω για να δω που θα πάει. Ωχ... καλόγρια, μοναχή... Δεν είναι συνώνυμες αυτές οι  λέξεις; Ή νομίζω η μία είναι λίγο πιο επίσημη από την άλλη. Πώς δεν το είχα σκεφτεί τόση ώρα; Είμαι όντως άχρηστος. Κατα πάσα πιθανότητα είναι η γυναίκα με τις πολλές μαγείες. Αλλά δεν φαίνεται και τόσο μεγάλη, το πολύ να είναι πενήντα χρονών. Κάτι δεν πάει καλά... Μπορεί να σχετίζεται όμως με την γυναίκα με τις πολλές μαγείες, η οποία σχετίζεται με τον Μαλχη. Σε αυτή την περίπτωση, αν έχει όντως κάποια σχέση με τον Μαλχη, δεν πρέπει να είναι και πολύ κοντινή. Είναι μεγάλο ρίσκο να την ακολουθήσω, αλλά είναι το μόνο πράγμα που μπορώ να κάνω, αφού ο στόχος μου δεν εμφανίστηκε πότε...

\chapter{}

«Λοιπόν παιδιά, εγώ θα σας καληνυχτίσω για σήμερα» είπε η καλόγρια. «Παρόλο που δεν κέρδισα πολλά χρήματα σήμερα, είχαμε μερικά αρκετά διασκεδαστικά χέρια»

«Καλά τα λες» της είπε ο άντρας που καθόταν δίπλα της. «Καλό βράδυ καλόγρια»

«Καλή συνέχεια» χαμογέλασε και έφυγε.

Δεν έπρεπε να την χάσω από το οπτικό μου πεδίο, ακόμη και αν αυτό σήμαινε πως θα έπρεπε να σηκωθώ στη μέση του επόμενου γύρου. Ευτυχώς για μένα, η καλόγρια στάθηκε για λίγο μπροστά στην έξοδο και έπιασε συζήτηση με την κοπέλα στην ρεσεψιόν.

«Θα παίξεις μουσάτε;»

«Α ναι, συγγνώμη. Όλα μέσα»

«Πάλι; Είναι η τρίτη φορά σε πέντε γύρους που στοιχηματίζεις όλα σου τα λεφτά» εκνευρίστηκε ο άντρας που καθόταν δίπλα στην καλόγρια.

«Εγώ πάω πάσο» είπε ο άντρας δίπλα μου.

«Εγώ όχι, τα ποντάρω όλα και εγώ, παρόλο που έχω λιγότερα»

«Και γω το ίδιο» είπε και ο τελευταίος άντρας.

«Για να δούμε τα χέρια σας» μας είπε ο κρουπιέρης. Οι δύο άντρες είχαν ένα ζευγάρι άσους και ένα ζευγάρι βαλέδες αντίστοιχα. Εγώ γύρισα τις κάρτες μου και τους έδειξα πως είχα 2 μπαστούνι και 3 καρό.

«Ώστε μπλόφαρες, ε;»

Δεν του απάντησα. Ήμουν συγκεντρωμένος στον στόχο μου, η οποία εκείνη τη στιγμή αποχαιρετούσε τους ανθρώπους στην είσοδο του καζίνο.

«Τι είναι αυτό;» άκουσα τον έναν από τους δύο άντρες να λέει όταν ο κρουπιέρης γύρισε τις πρώτες τρεις κάρτες. «Κοίτα να δεις που θα κερδίσει ο μουσάτος πάλι»

Δεν γύρισα να κοιτάξω καν τις κάρτες, γιατί θα έχανα τον μόνο άνθρωπο που θα μπορούσε να με φέρει πιο κοντά στον Μαλχη.

«Κοίτα δω, ο μουσάτος έχει κέντα. Δεν μπορούμε καν  να νικήσουμε τώρα, αφού και τα σύμβολα που τύχαμε δεν ταιριάζουν. Μπράβο φίλε, μας...»

«Συγγνώμη, αλλά πρέπει να φύγω» έπιασα όσες μάρκες μπορούσα να κρατήσω στην αγκαλιά μου και σηκώθηκα γρήγορα από το τραπέζι, καθώς είδα την καλόγρια να βγαίνει από το καζίνο.

«Αυτό που κάνεις είναι κακή ηθική. Δεν είναι σωστό να κερδίζεις ένα μεγάλο χρηματικό ποσό και ναι...»

Πριν τελειώσει την πρόταση του, είχα φύγει ήδη μακριά από το τραπέζι οδεύοντας προς την ρεσεψιόν, όπου θα έκανα την ανταλλαγή των μαρκών.

«Κερδίσατε βλέπω»

«Ναι ναι, αν μπορείτε κάντε την αντιστοίχηση όσο πιο γρήγορα γίνεται» είπα και κοίταξα έξω από την ανοιχτή πόρτα του καζίνο και είδα την καλόγρια να στρίβει δεξιά δυο στενά απέναντι από το σημείο που βρισκόμουν. Πήγαινε προς την εγκαταλελειμμένη εκκλησία. Μπίνγκο.

«Περιμένετε κάποιον;» με ρώτησε η κοπέλα.

«Ναι... περιμένω έναν φίλο μου να έρθει να με πάρει» είπα ψέματα για να μη φανώ ύποπτος.

«Αχ ευτυχώς, επειδή σας βλέπω κοιτάτε τόσην ώρα έξω και άρχισα να αγχώνομαι πως είστε κάποιος ληστής. Βλέπετε...»

«Δεν έχω χρόνο για διάλογο. Σας παρακαλώ, 24 μάρκες είναι όλες κι όλες, δώστε μου το αντίστοιχο ποσό για να φύγω»

«Εντάξει...» είπε η κοπέλα απογοητευμένη και μου έδωσε έξι κατοστάρικα. «Καλό σας βράδυ»

Δεν την χαιρέτησα πίσω. Έβαλα τα λεφτά στην τσέπη μου και έτρεξα να προλάβω την καλόγρια. Στρίβοντας στο στενό που έστριψε κι αυτή, την είδα να μπαίνει σε ένα μίνι μάρκετ, το οποίο βρισκόταν περίπου τέσσερα στενά μακριά από μένα. Εκείνη την στιγμή ανακουφίστηκα, καθώς δεν θα χρειαζόταν να τρέξω αλλό. Δυστυχώς για μένα, η τύχη μου έγινε ατυχία, καθώς όταν άρχισα να περπατάω πιο σιγά προς το μίνι μάρκετ, ένιωσα μερικές σταγόνες βροχής να χτυπούν στα μαλλιά μου.

«Πάλι;» αναρωτήθηκα. «Ελπίζω αυτή τη φορά να μη με χτυπήσει κεραυνός πάλι» έκανα ένα αστείο που μόνο εγώ μπορούσα να καταλάβω. «Κεραυνός, ε; Τι και αν... με χτύπησε όντως κεραυνός; Δεν μπορώ να θυμηθώ τι έγινε αφότου με αγνόησαν η Μπλαιρ, ο Βασι και ο Κουπ. Πώς μπόρεσα και περπάτησα με το που σηκώθηκα; Δεν βγάζει νόημα. Εκτός και αν... πλάκα θα είχε. Πλάκα θα είχε να έχω αποκτήσει μία αναβάθμιση στις δυνάμεις μου όπως κάποιος πρωταγωνιστής σε μία ταινία πριν πολεμήσει τον υπέρτατο κακό, δεν νομίζω όμως»

Όταν έφτασα έξω από το μίνι μάρκετ, κρύφτηκα μέσα σε ένα διπλανό στενό για να περιμένω τον καλόγρια. Από περιέργεια, τέντωσα το χέρι μου προς έναν κάδο σκουπιδιών και προσπάθησα, όπως έκανα μικρός, να βγάλω έναν κεραυνό από το χέρι μου, και όχι μόνο έβγαλα έναν κεραυνό, αλλά έναν μεγάλο κεραυνό, τόσο μεγάλο που με έστειλε μερικά μέτρα πίσω και ακούστηκε σε όλο το τετράγωνο. Εκείνη την στιγμή είχα μεικτά συναισθήματα, καθώς μπόρεσα να κάνω κάτι που όλη μου τη ζωή δεν μπορούσα, αλλά η στιγμή δεν ήταν κατάλληλη, καθώς θα μπορούσε να με είχε ακούσε η καλόγρια και να της κινήσω υποψίες.

Ευτυχώς για μένα, όταν η καλόγρια βγήκε από το μίνι μάρκετ, δεν είχε κανέναν προβληματισμό στο πρόσωπο της. Μάλλον η βροχή λειτούργησε τελικά σαν κάληψη μου, ίσως δεν είμαι τόσο άτυχος τελικά. Η καλόγρια κρατούσε πέντε σακούλες γεμάτες ψώνια, γεγονός που με παραξένεψε λίγο, καθώς ήταν βράδυ. Αποφάσισε να κάνει τα ψώνια της εβδομάδας από ένα μίνι μάρκετ ένα τυχαίο βράδυ; Και πώς μπορεί αυτό το λιγνό αθώο σώμα να κουβαλάει πέντε γεμάτες σακούλες, χωρίς να δείχνει κάποιο σημάδι κούρασης;

Την άφησα να προχωρήσει ένα στενό παραπέρα πριν συνεχίσω να την ακολουθώ. Η βροχή είχε δυναμώσει και η υπομονή μου τελείωνε. Η αδεξιότητα μου όμως θα εμφανιζόταν μία ακόμη φορά, καθώς όταν ξεκίνησα να περπατάω πίσω από την γυναίκα, μου ήρθε μία ισχυρή επιθυμία να φτερνιστώ. Όσο προσπαθούσα να κρατηθώ, τόσο μεγαλύτερη ήταν η ενόχληση που ένιωθα στην μύτη μου. Όταν δεν μπορούσα να κρατηθώ άλλο, φτερνίστηκα όσο πιο προσεκτικά μπορούσα, όμως ακόμη ένας κεραυνός μου ξέφυγε, από το στόμα μου  αυτή τη φορά, και με πέταξε μερικά μέτρα πίσω. Σηκώθηκα γρήγορα και κοίταξα προς το μέρος της καλόγριας και ανακουφίστηκα όταν την είδα να προχωράει ανενόχλητη μέσα στη βροχή, χωρίς να έχει γυρίσει το βλέμμα της προς την κατεύθυνση μου.

Συνέχισα να περπατάω από πίσω της, αλλά όταν έκανα το πρώτο βήμα, η ταχύτητα με την οποία κουνήθηκα μου φάνηκε περίεργη. Δεύτερο βήμα, το ίδιο. Τρίτο βήμα, περπατούσα αρκετές φορές πιο γρήγορα. Γιατί; Δεν κάνω κάτι διαφορετικό. Εκείνη τη στιγμή κοίταξα τα πόδια μου και ξαφνιάστηκα όταν είδα δύο... “παπούτσια” από κεραυνούς τριγύρω από τα κανονικά μου παπούτσια. Προσπάθησα να τρέξω προς την αντίθετη κατεύθυνση από την οποία προχωρούσαμε εγώ και η καλόγρια τόση ώρα και, πράγματι, η ταχύτητα μου είχε διπλασιαστεί. Έτρεξα έναν κύκλο του τετραγώνου μέσα σε δύο ή τρία δευτερόλεπτα. Λες μήπως να με χτύπησε όντως κεραυνός εκείνη τη στιγμή; Λες η μαγεία μου να έχει όντως αναβαθμιστεί τη στιγμή που τη χρειάζομαι περισσότερο σε όλη τη ζωή μου; Δύο αναβαθμίσεις όμως είναι υπερρεαλιστικό. Ούτε στα καλύτερα μου όνειρα δεν είχα τέτοια τύχη. Γιατί όμως εμφανίζεται τώρα αυτή η τύχη; Που ήταν τόσο καιρό;

Είδα από δύο στενά μακριά την καλόγρια να στρίβει αριστερά και να περπατάει στην πλατεία της εκκλησίας, πηγαίνοντας προς την πόρτα. Δεν αγχώθηκα όμως, με την καινούργια μου αναβάθμιση στην ταχύτητα μπορούσα να την προλάβω μέσα σε ένα δευτερόλεπτο. Έμεινα για λίγο πίσω και προσπάθησα να βγάλω τα νέα μου “παπούτσια” από κεραυνούς, σκύβοντας και ακουμπώντας με το δεξί μου χέρι το αντίστοιχο παπούτσι. Όταν ήρθε σε επαφή όμως το χέρι μου με τους κεραυνούς του παπουτσιού αυτού, ένιωσα ένα μικρό σε ένταση ηλεκτροσοκ να με διαπερνά.

«Δεν βγαίνουν έτσι...» ψιθύρισα στον εαυτό μου και σηκώθηκα από την λυγισμένη στάση που είχα πάρει. «Προφανώς, αφού είναι μαγεία...» σκέφτηκα πως θα έπρεπε κάπως να τα ελέγξω με το μυαλό μου ή με τα χέρια μου. Τέντωσα το χέρι μου προς τις πατούσες μου και προσπάθησα να συγκεντρωθώ. Έκλεισα τα μάτια μου και σκέφτηκα τους κεραυνούς να φεύγουν από γύρω από τα πόδια μου. Όταν ξανακοίταξα κάτω, έμεινα έκπληκτος από το μικρό μου αυτό επίτευγμα. Μπόρεσα να ελέγξω την μαγεία μου μετά από τόσα χρόνια. Τα “παπούτσια” μου είχαν εξαφανιστεί. Πώς; Δεν ξέρω ακριβώς ακόμη, αλλά εκείνη τη στιγμή προσπάθησα να τα ξαναφορέσω, καθώς για να έχω πλήρη έλεγχο της μαγείας μου θα έπρεπε να μπορώ να κάνω με επιτυχία και τις δύο αυτές ενέργειες. Τέντωσα το χέρι μου ξανά προς τα κάτω και φαντάστηκα τα “παπούτσια” αυτά να εμφανίζονται γύρω από τα πόδια μου. Δεν μπορούσα να πιστέψω στα μάτια μου. Τα “παπούτσια” είχαν πάρει μορφή για ακόμη μία φορά γύρω από τα κανονικά μου παπούτσια.

«Πάλι αποσπώμαι...» είπα και έστρεψα το βλέμμα μου στην πλατεία της εκκλησίας. «Ας πηγαίνουμε τώρα... Ναι πάμε... ωραία ξεκινάμε» έκανα μία μικρή συζήτηση με τον εαυτό μου και πήρα στάση εκκίνησης δρομέα στίβου. Προτού καταλάβω ότι ξεκίνησα να τρέχω, είχα ήδη φτάσει έξω από την πόρτα της εκκλησίας.

«Πολύ ωραία η ταχύτητα, αλλά... φοβάμαι μη με χτυπήσει πάλι κεραυνός...» είπα και έβγαλα τα “παπούτσια” αυτά. «Θα τα ξαναφορέσω μέσα στην εκκλησία αν είναι» είπα στον εαυτό μου και άνοιξα αργά την πόρτα, για να μην ακουστώ. Εκεί είδα την καλόγρια, στο κέντρο του χολ της εκκλησίας να κάθεται στα γόνατα της και να προσεύχεται, με τα ψώνια της ακουμπισμένα λίγα μέτρα πίσω της. Γύρω της, όπως θυμάμαι πως είχε περιγράψει η Ιφι, υπήρχαν πολλά διαφορετικά φώτα, τα οποία σχημάτιζαν έναν κύκλο στο πάτωμα της εκκλησίας. Η ένταση τους δεν ήταν και τόσο ισχυρή ώστε να έφταναν μέχρι τον ουρανό, όπως τα φώτα των πέντε κουκουλοφόρων στην αίρεση σε εκείνη την υπόθεση, αλλά ήταν σίγουρα ίδια ένταση με εκείνα της φωτογραφίας που μου είχε δείξει η Ιφι πριν δέκα χρόνια, την βραδιά πριν την μοιραία αποστολή. Προσπάθησα να ακούσω λίγο καλύτερα την προσευχή της καλόγριας... μοναχής... θα την λέω καλόγρια στις σκέψεις μου, αφού αυτό συνήθισα.

«...Ω πατέρα. Ω πατέρα που ζεις στους ουρανούς, Γιατί πρέπει οι αδύναμοι να είναι αυτοί που περνούν τα περισσότερα βάσανα. Τα τελευταία χρόνια έχω δει το κακό που υπάρχει έξω από το φως της χάρης Σου. Δώσε μου ένα σημάδι, πες μου πως ακούς τη φωνή μου. Τα παιδιά αυτά δεν μπορούν να ζήσουν έτσι για πολύ ακόμη. Δεν είναι σαν τα προηγούμενα μου παιδιά, όμως αυτό δεν σημαίνει πώς πρέπει και να ζουν όλη μέρα στο σκοτάδι, για αυτό, σε παρακαλώ, δώσε τους δύναμη, δώσε τους αγάπη, δώσε τους κίνητρο, όπως κάνεις τα τελευταία 22 χρόνια. Σε ευχαριστώ που με άκουσες και σήμερα και θα είμαι ευγνώμων πάντοτε για όσα μου έχεις δώσει σε αυτή τη ζωή...»

Τελειώνοντας την προσευχή της, η καλόγρια πήρε τις πέντε σακούλες της και άρχισε να κατεβαίνει τις σκάλες για το μοιραίο υπόγειο. Την άφησα για λίγο να πάρει ένα μικρό προβάδισμα, καθώς μπορούσα να την φτάσω μέσα σε κλάσματα δευτερολέπτου με την ταχύτητα μου. Έμεινα για λίγο έξω από την εκκλησία σκεπτόμενος.

«Η προσευχή της περιείχε στοιχεία από αυτά που μου είχε αναφέρει η Ιφι εκείνο το βράδυ, όμως ήταν υπερβολικά μικρότερη από τρεις ώρες. Λες να βιαζόταν; Τι λόγο έχει να βιάζεται; Μπορεί να την περιμένουν τα “παιδιά” της. Πόσα παιδιά έχει όμως για να πάρει πέντε γεμάτες σακούλες ψώνια το βράδυ; Μήπως... είμαι όντως άχρηστος. Τα παιδιά στα οποία αναφερόμαστε και εγώ και αυτή είναι τα παιδιά τα οποία χάθηκαν στις φωτιές πριν 22 χρόνια. Πως δεν το σκέφτηκα νωρίτερα, βγάζει απόλυτο νόημα. Α, και τα άλλα παιδιά είναι τα παιδιά που χάθηκαν στις πρώτες φωτιές νοσοκομείων, πριν 30 χρόνια. Όλα συνδέονται τώρα. Η καλόγρια που ακολουθούσα τόση ώρα είναι ο κύριος ύποπτος της υπόθεσης που είχε αναλάβει η ομάδα μου στην τελευταία μου αποστολή και τώρα, όπως ο πατέρας μου, έχω αναλάβει να τελειώσω αυτή την αποστολή μόνος μου. Δεν έχω καταλάβει ακόμη, όμως, πως συνδέεται ο Μαλχη με αυτή τη γυναίκα... Υποθέτω υπάρχει μόνο ένας τρόπος να το μάθουμε»

Μπαίνοντας στην εκκλησία, έσπευσα κατευθείαν προς τις σκάλες που οδηγούσαν στο υπόγειο. Όσο κατέβαινα, άκουγα την φωνή της καλόγριας, η οποία ήταν σαν να κάνει ομιλία σε κοινό. Ταυτόχρονα μπορούσα να διακρίνω από την ίδια κατεύθυνση πολλά ακόμη στόματα να μασουλάνε.

«Ακούσατε ότι ειπώθηκε η εντολή: “Θα αγαπάς τον πλησίον σου”, και θα μισείς τον εχθρό σου» ήταν οι πρώτες λέξεις που μπόρεσα να ακούσω καθαρά την καλόγρια να λέει «Εγώ, όμως, σας λέω: Να αγαπάτε τους εχθρούς σας, να ευλογείτε εκείνους που σας καταριούνται, να ευεργείτε εκείνους που σας μισούν, και να προσεύχεστε για εκείνους που σας βλάπτουν και σας κατατρέχουν...»

Φτάνοντας στο υπόγειο, δεν μπήκα στο δωμάτιο που βρίσκονταν οι καλόγρια και τα παιδιά της. Κάθισα πίσω από τη μισάνοιχτη πόρτα, καθώς με προβλημάτισαν τα λόγια της. “Να αγαπάτε τους εχθρούς σας”; Μετά από όσα έχουν κάνει; Ούτε κατα διάνοια. Στάθηκα λίγο ακόμη έξω από την πόρτα για να δω που θα κατέληγε ο λόγος της.

«Έτσι θα γίνεται παιδιά του ουράνιου Πατέρα σας, γιατί αυτός ανατέλλει τον ήλιο του για κακούς και καλούς, και στέλνει τη βροχή σε δίκαιους και άδικους. Γιατί αν αγαπήσετε μόνο όσους σας αγαπούν, ποια αμοιβή περιμένετε από το Θεό; Το ίδιο δεν κάνουν κι οι τελώνες; Κι αν χαιρετάτε μόνο τους φίλους σας, τι παραπάνω κάνετε από τους άλλους; Μήπως και οι τελώνες το ίδιο δεν κάνουν; Να γίνεται, λοιπόν, κι εσείς τέλειοι, όπως τέλειος είναι και ο Πατέρας σας ο ουράνιος» τελείωσε η καλόγρια και ακούστηκε ένας ήχος σαν να κλείνει κάποιο βιβλίο. «Λοιπόν παιδιά, τι έχετε να πείτε;»

«Εγώ έχω μία ερώτηση» είπα στερεώνοντας το πιστόλι μου στο πίσω μέρος του κεφαλιού της. Μέσα στο δωμάτιο η καλόγρια στεκόταν όρθια κρατώντας ένα ευαγγέλιο και γύρω της υπήρχαν δεκάδες ενήλικες, γύρω στην ηλικία των 20 ετών, οι οποίοι κάθονταν στο πάτωμα του δωματίου κρατόντας ο καθένας από μία σακούλα πατατάκια ή ποπ κορν, κανένας τους όμως δεν έκανε κάποια κίνηση για να με σταματήσει. Με κοιτούσαν όλοι με το φοβισμένο ύφος και περίμεναν από την καλόγρια να δράσει για να διώξει τον εισβολέα.

«Και ποια είναι η ερώτηση αυτή... Σπαρκ;»

Έμεινα άφωνος και χαλάρωσε για λίγο η λαβή μου, καθώς αιφνιδιάστηκα με τα λόγια της γυναίκας που σημάδευα.

«Πώς ξέρεις το όνομα... Τέλος πάντων» έπιασα το όπλο μου καλύτερα και την ξανασυμάδεψα. «Πες μου ποια είναι η σχέση σου με τον Μαλχη»

«Μήπως θα ήθελες να πάμε πάνω πρώτα;» μου πρότεινε η καλόγρια γυρνώντας άφοβα να με κοιτάξει, λες και δεν απειλούσα τη ζωή της εκείνη τη στιγμή.

«Πάνω;»

«Ναι. Δεν θέλω να γίνει κάποιο ατύχημα και να χάσει κανένα παιδί μου τη ζωή του» μου χαμογέλασε.

Δεν της απάντησα.

«Μην αγχώνεσαι, δεν πρόκειται να σου στήσουμε κάποια ενέδρα, αν σκέφτεσαι αυτό»

«Καλά... πάμε» είπα χωρίς να χαμηλώσω το όπλο μου.

«Λοιπόν, παιδιά, καθίστε φρόνιμα μέχρι να γυρίσω, όπως κάθεστε όταν πάω στο καζίνο» είπε η καλόγρια και άφησε το ευαγγέλιο σε έναν από τα “παιδιά”. Όλοι τους έγνεψαν θετικά με το κεφάλι ταυτόχρονα. «Πάμε, λοιπόν» προχώρησε η γυναίκα μπροστά μου και την ακολούθησα με το όπλο μου να μην ξεφεύγει από τον στόχο του.

Καθώς ανεβαίναμε, το μόνο που ακουγόταν ήταν τα βήματα μας πάνω στα μαρμάρινα σκαλοπάτια. Η καλόγρια δεν έδειχνε καμία πρόθεση να ξεκινήσει μία συζήτηση για να με αποσπάσει από τον στόχο μου, αν και δεν υπήρχε περίπτωση να αποσπαστώ τώρα που ήμουν τόσο κοντά στην πρώτη μικρή μου νίκη.

Φτάνοντας στο κύριο χολ, η γυναίκα προχώρησε μέχρι το σημείο στο κέντρο της εκκλησίας όπου την είχα δει να κάνει προσευχή και σταύρωσε τα χέρια της.

«Λοιπόν...» άκουσα μία γνώριμη αντρική φωνή από την κατεύθυνση της γυναίκας.

Πριν προλάβω να αντιδράσω, ρίζες μαγείας Φύσης με είχαν τυλίξει από τα πόδια μέχρι τον λαιμό, ρίχνοντας το πιστόλι μου ταυτόχρονα από το χέρι μου στο πάτωμα. Προσπάθησα να το ξαναπιάσω, αλλά οι ρίζες με είχαν ακινητοποιήσει τελείως από τον λαιμό και κάτω.

«Θα μιλήσεις;» άκουσα πάλι την αντρική φωνή από την κατεύθυνση της καλόγριας. Έστρεψα το βλέμμα μου προς το μέρος της, αλλά ο άνθρωπος που είδα όχι μόνο δεν ήταν η καλόγρια, αλλά ίσως ήταν το αντίθετο. Μπροστά μου στεκόταν ο Μαλχη με σταυρωμένα χέρια, ο οποίος περίμενε να του απαντήσω.

«Πώς...» ήταν η μόνη λέξη που μπόρεσα να αρθρώσω μέσα στην έκρηξη συναισθημάτων που είχα εκείνη τη στιγμή. Πώς γίνεται αυτό; Ο Μαλχη είναι... η καλόγρια; Μα πώς γίνεται αυτό;

«Μαγεία Αλλαγής Όψης» μου απάντησε.

Μα πώς δεν το σκέφτηκα. Ήξερα ότι υπάρχει τέτοια μαγεία, αφού η ίδια η μητέρα μου την είχα. Πώς δεν μου πέρασε ούτε μία φορά από το μυαλό εδώ και δέκα χρόνια; Είμαι όντως άχρηστος.

«Θα πεις κάτι ή θα κοιτιόμαστε; Εγώ σου απάντησα στην ερώτηση που έκανες, οπότε σειρά σου τώρα»

«Άσε με... κάτω» του πρότεινα. «...να λύσουμε τις διαφορές μας... με διάλογο»

Ένιωσα τη λαβή των ριζών του Μαλχη να χαλαρώνει. Εκείνη τη στιγμή, πριν προλάβει να αντιδράσει, έβαλα τα “παπούτσια” μου και έτρεξα προς το μέρος του. Με την φόρα μου μάζεψα, του έριξα μία μπουνιά στο πρόσωπο, η οποία τον εκτίναξε ως τον τοίχο. Αυτός, τραυματισμένος, προσπάθησε να με πυροβολήσει με το ένα του χέρι, ενώ ταυτόχρονα με το άλλο του χέρι θεράπευε την τραυματισμένη περιοχή του προσώπου του. Η αποφυγή των βλημάτων του ήταν εύκολη τώρα που είχα υπεράνθρωπη ταχύτητα, οπότε προσπάθησα να τον πλησιάσω καθώς απέφευγα τις σφαίρες του, αυτός όμως μετέτρεψε τα πόδια του σε ελατήρια και έκανε ένα υπεράνθρωπο άλμα προς την οροφή της εκκλησίας, από την οποία έβγαλε μερικές ρίζες και κρατήθηκε.

Ανίκανος να κάνω κάτι, κάθισα στο κέντρο της εκκλησίας και περίμενα την επόμενη του κίνηση. Θεώρησα πως αν είμαι στο κέντρο της εκκλησίας, θα έχω άπλετο χώρο να αποφύγω την επόμενη επίθεση του Μαλχη, ανεξάρτητα από το μέγεθος της. Το λάθος μου όμως το κατάλαβα όταν ξαφνικά είδα αριστερά μου να σηκώνονται ρίζες μαγείας Φύσης και δεξιά μου ένα κύμα μαγείας Φωτιάς. Ήμουν περιτριγυρισμένος από μαγεία. Δεν μπορούσα να κάνω τίποτα. Οι μυτερές ρίζες και το κύμα φωτιάς έφτασαν μέχρι το ύψος της οροφής. Είδα τον Μαλχη να κατεβάζει το χέρι που δεν χρησιμοποιούσε για να κρατηθεί από την οροφή και τα τείχει από μαγείες άρχισαν να πέφτουν προς το μέρος μου. 

Αυτό ήταν, ήμουν νεκρός. Δεν μπορούσα να κάνω κάτι για να σώσω τον εαυτό μου. Θα μπορούσα να είχα πετάξει έναν κεραυνό τόση ώρα προς την κατεύθυνση του Μαλχη, ή ακόμη να είχα πιάσει το όπλο μου με την ταχύτητα μου πριν προλάβει να σηκωθεί ο Μαλχη, αλλά δεν το σκέφτηκα εκείνη τη στιγμή. Είμαι... άχρηστος.

Κάθισα στο πάτωμα και αποδέχτηκε τη μοίρα μου. Κανείς δεν μπορεί να αποφύγει δύο τεράστια τείχη μαγείας να τον κλείνουν από όλες τις πιθανές γωνίες, σαν δύο χέρια που χτυπάνε παλαμάκια. Ήταν μοιραίο πως θα πέθαινα και εγώ από τον Μαλχη.

«Πατέρα έρχομαι...»

\chapter{}

Ανοίγοντας τα μάτια μου όμως, κατάλαβα πως δεν ήταν γραμμένο να τελειώσει εδώ η ιστορία μου. Το τείχος φωτιάς που ερχόταν καταπάνω μου από τα δεξιά είχε εξαφανιστεί και το στεγνό μαρμάρινο πάτωμα είχε βραχεί, λες και η βροχή έσβησε την φωτιά. Η τρύπα στην οροφή όμως δεν βρισκόταν στη δεξιά μεριά της εκκλησίας, αλλά στην αριστερή, από την οποία υποτίθεται δέχτηκα επίθεση από ένα τείχος μυτερών ριζών μαγείας Φύσης. Ο Μαλχη προφανώς είχε λάβει υπόψη του το περιβάλλον στο οποίο παλεύαμε, δεν υπήρχε περίπτωση να είχε κάνει τέτοιο λάθος. Τότε πως έζησα;

Έστρεψα το βλέμμα μου στα αριστερά και είδα μία σκηνή που μου θύμισε τις παλιές όχι τόσο καλές εποχές. Εναντίον του Γιν, σε αυτή τη μικρή μάχη, ο Χερμπ είχε διαπεραστεί από μία τεράστια ρίζα μαγείας φύσης. Αυτή τη φορά όμως, δεν ήταν μία, αλλά πολλές μικρότερες μυτερές ρίζες, οι οποίες δεν με είχαν χτυπήσει καθώς είχαν συναντήσει ένα εμπόδιο στη διαδρομή τους. Ο Βασι στεκόταν στην αριστερή μου μεριά και είχε μπλοκάρει τις ρίζες του Μαλχη, προστατεύοντας με από βέβαιο θάνατο, σε αντάλλαγμα με τη δική του ζωή.

«Γιατί...» σηκώθηκα γρήγορα και έπιασα το σώμα του πριν προλάβει να πέσει στο πάτωμα.

«Μην... ανησυχείς. Η Μπλαιρ και ο Κουπ...» έβηξε μία μικρή ποσότητα αίματος στο πάτωμα.

«Η Μπλαιρ και ο Κουπ;»

\textasciitilde ...θα έρθουν να σε βοηθήσουν όταν βρεθείς αντιμέτωπος με τη θνητότητά σου...\textasciitilde \  διάβασα κατά λάθος τις τελευταίες σκέψεις του Βασι καθώς έπεφτε στην αγκαλιά μου αδυνατώντας να μιλήσει. Τι σημαίνει αυτό; Γιατί χρησιμοποιεί ο Βασι μεγάλες λέξεις στην τρέχουσα κατάσταση που βρισκόμαστε;

Οι ρίζες του Μαλχη υποχώρησαν από το σώμα του Βασι, αποκαλύπτοντας τις δεκάδες τρύπες που είχαν δημιουργηθεί στο σώμα του πλέον νεκρού φίλου μου. Παρόλα αυτά, εκείνη τη στιγμή, ένιωσα τα χέρια του Βασι να προσπαθούν να κάνουν μία τελευταία κίνηση. Ο Βασι δεν είχε πεθάνει ακόμη, προσπαθούσε να μου δώσει τη μαγεία του, τη μαγεία Νερού που πενευόταν από μικρός πώς θα την χρησιμοποιούσε κάποια μέρα για να παλέψει στο πλευρό μου, και στην ουσία, στις τελευταίες του στιγμές σε αυτόν τον κόσμο αυτό έκανε. Χρησιμοποίησε την μαγεία του για να με σώσει.  Γιατί εμένα όμως; Αυτός ήταν πολύ καλύτερος στις μάχες από τον άχρηστο ανακριτή Σπαρκ. Γιατί να μη με αφήσει να πεθάνω και να μην επιτεθεί αυτός στον Μαλχη; Και η Μπλαιρ και ο Κουπ που κολλάνε; Μήπως θα έρθουν να με βοηθήσουν όταν πάω όντως να πεθάνω;

Ο Μαλχη χρησιμοποίησε άλλη μία φορά τα ελατήρια στα πόδια του και προσγειώθηκε στο πάτωμα. 

«Ευτυχώς που έπεισα την Ιφι πριν τόσα χρόνια να μου δώσει τη μαγεία της. Με έχει σώσει αρκετές φορές τα τελευταία χρόνια...»

Δεν γύρισα καν να τον κοιτάξω. Με τον Βασι στα χέρια μου να προσπαθεί να τελειώσει την χειρονομία μεταφοράς μαγείας, τα συναισθήματα μου είχαν φτάσει ήδη στο ζενίθ.

Τελειώνοντας τις κινήσεις με τα χέρια του, ο Βασι προσπάθησε να σηκώσει το κεφάλι του αλλά πριν προλάβει ένιωσα όλους του μυς στο σώμα του να χαλαρώνουν και να πέφτουν πάνω στο αδύναμο λιγνό μου στήθος. Μαζί με το βάρος του Βασι, ένιωσα και μία περίεργη αίσθηση στο αίμα μου, σαν να είχε αυξηθεί η κυκλοφορία του αίματος, αλλά ο ρυθμός με τον οποίο χτυπούσε η καρδιά μου είχε μείνει ίδιος. Κοίταξα τα χέρια μου και αυτή τη φορά δεν είχα γύρω μου μόνο μικρές σπίθες κεραυνών, αλλά και ένα στρώμα από σταγόνες νερού. Είναι θαύμα πως δεν είχα πάθει ηλεκτροπληξία εκείνη τη στιγμή που με περιτριγύριζαν δύο τόσο αντίθετες μαγείες, καθώς κάθε φορά που οι δύο μας μαγείες ακουμπούσαν η μία την άλλη, το μέγεθος της σπίθας κεραυνού αυξανόταν τρομακτικά.

«Σπαρκ, για πες, τι είναι αυτές οι μαγείες γύρω σου; Μήπως είσαι λιγουλάκι θυμωμένος;»

Τι λέει; Προσπαθεί να μου αυξήσει την ένταση των συναισθημάτων που ήδη οργιάζουν στο μυαλό μου; Αν είναι αυτός ο σκοπός του, το πέτυχε. Οι κεραυνοί γύρω μου, όσο μεγάλωνε ο θυμός, μου τόσο μεγαλύτεροι και πιο τρομακτικοί γίνονταν, συνοδευόμενοι από αντίστοιχους θορύβους που μπορεί να ακούσει κανείς μόνο σε περιόδους καταρρακτώδης βροχής.

«Δεν μιλάς, ε; Λογικό, αφού...»

Έβαλα τα “παπούτσια” ταχύτητας μου και έσπευσα προς το μέρος του με τις γροθιές μου οπλισμένες, οδηγούμενος αυτή τη φορά μόνο από τα συναισθήματα μου, ανίκανος να σκεφτώ καθαρά. Αυτός, όμως, απέφυγε το χτύπημα μου κάνοντας μόνο ένα βήμα προς τα δεξιά και με παγίδευσε άλλη μία φορά με τις ρίζες της μαγείας Φύσης του, ακινητοποιώντας με αυτή τη φορά από τα πόδια μέχρι και το κεφάλι

«Δεν πιστεύω να νόμιζες πώς η ίδια επίθεση θα δούλευε δεύτερη φορά στον αρχηγό της Ειδικής Υπηρεσίας» είπε και με πυροβόλησε κοντά στον δεξιό πνεύμονα, αποφεύγοντας να με τραυματίσει επίτηδες θανάσιμα. 

Βήχοντας μία μικρή ποσότητα αίματος, ένιωσα τη λαβή των ριζών του να χαλαρώνει και να με αφήνει άλλη μία φορά να πατήσω στα πόδια μου. Αυτή η μικρή ευχαρίστηση μου δεν κράτησε πάνω από μερικά δευτερόλεπτα, καθώς ο Μαλχη με πυροβόλησε στα πόδια, αφαιρώντας μου την επιλογή να ξαναφορέσω τα “παπούτσια” μου και να του ξανα επιτεθώ με ταχύτητα.

«Ωραία, τώρα μάλλον είσαι τελείως άχρηστος» μου είπε ο Μαλχη περπατώντας προς το μέρος μου. «Σαν πρωταγωνιστής μιας υπερρεαλιστικής ιστορίας επέζησες την πρώτη ενέδρα, μετά δέκα χρόνια στη φυλακή, ακόμη και την επίθεση που σου έκανα περι τριγυρίζοντας σε με μαγεία. Τώρα όμως αυτή η αόρατη πανοπλία που φοράς που σε σώζει από ακραία γεγονότα για να συνεχίσει η ζωή σου έσπασε και η στιγμή του ηρωισμού σου τελείωσε»

«Γιατί;...» τον ρώτησα από την πλέον ξαπλωμένη πλαγίως θέση που είχα πάρει, κρατώντας το σημείο στο στήθος μου όπου με πυροβόλησε με την πρώτη του πετυχημένη σφαίρα.

«Γιατί η περιέργεια σκότωσε τη γάτα Σπαρ, για αυτό. Αν σαν καλό παιδί μου έδινες τη μαγεία σου εκείνο το βράδυ, χωρίς να διαβάσεις τις σκέψεις μου, δεν θα βρισκόμασταν σε αυτό το σημείο τώρα»

«Δεν... εννοούσα αυτό... αλλά η Ιφι και ο Φυτ... γιατί»

«Ιφι και Φυτ;» είπε ο Μαλχη προβληματισμένος, κοιτώντας την κάννη του πιστολιού του. «δεν ξέρω για ποιους μιλάς»

«Εσύ δεν ανέφερες πριν λίγο την Ιφι;»

Ο Μαλχη με πυροβόλησε στο στομάχι και έβγαλα μία κραυγή πόνου ισάξια σε ένταση με το κάλεσμα ενός ελέφαντα, ή τουλάχιστον έτσι νόμισα εκείνη τη στιγμή.

«Ποιος είσαι εσύ να μου πεις τι είπα και τι δεν είπα; Ένας χαζός και άχρηστος πράκτορας είσαι Σπαρκ. Ούτε καν τι μαγεία έχεις δεν ξέρεις»

«Τι... εννοείς;»

«Νομίζεις πώς είσαι έξυπνος; Πώς χρησιμοποιείς τους κεραυνούς σου τόσο καιρό για να διαβάσεις τις σκέψεις αλλονών; Ε, λοιπόν αν είχες την παραμικρή περιέργεια θα καταλάβαινες πώς κάτι πάει λάθος. Γιατί άρχισες να ακούς τις σκέψεις των άλλων όταν πέθανε ο πατέρας σου; Δεν σου έχει περάσει ούτε μία φορά από το μυαλό;»

Δεν μπορούσα να του απαντήσω στην κατάσταση που βρισκόμουν.

«Ο πατέρας σου Σπαρκ, είχε μεν μαγεία Κεραυνών, αλλά είχε κληρονομήσει από την αδερφή του μαγεία Σκέψεων, την οποία χρησιμοποιούσε όμοια με εσένα, αλλά καταλήξατε και οι δυο σας να έχετε το ίδιο τέλος. Δεν μπόρεσα να πείσω τον πατέρα σου να μου δώσει τις μαγείες του, οι οποίες είναι από τις λίγες που δεν έχω καταφέρει να αποκτήσω, για αυτό προσπάθησα να σε κάνω να μου δώσεις εσύ τις μαγείες σου. Η εθελοτυφλία σου, πάντως, με εκπλήσει. Γιατί δεν αναρωτήθηκες ποτέ κάποια παραπάνω πράγματα για την μαγεία; Από που προήλθε, πώς κληρονομείται, πόσες διαφορετικές μαγείες υπάρχουν, αλήθεια δεν αναρωτήθηκες ποτέ;»

Δεν ήταν πώς δεν είχα αναρωτηθεί ποτέ. Από μικρός είχα περιέργεια να μάθω για τη μαγεία, αλλά όταν μεγάλωσα, άρχισα να βλέπω τη μαγεία ως μαύρο κουτί, το οποίο τι χρησιμοποιώ σε ειδικές περιπτώσεις για να ολοκληρώσω συγκεκριμένες δουλειές. Αυτό γιατί είχα την εντύπωση πως η ανθρωπότητα δεν έχει τις απαραίτητες γνώσεις ακόμη για τη μαγεία, αλλά απ’ ό,τι φαίνεται, αφού ο παππούς του Μαλχη κιόλας είχε γράψει ολόκληρο βιβλίο για την μαγεία, γνωρίζαμε περισσότερα πράγματα για τη μαγεία από όσα πίστευα. Που να το εξηγήσω τώρα όλο αυτό στον Μαλχη όμως τώρα, ειδικά στην κατάσταση που βρίσκομαι.

«Για... τι;» του είπα.

«Γιατί δεν σε αποτελειώνω ή γιατί συλλέγω μαγείες; Είναι τελείως διαφορετικές ερωτήσεις»

«Το δευ...»

«Εσύ γιατί νομίζεις» είπε και με πυροβόλησε, αυτή τη φορά στον αριστερό πνεύμονα. «Όλα ξεκίνησαν όταν...» άρχισε να μου εξηγεί ο Μαλχη, αλλά η ένταση της φωνής του χαμήλωνε και η όραση μου θόλωνε με το δευτερόλεπτο. Δεν μπορούσα να διακρίνω τις λέξεις που έβγαιναν από το στόμα του, έβλεπα απλά μία θολή εικόνα ενός ανθρώπου να περπατάει σε κύκλους μπροστά μου κουνώντας τα χέρια του σαν Ιταλός.

«...καλα, αφού δε με ακούς, αντίο» ήταν οι τελευταίες λέξεις που κατάφερα να διακρίνω στην κατάσταση που βρισκόμουν. Τελειώνοντας την πρόταση του αυτή, άκουσα το όπλο του να πυροβολεί μία τελευταία φορά και τα φώτα για μένα έσβησαν.

\chapter{}

Ώστε έτσι είναι η μετά θάνατον ζωή. Δεν μπορώ να δω τίποτα, υπάρχει σκοτάδι παντού. Είμαι απλά μόνος μου με τις σκέψεις μου. Μου θυμίζει τα χρόνια που πέρασα στη φυλακή, απλά τότε τουλάχιστον είχα κι άλλες φωνές απ’ έξω από το κελί μου να μου κρατούν συντροφιά, αλλά και φως από το παράθυρο του κελιού να μου δείχνει τι ώρα είναι, τι εποχή είναι, τι... τι με νοιάζει τώρα αυτό; Αφού είμαι νεκρός, ή τουλάχιστον... έτσι νομίζω. Δεν μπορώ να φανταστώ πως είναι δυνατόν να επιζήσω μία τέτοια κατάσταση.

“Να αγαπάτε τους εχθρούς σας”; Τι φράση και αυτή. Ο Μαλχη βέβαια δεν την τύρισε, αφού με σκότωσε εν ψυχρώ, χωρίς δεύτερη σκέψη. Βασικά... με άφησε για λίγο να ζήσω, έστω και τραυματισμένος. Τι διπολικός άνθρωπος. Δεν μπόρεσα να καταλάβω τελικά ποιος ήταν ο στόχος του. Δεν κατάφερα να γίνω ο ήρωας πρωταγωνιστής της ιστορίας μου, που μετά από όλες τις δυσκολίες καταφέρνει να νικήσει τον κακό. Και γιατί; Επειδή είμαι άχρηστος. Θα μπορούσα άνετα, όσο ο Μαλχη ήταν στο ταβάνι και θεράπευε τον εαυτό του να είχε πιάσει το όπλο μου από το πάτωμα και να τον είχα πυροβολήσει, ή ακόμη να του είχα ρίξει έναν κεραυνό με την μαγεία μου. Όλα αυτά είναι απλά κριτικές που κάνω στον εαυτό μου. Δεν σημαίνουν τίποτα στην τελική. Εκείνη τη στιγμή ήμουν όσο άχρηστος όσο πάντα, δεν μπορούσα να σκεφτώ καθαρά, με καθοδηγούσε το μίσος και η εκδίκηση. Και γιατί; Αν σκότωνα τον Μαλχη θα ολοκληρωνόμουν σαν άνθρωπος; Αν σκότωνα τον εχθρό μου θα μπορούσα να επιστρέψω στην παλιά μου ζωή; Σαχλαμάρες, μόνο χειρότερα θα γίνονταν τα πράγματα. Τι σκεφτόμουν πραγματικά.

“Δεν ξέρω ποιος είμαι” και “Δεν ξέρω που είμαι”. Αυτές είναι λέξεις που θα μπορούσαν να περιγράψουν την κατάσταση μου αυτή τη στιγμή, καθώς ούτε ποιος είμαι ξέρω, ούτε που είμαι ξέρω. Μπορεί να είμαι μία τυχαία ψυχή που οδεύει προς τον παράδεισο ή την κόλαση. Μπορεί να είμαι ένα κοτοπουλάκι και περιμένω να εκκολαφθεί το αυγό μου. Μπορεί απλά να συμβαίνει ακριβώς αυτό που βλέπω: τίποτα. Ποιος ξέρει... και ποιος νοιάζεται πλέον.

Ο Κουπ με την Μπλαιρ; Δεν ήρθαν ποτέ να με βοηθήσουν. Η ζωή μου τελείωσε και ήμουν μόνος, χωρίς φίλους, χωρίς οικογένεια, χωρίς... να έχω κάποιον να μιλήσω για όσα περνούσα. Μόνο ο Βασι ήταν εκεί για μένα τις τελευταίες μου στιγμές. Με έβγαλε από τη φυλακή και με έσωσε από βέβαιο θάνατο, παρόλο που πέθανα έτσι κι αλλιώς μετά από πέντε λεπτά. Γιατί να το κάνει αυτό; Είχαμε να μιλήσουμε πάνω από δέκα χρόνια και μου έσωσε τη ζωή λες και ήμασταν αδέρφια. Βέβαια, όταν ήμασταν μικροί τον ένιωθα σαν μεγάλο αδελφό, αλλά αδίκησε τον εαυτό του δίνοντας τη ζωή του για έναν άνθρωπο σαν κι εμένα, την οποία έχασα ακριβώς μετά την ηρωική του πράξη. Είμαι όντως άχρηστος.

Όλα αυτά όμως δεν μετράνε στη μετά θάνατον ζωή που βρισκόμαστε και οι δυο μας αυτή τη στιγμή. Βασικά, λες να μετράνε; Αφού η καλόγρια, δηλαδή... ο Μαλχη, είχε πει πως αν αγαπάμε τους εχθρούς μας γινόμαστε σαν τον ουράνιο Πατέρα μας ή κάτι τέτοιο, οπότε αν... έπαιρνα το μέρος του Μαλχη... δεν βγάζει νόημα. Γιατί να το κάνω αυτό; Να βοηθήσω έναν εγκληματία, η οποία δεν είναι τόσο βαριά λέξη για αυτόν, να φτάσει τον στόχο του; Μπορεί να αναφερόταν σε πιο απλές περιπτώσεις όπου δύο άνθρωποι τσακώνονται και μένουν για πάντα εκνευρισμένοι ο ένας με τον άλλον για μικροπράγματα, τότε το καταλαβαίνω. Αλλά δεν μπορώ παρά να σκέφτομαι αν όντως έπρεπε να πάρω το μέρος του Μαλχη. Δεν ξέρεις, μπορεί τελικά να είχε κάποιον καλό σκοπό συλλέγοντας μαγείες, απλά ο τρόπος με τον οποίο τις σύλλεγε να ήταν λίγο... περίεργος. Άμα είχα βοηθήσει, όμως, τον Μαλχη λες να είχα πάει στον παράδεισο; Να ήμουν τιμωρημένος σε μία αιωνιότητα μοναχικότητας, συνοδευόμενος μόνο από τις σκέψεις μου. Αυτή η σκέψη θα με στοιχειώνει όσο καιρό είμαι παγιδευμένος εδώ.

Δεν ήμουν και μεγάλος οπαδός της θρησκείας, μπορεί να παίζει και αυτό ρόλο. Αλλά μου λες ότι αν πήγαινα, πχ, εκκλησία κάθε Κυριακή και προσευχόμουν κάθε μέρα θα κατέληγα σε άλλο μέρος μετά τον θάνατο μου; Τι ανοησίες είναι αυτά. Πώς μπορεί να ξέρει ένας θνητός τι θα γίνει όταν η θνητότητα του λήγει; Μπορεί να έχουν δίκιο, ποιος ξέρει, απλά η προσωπική μου εμπειρία μέχρι στιγμής διαφωνεί... τι είναι αυτό; Μία μύγα; Τι κάνει μία μύγα εδώ; Ωπ, έφυγε από το... “οπτικό πεδίο” μου, θα έλεγα αν ήμουν ζωντανός, αλλά δεν είμαι. Αλλά για να υπάρχει κάτι άλλο ζωντανό εδώ που βρίσκομαι σημαίνει ότι κάπου όντως βρίσκομαι. Δεν μπορώ να δω το σώμα μου, δεν μπορώ να ξεχωρίσω το μπροστά και το πίσω, δεν μπορώ να κάνω τίποτα. Νιώθω σαν να μη μπορώ να κουνήσω το σώμα μου όπως το θέλω εγώ, λες και είμαι σε όνειρο.

Κι άλλος ήχος; Ένα βέλασμα κατσίκας; Δεν μπορώ να γυρίσω για να δω αν υπάρχει όντως κάποια κατσίκα, αλλά σίγουρα άκουσα αυτό το χαρακτηριστικό βέλασμα που κάνει μία κατσίκα όταν θέλει να στρέψει την προσοχή στον εαυτό της... ή μήπως ήταν πρόβατο; Πάντα τα μπέρδευα αυτά... τι είναι αυτό; Φρένα αυτοκινήτου; Φωνές;

«ΡΙΝΤ ΤΙ ΚΑΝΕΙΣ;»

«ΑΑΑΑΑ»

«ΜΗ ΦΩΝΑΖΕΤΕ»

*μπουμ*

Ριντ; Μπαμπά είσαι εκεί; Που συμβαίνουν όλα αυτά; Βλέπω μόνο ένα σκοτάδι. Το αυτοκίνητο ακούστηκε να εκρήγνυται. Είναι όλοι καλά; Τι ρωτάω, λες και μπορώ να τους βοηθήσω στην κατάσταση που βρίσκομαι.

«Αργήσαμε πολύ, οι τέσσερις τους δεν έχουν παλμό. Ο μόνος που μάλλον ζει είναι αυτός εδώ με το βιβλίο»

Άλλες φωνές; Διαφορετικές, πιο μπάσες από τις προηγούμενες. Μπορεί να ήρθαν να τους βοηθήσουν.

«Δεν θέλω να χάσω και τους πιο κοντινούς μου φίλους»

Κουπ;

«Πρέπει να ζήσεις Σπαρκ...»

Μαμά;

«Δεν πρέπει να χάσω και εσένα Σπαρκ. Σε παρακαλώ, ζήσε... είσαι ο μόνος άνθρωπος που έχω»

Μπλαιρ; Τι γίνεται; Δεν μπορώ να δω τιπο... ένα φως; Με πλησιάζει... Τι πρέπει να κάνω; Είστε στην άλλη μεριά; Σας παρακαλώ... ας μου απαντήσει κάποιος...

\chapter{}

Περνώντας μέσα από το φως, βρέθηκα ξαπλωμένος σε ένα κρεβάτι νοσοκομείου. Ξαφνιασμένος σηκώθηκα, αλλά όλοι μου οι μυς τραντάχτηκαν και έπεσα πάλι πίσω στο μαξιλάρι. Προσπάθησα να επεξεργαστώ τον χώρο από την ανάσκελη θέση που είχα βρεθεί. Έξω ήταν ακόμη βράδυ, με ένα πλήρες φεγγάρι να βασιλεύει στον νυχτερινό ουρανό. Στα δεξιά μου η μητέρα μου κοιμόταν σε ένα σκαμπό, ακουμπώντας το κεφάλι και τον ώμο της στην γωνία του δωματίου.

«Μαμά» προσπάθησα να την ξυπνήσω.

Ανοίγοντας αργά τα μάτια της, η μητέρα μου ξαφνιάστηκε και σηκώθηκε από το κάθισμα της για να πέσει στην αγκαλιά μου.

«Ω, Σπαρκ» είπε ξαπλωμένη μπρούμυτα στο κρεβάτι, σίγκοντας με το δεξί της χέρι το πάπλωμα μου. Στα μάγουλα της κυλούσαν δάκρυα χαράς συνοδευόμενα από ένα πλατύ χαμόγελο στα χείλη της.

«Μαμά;»

«Δεν ξέρεις πόσο καιρό ελπίζω πώς θα γυρίσεις πίσω...» είπε η μητέρα μου προσπαθώντας να διώξει τα δάκρυα από το πρόσωπο της. «Ήξερα ότι ήταν κακή ιδέα αυτό το ταξίδι, αλλά...»

«Ποιο ταξίδι;»

Με κοίταξε προβληματισμένη. Εκείνη τη στιγμή φάνηκε λες και της πέρασε μία πολύ συγκεκριμένη σκέψη από το μυαλό αλλά δεν την ανέφερε, ξέσπασε απλά πάλι σε ελιγμούς πάνω μου.

«Μαμά, ποιο...»

«Μην ανησυχείς» με διέκοψε «όλα είναι υπό έλεγχο» μου απάντησε και σηκώθηκε για να βγει έξω από το δωμάτιο. 

Πώς είμαι ακόμη ζωντανός; Πώς και με θυμήθηκε η μητέρα μου αυτή τη στιγμή; Τι μου συμβαίνει; Προσπάθησα να σηκωθώ για να την ακολουθήσω, αλλά μπόρεσα μόνο να αλλάξω τη στάση του σώματος μου στο κρεβάτι, από ξαπλωμένος σε καθιστός. Το σώμα μου δεν μου επέτρεπαν να περπατήσω, καθώς όταν προσπάθησα να πατήσω το πόδι μου στο πάτωμα, ένας βασανιστικός πόνος προκαλούταν στις γάμπες και τους αστραγάλους μου. Όχι πως αν σηκωνόμουν θα μπορούσα να πάω και πολύ μακριά, καθώς ήμουν μπλεγμένος σε κάτι περίεργα καλώδια που συνδέονταν με την καμπή των αγκώνων μου, αλλά και το κρανίο μου. 

Ευτυχώς δεν χρειάστηκε να περιμένω πολύ ώρα. Η μητέρα μου επέστρεψε μέσα σε λίγα λεπτά με έναν γιατρό, ο οποίος όταν με πρωτοείδε χαμογέλασε και άρχισε να γράφει κάτι στο λάπτοπ του.

«Ωραιά. Τώρα που ο γιος σας είναι ξύπνιος θα χρειαστεί να κάνουμε κάποιους ελέγχους για να δούμε σε τι επίπεδο βρισκόμαστε»

«Κάντε ό,τι χρειάζεται, αρκεί το παιδί μου να είναι καλά» είπε η μητέρα μου και προχώρησε πάλι προς την πόρτα του δωματίου.

«Μισό» την σταμάτησα πριν προλάβει να βγει. «Να σε ρωτήσω κάτι;»

«Ό,τι θες παιδί μου»

«Που... είμαι;»

«Μην ανησυχείς, σου είπα, όλα θα πάνε καλά» μου χαμογέλασε και έφυγε από το δωμάτιο, κλείνοντας μας την πόρτα.

«Λοιπόν Σπαρκ» μου είπε ο γιατρός και μου έδειξε στο λάπτοπ του μία εικόνα του πατέρα μου, την ίδια μάλιστα εικόνα που υπήρχε ως κορνίζα στα κεντρικά της Ειδικής Υπηρεσία, απλά χωρίς το όνομα και τον τίτλο του στο κάτω δεξί μέρος της εικόνας. «Αναγνωρίζεται αυτόν τον άντρα;»

«Ναι, είναι ο πατέρας μου, Ριντ Θοτς»

«Πολύ ωραία. Αυτόν τον άντρα;» συνέχισε ο γιατρός και μου έδειξε αυτή τη φορά μία εικόνα του Χερμπ, πάλι την ίδια που υπήρχε και στον διάδρομο πίσω στην Ειδική Υπηρεσία.

«Ναι... νομίζω τον λένε Χερμπ»

«Πολύ ωραία. Αυτόν τον άντρα;» μου έδειξε μία εικόνα του Ρομπ, την οποία δεν είχα ξαναδεί.

Πλάκα πλάκα δεν θυμάμαι να βλέπω κορνίζα του Ρομπ στην Ειδική Υπηρεσία. Λες να μην του είχαν κρεμάσει κορνίζα;

«Δεν τον θυμάστε;»

«Όχι τον θυμάμαι... Ρομπ;»

«Πολύ ωραία. Και τέλος αυτόν τον άντρα;» μου έδειξε μία εικόνα του Βασι.

«Βασι;»

«Πάρα πολύ ωραία, κύριε Θοτς. Μήπως θυμάστε και ποια είναι η γυναίκα που βγήκε από το δωμάτιο πριν λίγο;»

«Η... μητέρα μου;»

«Εξαιρετικά» είπε και ακούμπησα το λάπτοπ του σε ένα τραπεζάκι λίγο πιο δίπλα. «Η μνήμη σας δεν έχει χαθεί τελείως, οπότε η ανάρρωση σας δεν πρέπει να μας πάρει πάνω από έξι μήνες»

«Συγγνώμη, αλλά... γιατί είμαι στο νοσοκομείο;»

«Ήσασταν σε κώμα τους τελευταίους... έντεκα μήνες, λίγο παραπάνω από έντεκα μήνες, βασικά. Ευτυχώς, θυμάστε τα κοντινά σας πρόσωπα, οπότε δεν χρειάζεται να σας βοηθήσουμε σε αυτόν τον τομέα. Έχετε προσπαθήσει να σηκωθείτε;»

«Ναι...»

«Και πώς πήγε;»

«Δεν... μπόρεσα...»

«Άρα θα χρειαστεί να μάθουμε να περπατάμε από την αρχή, ε; Δεν πειράζει, δεν είναι έξτρα βάρος. Μήπως μπορείτε να θυμηθείτε τις τελευταίες σας στιγμές πριν χάσετε τις αισθήσεις σας; Ίσως τον λόγο που μπήκατε σε κώμα;»

«Λογικά... από την μάχη μου με τον Μαλχη. Με είχε πυροβολήσει...»

«Μάχη;» με διέκοψε ο γιατρός κοιτώντας με προβληματισμένος.

«Τι έγινε;»

«Τίποτα...» είπε ο γιατρός και έκλεισε το λάπτοπ του. «Θα έρθουν σε λίγο δύο νοσοκόμες να σου φτιάξουν τους ορούς σου. Ελπίζω να πάνε όλα καλά και σε έξι μήνες περίπου θα μπορέσεις να συνεχίσεις τη ζωή σου κανονικά» είπε ο γιατρός και άνοιξε την πόρτα του δωματίου.

«Τι έγινε γιατρέ;» τον ρώτησε η μητέρα μου.

«Το παιδί σας φαίνεται να είναι μια χαρά. Θα καλέσω δύο νοσοκόμες για να...»

«Μπορούν να έρθουν επισκέψεις αύριο το πρωί;» τον διέκοψε η μητέρα μου.

«Εννοείτε αυτή την κοπέλα και το άλλο το ψιλό αγόρι που έρχονταν κάθε μέρα πριν κάτι μήνες;»

«Ίσως και όλους τους φίλους του γιου μου.»

«Δεν πιστεύετε πως θα του φανεί λίγο ξαφνικό;»

«Όχι»

«...μόνο όχι;»

«Καλά. Αφού σας θυμάται, λογικά θα θυμάται και τα υπόλοιπα πρόσωπα που γνωρίζει, οπότε...»

«Οπότε;»

«Απλά...»

«Απλά τι;»

«Θα σας πω άλλη στιγμή» είπε ο γιατρός, ρίχνοντας μία πλαϊνή ματιά προς το μέρος μου και μας άφησε μόνους μας.

«Ακούς Σπαρκ, οι φίλοι σου θα έρθουν αύριο... δηλαδή σήμερα το πρωί. Δεν είσαι χαρούμενος;»

Δεν της απάντησα, έμεινα απλά καθιστός στο κρεβάτι μου με τα χέρια διπλωμένα και ακουμπισμένα στα μπούτια μου.

Οι φίλοι μου; Ο μόνος πραγματικός μου φίλος που μπορώ να θυμηθώ είναι ο Βασι. Έχω κι άλλους φίλους; Μήπως εννοεί τους συναδέλφους μου από την Ειδική Υπηρεσία; Αυτοί όχι μόνο δεν είναι φίλοι μου. Πώς μπορεί κάποιος στη θέση μου να τους θεωρεί φίλους, μετά από όσα έχουν συμβεί; Λες να μην γνωρίζει η μητέρα μου τι έχει γίνει; Να μην ξέρει για τον Μαλχη; Αυτή είναι η μόνη εξήγηση που μπορώ να σκεφτώ.

«Σπαρκ, όλα καλά;»

«Ναι, συγγνώμη μαμά, αφαιρέθηκα...»

«Γεια σας, ήρθαμε να σας βοηθήσουμε» είπε μία από τις δύο νοσοκόμες καθώς έμπαιναν στο δωμάτιο.

«Τόσο γρήγορα;» την ρώτησε η μητέρα μου.

«Είναι βράδυ και η κίνηση στο νοσοκομείο δεν είναι τόσο μεγάλη, οπότε οι μικροδουλειές γίνονται γρήγορα» της απάντησε χαμογελαστή η ίδια νοσοκόμα και ήρθε προς το μέρος μου.

«Καλά... Σπαρκ σε αφήνω να ξεκουραστείς μέχρι το πρωί και θα έρθω με το αγαπημένο σου πρωινό, εντάξει;»

«Εντάξει...»\\\\

«Ευχαριστώ για το πρωινό μαμά. Είχα να φάω καιρό φαγητό φτιαγμένο από τα χέρια σου»

«Αυτό ξαναπές το...» μου χαμογέλασε η μητέρα μου και πήρε το τάπερ με την ομελέτα που μου είχε φτιάξει. Η μητέρα μου είχε επιστρέψει νωρίς το πρωί για να με ξυπνήσει και να με ταΐσει, απλά δεν ξέρω κατά πόσο ήταν “σωστό” να τρώω φαγητό ενώ έχω ορούς στα χέρια μου. Η μητέρα μου, όμως, δεν μου άφηνε άλλη επιλογή.

«Μαμά...»

«Ναι Σπαρκ;»

«Πότε θα έρθουν αυτοί οι... “φίλοι” μου;»

«Τους είπα όποτε να έρθουν όποτε μπορούν. Είναι ακόμη 7 παρά η ώρα, μπορεί να μην έχουν δει καν τα μηνύματα μου. Μην αγχώνεσαι πάντως, όλα θα πάνε καλά τώρα που ξύπνησες»

«Γιατί το λες αυτό συνέχεια;»

«Προσπαθώ να σε...»

Πριν προλάβει να τελειώσει η μητέρα μου την πρόταση της, είδα με την άκρη του ματιού μου την πόρτα του δωματίου να ανοίγει από μία κοπέλα με πράσινα μάτια και καστανά κυματιστά μαλλιά.

«Ιφι;» είπα ξαφνιασμένος.

Αυτή έτρεξε προς το μέρος μου και έπεσε στην αγκαλιά μου, ξεσπώντας σε ελιγμούς. Δεν ήξερα πως να αντιδράσω. Τελικά με θυμόταν; Γιατί να με θυμηθεί τώρα; Μήπως ήταν κάποιο κόλπο του Μαλχη που τους είχε κάνει να με ξεχάσουν;

«Ιφι...»

«Συγγνώμη Σπαρκ... συγγνώμη...» μου είπε η Ιφι σφίγγοντας την αγκαλιά της και ακουμπώντας το κεφάλι της απαλά στο δικό μου. Ένιωθα τα δάκρυα της να κυλούν στον ώμο μου και να κατεβαίνουν την πλάτη μου, πέφτοντας τελικά στο μαξιλάρι.

«Γιατί κλαις;» τον ρώτησα θέλοντας ταυτόχρονα να υποχωρήσω από την αγκαλιά για να την κοιτάξω στα ανοιχτόχρωμα μάτια της. 

Αυτή, όμως, δεν με άφησε να την διώξω από κοντά μου. Στην αρχή προβληματίστηκα, καθώς η αντίδραση της εκείνη τη στιγμή ήταν κυριολεκτικά η αντίθετη σε σχέση με την τελευταία φορά που την είχα συναντήσει. Μήπως μου έλεγε συγγνώμη για αυτό ακριβώς το συμβάν; Ανήμπορος να αναλύσω την κατάσταση, απλά αγκάλιασα την Ιφι και κοίταξα την μητέρα μου, η οποία μου χαμογέλασε πίσω.

Στρέφοντας το βλέμμα μου προς το παράθυρο είδα έναν ακόμη επισκέπτη στο δωμάτιο μου, ο οποίος κοιτούσε τα αμάξια που περνούσαν στη λεωφόρο κάτω από το νοσοκομείο.

«Φυτ;» είπα ακόμη πιο προβληματισμένος.

Αυτός γύρισε το βλέμμα του για να μου δείξει πώς με είχε ακούσει, αλλά σύντομα ξανακοίταξε έξω από το παράθυρο όπως και πριν, βγάζοντας μία βαριά ανάσα απογοήτευσης.

«Φυτ όλα καλά;»

Δεν μου απάντησε λεκτικός, κούνησε μόνο το κεφάλι του γνέφοντας αρνητικά με ένα απογοητευμένο ύφος στο πρόσωπο του.

«Τι έγινε Φυτ;»

Δεν μου απάντησε, ούτε καν με τη γλώσσα του σώματος. Αυτοί οι δύο είναι η κοπέλα και το ψιλό αγόρι που είχε αναφέρει ο γιατρός; Πώς γίνεται αυτό; Πώς γίνεται τη μία μέρα να μη με θυμούνται, την επόμενη να έρχονται κλαίγοντας στο δωμάτιο μου και να μου ζητάνε συγγνώμη; Κάτι δεν πάει καλά. Μήπως δεν είμαι ο Σπαρκ που γνωρίζω; Μήπως έχω περάσει κάποιο είδους μετενσάρκωση; Όχι, δεν βγάζει νόημα, αφού όλοι εδώ είναι γνωστοί μου και έχουν τις ίδιες σχέσεις με έμενα όπως πριν. Λες να βλέπω κάποιο “όνειρο" στην μετά θάνατον ζωή; Μόνο έτσι εξηγείται.

«Ω Σπαρκ...» άκουσα μία καινούργια, αλλά γνώριμη γυναικεία φωνή από την πόρτα.

«Μπλαιρ;» είπα και κοίταξα προς το μέρος της, όπου είδα την Μπλαιρ να έχει καταρρεύσει αγκαλιάζοντας στον ψηλό και χαζοχαρούμενο όπως πάντα Κουπ, ο οποίος αυτή τη φορά είχε ένα χαμόγελο ανακούφισης στο πρόσωπο του. Οι δυο τους φαίνονταν σαν να προσπαθούσαν να κρατήσουν τα δάκρυα τους και να μην κλάψουν, η Μπλαιρ όμως υπέκυψε στα συναισθήματα της και άρχισε να κλαίει πάνω στην ποδοσφαιρική φανέλα που φορούσε ο Κουπ.

«Τι... γίνεται;» ρώτησα προβληματισμένος.

«Συγγνώμη Κουπ, να περάσω;» άκουσα μία φωνή που με είχε τραυματίσει τα τελευταία χρόνια. Είχε μπει στο δωμάτιο το μόνο άτομο που δεν έπρεπε να εμφανηστεί.

«Εσύ...» εκνευρίστηκα και έσπρωξα την Ιφι στο πάτωμα, εκτινάζοντας ταυτόχρονα τους σωλήνες που είχα στα χέρια και στη μύτη μου στο πλάι για να σηκωθώ και να τον πλησιάσω.

Όταν σηκώθηκα ο πόνος που ένιωθα ήταν απερίγραπτος. Τα πόδια μου ήταν υπερβολικά αδύναμα για να με κρατήσουν όρθιο, αλλά στερεώνοντας τον εαυτό μου στο κρεβάτι και στον τοίχο μπόρεσα να τον πλησιάσω.

«Εσύ...»

«Τι έκανα;» προβληματίστηκε ο Μαλχη.

«Τι έκανες;» είπα εκνευρισμένος και του έριξα μία μπουνιά στο αριστερό του μάγουλο, ρίχνοντας τον στο γρανιτένιο πάτωμα του διαδρόμου.

«Τι έπαθες Σπαρκ; Είσαι καλ...»

Πριν προλάβει να τελειώσει την πρόταση του είχα πέσει πάνω του.

«Σταμάτα. Τι κάνεις;» είπε κρατώντας τα χέρια μου για να μην μπορέσω να τον χτυπήσω ξανά. Εκείνη τη στιγμή, τράβηξα τους καρπούς μου ψιλά, σηκώνοντας τον Μαλχη από το πάτωμα και φέρνοντας τον μπορστά στο πρόσωπό μου, όπου του έριξα μία κουτουλιά η οποία του άνοιξε το κεφάλι και τον έστειλε πάλι ανάσκελα στο πάτωμα του νοσοκομείου.

«Σπαρκ σταμάτα. Τι έχεις πάθει;» προσπάθησε να με συγκρατήσει ο Κουπ πιάνοντας το αριστερό μου χέρι, αλλά εγώ τον κλώτσησα στο στομάχι για να τον απομακρύνω και έστρεψα το βλέμμα μου πάλι στον Μαλχη. Πριν προλάβω όμως να ξαναπάω πάνω του, η μητέρα μου με είχε πιάσει από τη μέση για να με σταματήσει.

«Ήρεμα τώρα Σπαρκ, όλα θα πάνε... καλά» είπε ενώ προσπαθούσε με όλη της τη δύναμη να με τραβήξει προς τα πίσω και να με τοποθετήσει πίσω στο κρεβάτι.

«Τι εννοείς ήρεμα; Πώς μπορεί κάποιος να είναι ήρεμος σε αυτή τη κατάσταση;»

«Σπαρκ... ήρεμα» άκουσα τη φωνή της μητέρας μου πίσω από την πλάτη μου ενώ με έσερνε προς το κρεβάτι κλαίγοντας.

«Μα...»

«Δεν θέλω να ακούσω τίποτα» είπε και με πέταξε στο κρεβάτι. «Μαλχη είσαι καλά;»

«Ναι, μια χαρά είμαι»

«Και το αίμα που τρέχει από το κεφάλι σου;»

«Τρέχει αίμα από το κεφάλι μου;» είπε ο Μαλχη και έλεγξε με το δεξί του χέρι το καρούμπαλο που του είχα δημιουργήσει με την κουτουλιά μου. «Κοίτα δω, μη βρίσω το σπανάκι μου... Καλά, απ’ ό,τι φαίνεται ο Σπαρκ δεν είναι και πολύ χαρούμενος να με δει, οπότε καλύτερα να φεύγω μη φάω κι άλλη κουτουλιά...» είπε ο Μαλχη καθώς σηκωνόταν και αποχώρησε. Χωρίς να πει τίποτα, η Ιφι σηκώθηκε και αυτή από το πάτωμα, που την είχα ρίξει προηγουμένως,  και τον ακολούθησε στον διάδρομο του νοσοκομείου.

Έστρεψα το βλέμμα μου στους υπόλοιπους επισκέπτες μου, οι οποίοι με κοιτούσαν όλοι φοβισμένοι και προβληματισμένη με το σκηνικό που εξελίχθηκε πριν λίγα δευτερόλεπτα. Ο Φυτ, ίσως ο λιγότερο προβληματισμένος, αποχώρησε και αυτός από το δωμάτιο χωρίς να βγάλει άχνα.

«Που πας Φυτ;» τον ρώτησε η Μπλαιρ.

Αυτός δεν της απάντησε.

«Τι έγινε εδώ;» εμφανίστηκε από το πουθενά ο γιατρός που με είχε αναλάβει το βράδυ όταν ξύπνησα. «Γιατί με κοιτάτε όλοι σας περίεργα; Θα μου πει κάποιος τι έγινε εδώ; Γιατί έχει αίμα το πάτωμα; Γιατί ο ασθενής έχει γρατζουνιές στα χέρια του;» εκνευρίστηκε ο γιατρός επειδή δεν του απαντούσε κανείς.

«Να, βλέπετε...» ξεκίνησε η μητέρα μου.

«Ο ασθενής σας τραυμάτισε έναν φίλο μας στο κεφάλι» άκουσα την φωνή του Φυτ από τον διάδρομο.

«Τι λες τώρα. Άρα περπατάει;»

«Απ’ ό,τι φαίνεται...» άκουσα τη φωνή του Φυτ να απομακρύνεται.

«Γιατί φεύγεις; Δεν θα μου εξηγήσεις τι έγινε;»

Ο Φυτ δεν του απάντησε.

«Καλά, δεν πειράζει...» είπε απογοητευμένος ο γιατρός. «Εγώ όμως κυρία Θοτς σας είπα πως είναι νωρίς για επισκέψεις, αλλά εσείς δεν με ακούσατε...»

«Ποιος θα μπορούσε να προβλέψει ότι κάτι τέτοιο θα γινόταν;» του απάντησε η μητέρα μου.

«Κυρία μου, ο άνθρωπος είναι έντεκα μήνες και κάτι σε κώμα, πώς νομίζετε πως θα αντιδρούσε;» την ρώτησε εκνευρισμένος ο γιατρός.

«Όχι έτσι πάντως...» του απάντησε η μητέρα μου και σταύρωσε τα χέρια της.

«Όχι άλλες επισκέψεις, εκτός της άμεσης οικογένειας του ασθενή» επισήμανε ο γιατρός και κοίταξε την Μπλαιρ και τον Κουπ. «Συγγνώμη, αλλά δεν μπορούμε να ρισκάρουμε την υγεία του ασθενή»

«Μα γιατί;» προβληματίστηκε η Μπλαιρ.

«Νομίζω πώς είναι προφανής ο λόγος» της απάντησε ο γιατρός εκνευρισμένος.

«Μα...»

«Εντάξει γιατρέ, δεν πειράζει. Καλή συνέχεια...» την διέκοψε ο Κουπ και αποχώρησαν από το δωμάτιο.

«Σας ευχαριστώ για την κατανόηση... Λοιπόν, οι γρατζουνιές στα χέρια σου μάλλον είναι από τους ορούς. Θα φωνάξω πάλι τις νοσοκόμες να σου ξαναφορέσουν τους σωλήνες και τα καλώδια»

«Εγώ τι να κάνω;» τον ρώτησε η μητέρα μου.

«Μπορείτε να καθίσετε μαζί με τον γιο σας. Τώρα που είναι ξύπνιος υποθέτω θα θέλει και λίγη συντροφιά»

«Έχετε δίκιο... Α και κάτι ακόμη. Πόσο πιστεύετε πως θα διαρκέσει η διαδικασία της ανάρρωσης του εδώ στο νοσοκομείο; Θέλω να ξέρω για να πάρω τις κατάλληλες άδειες»

«Αφού ο ασθενής απ’ ό,τι κατάλαβα περπατάει...»

«Καλά, δεν περπατάει...» τον διέκοψε η μητέρα μου «Στερεονόταν από τους τοίχους και περπατούσε λίγο πλάγια σαν ζόμπι, οπότε δεν νομίζω πως είναι έτοιμος να περπατήσει σε πεζοδρόμιο ακόμη»

«Α, εντάξει κατάλαβα, τότε η ανάρρωση του εδώ στο νοσοκομείο θα του πάρει γύρω στους πέντε μήνες, ίσως λίγο λιγότερο, όμως πρέπει να συνεχίσετε να τον φροντίζετε μέχρι να ενταχθεί πλήρως ξανά στην κοινωνία»

«Αυτό είναι αυτονόητο γιατρέ...» του απάντησε η μητέρα μου με ένα ψεύτικο χαμόγελο στα χείλη της.

«Μπορώ να ρωτήσω και εγώ κάτι;» είπα διστακτικά.

«Παρακαλώ, πείτε μου»

«Πόσο... έχουμε σήμερα;»

«8 Αυγούστου. Γιατί ρωτάς;»

«Τίποτα, απλά ήθελα να μάθω...»

\chapter{}

«Είσαι έτοιμος Σπαρκ;»

«Έτσι νομίζω»

«Ωραία, πάμε σπίτι. Καλή συνέχεια γιατρέ, ευχαριστούμε για όλα»

Ο γιατρός δεν της απάντησε λεκτικά, άφησε μόνο ένα χαμόγελο προς το μέρος μας και μας άνοιξε την πόρτα του δωματίου.

«Πώς πέρασε έτσι ο καιρός ρε Σπαρκ. Θυμάμαι σαν να ήταν χθες το βράδυ που ξύπνησες και τώρα μπορείς και περπατάς χωρίς βοήθειες. Είναι πραγματικά θαύμα»

«Όντως αυτοί οι πέντε μήνες πέρασαν πολύ γρήγορα...»

«Εντάξει δεν ήταν και πέντε μήνες... ίσως κάτι λιγότερο. Ο χρόνος όμως κυλάει πιο γρήγορα όταν έχεις καλή παρέα, ε; Δεν συμφωνείς;»

«Ναι...». Τι εννοεί καλή παρέα; Τις περισσότερες μέρες ήμασταν απλά στο δωμάτιο του νοσοκομείου και κοιτούσαμε τους τοίχους που ξεφλουδίζουν. Ακόμη δεν έχω καταλάβει πως έχω καταλήξει σε αυτή την κατάσταση και με προβληματίζει. Γιατί ζω ακόμη; Γιατί είναι όλοι με το μέρος μου πλέον; Γιατί ο Μαλχη άλλαξε; Πολλές αναπάντητες ερωτήσεις που η μητέρα μου δύσκολα θα ξέρει την απάντηση. Καλύτερα να κρατήσω αυτές τις σκέψεις για τον εαυτό μου και να ερευνήσω την υπόθεση μόνος μου.

«Ωχ, κοίτα. Οι φίλοι σου Σπαρκ»

«Οι φί...λοι μου;» έστρεψα το βλέμμα μου στην πόρτα της εισόδου του νοσοκομείου, όπου είδα μία γνώριμη κοτσίδα δίπλα σε ένα ψηλό άντρα που κρατούσε μια ανθοδέσμη να συζητάνε με μία νοσοκόμα.

«Ωχ να τος. Σας ευχαριστούμε πολύ, τον βρήκαμε. Σπαρκ έλα εδώ»

«Για..τί;»

«Για να σου δώσουμε αυτά τα λουλούδια» είπε η Μπλαιρ και μου έδωσε τα λουλούδια που κρατούσε ο Κουπ.

«Τι;»

«Να, βλέπεις, αφού σήμερα βγήκες από το νοσοκομείο, είπαμε να σου κάνουμε ένα...»

«Γιατί;»

«Τι εννοείς γιατί;»

«...» έκανα μία παύση για να σκεφτώ και αποφάσισα να μην αρχίσω συζήτηση στη μέση του νοσοκομείου. «Δεν θέλω λουλούδια»

«Γιατί;»

«Τι εννοείς γιατί;» την πλησίασα εκνευρισμένος και αυτή έκανε ένα βήμα πίσω.

«Άστο Μπλαιρ, ίσως θέλει λίγο χρόνο να σκεφτεί, όπως πάντα...» τράβηξε ο Κουπ προς το μέρος του τη Μπλαιρ.

«Καλά, αλλά θυμάσαι ότι σας είχα πει πως έχω κάτι σημαντικό να σας ανακοινώσω, έτσι;» με ρώτησε η Μπλαιρ. «Είσαι ο μόνος που δεν το ξέρει πλέον, αφού δεν ήσουν κοντά μας τόσο καιρό. Ελπίζω να μην αποχωρήσεις τελείως αυτή τη φορά...»

«Δεν... καταλαβαίνω...»

«Άστο Μπλαιρ, είπαμε θέλει τον χρόνο του»

«Μα...» έκανε μία μικρή παύση η Μπλαιρ για να με κοιτάξει. «Καλά... αλλά μέσα στη βδομάδα θα ξαναμιλήσουμε Σπαρκ, εντάξει;»

«Δεν... υπόσχομαι τίποτα...»

«Πω, σταμάτα να είσαι γκρίζος όλη την ώρα. Ναι ή όχι;»

Δεν της απάντησα.

«Καλά, το εκλαμβάνω σαν “Ναι”. Πάμε Κουπ, οδηγάς εσύ γιατί εκνευρίστηκα»

«Καλή συνέχεια Σπαρκ και κυρία Θοτς» μας αποχαιρέτησε ο Κουπ.

«Καλή συνέχεια παιδιά! Έλα Σπαρκ, πάμε και εμείς στο αυτοκίνητο» είπε η μητέρα μου και με έπιασε από το χέρι σαν να ήμουν μικρό παιδί.\\\\

«Ωραία, κάτσε να ξεκλειδώσω και θα σου πω» είπε η μητέρα μου και σήκωσε το πατάκι για να πάρει το κλειδί. 

Το σπίτι είναι ακριβώς όπως το θυμάμαι, άσπρο στο εξωτερικό, μονώροφο, με μία μικρή αυλή, το κλειδί ήταν κάτω από το πατάκι. Αυτές οι μικρές λεπτομέρειες μου δείχνουν πώς ίσως είμαι ακόμη ζωντανός και ίσως δεν έχω περάσει στη μετά θάνατον ζωή. Πώς γίνεται όμως αυτό; Ακόμη δεν μπορώ να καταλάβω, λες η μητέρα μου να ξέρει;

«Λοιπόν Σπαρκ, μπες μέσα και βγάλε τα παπούτσια σου. Αν και ο γιατρός μου είπε να μην σε αφήσω μόνο σου για μερικούς μήνες, πάω μία γρήγορη να πάρω κάτι πράγματα από το σούπερ μάρκετ τώρα που είναι ακόμη ανοιχτό και θα ‘ρθω να φτιάξω φαγητό» μου είπε η μητέρα μου και έκλεισε την πόρτα πίσω μου καθώς έμπαινα μέσα στο σπίτι. «Ελπίζω να μην κάνεις τίποτα όσο λείπω...» την άκουσα να ψιθυρίζει όσο κλείδωνε την πόρτα του σπιτιού.

Ακόμη και στο εσωτερικό, το σπίτι μας ήταν ολόιδιο με τις αναμνήσεις που είχα. Έβγαλα τα παπούτσια μου και προχώρησα προς το σαλόνι. Είδα τον καναπέ και το τραπεζάκι όπου είχα προσπαθήσει να πάρω τη ζωή μου. Ο τοίχος απέναντι από τον καναπέ είχε ένα σημάδι. Λες να είναι από τότε που πέταξα το κινητό μου; Πήγα προς την κουζίνα και ήταν όπως ακριβώς την θυμόμουν, απλά ο νεροχύτης ήταν γεμάτος. Η κυκλική τραπεζαρία όπου καθόμουν και διάβαζα είχε τις τρεις, όπως πάντα, καρέκλες γύρω της οι οποίες μου θύμισαν αναμνήσεις του πατέρα μου να μου μαγειρεύει. Δεν ήταν και ο καλύτερος μάγειρας, αλλά τουλάχιστον είχα φαγητό στο πιάτο κάθε μέρα, νομίζω το έχω ξαναπεί αυτό... Τι λέω, αφού είμαι μόνος μου με τις σκέψεις μου, και να το έχω ξαναπεί δεν έχει σημασία, δεν θα με κρίνει κάποιος που είπα το ίδιο πράγμα δύο φορές.

Περπάτησα προς το δωμάτιο μου για να δω αν ακόμη και αυτό ήταν όπως το θυμόμουν. Σταμάτησα στα μέσα της διαδρομής για να κοιτάξω τον εαυτό μου στον καθρέφτη. Η αντανάκλαση μου δεν μου προκάλεσε και τόση περιέργεια. Φαινόμουν όπως ακριβώς με θυμάμαι, μουσάτος, λεπτός, με ένα βλέμμα απογοήτευσης. Ίσως είχα λιγότερα σημάδια ηλικίας, αλλά μπορεί να ήταν και η ιδέα μου.

Μπαίνοντας στο δωμάτιο μου, άνοιξα κατευθείαν το κομοδίνο μου και...

«Προφανώς και είναι εδώ... γιατί όμως;» είπα και σήκωσα το πιστόλι. «Ωχ» είδα κάτι που με προβλημάτισε κάτω από το πιστόλι. «Πώς γίνεται αυτό;» είπα και έπιασα με το ελεύθερο μου χέρι το κινητό μου που βρισκόταν κάτω από	πιστόλι. «Αφού το είχα... δουλεύει κιόλας. Κάτι πάει λάθος...»

Παρά τον προβληματισμό μου, έβαλα το κινητό στην τσέπη μου και συνέχισα να ψάχνω για να βρω τον γεμιστήρα του όπλου μου.

«Δεν είναι εδώ, λες... ναι ήταν ήδη γεμάτο. Γιατί δεν τσέκαρα αυτό πρώτα;... είμαι όντως άχρηστος» είπα και ξαναέβαλα τον γεμάτο σφαίρες γεμιστήρα πίσω στο πιστόλι. «Αλλά, ποιο είναι το κίνητρο μου τώρα. Δεν έχω κάποιο κίνητρο. Από που να ξεκινήσω;»

Επέστρεψα στο σαλόνι και κάθησα στον καναπέ, αφήνοντας το πιστόλι μου στο τραπεζάκι και στρέφοντας το βλέμμα μου στο σημάδι στον τοίχο. Κοίταξα άλλη μία φορά το κινητό μου.

«Πώς... δεν καταλαβαίνω. Μία ζωή αυτό λέω. Δεν έχω και άδικο όμως, αφού έχω ζήσει μία ζωή γεμάτη ανεξήγητες προδοσίες, ή τουλάχιστον... έτσι νιώθω. Γιατί όμως είναι όλοι τώρα με το μέρος μου; Όταν λέω όλοι, το εννοώ όλοι. Πώς γίνεται ο Μαλχη να είναι με το μέρος μου μετά από όσα έχει κάνει; Τι σημαίνει καν αυτό που λέω; “Είναι με το μέρος μου”, ή όπως το λέει η μητέρα μου “είναι φίλοι μου”. Κανείς τους δεν είναι πραγματικός φίλος μου. Ο μόνος φίλος που είχα πραγματικά ήταν ο Βασι, ο Βασι... ποιο είναι το επίθετο του πάλι; Πώς γίνεται να μην θυμάμαι το επίθετο ενός από των πιο σημαντικών ανθρώπων στη ζωή μου; Είμαι όντως άχρηστος»

Έπιασα το πιστόλι μου για να το επεξεργαστώ.

«Και ξέρεις γιατί είμαι όντως άχρηστος; Δεν διάβασα τις σκέψεις καθενός στο νοσοκομείο. Γιατί; Πώς δεν το σκέφτηκα; Μήπως φταίει αυτό που είπα πριν, το γεγονός ότι δεν έχω κάποιο κίνητρο. Ποιος ο λόγος να ζω καν; Τόσο καιρό είμαι ένα βάρος για τους ανθρώπους με τους οποίους συναναστρέφομαι. Ακόμη και όταν πήγα μόνος μου για την τελευταία μου αποστολή, κατέληξε άλλος ένας γνωστός μου να χάνει τη ζωή του. Χωρίς λόγω κιόλας, αφού μετά από λίγα λεπτά πέθανα και γω... ή μάλλον, δεν πέθανα, αλλά...» ξέσπασα σε δάκρυα. Εκείνη τη στιγμή είδα μία μύγα να προσγειώνεται στο τραπεζάκι.

«Τι είναι καν η ζωή» είπα, σκοτώνοντας τη μύγα με την παλάμη μου. «Αν είναι τόσο εύκολο να την πάρεις από κάποιον, γιατί να τόσο πολύτιμη; Ίσως αυτό την κάνει πολύτιμη, το γεγονός ότι κάποια στιγμή μπορεί να πεθάνεις χωρίς να το περιμένεις» είπα και σκούπισα τα δάκρυα μου. «Τουλάχιστον όμως τα έντομα έχουν κάποιο λόγο ύπαρξης, είναι τροφή για τα ζώα πάνω τους στην τροφική αλυσίδα. Ο άνθρωπος τι λόγο ύπαρξης έχει; Τι προσφέρω εγώ στο ζωικό βασίλειο που ονομάζουμε κοινωνία; Αν προσφέρω όντως κάτι, γιατί ο Θεός είναι τόσο κακώς μαζί μου; Αν η ζωή μου ήταν τόσο πολύτιμη, γιατί να μην μπορέσω να νιώσω κάποια ευχαρίστηση αυτά τα χρόνια; Τι είναι η ευχαρίστηση καν; Έχω τόσο πολύ καιρό να νιώσω αυτό το συναίσθημα που έχω ξεχάσει πως είναι, από που προέρχεται, πώς το αποκτώ. Ίσως φταίει που δεν έχω φίλους, ίσως... εγώ είμαι το πρόβλημα...» κοίταξα την κάννη του πιστολιού μου.

«Κάνει κρύο σήμερα πατέρα. Δεν ξέρω αν με ακούς ή που βρίσκεσαι. Δεν ξέρω καν που βρίσκομαι εγώ πλέον, δεν ξέρω καν ποιος είμαι πλέον. Μικρός ήμουν γεμάτος χαρά, ήμουν μία ήσυχη ψυχή που περιπλανιόταν στην πόλη χωρίς να σκέφτομαι. Όταν απέκτησα συνείδηση και άρχισα να αναρωτιέμαι πράγματα για την ζωή, κατάλαβα τους προβληματισμούς σου. Δεν μπόρεσα να τους λύσω όμως. Συνέβησαν πολλά, ή τουλάχιστον μου φαίνεται πως συνέβησαν πολλά. Μπορεί από τις οπτικές των άλλων να μην έχει συμβεί τίποτα και απλά να κλαίω μόνος μου, αλλά πατέρα...» έβαλα το πιστόλι στο κρανίο μου για μία τελευταία, ελπίζω, φορά «...ελπίζω να με ακούς πατέρα»

\chapter{}

*Ντριν ντριν... ντριν ντριν... ντριν*

«Πω... Σαν να μην θέλει το σύμπαν να τελειώσει αυτή η ιστορία. Παρακαλώ;»

«Έλα Σπαρκ. Συγγνώμη για ό,τι έγινε το πρωί. Είναι καλή ώρα να μιλήσουμε ή κάνεις κάτι;»

Κατέβασα το όπλο από το κεφάλι μου.

«Σπαρκ;»

Τι να της απαντήσω; Δεν έχω καμία όρεξη να μιλήσω μαζί της αυτή τη στιγμή.

«Σπαρκ;»

Και να μιλήσουμε τι θα γίνει; Θα καταλήξει να είναι άλλη μία συζήτηση χωρίς νόημα. Δεν έχω πλέον κίνητρο για να ζήσω, δεν μπορείτε απλά να... με αφήσετε ήσυχο;

Σαν να με άκουσε, η Μπλαιρ μου έκλεισε το τηλέφωνο στα μούτρα χωρίς να πει κάτι παραπάνω. Εκείνη τη στιγμή άκουσα την πόρτα να ξεκλειδώνει. ‘Εκρυψα γρήγορα το όπλο κάτω από τα μαξιλάρια του καναπέ για να μην υποψιαστεί κανείς τίποτα.

«Ξέρεις, δεν είναι και τόσο καλή ιδέα να αφήνετε το κλειδί κάτω από το πατάκι σας» μου είπε η Μπλαιρ. «Μισό να βγάλει τα παπούτσια του και ο Κουπ και... Τι κάνεις εκεί Σπαρκ;»

«Τίποτα... σίγουρα δεν κρύβω κάτι, χαχα...»

«Συγγνώμη που μπήκαμε τόσο ξαφνικά, αλλά δεν μπορούσαμε να κάνουμε κάτι άλλο, αφού στο τηλέφωνο δεν απαντούσες και...»

«Μπορούσατε να... με αφήσετε ήσυχο»

«Γιατί τόση κατάθλιψη ξαφνικά; Είσαι πλέον μαζί μας, τι άλλο θες;»

«Δίκιο έχει Σπαρκ» πετάχτηκε ο Κουπ από το πουθενά και κάθισε στο τραπεζάκι.

«Και αυτό το περιστατικό όταν ξύπνησες; Γιατί τόσα νεύρα προς εμάς και ιδιαίτερα προς τον Μαλχη; Όλα καλά;» με ρώτησε η Μπλαιρ καθώς έπαιρνε θέση δίπλα στον Κουπ στο τραπεζάκι μπροστά στον καναπέ.

«Πού...» σταμάτησα στη μέση της πρότασης μου, επειδή πίστευα πώς θα έλεγα κάτι χαζό.

«”Που” τι; Πες το και ας είναι χαζομάρα» μου απάντησε η Μπλαιρ σαν να διάβαζε τις σκέψεις μου.

«Πού... είμαι;»

«Είσαι σόος και υγιής δίπλα στους καλύτερους σου φίλους» μου απάντησε η Μπλαιρ χωρίς δεύτερη σκέψη.

«Καλύτεροι μου φίλοι;»

«Ναι, είτε το πιστεύεις είτε όχι.»

«Πώς, μα...»

«Τι έγινε Σπαρκ;» η Μπλαιρ ακούμπησε το χέρι της στον ώμο μου συμπονετικά. 

«Ποιος... είμαι; Τι κάνω εδώ;»

Ο Κουπ και η Μπλαιρ κοιτάχτηκαν για μία στιγμή προβληματισμένοι από τις ερωτήσεις μου.

«Είσαι ο πιο τυχερός άνθρωπος του κόσμου, είτε το πιστεύεις, είτε όχι»μου απάντησε η Μπλαιρ. «Το προηγούμενο καλοκαίρι, πριν έναν χρόνο και κάτι μήνες δηλαδή, είχατε οργανώσει να πάτε ταξίδι με τον πατέρα, τον αδελφό, τον ξάδελφο και... τον θείο σου. Δυστυχώς, όποιος οδηγούσε έκανε μία μικρή πατάτα και το αμάξι σας εξεράγει. Ευτυχώς για σας, ο Κουπ σας είχε πάρει τηλέφωνο εκείνη τη στιγμή καθώς οι γονείς του είχαν μόλις χτυπηθεί από ένα φορτηγό εκείνο το πρωί και ήθελε να σας το πει»

«Δεν πιστεύεις ότι του τα λες λίγο μαζεμένα;» την ρώτησε ο Κουπ.

«Τι εννοείς;»

«Ίσως να του λέγαμε τα γεγονότα πιο... διάσπαρτα. Μην πάθει καμία κρίση και...»

«Δεν πειράζει. Εμένα αυτό είναι το στυλ μου, τι θες τώρα;»

«Δεν καταλαβαίνω...» τους διέκοψα. «Και πώς...»

«Πώς επέζησες;» με διέκοψε αυτή τη φορά η Μπλαιρ. «Όταν ο Κουπ άκουσε την έκρηξη από την κλήση πήρε κατευθείαν την πυροσβεστική και τους είπε που ήσασταν. Ευτυχώς, σας βρήκαν μέσα στη φωτιά, αλλά δυστυχώς... ήσουν ο μόνος που είχε παλμό. Οι άλλοι τέσσερις είχαν πεθάνει»

«Δυστυχώς;» προβληματίστηκα με τον τρόπο που διάλεξε να εκφράσει την πρόταση της.

«Όχι δεν το εννοώ με κακή πρόθεση, απλά... δεν πιστεύεις ότι είναι κρίμα που πέθαναν τόσα μέλη της κοντινής σου κιόλας οικογένειας μέσα σε...»

«Δεν καταλαβαίνω... αφού ο πατέρας μου είχε πεθάνει πολύ καιρό τώρα. Πώς γίνεται...»

«Τι λες μωρέ;» με διέκοψε πάλι η Μπλαιρ. «Αφού μας έλεγες πριν πάτε διακοπές πώς γίνεται... Αχ, τώρα με μπέρδεψες και εμένα...»

«Κάτσε λίγο...» μας τράβηξε την προσοχή ο Κουπ «...δεν θυμάσαι τον πατέρα σου;»

«Όχι τον θυμάμαι...» του απάντησα «...αλλά έχω μόνο αναμνήσεις από όταν ήμουν πολύ μικρός, αφού σας είπα πέθανε όταν ήμουν οκτώ»

«Χμ... για πες μας, τι άλλο θυμάσαι; Πες τα μας όλα από την αρχή»

«Όλα;»

Μου έγνεψε καταφατικά με το κεφάλι του ο Κουπ.

«Τι άλλο θυμάμαι, ε; Θυμάμαι... μερικές μέρες από τη σχολική μου ζωή οι οποίες είναι κοντά χρονολογικά με τον θάνατο του πατέρα μου. Μετά... θυμάμαι να σας γνωρίζω την πρώτη μέρα που μπήκαμε με τις αξιολογίσεις μας στην Ειδική Υπηρεσία. Τι άλλο... τις αποστολές μου θυμάμαι, σε μερικές ήμασταν όλοι, σε μερικές έλειπες εσύ Κουπ γιατί δεν ήσουν ακόμη προφανώς μέλος της ομάδας...»

Οι δυο τους κοιτάχτηκαν προβληματισμένοι.

«Γιατί κοιταζόσαστε έτσι;»

«Τίποτα συνέχισε» μου απάντησε ο Κουπ.

«Ε μετά θυμάμαι τον Μαλχη να με προδίδει και να με βάζει φυλακή για 100 χρόνια, αλλά εσείς και ο Βασι πληρώσατε για να με...»

«Όπα όπα. Πότε πήγες φυλακή;»

«Δεν νομίζω πως είναι αυτό το πρόβλημα Μπλαιρ. Σπαρκ, θυμάσε τίποτα άλλο... περίεργο;»

«Τι άλλο μπορεί να είναι περίεργο; Ότι πέθανα και ξαναγεννήθηκα ως ο ίδιος άνθρωπος; Ότι η μαγεία μου στις τελευταίες μου ώρες ξαφνικά έγινε πιο ισχυρή;»

«Μαγεία;» είπαν και οι δυο τους ταυτόχρονα.

«Ναι... τι;»

«Τι εννοείς μαγεία;»

«Νόμιζα πως είναι αυτονόητο, αφού... έχετε και εσείς μαγεία. Όχι;»

Κοιτάχτηκαν για πολλοστή φορά μεταξύ τους προβληματισμένοι.

«Κοίτα Σπαρκ, είναι προφανές ότι είσαι θύμα κάποιου είδους απώλειας μνήμης. Δεν είμαι γιατρός, αλλά όσο ήσουν σε κώμα, ίσως έβλεπες κάποιου είδους παραισθήσεις σαν να βρισκόσουν σε όνειρο» μου είπε ο Κουπ. Αλλά από περιέργεια, για πες μας τι ακριβώς ήταν αυτές οι αποστολές και πώς ακριβώς σε πρόδωσε ο Μαλχη»

«Δεν... νομίζω πώς θέλετε να τα ακούσετε...»

«Θέλουμε» με διέκοψε πάλι η Μπλαιρ. «Εξάλλου, για αυτό είμαστε εδώ»

«Καλά...»\\\\

«Ωραίος εφιάλτης πάντως» σχολίασε η Μπλαιρ. «Πολύ δημιουργικός θα έλεγε κάποιος. Αλλά δεν μου είπες, τι μαγεία είχα εγώ;»

«Δεν είναι ώρα για τέτοια Μπλαιρ. Η ιστορία σου πάντως μοιάζει πάρα πολύ με την πραγματικότητα» είπε ο Κουπ και σηκώθηκε από το τραπεζάκι.

«Τι εννοείς;» ρώτησα.

«Καλά δεν είναι και ένα προς ένα αντιστοιχία, αλλά θυμάμαι τις πρώτες μέρες στο γυμνάσιο να σας προσεγγίζω εσένα και τον Βασι και να σας ζητάω να γίνουμε φίλοι. Στη συνέχεια να γνωρίζουμε την Μπλαιρ και τον Μαλχη, λίγο αργότερα τον Φυτ και ακόμη πιο μετά την Ιφι. Αλλά δεν είναι αυτό που με...»

«Άρα είμαστε όντως... φίλοι;»

«Ναι, γιατί να μην είμαστε; Ο Βασι βέβαια ήταν αδελφός σου, για αυτό...»

«ΤΙ; Ο Βασι είναι αδελφός μου;» ξαφνιάστηκα.

«Αυτό σε προβλημάτισε;» σχολίασε η Μπλαιρ

«Μη θεωρείς αυτά που ξέρουμε εμείς αυτονόητα Μπλαιρ, αφού δεν θυμάται τίποτα ο Σπαρκ»

«Θυμάται τα κοντινά του πρόσωπα...»

«Ναι, ίσως υποσυνείδητα, αλλά δεν φαίνεται να θυμάται τις σχέσεις του με αυτά τα πρόσωπα»

«Καλά... συνεχίστε»

«Από που να συνεχίσω; Μου έκοψες τη ροή τώρα...»

«Να ρωτήσω κάτι;» τους διέκοψα.

«Για πες Σπαρκ»

«Άρα εμείς είμαστε ίδια ηλικία;»

«Αμέ, όλοι 21 είμαστε. Ακόμη και ο Μαλχη που το μουστάκι του τον κάνει να φαίνεται δεκαπέντε χρόνια μεγαλύτερος είναι 21. Ο Βασι ήταν ο μόνος στην παρέα που ήταν δύο χρόνια μεγαλύτερος»

«Πώς... τέλος πάντων»

«Τι έγινε Σπαρκ;»

«Θυμάμαι μερικά ακόμη πρόσωπα. Μπορώ να σας ρωτήσω αν επιτρέπεται;»

«Ρώτα Σπαρκ, μη ντρέπεσαι»

«Ο Χερμπ και ο Ρομπ... τι σχέση έχω με αυτούς;»

«Αυτό θα το απαντήσω εγώ» είπε η Μπλαιρ πριν προλάβει να μιλήσει ο Κουπ. «Αυτοί οι δύο είναι...»

«Ο πατέρας σου και ο αδελφός σου;»

«Πώς;» προβληματίστηκε η Μπλαιρ με την ερώτηση μου. «Πώς το ήξερες;»

«Να, στις αναμνήσεις μου είχες μία παραπάνω έλξη προς τον Χερμπ και τον Ρομπ, ειδικά όταν... πέθαναν»

«Α... Ακραία σύμπτωση πάντως, γιατί αυτοί οι δύο είναι όντως ο πατέρας και ο αδελφός μου. Δηλαδή... αυτό ήθελα να σας πω πριν φύγετε για το ταξίδι αλλά δεν πρόλαβα»

«Τι εννοείς;»

«Να... ο θείος σου ο Χερμπ δεν είναι ο πατέρας μου, με την έννοια ότι...»

«Κάτσε τι;» την διέκοψα. «Ο Χερμπ είναι θείος μου;»

«Ναι βρε, αυτό δεν σου είπα και πριν. Αφού είπαμε είχατε πάει ταξίδι ο πατέρας σου, ο κύριος Ριντ δηλαδή, ο Βασι, ο θείος σου ο Χερμπ και ο Ρομπ. Μη μου πεις ότι δεν θυμάσαι ούτε ότι ο Ρομπ είναι ξάδελφος σου;»

«Ο Ρομπ είναι ξάδελφος μου;»

«Μπλαιρ, σου είπα ότι τα αυτονόητα για εμάς δεν είναι αυτονόητα για αυτόν...» της είπε ο Κουπ.

«Μπορείτε να σταματήσετε να με διακόπτετε; Εκεί που ήθελα να καταλήξω είναι ότι όταν ήμουν μωρό, οι πραγματικοί μου γονείς, οι οποίοι στην πραγματικότητα δεν ξέρω ποιοι είναι, δηλαδή... ξέρω, αλλά είναι περίπλοκη κατάσταση και δεν θέλω να μπω σε παραπάνω λεπτομέρειες. Δεν ήθελαν, τέλος πάντων, να με κρατήσουν και με θεώρησαν ως ένα λάθος που δεν έπρεπε να γεννηθεί. Ο Χερμπ όμως με υιοθέτησε και με μεγάλωσε ως δική του κόρη, παρόλες τις εμπλοκές που είχε περάσει εκείνη τη περίοδο. Δεν σου το είχα πει τόσο καιρό γιατί δεν έτυχε ποτέ να μιλήσουμε για τα οικογενειακά μου, αλλά ελπίζω να μην σε ξάφνιασα πολύ...»

«Ωχ...» μου ήρθε μία ανάμνηση με τον Χερμπ να μου εξηγεί ακριβώς αυτό που έλεγε η Μπλαιρ, αλλά δεν μπορούσα να διακρίνω τον χρόνο και το μέρος που η ανάμνηση αυτή λάμβανε χώρα.

«Τι έγινε Σπαρκ;»

«Τίποτα, απλά... έχω αρχίσει να συνδέω τα κομμάτια σιγά σιγά... Λάθος, ε;» είπα και κοίταξα κάτω από το μαξιλάρι. «Κάτσε όμως... θυμάμαι πώς, εκτός από τα κοντινά μέλη της οικογένειας μου πέθαναν και άλλοι δυό στις “αναμνήσεις” μου»

«Ποιος άλλος;» με ρώτησε ο Κουπ.

«Η... Ιφι και ο Φυτ. Αλλά αυτοί οι δύο για κάποιο λόγο μετά από κάποια χρόνια ήταν... ζωντανοί;»

«Χμ... ίσως αυτό που σκεφτόμουν πριν σχετίζεται με αυτό που λες» ο Κουπ άρχισε να κάνει κύκλους στο σαλόνι.

«Τι σκεφτόσουν πριν;» τον ρώτησε η Μπλαιρ.

«Ότι ο Σπαρκ όσο ήταν σε κώμα αυτούς τους έντεκα μήνες μας άκουγε»

«Α, θα μπορούσε»

«Τι εννοείται, δεν... κατάλαβα. Πώς προέκυψε αυτό το συμπέρασμα;»

«Να, γιατί... σου αναφέραμε ότι η Ιφι είναι η κοπέλα σου, έτσι;»

«Η Ιφι ήταν η κοπέλα μου;» φώναξα ξαφνιασμένος.

«Α, άρα δεν στο είχαμε αναφέρει...»

«Όχι, θα το θυμόμουν»

«Τέλος πάντων, δεν ήθελα να πω αυτό. Ήθελα να σου πω πως...» κοιτάχτηκαν για λίγο οι δυο τους πριν συνεχίσει την πρόταση του ο Κουπ.

Η Μπλαιρ του έγνεψε καταφατικά.

«...πως όταν έμαθε τα νέα για σένα η Ιφι... δεν ήθελε να είναι πια μαζί σου, επειδή σε θεωρούσε ήδη νεκρό. Την πρώτη μέρα κιόλας σκεφτόταν να προσεγγίσει ερωτικά τον Φυτ, πράγμα το οποίο έκανε, αλλά η “επίσημη σχέση” τους δεν κράτησε πάνω από έναν μήνα. Μερικούς μήνες μετά... η Ιφι έκανε το ίδιο με τον Μαλχη»

«Και πώς συνδέονται όλα αυτά με το γεγονός ότι μπορεί να σας άκουγα;» ρώτησα προβληματισμένος.

«Κάτσε θα φτάστουμε και εκεί. Τη μέρα που μας ανακοίνωσαν η Ιφι και ο Μαλχη ότι ήταν πλέον μαζί, είχαμε μία μικρή διαφωνία, ας το πούμε, στο δωμάτιο σου όσον αφορά το σκεπτικό της Ιφι και το πόσο εύκολα την δέχτηκαν ο Φυτ και ο Μαλχη»

«Και σα να μην έφτανε αυτό...» πετάχτηκε η Μπλαιρ. «...η Ιφι, ενώ τα είχε με τον Μαλχη, έκανε πάλι κάτι κρυφά με τον Φυτ»

«Όντως;» γέλασε ο Κουπ. «Πάλι καλά τότε που δεν τους κάνουμε πλέον παρέα. Ας τους αφήσουμε να τα λύσουν μόνοι τους. Εκεί που ήθελα να καταλήξω είναι ότι ο θάνατος της Ιφι και του Φυτ και το μίσος σου προς τον Μαλχη στο όνειρο σου μπορεί να συνέβισαν συμβολικά επειδή μας άκουγες... ή μπορεί να είναι μία σύμπτωση, ποιος ξέρει»

«Τι;» μπλέχτηκαν τα κομμάτια του παζλ. «Πάμε λίγο πίσω. Δεν μου είπατε πως είμαστε όλοι μία παρέα;»

«Ναι...» μου απάντησε η Μπλαιρ «...ίσως να μην θέλεις να τους ξανα πλησιάσεις. Μετά ειδικά από όλα όσα έκαναν πίσω από την πλάτη σου και την πλάτη μας»

«Γιατί τι άλλο έκαναν;»

«Όχι δεν έκαναν κάτι άλλο, απλά λέω πως αυτά είναι ήδη πολλά. Αν ήμουν στη θέση σου ίσως αυτοκτονούσα κιόλας, χαχα»

«Θα αυτοκτονούσες, ε;» είπα και κοίταξα το μαξιλάρι που είχε από κάτω το πιστόλι μου.

«Όχι όχι, για πλάκα το είπα. Μη μου πεις ότι το πίστεψες...» μου είπε κουνώντας τα χέρια της μπροστά μου. «Γιατί κοιτάς όμως συνέχεια εκείνο το μαξιλάρι δίπλα σου;»

«Όχι μη!...»

«Σπαρκ... τι είναι αυτό;» με κοίταξε φοβισμένη, πιάνοντας το πιστόλι μόνο με δύο δάχτυλα από την άκρη και σηκώνοντας το για να το δει κι ο Κουπ.

«Ένα... πιστόλι;» 

«Όχι χαζούλη, εννοώ γιατί το έχεις εδώ; Μας είχες πει ότι συνήθως το κρύβεις στο κομοδίνο σου. Μη μου πεις ότι...»

«Όχι δεν... βασικά... είναι ακριβώς αυτό που σκέφτεσαι»

«Δεν καταλαβαίνω» πετάχτηκε ο Κουπ. «Καταρχήν γιατί έχει πιστόλι στο σπίτι του ο Σπαρκ; Είπαμε, είσαι αστυνομικός, αλλά δεν είμαι σίγουρος αν επιτρέπεται να έχει όπλο στο σπίτι του»

«Αυτό σε πείραξε;» είπε εκνευρισμένη η Μπλαιρ και πέταξε το πιστόλι στο πάτωμα. «Εδώ ο άνθρωπος πήγε να αυτοκτονήσει χωρίς λόγο και εσένα σε προβληματίζει αν επιτρέπεται να έχουμε όπλα στα σπίτια μας;»

«Ναι, γιατί;»

«Μη τον ακούς Σπαρκ. Κοίτα με» σηκώθηκε η Μπλαιρ και έβαλε τις παλάμες της στο πρόσωπό μου και μου έστρεψε το βλέμμα στα μάτια της. «Γιατί;»

«Γιατί... πίστευα πώς δεν αξίζω να ζω...»

«Τι σαχλαμάρες λες; Πες μου ότι ακόμη το πιστεύεις αυτό, μετά από όλη αυτή τη συζήτηση που έχουμε κάνει»

«Ναι... το πιστεύω. Γιατί να είμαι τόσο άτυχος; Η οικογένεια μου έχει πεθάνει, η κοπέλα μου με παράτησε, οι φίλοι μου με πρόδωσαν και... πιστεύω δεν έχω πια κίνητρο να ζήσω...»

«Σπαρκ... μη λες τέτοια λόγια.»

«Μετά από όλα όσα έχω ακούσει δεν μπορώ παρά να πω και κάτι άλλο...»

«Μα ακόμη και να μη σε αγαπάει κανένας, πράγμα που δεν ισχύει παρεμπιπτόντως αφού έχεις εμάς, ο Θεός θα σε αγαπάει, θα είναι πάντα μαζί σου να σε βοηθάει»

«Γιατί έγινες τόσο θρήσκα ξαφνικά;» γέλασε ο Κουπ.

«Ο Θεός; Πώς γίνεται να με αγαπάει και να με βάζει να περνάω τόσα δεινά;»

«Κοίτα με Σπαρκ. Όταν ακούς τις λέξεις “ο Θεός σε αγαπάει”, σκέφτεσαι “Ναι, όλους τους αγαπάει, ποιο το νόημα;”. Έχεις σταματήσει ποτέ, όμως, να σκεφτείς ότι δεν χρειαζόταν να υπάρχεις; Η ύπαρξη σου δεν είναι απαραίτητο μέρος της δημιουργίας του Θεού, ούτε η δική μου είναι. Οπότε ναι, ο Θεός σε αγαπάει, αλλά κοίτα, δεν ήταν απαραίτητο να ΥΠΑΡΧΕΙΣ. Υπάρχει μόνο ένας λόγος που ΥΠΑΡΧΕΙΣ και αυτός είναι επειδή ο Θεός, από όλη την αιωνιότητα, κοίταξε αυτό το σύμπαν και είπε: “Δεν θέλω ένα σύμπαν... χωρίς εσένα μέσα του”»

«Ω, τι ωραία λόγια...»

«Βούλοσ’ το Κουπ» είπε εκνευρισμένη η Μπλαιρ. «Ατυχίες συμβαίνουν σε όλους μας, όπως εσύ έχασες την οικογένεια σου, και εγώ έχασα την μοναδική οικογένεια που είχα, και ο Κουπ έχασε την μοναδική οικογένεια που είχε. Δεν μπορούμε όμως να αφήνουμε αυτές τις ατυχίες να μας ρίχνουν τόσο πολύ ψυχολογικά σε επίπεδο που βρισκόμαστε αντιμέτωποι με τη θνητότητά μας»

«Αυτό όμως σημαίνει πως... ο Θεός δεν αγαπούσε τους υπόλοιπους που πέθαναν;»

«Πάλι βλακείες λες, βγάλτο από το μυαλό σου. Είσαι, όπως είπα πριν, ο πιο τυχερός άνθρωπος του κόσμου, είτε το πιστεύεις, είτε όχι»

«Ο πιο τυχερός άνθρωπος του κόσμου; Γιατί συνεχίζεις και το πιστεύεις αυτό;»

«Αχ, άστο» κάθισε εκνευρισμένη με εμένα πάλι στο τραπεζάκι. «Κουπ πες του κάτι, μη κάνεις κύκλους μόνος σου»

«Δεν έχω κάτι να πω»

«Μπορώ... να ρωτήσω εγώ κάτι;»

«Για πες Σπαρκ» μου χαμογέλασε ο Κουπ.

«Είμαι... χρήσιμος;»

«Τι εννοείς;» είπαν ταυτόχρονα.

«Δηλαδή... τι... πρέπει να κάνω για να...» δεν μπόρεσα να κρατήσω τα δάκρυα μου άλλο πίσω.

«Έλα Σπαρκ, μην κλαις» με πήρε αγκαλιά η Μπλαιρ και με χτύπησε στην πλάτη για να με ανακουφίσει. «Είσαι χρήσιμος, ο καθένας μας είναι “χρήσιμος” με τον δικό του τρόπο. Τώρα δεν ξέρω πως ακριβώς το εννοείς, αλλά είσαι από τους πιο έξυπνους ανθρώπους που ξέρω. Το μόνο ελάττωμα σου είναι ότι υπεραναλύεις τα πάντα στο μυαλό σου. Δεν χρειάζεται, σταμάτα το αυτό. Οι σκέψεις σου είναι ο μεγαλύτερος σου εχθρός σε αυτή τη ζωή, ειδικά όταν είναι αρνητικές σκέψεις. Είδες που σε οδήγησε η υπερανάλυση των σκέψεων σου»

«Ναι...»

«Μπορείς να μου κάνεις μία υπόσχεση;» μου πρότεινε η Μπλαιρ υποχωρώντας από την αγκαλιά. «Θέλω να χαμογελάς από εδώ και πέρα, είτε βρίσκεσαι με εμάς είτε όχι. Μην αφήνεις τις σκέψεις σου να σου καταστρέφουν τη ζωή. Το παρελθόν είναι παρελθόν, είτε αυτό είναι πλαστό, είτε πραγματικό παρελθόν. Όπως μου είχε πει ένας φίλος, δεν μπορούμε να αλλάξουμε το παρελθόν. Μπορούμε μόνο να ανατρέξουμε τις αναμνήσεις που κάναμε και να κρίνουμε τον εαυτό μας για τα λάθη μας. Για αυτό, αντί να σκέφτεσαι “τι θα μπορούσα να είχα κάνει αλλιώς”, ζήσε στο παρόν. Προσπάθησε να κάνεις τις καλύτερες επιλογές τώρα έτσι, ώστε ο μελλοντικός σου εαυτός να μην σε κρίνει, αλλά να ευχαριστιέται τις αναμνήσεις του. Άρχισε να περνάς ποιοτικό χρόνο με τα άτομα που αγαπάς και σε αγαπάνε πίσω. Όσον αφορά τις αναμνήσεις σου... θα σε βοηθήσουμε να τις ανακτήσεις σιγά σιγά όσο μπορούμε. Είμαστε φίλοι εξάλλου εδώ και 9 χρόνια, αν και εσένα μπορεί να σου φαίνονται 8, αφού ήσουν σε κώμα... αλλά κατάλαβες»

Ο Κουπ κάθισε επιτέλους πάλι στο τραπεζάκι δίπλα στην Μπλαιρ. Εμένα όμως μου ήρθε μία ιδέα και σηκώθηκα από τον καναπέ.

«Που πας Σπαρκ;»

«Για μια βόλτα»

«Και θα μας αφήσεις μόνους μας στο σπίτι σου;»

Δεν της απάντησα, άνοιξα την εξώπορτα και ξεκίνησα να βάζω τα παπούτσια μου.

«Καλά, δεν θα σε σταματήσουμε» συνέχισε η Μπλαιρ. «Απλά κάνει κρύο, ζακέτα να πάρεις» είπε και ένιωσα ένα ρούχο να προσγειώνεται στην πλάτη μου.

Γύρισα πίσω μου να δω ποιος μου το πέταξε, αλλά στην προσπάθεια μου σκόνταψα στο μισοφορεμένο μου παπούτσι. Οι δύο τους ξέσπασαν σε γέλια βλέποντας με να χτυπάω. Δεν ξέρω γιατί, αλλά μου βγήκε και εμένα ένα νευρικό γέλιο εκείνη τη στιγμή.

«Αυτό ακριβώς εννοώ» μου είπε η Μπλαιρ.

«Τι;»

«Άστο, δεν πειράζει. Πήγαινε τώρα εσύ τη βόλτα σου, απλά μην αργήσεις. Δεν θέλουμε να αγχωθεί πολύ τη μητέρα σου»

\chapter{}

«Ώστε οι αναμνήσεις μου είναι ψεύτικες, ε; Ό,τι νόμιζα πως είναι πραγματικό ήταν απλά στο μυαλό μου. Μία ψεύτικη εικόνα της πραγματικότητας πλασμένη με μία δόση υποκειμενικότητας από πλευράς μου. Όλες μου οι αναμνήσεις είναι απλά λάθος σκέψεις που έκανα για τα άτομα στη ζωή μου. Αν και... αυτές οι σκέψεις μου δεν ήταν τελείως άσχετες με την πραγματικότητα.

Πάλι όμως, θα έπρεπε να το είχα καταλάβει ότι δεν βρισκόμουν στην πραγματικότητα, τα σημάδια ήταν όλα εκεί. Οι παράλογες προδοσίες, η ξαφνική εμφάνιση πραγμάτων από το πουθενά, ο μακρύς ανεξήγητος διάδρομος στην Ειδική Υπηρεσία, για να μη μιλήσω για τη μαγεία... δεν μπορούσαν να μου είχαν δοθεί περισσότερα σημάδια. Ας μη μιλήσω επίσης για τη ταχεία εξέλιξη των πραγμάτων. Μέσα σε υποτίθετε 10 μέρες είχα μπει στην υποτιθέμενη Ειδική Υπηρεσία και είχα σχεδόν ολοκληρώσει τρεις αποστολές, οι συνεργάτες μου από την πρώτη τους αποστολή μαζί μου φέρονταν λες και είμαστε κολλητοί για χρόνια, τα δέκα χρόνια στην φυλακή πέρασαν λες και ήταν δέκα λεπτά. Πρέπει να είμαι πιο παρατηρητικός»

Καθώς το είπα αυτό, πέρασα έξω από έναν δρόμο που μου θύμισε κάτι. Δεν μπορούσα να θυμηθώ τι ή γιατί αυτός ο δρόμος που προκάλεσε περιέργεια, μέχρις ότου έφτασα έξω από ένα πολυόροφο κτίριο με μία γυάλινη περιστρεφόμενη πόρτα. Τότε ήταν που θυμήθηκα γιατί αυτός ο δρόμος μου είχε κινήσει το ενδιαφέρον.

«Ωχ, η... Ειδική Υπηρεσία; Ώστε είναι πραγματικό κτίριο, ε;» είπα και προσπάθησα να κοιτάξω μέσα από την γυάλινη πόρτα για να δω αν το εσωτερικό της κτιρίου ήταν όπως το θυμόμουν. Πράγματι, από τους ανθρώπους μέχρι την εσωτερική αρχιτεκτονική του κτιρίου, ή τουλάχιστον όσα μπορούσα να διακρίνω από το σημείο που βρισκόμουν, ήταν όμοια με όσα θυμόμουν, σε πολύ μικρότερη κλίμακα βέβαια.

Μπαίνοντας μέσα, πλησίασα το γραφείο της γραμματέα. Αυτή προβληματίστηκε, αλλά δεν με ρώτησε κάτι. Προχώρησα προς τον υποτιθέμενο μακρύ διάδρομο, αλλά το μόνο που είδα ήταν ένας μικρός διάδρομος με τέσσερις πόρτες, μία στα δεξιά στη μέση του διαδρόμου και άλλες τρεις στο τέλος του διαδρόμου.

«Μπορώ να σας βοηθήσω;» με ρώτησε η γραμματέας βλέποντας με προβληματισμένο μπροστά στο διάδρομο. «Έχετε κάποιο ραντεβού; Αν και θα ήταν περίεργο να έχετε ραντεβού τέτοια μέρα...»

«Ε...»

«Αν δεν έχετε, παρακαλώ περάστε από εδώ και θα σας βοηθήσω εγώ. Αν έχετε, πείτε μου τι ακριβώς ψάχνετε» μου εξήγησε η γραμματέας. «Γιατί συμφώνησα να έρθω σήμερα για να αναλάβω τα έκτακτα ραντεβού, ρε φίλε» ψιθύρισε εκνευρισμένη με τον εαυτό της.

«Που... είμαι; Τι είναι αυτό το μέρος;»

«Τι είναι αυτό το μέρος; Κύριε μου έχουμε μία ταμπέλα απ’ έξω στην οποία αναγράφεται “Αστυνομικό Τμήμα”, που αλλού νομίζετε πως βρίσκεστε; Μήπως είστε υπό την επήρεια κάποιας ουσίας;»

«Αστυνομικό τμήμα, ε;» ο Κουπ μου είχε πει πως είμαι αστυνομικός. Αυτό θα εξηγούσε πολλά πράγματα, στις “αναμνήσεις” μου. «Να σας ρωτήσω κάτι;»

«Παρακαλώ»

«Μήπως δουλεύει εδώ κάποιος... Σπαρκ Θοτς;»

«Ναι, αλλά γιατί ρωτάτε; Είστε συγγενής;»

«Βασικά... εγώ είμαι»

«Τι λες, όντως; Μα άκουσα πώς είσαι...» η γραμματέας σηκώθηκε για να με κοιτάξει λίγο καλύτερα. «Να πω την αλήθεια μοιάζεις, αλλά ο Σπαρκ δεν είχε μούσι, ούτε ήταν τόσο λεπτός»

«Ίσως... αυτό μπορεί να σε πείσει...» είπα και έβγαλα το κινητό μου από την τσέπη μου.

«Το κινητό σου; Πώς θα με πείσει αυτό;»

Άνοιξα τις επαφές μου και για καλή μου τύχη είχα αρκετά γνωστά ονόματα και αρκετά όχι τόσο γνωστά ονόματα. Ίσως τα άγνωστα ήταν συνάδελφοι μου που δεν μπορούσα να τους θυμηθώ αυτή τη στιγμή. Συνέχισα να ψάχνω όμως μέχρις ότου να βρω την επαφή της γραμματέα.

«Μήπως μπορώ να σας ρωτήσω... πώς σας λένε;»

«Λιλι»

«Λιλι, Λιλι...»

*ντριν ντριν... ντριν ντριν...*

«Μίσο λεπ... ωχ... είσαι όντως εσύ. Σπαρκ... όλοι νομίζαμε πώς είχες πεθάνει. Ωχ θα μου φύγουν δάκρυα χαράς...»

«Δάκρυα χαράς;»

«Δεν μπορώ, έλα εδώ να σε κάνω μία αγκαλιά» είπε και σηκώθηκε, ανοίγοντας τα χέρια της διάπλατα.

«Συγγνώμη, αλλά... βιάζομαι. Ίσως να τα πούμε κάποια άλλη στιγμή» είπα αλλά δεν είχα στην πραγματικότητα κάποιο σοβαρό λόγο. Δεν ήθελα απλά να κάτσω κι άλλο εκεί πέρα γιατί η παρουσία μου θα μπορούσε να ξεκινήσει κι άλλες τέτοιες συζητήσεις με πρακτικά αγνώστους για μένα ανθρώπους στο αστυνομικό τμήμα, τις οποίες ήθελα να αποφύγω.

«Βιάζεσαι; Για ποιο λόγο; Δεν ήσουν στο νοσοκομείο τόσο καιρό;»

«Ναι, αλλά... ας πούμε ότι έχω τους λόγους μου»

«Καλά... όπως πάντα μένεις μυστήριος. Εντάξει λοιπόν, απλά την επόμενη φορά να μας ενημερώσεις πριν έρθεις. Μην εμφανιστείς μία τυχαία μέρα για δουλειά χωρίς προειδοποίηση» είπε και γέλασε μόνη της η γραμματέας.

«Ναι... τέλος πάντων» είπα και γύρισα προς την είσοδο.

«Αντίο Σπαρκ!»

Δεν της απάντησα. Βγήκα από το κτίριο σκεπτόμενος. Αφού το κτίριο της “Ειδικής Υπηρεσίας” υπήρχε όντως, θα μπορούσαν και άλλα μέρη που θυμόμουν να είναι πραγματικά, όπως το καζίνο, η εγκαταλελειμμένη εκκλησία ή τα χωριά έξω από την πόλη.

«Υπάρχει μόνο ένας τρόπος να μάθω» είπα και ξεκίνησα να προχωράω προς το καζίνο. «”Δάκρυα χαράς”, ε; Γιατί όμως; Τι μπορεί να έχω κάνει για να προκαλώ τέτοια συναισθήματα σε κάποιον που ξέρω τόσο ελάχιστα; Τι είναι η χαρά; Πώς μπορώ να περάσω ευχάριστα τον χρόνο μου; Τι κάνει έναν άνθρωπο χαρούμενο; Δεν μπορώ να καταλάβω. Κάνω κάτι λάθος; Έχω κάποιο σκεπτικό; Θεέ, αν όντως με αγαπάς, όπως είπε η Μπλαιρ, δώσε μου ένα σημάδι, σε παρακαλώ»

Εκείνη τη στιγμή η προσοχή μου στράφηκε σε τρία παιδιά που έτρεχαν, ένα εκ των οποίων σκόνταψε. Τα άλλα δύο άρχισαν να γελάνε με τον πόνο του άλλου παιδιού, όμως το τρίτο παιδί άρχισε να γελάει μαζί τους. Μου θύμισαν το σκηνικό που συνέβη πριν φύγω από το σπίτι μου. Δεν νομίζω όμως πως ένα συναίσθημα που σε γεμίζει τόσο πολύ να αποκτάται τόσο εύκολα. Κάνω λάθος;

Φτάνοντας στην αυλή της εγκαταλελειμμένης εκκλησίας, ένιωσα μία μικρή απογοήτευση με τον εαυτό μου. Αλλά γιατί; Επειδή δεν μπόρεσα στις ψεύτικες αναμνήσεις μου να ολοκληρώσω την αποστολή μου; Γιατί είμαι τόσο σκληρός με τον εαυτό μου; Ήμουν πάντα έτσι; Στις “αναμνήσεις” μου ναι, στην πραγματικότητα... δεν μπορώ να ξέρω. Η Μπλαιρ είχε πει πως τείνω να υπεραναλύω τα πάντα στο μυαλό μου, οπότε μάλλον είμαι κάποιου είδους τελειωμανής; Αυτό θα εξηγούσε πολλά. Τώρα πλέον όμως, δεν ξέρω... ποιος είμαι. Τουλάχιστον έχω μέχρι στιγμής μία ιδέα για το που είμαι...

«Πάμε προς το καζίνο τώρα, να δούμε αν υπάρχει και αυτό» είπα και έσπευσα προς το καζίνο. Δεν είχα καμία πρόθεση να παίξω κάποιο παιχνίδι, ήθελα μόνο να επιβεβαιώσω την ύπαρξη του.

Γυρνώντας το βλέμμα μου από την αυλή της εκκλησίας προς το στενό στο οποίο ήθελα να στρίψω, είδα έναν μοναχικό συνοφρυωμένο άντρα ένα στενό παραπέρα να περπατάει προς το μέρος μου. Με την άκρη του ματιού μου είδα ένα χαρούμενο ζευγαράκι στο απέναντι στενό να πλησιάζει από τα δεξιά τον άντρα. Όταν και οι τρεις τους φτάσανε στη γωνία, υπήρξε μία ανεπαίσθητη αλλά ακούσια ακούσια σύγκρουση μεταξύ των δύο αρσενικών. Ο χαρούμενος άντρας άρχισε να ζητάει συγγνώμη, διατηρώντας το χαμόγελο στο πρόσωπό του, ενώ ο συνοφρυωμένος άντρας τους αγνόησε και συνέχισε να περπατάει προς το μέρος μου.

Περπατώντας προς το καζίνο είδα μία ξανθιά κοπέλα να κάθεται στα σκαλοπάτια μιας πολυκατοικίας στηρίζοντας με τα χέρια το κεφάλι της, καλύπτοντας με τις παλάμες το πρόσωπο της, λες και δεν ήθελε να δει κανείς το πρόσωπο της. Πλησιάζοντας την και κοιτάζοντας την διακριτικά κατάλαβα πως αυτή η κοπέλα έκλαιγε. Μου πέρασε από το μυαλό η σκέψη να κάτσω δίπλα της και να της μιλήσω, αλλά την απέρριψα καθώς μπορεί να φαινόμουν περίεργος. Συνέχισα να περπατώ λες και δεν είχα δει τίποτα.

«Έλα, μη κλαις, ήρθα» άκουσα μία γυναικεία φωνή από πίσω μου. Γυρνώντας είδα μία ακόμη ξανθιά κοπέλα να στέκεται δίπλα στην λυπημένη κοπέλα από πριν και να προσπαθεί να την καθησυχάσει σηκώνοντας την και αγκαλιάζοντας την. Εκείνη τη στιγμή, η πρώτη κοπέλα, παρότι συνέχιζε να έκλαιγε, είχε ένα χαμόγελο στα χείλη της λες και όλα λύθηκαν επειδή η δεύτερη κοπέλα εμφανίστηκε.

Πώς γίνεται αυτό; Πώς μπορεί κάποιος να αλλάζει τη διάθεση του τόσο γρήγορα βλέποντας μόνο έναν άλλον άνθρωπο; Αυτή είναι η χαρά που αποζητώ; Δεν μπορώ να καταλάβω.

Φτάνοντας έξω από το καζίνο, είδα μία παρέα από πέντε άντρες να συνεννοούνται μεταξύ τους με λίγες σχετικά λέξεις, αλλά να έχουν ξεκαρδιστεί μόνοι τους έξω από την πόρτα του καζίνο. Τους πλησίασα για να κρυφακούσω τη συζήτηση τους.

«Και ρε, ο μπρο τα έβαλε όλα στο κόκκινο» είπε ένας από τους άντρες και άρχισαν όλοι να γελάνε λες και υπήρχε κάτι αστείο στην πρόταση που μόλις ειπώθηκε.

«Καλά τι μιλάς εσύ, αφού κι εσύ όλα τα έχασες»

«Ναι αλλά στο κόκκινο; Αν τα είχες βάλει στο μαύρο όπως σου λέγαμε όλοι θα είχες κερδίσει»

«Τουλάχιστον χάσαμε όλοι μαζί, αυτό κάνει την εμπειρία μας ευχάριστη»

«Ναι ισχύει, αφού τώρα όλοι είμαστε άφραγκοι, χαχα»

Ακόμη δεν μπορώ να καταλάβω. Τι το αστείο έχει να χάνεις τα λεφτά σου; Πώς μπορεί να σου προκαλεί αυτό το γεγονός τόση ευχαρίστηση; Ήθελα να πεταχτώ και να τους διακόψω για να τους ρωτήσω αυτό ακριβώς, αλλά απέρριψα αυτή μου τη σκέψη γιατί φοβήθηκα πάλι μην φανώ παραπάνω περίεργος απ’ ό,τι ήδη ήμουν.

«Που αλλού θα μπορούσα να πάω; Στο χωρίο του Γιν; Μπα, είναι πολύ μακριά. Στο σπίτι των αιρετικών; Γιατί όχι. Είναι μισή ώρα οδικώς απ’ ότι θυμάμαι, που σημαίνει... Πόσην ώρα είναι αυτό με τα πόδια; Πω, είναι και έξω από την πόλ. Τέλος πάντων, ας ξεκινήσω τώρα να πηγαίνω προς τα εκεί γιατί είμαι περίεργος»

Η προσοχή μου στράφηκε για άλλη μία φορά τους άντρες έξω από καζίνο πριν ξεκινήσω να περπατάω. Είχαν απομακρυνθεί λίγο και δεν μπορούσα να τους ακούσω. Παρατηρούσα απλά την εικόνα μιας παρέας πέντε ατόμων να ξεκαρδίζονται με το παραμικρό. Μία παρέα... Λες αυτή να είναι η λέξη κλειδί; “Παρέα”. Όποιον χαρούμενο έχω δει μέχρι στιγμής είχε συντροφιά στο πλάι του. Όποιον λυπημένο έχω παρατηρήσει κατά τη διάρκεια αυτής της περιπλάνησης μου ήταν μοναχικός. Μήπως, αυτό είναι το μυστικό; Η παρέα. Γιατί σου προκαλεί ευχαρίστηση όμως η παρέα; Τι το ιδιαίτερο έχει η συνύπαρξη με άλλους ανθρώπους που κάνει ένα άτομο χαρούμενο; Δεν μπορώ να καταλάβω. Ξεκίνησα να περπατάω μόνος μου την διαδρομή που θυμόμουν πως κατέληγε στο σπίτι των αιρετικών.

Η Μπλαιρ είπε πως το μόνο μου ελάττωμα είναι πως υπεραναλύω τα πάντα στο μυαλό μου. Απ’ ό,τι φαίνεται δεν έχει και άδικο. Μέχρι στιγμής αυτή η ανάλυση δεν με έχει οδηγήσει κάπου. Όλο τελειώνω τις σκέψεις μου λέγοντας σον εαυτό μου ότι δεν καταλαβαίνω. Η παρέα όμως πώς μπορεί να είναι το κλειδί για την χαρά; Μπορεί να μην υπάρχει κάποιος προφανής λόγος, ίσως χρειάζεται να το πάρω σαν αξίωμα ότι “παρέα” συνεπάγεται “ευχάριστος χρόνος”. Πλάκα πλάκα, ούτε στις “αναμνήσεις” μου, ούτε στις μέρες μου στο νοσοκομείο αποδεχόμουν με ανοιχτά χέρια την συντροφιά άλλων ανθρώπων. Παρόλο που “δουλεύαμε” ως ομάδα στις “αναμνήσεις” μου, συνήθως απέφευγα το κοινωνικό κομμάτι εκτός δουλειάς. Όσες φορές μας είχε πει ο Μαλχη να πάμε στο καζίνο ή ο Κουπ με τον Φυτ να πάμε για κανέναν καφέ, εγώ πάντα αρνιόμουν και προτιμούσα να μείνω πίσω μόνος μου. Στο νοσοκομείο, όταν δεν με εξέταζε κάποιος γιατρός ή δεν έκανα κάποια άσκηση για να βελτιώσω το περπάτημα μου, καθόμουν απλά στο κρεβάτι και κοίταζα τους τοίχους με τη μητέρα μου. Ναι μεν πάντα είχα παρέα, αλλά δεν την αξιοποιούσα. Θεωρούσα τους ανθρώπους γύρω μου δεδομένους και προχωρούσα χωρίς να κάνω κάτι παραπάνω από την απλή συνύπαρξη στο ίδιο χώρο. Ίσως αυτό είναι άλλο ένα ελάττωμα μου. Αν βελτιωθώ σε αυτό θα... γίνω χαρούμενος;

Ας ξεκινήσω πρώτα μειώνοντας την υπερανάλυση των σκέψεων μου. Για να το πετύχω αυτό, δεν θα ξανα αφήσω τον εαυτό μου να αρχίσει αυτό το ασταμάτητο εσωτερικό μουρμουρητό που κάνω τόση ώρα, τουλάχιστον μέχρις ότου να φτάσω στο σπίτι των αιρετικών, ή ό,τι και να υπάρχει σε εκείνο το μέρος. Όταν φτάσω, θα δω τι θα κάνω με το άλλο ελάττωμα μου.

\chapter{}

«Πρέπει να έχει περάσει κανένα δίωρο που περπατάω στην ησυχία. Ούτε σκέψεις, ούτε άσκοποι μονόλογοι, ούτε τίποτα. Είναι ίσως η πρώτη φορά που μιλάω στον εαυτό μου από τότε που αποχώρησα από την πόλη. Έχω να πω πως ίσως έχω χαθεί, αλλά πιστεύω πως η διαδρομή πίσω είναι παιχνιδάκι. Το μόνο που έχω να κάνω είναι να ακολουθήσω τον δρόμο που περπατάω και κατά πάσα πιθανότητα θα φτάσω στην πόλη. Τι είναι όμως αυτή η πεδιάδα; Πρώτη φορά έρχομαι από εδώ. Ούτε στις ψεύτικες μου αναμνήσεις δεν θυμάμαι τέτοιο μέρος. Χορτάρι και άσφαλτος μέχρι τον ορίζοντα προς τα πίσω και μπροστά μου βλέπω εδώ αρκετή ώρα έναν μικρό λόφο με ένα μικρό κτίριο σαν σπιτάκι πάνω του. Είναι το ίδιο με το σπιτάκι των αιρετικών που θυμάμαι; Όχι, αλλά οι ελπίδες μου δεν πέφτουν. Τώρα που έχω φτάσει στους πρόποδες του λόφου δεν αξίζει να τα παρατήσω.

Πολύ μεγάλη διαδρομή πάντως. Η ησυχία την έκανε ακόμη μεγαλύτερη. Λες αν είχα παρέα να πέρναγε πιο γρήγορα ο χρόνος, όπως είπε και η μητέρα μου όταν βγήκα από νοσοκομείο; Τώρα είναι πολύ αργά για να μάθουμε Σπαρκ, προχώρα και σταμάτα να μιλάς. Εντάξει Σπαρκ»

Ανεβαίνοντας τον λόφο, άρχισα να καταλαβαίνω περισσότερο τον κρύο καιρό που επικρατούσε. Παρόλο που ήταν ηλιόλουστη μέρα, η υγρασία και ο ψυχρός άνεμος σίγουρα κατέβαζαν την αισθητή θερμοκρασία κάτω από τους δέκα βαθμούς κελσίου. Τα πόδια μου έτρεμαν, καθώς είχα περπατήσει ήδη υπερβολικά πολλά χιλιόμετρα και δεν ήμουν ακόμη ούτε στα μισά της διαδρομής μου, αφού έπρεπε και να επιστρέψω πίσω στην πόλη.

Φτάνοντας στην κορυφή του λόφου, δεν προχώρησα προς το μικρό σπιτάκι που βρισκόταν εκεί. Σταμάτησα για λίγο να απολαύσω τη θέα. Μία τεράστια λίμνη βρισκόταν στην άλλη μεριά του λόφου, η οποία περιτριγυριζόταν από ένα δάσος που εκτεινόταν μέχρι εκεί που δεν μπορεί να δει το μάτι. Κάθησα για λίγο στο χορτάρι στην κορυφή του λόφου και παρατήρησα τα ζώα γύρω και μέσα στη λίμνη. Πάπιες που κολυμπούσαν, σκίουροι στα δέντρα, μυρμήγκια δίπλα στα πόδια μου, όλα τα ζώα λειτουργούσαν σε ομάδες. Σε τι επίπεδο όμως αυτά τα ζώα μπορούν να νιώθουν ευχαρίστηση; Μπορούν να είναι χαρούμενα; Αν ναι, γιατί; Πώς συνδέεται η χαρά τους με την ανθρώπινη χαρά;

«Άρχισα πάλι τις σκέψεις. Ίσως οι σκέψεις είναι όντως ο χειρότερος μου εχθρός. Ο Μαλχη όμως στις αναμνήσεις μου είχε πει να αγαπάμε τους εχθρούς μας. Κατά πόσο αυτό ήταν αλήθεια... δεν ξέρω, και δεν μπορώ να τον ρωτήσω κιόλας. Τι, να τον πάρω έτσι στο άσχετο τηλέφωνο αφότου πριν τέσσερις και κάτι μήνες τον είχα πλακώσει και σε αυτό το χρονικό διάστημα δεν του έχω μιλήσει καν; Είναι λίγο περίεργο όπως και να το κοιτάξει κανείς»

«Συγγνώμη» άκουσα μία αντρική φωνή να με πλησιάζει από την κατεύθυνση του σπιτιού. 

Γύρισα το βλέμμα μου προς την κατεύθυνση του και είδα έναν νέο άντρα, πιο νέο από μένα, ίσως να ήταν και έφηβος. Ο άντρας αυτός φορούσε μία μαύρη ποδιά και ένα ταιριαστό καπέλο που έγραφαν πάνω με άσπρα γράμματα “Καφέ της Λίμνης”.

«Συγγνώμη κύριε, είστε μόνος σας;» ήρθε δίπλα μου ο άντρας

«Τι;»

«Είστε μόνος σας;» με ξανα ρώτησε ο νέος άντρας και υπήρξε μία στιγμή ησυχίας

«Δεν... ξέρω»

«Τέτοια μέρα και είστε μόνος σας; Είναι γιορτινές μέρες, δεν θα έπρεπε να είστε με την οικογένεια σας; Τι κάνετε εδώ πέρα μόνος σας;»

Δεν του απάντησα.

«Δεν... θα έπρεπε να είστε με την οικογένεια σας;»

«Συγγνώμη, αλλά... τι μέρα είναι;»

«Παραμονή Χριστουγέννων. Δεν το γνωρίζατε;»

Άνοιξαν τα μάτια μου διάπλατα. Εκείνη τη στιγμή μπήκα σε πολλές σκέψεις. Λες για αυτό να μας άφησαν πρόωρα από το νοσοκομείο; Μπα, δεν νομίζω να έπαιξε ρόλο αυτό. Η Μπλαιρ και ο Κουπ όμως; Γιατί ήρθαν συγκεκριμένα αυτή τη μέρα σπίτι μου με πρόθεση να συζητήσουμε; Δεν έχουν οικογένειες και αυτοί; Δεν... έχουν...

«Κύριε; Τελικά είστε μόνος σας ή...»

Εκείνη τη στιγμή χτύπησε το κινητό μου και του διέκοψε την ροή του λόγου του.

«Ο Κουπ;... Παρακαλώ»

«Έλα Σπαρκ, που είσαι; Έχουν περάσει πέντε ώρες από τότε που έφυγες, θα νυχτώσει σε κανένα μισάωρο. Είσαι καλά; Ακούγεσαι κουρασμένος»

«Όχι όχι... καλά είμαι»

«Ωραία, πες μας που είσαι για να έρθουμε να σε πάρουμε»

«Που... είμαι;» κοίταξα τον νέο άντρα. «Στο “Καφέ της Λίμνης”»

«Πω, πώς έφτασες μέχρι εκεί με τα πόδια. Αυτό είναι πολύ έξω από την πόλυ. Σε κανένα σαραντάλεπτο θα φτάσουμε το γρηγορότερο αν το πατήσω λίγο. Μην κουνηθείς από εκεί, ερχόμαστε» είπε ο Κουπ και μου το έκλεισε.

«Θέλετε να έρθετε μέσα στο καφενείο τουλάχιστον για να μην είστε μόνος σας έξω στο κρύο; Ή δεν είστε μόνος σας;»

«Όχι δεν... απλά... ναι αλλά...» πήρα μία βαθιά ανάσα και τον κοίταξα στα μάτια καθώς ένα κρύο αεράκι με χτύπησε απαλά στην πλάτη μου. «Πάμε απλά μέσα, αφού είπες και μόνος σου κάνει κρύο εδώ έξω»\\\\

«Καλησπέρα σας. Εδώ δεν είναι το Καφέ της Λίμνης;»

«Ναι, καλώς ήρθατε. Μήπως είστε ο φίλος του κυρίου εκεί πέρα;»

«Αχ ναι, ευχαριστώ πολύ. Βλέπω παρήγγειλε και καφέ, ε; Εκπλήσσομαι που θυμάται τι καφέ πίνει»

«Δεν παρήγγειλε ο ίδιος. Του φτιάξαμε απλά ένα ζεστό ρόφημα γιατί φαίνεται λες και έχει περάσει πολλά ο άνθρωπος»

«Μήπως θέλετε να σας πληρώσω εγώ;»

«Όχι φιλέ, κράτα τα. Δες το σαν χριστουγεννιάτικο δώρο»

«Α... ευχαριστούμε τότε, για όλα, όχι μόνο για τον καφέ. Σπαρκ!»

«Άσε θα πάω να τον σηκώσω εγώ. Κάθετε σε εκείνη την καρέκλα εδώ και τρία τέταρτα κοιτάει την λίμνη. Δεν ξέρω τι σκέφτεται αλλά ελπίζω να πάνε όλα καλά» είπε στον Κουπ ο νέος άντρας και τον άκουσα να με πλησιάζει. «Κύριε... Σπαρκ; Ήρθαν οι...»

«Όχι δεν είμαι» τον διέκοψα.

Ο νέος φάνηκε προβληματισμένος με τα λόγια μου.

«Δεν είμαι μόνος μου»

Είδα τον προβληματισμό στο πρόσωπο του νέου να μεταμορφώνεται σε ευχαρίστηση όταν κατάλαβε σε τι αναφερόμουν.

«Χαίρομαι» μου απάντησε μονολεκτικά.

«Τι λέτε, δεν καταλαβαίνω;»

«Άστο Κουπ, πάμε στο αυτοκίνητο»

«Καλά... Καλή συνέχεια φίλε. Ευχαριστούμε και πάλι»

«Καλή συνέχεια παιδιά και καλά χριστούγεννα»

Αποχαιρέτησα και εγώ με ένα απλό χαμόγελο τον νέο και το βλέμμα μου στράφηκε στο τραπέζι που καθόμουν. Είχε δίκιο ο νέος, καθόμουν τόση ώρα στην καρέκλα, όμως είχα βγάλει όλες μου τις σκέψεις από το κεφάλι μου. Είχε δίκιο η Μπλαιρ, ούτε ο Μαλχη, ούτε η Ιφι, ούτε κανένας άνθρωπος δεν ήταν εχθρός μου. Ο πραγματικός μου εχθρός ήταν οι σκέψεις μου, οι λάθος κρίσεις μου, η υπεραναλύσεις που δεν οδηγούσαν πουθενά, οι αρνητικές σκέψεις που μου έφερναν αυτοκτονικές τάσεις. Δεν μπορώ όμως να σταματήσω να σκέφτομαι, αυτό θα ήταν χαζό. Εξάλλου αυτή τη στιγμή πάλι μονολογώ με τον εαυτό μου. “Να αγαπάς τον εχθρό σου” είχε πει κάποιος στο όνειρο μου, όμως αναφερόταν στη γενική περίπτωση ή αναφερόταν σε εμένα;

«Σπαρκ;»

«Ναι;»

«Πάλι αφαιρέθηκες;» γέλασε ο Κουπ. «Μου είπες να πάμε στο αυτοκίνητο και με άφησες να φύγω μόνος μου»

«Α... συγγνώμη» είπα και τον ακολούθησα. «Η Μπλαιρ;»

Φτάνοντας στο αυτοκίνητο, είδα την Μπλαιρ να κάθεται στη θέση του συνοδηγού και να κάνει τη χαρακτηριστική της κίνηση με το δεξί της χέρι και να κουνάει το στόμα της λες και μουρμούριζε κάτι στον εαυτό της.

«Σολφέζ, ε;» της είπα μόλις άνοιξα την πόρτα για να κάτσω στο πίσω κάθισμα.

«Ναι... πώς το ήξερες; Αφού ήσουν σε κώμα όταν ξεκίνησα μαθήματα μουσικής, δεν πρόλαβα...»

«Είχα μία προαίσθηση» είπα και κάθισα.

Ο Κουπ έβαλε μπρος τη μηχανή και η Μπλαιρ άπλωσε το χέρι της γρήγορα προς το ραδιόφωνο για να το ανοίξει.

«Όχι πάλι εσύ μουσική» αναστέναξε ο Κουπ.

«Γιατί; Δεν σου αρέσουν τα τραγούδια μου;»

«Όχι, απλά... βαριέμαι»

«Τι σημαίνει βαριέμαι; Βαριέσαι να ακούσεις μουσική; Πώς γίνεται αυτό;»

«Καλά καλά, βάλε ότι θες» της είπε ο Κουπ επειδή δεν ήθελε να τσακωθεί και ξεκίνησε να οδηγάει.

«Τι σας πήρε τόση ώρα τέλος πάντων;»

«Δεν κάναμε καν τόση ώρα, τι λες; Εσένα τι σου πήρε τόση ώρα να πάρεις καφέ πριν φύγουμε;»

«Τι συγκρίνεις τώρα;» είπε εκνευρισμένη η Μπλαιρ καθώς άλλαζε τους σταθμούς μέχρι να βρει κάτι που της άρεσε.

«Δεν συγκρίνω τίποτα, απλά λέω...»

«Να... ρωτήσω μερικά πράγματα;» τους διέκοψα.

«Ναι Σπαρκ, ό,τι θες» μου είπε η Μπλαιρ με ένα χαμόγελο στα χείλη της, λες και δεν είχε τσακωθεί πριν μερικά δευτερόλεπτα με τον Κουπ.

«Καταρχάς, ξέρετε κανέναν Πιτ;»

«Πιτ;» ξαφνιάστηκε η Μπλαιρ. «Όχι δεν νομίζω. Γιατί ρωτάς;»

«Γιατί στο όνειρο μου την περίοδο που ήμουν φυλακή ένας Πιτ με... “φρόντιζε”; Δεν ξέρω πως να το εξηγήσω»

«Α τώρα που είπες φυλακή ναι κάτι μου θυμίζει» μου απάντησε η Μπλαιρ και σταμάτησε να αλλάζει σταθμούς «Ου ABBA, μου αρέσει αυτό»

«Τι εννοείς σου θυμίζει κάτι;»

«Να, όταν σε είχαν βρει ανάμεσα στα συντρίμμια του αυτοκινήτου, κρατούσες ένα βιβλίο που έλεγε για την ζωή στη φυλακή και ο πρωταγωνιστής ήταν ένας φύλακας που λεγόταν Πιτ νομίζω, δεν είμαι σίγουρη όμως. Είχα ξεκινήσει να το διαβάζω μια μέρα όσο καθόμασταν στο νοσοκομείο δίπλα σου, αλλά να πω την αλήθεια δεν το τελείωσα πότε»

«Ούτε καν που το άρχισες. Δέκα σελίδες διάβασες και έλεγες ότι κουράστηκες» σχολίασε ο Κουπ.

«Τι λες μωρέ, έντεκα διάβασα»

«Ακριβώς»

Η Μπλαιρ έκανε μία μικρή παύση και γέλασε.

«Στο νοσοκομείο δηλαδή... εσείς ήσασταν που με επισκεπτόσασταν;»

«Ναι ποιος άλλος να ήταν; Εμείς και η μητέρα σου προφανώς. Οι άλλοι τρεις ερχόντουσαν στην αρχή και μετά σε παράτησαν»

Άλλη μία φορά που έκανα λάθος στις σκέψεις μου. Νόμιζα πώς ο Φυτ και η Ιφι ήταν αυτοί που με επισκέπτονταν κάθε μέρα αλλά στην πραγματικότητα, είχα πολύ καλύτερη παρέα όσο ήμουν αναίσθητος.

«Γιατί ρωτάς;»

«Α... χωρίς λόγο... Τώρα που είναι στο θέμα μας όμως, τι έγινε εκείνη τη μέρα;»

«Ποια μέρα;» είπαν και οι δύο τους ταυτόχρονα.

«Όταν ξύπνησα; Που φύγατε όλοι; Έγινε κάτι μετά;»

«Όχι κάτι ιδιαίτερο» μου απάντησε ο Κουπ. «Η Μπλαιρ απλά έκραξε άλλη μία φορά την Ιφι, ο Μαλχη...»

«Δεν την έκραξα...» τον έκοψε η Μπλαιρ. «Καλά... ναι βασικά την έκραξα»

«Βλέπεις» είπε ο Κουπ κοιτάζοντας της και σαν να είχε ειπωθεί κάτι αστείο γελάσαν πάλι οι δυο τους μόνοι τους. «Ναι αυτό βασικά. Ο Μαλχη πήγε να τζογάρει μετά και η Ιφι με τον Φυτ μας παράτησαν και αυτοί γιατί είχαν έρθει μαζί. Εμείς κάτσαμε... Ωχ ρε Μπλαιρ, πώς δεν το είχα σκεφτεί πριν μου το πεις. Αμάν... έχω πραγματικά μηδέν παρατηρητικότητα»

Ξανακοιτάχτηκαν και άρχισαν πάλι να γελάνε χωρίς λόγο. Το γέλιο τους, σαν κολλητική αρρώστια, άρχισε να με επηρεάζει και εμένα και πριν το καταλάβω είχα αρχίσει και εγώ να γελάω χωρίς να ξέρω γιατί.

«Ξέρεις όμως τι άλλο είναι αστείο;» του είπε η Μπλαιρ.

«Τι;»

«Είναι... αχ, το ξέχασα»

«Εγώ δεν έχω παρατηρητικότητα και εσύ μνήμη απ’ ό,τι φαίνεται»

«Τι λες βρε...» είπε η Μπλαιρ και τον έσπρωξε παιχνιδιάρικα στον ώμο.

«Να ρωτήσω και κάτι ακόμη;»

«Ναι ρε Σπαρκ, δεν χρειάζεται να το ανακοινώνεις κάθε φορά» μου είπε η Μπλαιρ.

«Μήπως... είστε μαζί;»

«Τι λες; Εγώ; Με αυτό το ζωντόβολο; Ούτε να με πλήρωρες»

Ο Κουπ δεν φάνηκε να προσβλήθηκε καθόλου από τα αρνητικά λόγια της Μπλαιρ.

«Πώς σου ήρθε αυτή η σκέψη;» με ρώτησε προβληματισμένη.

«Ξέρω γω, πώς του ήρθε;» σχολίασε ο Κουπ.

«Τι λέτε; Είναι κάποιο αστείο που δεν ξέρω; Τι είπατε μέσα στο καφέ;»

«Πφ» του έφυγε ένα μικρό γελάκι. «Τίποτα»

«Α, για αυτό αργήσατε, ε;»

«Ναι ρε, σίγουρα...» είπε ο Κουπ και μου ξέφυγε και εμένα ένα μικρό γελάκι ακούγοντας τους να τσακώνονται.

«Με κοροϊδεύεις, ε; Μήπως θες να σε χτυπήσω όπως είχε χτυπήσει η μητέρα του Σπαρκ τον τοίχο»

«Ω, φοβάμαι τώρα»

«Θα έπρεπε»

«Κάτσε» τους διέκοψα. «Ποιος; Ποιανού τον τοίχο;»

«Α ναι, δεν στο είπα. Όταν γύρισε η μητέρα σου από το επαγγελματικό της ταξίδι, πήγαμε σπίτι σου για να της κάνουμε παρέα και είχε ο τοίχος στο σαλόνι είχε ένα σημάδι, οπότε ρώτησα τη μητέρα σου αν αυτό ήταν πάντα εκεί και αυτή είπε πως είχε χτυπήσει τον τοίχο εκνευρισμένη. Γιατί ρωτάς; Νόμιζα πως δεν το είχες παρατηρήσει»

«Ίσα ίσα... με προβλημάτισε αρκετά αυτό το σημάδι. Επειδή υπήρχε και στο όνειρο μου αυτό το σημάδι στον τοίχο, νόμιζα για λίγο πως βρισκόμουν ακόμη στο όνειρο, αλλά...»

«Αλλά ευτυχώς για σένα δεν είσαι. Αυτό δεν θα έλεγες;» γύρισε η Μπλαιρ και μου χαμογέλασε.

«Ναι...»

«Ου, κι άλλο κομμάτι ABBA» είπε η Μπλαιρ και άρχισε να σιγοτραγουδάει χωρίς να λέει ακριβώς τους στίχους του τραγουδιού.

«Ποιος το τραγουδάει αυτό το τραγούδι είπαμε;» την ρώτησε ο Κουπ.

«Οι ABBA, το είπα κυριολεκτικά πριν...»

«Καλύτερα να τους αφήσουμε να το τραγουδήσουν τότε...» είπε ο Κουπ και άφησε άλλο ένα μικρό γελάκι που με παράσυρε και εμένα να γελάσω.

«Άσε μας, ούτε να τραγουδήσω δεν μπορώ...» είπε η Μπλαιρ και δυνάμωσε το ράδιο καθώς ταυτόχρονα έστριβε προς το παράθυρο. «Καλή ατάκα πάντως, θα τη θυμάμαι για όταν αρχίσεις εσύ να τραγουδάς»

«Πότε δηλαδή;»

«Δεν ξέρω, όλο και κάποια στιγμή δεν θα τραγουδήσεις;»

«Μπα» είπε ο Κουπ διατηρώντας το χαμόγελο που είχε τόση ώρα στο πρόσωπο του.

Έστρεψα το βλέμμα μου προς την Μπλαιρ που καθόταν μπροστά μου και την είδα να χαμογελάει παρά όλα τα σχόλια που είχαν ειπωθεί.

«Γιατί γελάτε; Τώρα δεν τσακωνόσασταν;»

«Τι; Δεν τσακωνόμαστε» είπε ο Κουπ. «Φίλοι είμαστε, πειράζουμε λίγο ο ένας τον άλλον. Αν ήταν να τσακωνόμασταν με το παραμικρό πράγμα που έλεγε ο καθένας μας θα είχαμε ξεκόψει πολύ καιρό»

«Τι;» προβληματίστηκα.

«Μη σε νοιάζει. Σε ακούω τόσην ώρα γελάς, που σημαίνει ότι απολαμβάνεις το σόου»

«Ποιο σόου μωρέ» τον έσπρωξε πάλι παιχνιδιάρικα η Μπλαιρ, αφήνοντας και αυτή ένα μικρό γελάκι, παρασύροντας και τον Κουπ σε γέλια.

Πάλι γελάνε χωρίς λόγο. Είναι ακριβώς αυτό που σκεφτόμουν όταν ήμουν ακόμη στην πόλη: όσοι έχουν παρέα είναι χαρούμενοι, γελούσαν φαινομενικά χωρίς κάποιο λόγο, χωρίς να έχει γίνει ή να έχει ειπωθεί κάτι αστείο. Όσοι ήταν μόνοι τους βρίσκονταν στην ίδια κατάσταση που βρισκόμουν και εγώ, σκεπτόμενοι, χαμένοι στον κόσμο τους τον οποίο δεν τον μοιράζονταν με κάποιον άλλο γιατί δεν μπορούσαν, δεν είχαν εκείνη τη στιγμή κάποιον να μιλήσουν. Εκείνη η κοπέλα συγκεκριμένα θυμάμαι πώς όταν είχε έρθει η φίλη της άλλαξε κατευθείαν έκφραση, επειδή ακριβώς είχε κάποιον επιτέλους που μπορούσε να μιλήσει, κάποιον να την βοηθήσει, να την κάνει να γελάσει, να περάσει μαζί της τις όποιες δυσκολίες.

Εγώ από την άλλη τι έκανα; Στο όνειρο μου αλλά και στην πραγματικότητα είχα κλειστεί τελείως από τους υπόλοιπους νομίζοντας πως θα μπορούσα να περάσω τις δυσκολίες μου, υπαρκτές ή ανύπαρκτες, μόνος μου. Και αυτό με οδήγησε στην χειρότερη λύση. Βλέποντας όμως επιτέλους τους δύο καλύτερους μου φίλους να είναι δίπλα μου, να με βοηθάνε χωρίς να επιθυμούν κάτι άμεσα ως αντάλλαγμα, να χαμογελάνε παρόλο που κανένας από τους δύο δεν έχει πλέον οικογένεια και ο καλύτερος τους φίλος έχει απώλεια μνήμης, δεν μπορώ παρά να σκέφτομαι ότι ήμουν λάθος. Αν έπαιρνα τη ζωή μου... δεν θέλω να σκεφτώ τι θα μπορούσε να συμβεί στη μητέρα μου ή στους φίλους μου.

Τουλάχιστον πατέρα, αν με ακούς, και αν όντως είχες πει αυτές τις λέξεις, επιτέλους μπορώ να πω πως βρήκα απάντηση. Για το “που είμαι”, έχω να πω πως βρίσκομαι σε καλά χέρια, οπότε μην αγχώνεσαι για αυτό. Ακόμη και αν δεν με βοηθήσουν να ανακτήσω τις αναμνήσεις μου, θα μπορέσω να δημιουργήσω καινούργιες, επειδή έχω στο πλευρό μου την καλύτερη συντροφιά που θα μπορούσα να έχω, καθώς όσον αφορά το “ποιος είμαι”...

«Μπλαιρ...»

«Ναι Σπαρκ;»

«Είμαι όντως ο πιο τυχερός άνθρωπος του κόσμου»


\backmatter

%\chapter{Επίλογος}
%Το τέλος.
% ΤΟΔΟ: δεν ξέρω τι να το κάνω αυτό

\end{document}
